
\documentclass{article}
\usepackage[utf8]{inputenc}
\usepackage[T1]{fontenc}
\usepackage{lmodern}
\usepackage{hyperref}

\begin{document}

\section*{Résumé}
UNIVERSITÉ DE NICE - SOPHIA ANTIPOLIS
ÉCOLE DOCTORALE STIC
SCIENCES ET TECHNOLOGIES DE L’INFORMATION
ET DE LA COMMUNICATION
T H È S E
pour obtenir le titre de
Docteur en Sciences
de l’Universit\'{e} de Nice - Sophia Antipolis
Mention : INFORMATIQUE
Pr\'{e}sent\'{e}e et soutenue par
Andr\'{e} KALAWA
Migration des applications vers les tables
interactives par recherche d’\'{e}quivalences
Th\`{e}se dirig\'{e}e par Michel RIVEILL
pr\'{e}par\'{e}e au sein du laboratoire I3S, Equipe RAINBOW
soutenue le 30 ao\^{u}t 2013
Jury :
Rapporteurs :
Daniel HAGIMONT
Professeur, INPT/ENSEEIHT de Toulouse
Jacky ESTUBLIER
Professeur, LIG Grenoble
Directeur :
Michel RIVEILL
Professeur, Universit\'{e} de Nice Sophia Antipolis
Co Encadrante :
Audrey OCCELLO
Docteur, Universit\'{e} de Nice Sophia Antipolis
Pr\'{e}sident :
Mireille BLAY-FORNARINO
Professeur, Universit\'{e} de Nice Sophia Antipolis
2
R\'{e}sum\'{e}
Dans le domaine du g\'{e}nie logiciel pour les Interactions Homme Machine (IHM), la migration des interfaces
utilisateurs (UI) est un moyen pour r\'{e}utiliser des applications sur des plateformes ayant des modalit\'{e}s d’interactions diff\'{e}rentes des environnements de d\'{e}part. Les approches existantes de migration des UI sont manuelles
dans le cadre des approches sp\'{e}cifiques, elles sont automatiques dans le cadre des services d’adaptation des UI
aux contextes d’usage, ou elles sont semi automatiques dans le cadre d’une migration flexible dirig\'{e}e par un
concepteur.
Dans cette th\`{e}se nous nous int\'{e}ressons \`{a} la migration semi automatique des UI vers une cible comme une
table interactive dans l’objectif de transformer des UI Desktop en UI qui favorisent la collaboration et l’utilisation des objets tangibles. Les tables interactives sont des plateformes qui disposent des instruments d’interactions permettant de d\'{e}crire des UI tangibles et multi-utilisateurs. En consid\'{e}rant que le noyau fonctionnel
(NF) des applications de d\'{e}part peut \^{e}tre r\'{e}utilis\'{e}s sur les cibles sans changement, les UI des applications sont
caract\'{e}ris\'{e}es par la dimension des dialogues entre les utilisateurs et le syst\`{e}me, la dimension de la structure et
du positionnement des \'{e}l\'{e}ments graphiques et la dimension du style des \'{e}l\'{e}ments visuels. La migration d’une
UI dans ces conditions consiste \`{a} transformer ou \`{a} recr\'{e}er les diff\'{e}rentes dimensions d’une UI de d\'{e}part pour
la cible tout en consid\'{e}rant les crit\`{e}res de conception des UI pour les tables interactives.
Nous proposons dans cette th\`{e}se un mod\`{e}le d’interactions abstraites pour \'{e}tablir les \'{e}quivalences entre les
dialogues et la structure des UI ind\'{e}pendamment des modalit\'{e}s d’interactions des plateformes source et cible.
Les primitives d’interactions et la structure des composants graphiques permettent de d\'{e}crire des op\'{e}rateurs
d’\'{e}quivalences pour retrouver et classer les \'{e}l\'{e}ments graphiques \'{e}quivalents en prenant en compte les guidelines des tables interactives. Nous proposons aussi des r\`{e}gles de substitution et de concr\'{e}tisation pour accro\^{i}tre
l’accessibilit\'{e} des \'{e}l\'{e}ments graphiques et favoriser l’utilisation des objets tangibles.
Mots cl\'{e}s :
migration des interfaces utilisateurs, \'{e}quivalences des plateformes, modalit\'{e}s d’interactions,
crit\`{e}res de conception, guidelines
Abstract
In software engineering, in the field of human computer interaction (HCI), the migration of user interface
(UI) is a way to reuse existing applications on platforms with different interactions modalities. The existing
approaches for UI migration can be manual (for specific applications), they can be automatic (for services
which adapt UI based on context aware), or they can be mix of the previous - semi automatic (providing a
flexible migration process driven by the person in charge).
This thesis proposes a semi automatic process for migration of UI from a desktop to interactive table for the
purpose of transforming the UI of desktop to support further collaboration and usage of tangible objects. The
interactive tables are platforms with interactions instruments which allow the describtion of tangible and multi
users UIs. Considering that the functional core (FC) of source applications can be reused on target platform
without transformation, any UI can be characterized with three dimensions : the first dimension concerns the
dialogues between the users and the system, the second dimension concerns the structure and the layout of
graphical components, and the third dimension concerns the visual style of graphical elements. In this context,
the problematic regarding the UI migration is how to transform or re inject these different dimensions of source
UI into the target, while considering the UI design criteria for interactive tables.
This thesis proposes an abstract interactions model for establishing equivalences (independent of modalities
of interactions) between the source and the dialogue and structure of the target. The primitives of interaction and
the structure of graphical components are used to describe equivalence operators to find and to rank equivalent
elements on interactive tables. Furthermore, this thesis proposes substitution and concretization rules to increase
the accessibility of graphical elements and to facilitate the usage of tangible objects. The ranking process and
the transformation rules are based on guidelines for UI migration to interactive tables which are interpreted
form design criteria.
Keywords :
user interface migration, \'{e}quivalence of plateform, design criteria, guidelines
ii
Table des mati\`{e}res
I
Introduction
3
1
Introduction
5
1.1
Contexte . . . . . . . . . . . . . . . . . . . . . . . . . . . . . . . . . . . . . . . . .
5
1.2
Enjeux de la migration
. . . . . . . . . . . . . . . . . . . . . . . . . . . . . . . . .
6
1.3
Contribution de la th\`{e}se . . . . . . . . . . . . . . . . . . . . . . . . . . . . . . . . .
7
1.4
Plan du manuscrit . . . . . . . . . . . . . . . . . . . . . . . . . . . . . . . . . . . .
7
2
Migration des applications
9
2.1
Motivations . . . . . . . . . . . . . . . . . . . . . . . . . . . . . . . . . . . . . . .
9
2.2
Sc\'{e}nario \textbackslash{}& Probl\`{e}mes
. . . . . . . . . . . . . . . . . . . . . . . . . . . . . . . . .
10
2.2.1
Cas d’une application desktop . . . . . . . . . . . . . . . . . . . . . . . . .
10
2.2.2
Probl\`{e}mes li\'{e}s \`{a} la migration du dialogue . . . . . . . . . . . . . . . . . . .
12
2.2.3
Probl\`{e}mes li\'{e}s \`{a} la migration de la structure et du positionnement
. . . . . .
13
2.2.4
Probl\`{e}mes li\'{e}s \`{a} la migration du style . . . . . . . . . . . . . . . . . . . . .
15
2.2.5
Espace des probl\`{e}mes li\'{e}s la migration vers les tables interactives . . . . . .
16
2.3
P\'{e}rim\`{e}tres de la migration des UI vers les tables interactives
. . . . . . . . . . . . .
17
2.3.1
Hypoth\`{e}se de travail . . . . . . . . . . . . . . . . . . . . . . . . . . . . . .
17
2.3.2
Pistes de migration . . . . . . . . . . . . . . . . . . . . . . . . . . . . . . .
18
2.4
Objectifs . . . . . . . . . . . . . . . . . . . . . . . . . . . . . . . . . . . . . . . . .
19
II
Domaine d’\'{e}tude
21
3
Tables interactives et Migration des UI
23
3.1
Mod\`{e}le d’interactions d’une table interactive
. . . . . . . . . . . . . . . . . . . . .
23
3.1.1
Les dispositifs mat\'{e}riels d’interactions d’une table interactive
. . . . . . . .
24
3.1.2
Biblioth\`{e}ques graphiques des tables interactives . . . . . . . . . . . . . . . .
27
3.1.3
Modalit\'{e}s d’interactions . . . . . . . . . . . . . . . . . . . . . . . . . . . .
29
3.1.4
Mod\`{e}le d’interactions abstraites pour la migration des UI . . . . . . . . . . .
31
3.2
Principes de conception des UI pour les tables interactives
. . . . . . . . . . . . . .
33
3.2.1
Propri\'{e}t\'{e}s caract\'{e}ristiques des tables interactives . . . . . . . . . . . . . . .
33
3.2.2
Crit\`{e}res ergonomiques de conception des UI
. . . . . . . . . . . . . . . . .
35
3.2.3
Recommandations pour la migration des UI vers les tables interactives . . . .
38
3.3
Synth\`{e}se . . . . . . . . . . . . . . . . . . . . . . . . . . . . . . . . . . . . . . . . .
43
4
Approches de migration des UI
45
4.1
Introduction . . . . . . . . . . . . . . . . . . . . . . . . . . . . . . . . . . . . . . .
45
4.2
Approches sp\'{e}cifiques de migration des UI
. . . . . . . . . . . . . . . . . . . . . .
46
4.2.1
Migration manuelle . . . . . . . . . . . . . . . . . . . . . . . . . . . . . . .
46
4.2.2
Portage des applications existantes sur des tables interactives . . . . . . . . .
48
4.3
Approches de migration bas\'{e}es sur les mod\`{e}les de l’UI . . . . . . . . . . . . . . . .
51
4.3.1
Approches de migration automatiques des UI . . . . . . . . . . . . . . . . .
52
4.3.2
Approche semi automatique de migration des UI . . . . . . . . . . . . . . .
57
iii
iv
TABLE DES MATIÈRES
4.4
Synth\`{e}se et objectifs
. . . . . . . . . . . . . . . . . . . . . . . . . . . . . . . . . .
62
4.4.1
Synth\`{e}se
. . . . . . . . . . . . . . . . . . . . . . . . . . . . . . . . . . . .
62
4.4.2
Objectifs
. . . . . . . . . . . . . . . . . . . . . . . . . . . . . . . . . . . .
64
III
M\'{e}thodes
65
5
Mod\'{e}lisation des UI
67
5.1
Introduction . . . . . . . . . . . . . . . . . . . . . . . . . . . . . . . . . . . . . . .
67
5.2
Primitives d’interactions
. . . . . . . . . . . . . . . . . . . . . . . . . . . . . . . .
68
5.2.1
Les primitives d’interactions en entr\'{e}e . . . . . . . . . . . . . . . . . . . . .
68
5.2.2
Les primitives d’interactions en sortie . . . . . . . . . . . . . . . . . . . . .
72
5.3
Mod\`{e}les de composants graphiques
. . . . . . . . . . . . . . . . . . . . . . . . . .
73
5.3.1
Un mod\`{e}le de types de composants graphiques . . . . . . . . . . . . . . . .
75
5.3.2
Un mod\`{e}le d’instance d’une UI
. . . . . . . . . . . . . . . . . . . . . . . .
81
5.3.3
Synth\`{e}se des mod\`{e}les abstraits . . . . . . . . . . . . . . . . . . . . . . . . .
90
5.4
Op\'{e}rateurs d’\'{e}quivalences
. . . . . . . . . . . . . . . . . . . . . . . . . . . . . . .
91
5.4.1
Op\'{e}rateurs d’\'{e}quivalences . . . . . . . . . . . . . . . . . . . . . . . . . . .
91
5.4.2
Op\'{e}rateurs d’\'{e}quivalences \textbackslash{}& types de donn\'{e}es
. . . . . . . . . . . . . . . .
95
5.5
Synth\`{e}se . . . . . . . . . . . . . . . . . . . . . . . . . . . . . . . . . . . . . . . . .
98
6
M\'{e}canismes de migration des UI
99
6.1
Introduction . . . . . . . . . . . . . . . . . . . . . . . . . . . . . . . . . . . . . . .
99
6.2
Prise en compte des guidelines . . . . . . . . . . . . . . . . . . . . . . . . . . . . .
100
6.2.1
Conception assist\'{e}e des UI . . . . . . . . . . . . . . . . . . . . . . . . . . .
101
6.2.2
Interpr\'{e}tation des guidelines pour la migration
. . . . . . . . . . . . . . . .
102
6.2.3
Utilisation effective des guidelines . . . . . . . . . . . . . . . . . . . . . . .
104
6.3
Transformations du mod\`{e}le de l’UI source . . . . . . . . . . . . . . . . . . . . . . .
108
6.3.1
Transformation des groupes d’\'{e}l\'{e}ments graphiques . . . . . . . . . . . . . .
108
6.3.2
Transformation d’un \'{e}l\'{e}ment . . . . . . . . . . . . . . . . . . . . . . . . . .
119
6.4
Classement des \'{e}l\'{e}ments \'{e}quivalents . . . . . . . . . . . . . . . . . . . . . . . . . .
124
6.4.1
Conformit\'{e} des composants graphiques aux guidelines . . . . . . . . . . . .
124
6.4.2
Charge de travail . . . . . . . . . . . . . . . . . . . . . . . . . . . . . . . .
127
6.4.3
Algorithme de classement
. . . . . . . . . . . . . . . . . . . . . . . . . . .
129
6.5
Synth\`{e}se . . . . . . . . . . . . . . . . . . . . . . . . . . . . . . . . . . . . . . . . .
133
IV
Exp\'{e}rimentations
135
7
Validations du processus
137
7.1
Introduction . . . . . . . . . . . . . . . . . . . . . . . . . . . . . . . . . . . . . . .
137
7.2
Impl\'{e}mentation du processus de migration . . . . . . . . . . . . . . . . . . . . . . .
138
7.2.1
Abstraction de l’UI source . . . . . . . . . . . . . . . . . . . . . . . . . . .
139
7.2.2
Proposition d’une version des UI migr\'{e}es . . . . . . . . . . . . . . . . . . .
143
7.2.3
Personnalisation des UI propos\'{e}es . . . . . . . . . . . . . . . . . . . . . . .
145
7.2.4
G\'{e}n\'{e}ration de l’UI finale . . . . . . . . . . . . . . . . . . . . . . . . . . . .
148
7.2.5
Remarques . . . . . . . . . . . . . . . . . . . . . . . . . . . . . . . . . . .
149
7.3
Évaluation . . . . . . . . . . . . . . . . . . . . . . . . . . . . . . . . . . . . . . . .
149
TABLE DES MATIÈRES
1
7.3.1
Crit\`{e}res d’\'{e}valuation de l’approche propos\'{e}e
. . . . . . . . . . . . . . . . .
150
7.3.2
R\'{e}sultats
. . . . . . . . . . . . . . . . . . . . . . . . . . . . . . . . . . . .
151
7.3.3
Interpr\'{e}tation des r\'{e}sultats . . . . . . . . . . . . . . . . . . . . . . . . . . .
154
V
Conclusions et Perspectives
157
8
Conclusions et perspectives
159
8.1
Synth\`{e}se . . . . . . . . . . . . . . . . . . . . . . . . . . . . . . . . . . . . . . . . .
159
8.2
Perspectives . . . . . . . . . . . . . . . . . . . . . . . . . . . . . . . . . . . . . . .
161
8.2.1
A moyen terme . . . . . . . . . . . . . . . . . . . . . . . . . . . . . . . . .
161
8.2.2
A long terme . . . . . . . . . . . . . . . . . . . . . . . . . . . . . . . . . .
161
Appendices
169
A Application des mod\`{e}les de l’UI et des PI
171
A.1
Comment d\'{e}crire une instance du mod\`{e}le de type de composants graphiques ? . . . .
171
A.1.1
Identification d’un Widget \`{a} partir d’un Assembly
. . . . . . . . . . . . . .
171
A.1.2
Identification des comportements
. . . . . . . . . . . . . . . . . . . . . . .
171
A.1.3
Remarque . . . . . . . . . . . . . . . . . . . . . . . . . . . . . . . . . . . .
172
A.2
Abstraction des UI finales . . . . . . . . . . . . . . . . . . . . . . . . . . . . . . . .
172
A.2.1
Identification des UIComponent \`{a} partir d’un fichier XAML . . . . . . . . .
172
A.2.2
Identification des Containers \`{a} partir d’un fichier XAML . . . . . . . . . . .
173
A.2.3
Identification de la classe Content \`{a} partir d’un fichier XAML . . . . . . . .
173
A.2.4
Identification de la classe ImplementedEvent \`{a} partir d’un fichier XAML
. .
174
A.2.5
Exemple d’UIStructure XAML
. . . . . . . . . . . . . . . . . . . . . . . .
174
A.2.6
Exemple de UIBehavior C\textbackslash{}#
. . . . . . . . . . . . . . . . . . . . . . . . . .
176
A.2.7
Exemple d’AbstractView C\textbackslash{}# . . . . . . . . . . . . . . . . . . . . . . . . . .
177
A.3
Primitives d’interactions
. . . . . . . . . . . . . . . . . . . . . . . . . . . . . . . .
180
A.3.1
Liste des Widgets . . . . . . . . . . . . . . . . . . . . . . . . . . . . . . . .
180
2
TABLE DES MATIÈRES
Premi\`{e}re partie
Introduction
3
CHAPITRE 1
Introduction
1.1
Contexte
Le nombre grandissant de plateformes et surtout leurs tr\`{e}s grandes vari\'{e}t\'{e}s de dispositifs d’interactions dont ils disposent telles que les smartphones, les tablettes, les terminaux tactiles ou les tables
interactives ont g\'{e}n\'{e}ralis\'{e} l’usage des applications ayant des modalit\'{e}s d’interactions [NC93] beaucoup plus sophistiqu\'{e}es qu’une simple utilisation du clavier et de la souris. Les tables interactives par
exemple offrent la possibilit\'{e} de mettre en œuvre des applications r\'{e}ellement multi-utilisateurs m\^{e}lant
interactions tactiles et objets tangibles 1 sur une surface de pr\^{e}t de 1m2.
Le domaine du g\'{e}nie logiciel pour les Interactions Homme Machine (IHM) propose plusieurs
approches pour la conception d’une application. En particulier la tendance actuelle est de concevoir les Interfaces Utilisateurs (UI) selon les mod\`{e}les du Framework de R\'{e}f\'{e}rence Cameleon
(CRF) [CCB+02] identifiant diff\'{e}rents niveau d’interface de la plus abstraite, ind\'{e}pendante des dispositifs d’interactions disponibles sur la plateforme \`{a} un niveau concret permettant la mise en œuvre
de l’UI selon la modalit\'{e} d’interaction souhait\'{e}e disponible sur la plateforme cible. Si les UI de l’application sont construites selon cette approche, alors la migration des applications entre terminaux
ayant des modalit\'{e}s diff\'{e}rents a d\'{e}j\`{a} \'{e}t\'{e} partiellement abord\'{e} en particulier dans [Pat11]. Notre travail
aborde la possibilit\'{e} de faire migrer entre plateformes ayant des modalit\'{e}s d’interactions diff\'{e}rentes
des applications qui n’ont pas \'{e}t\'{e} pr\'{e}vue selon les diff\'{e}rents mod\`{e}les de r\'{e}f\'{e}rence de Cameleon.
La migration des UI est une activit\'{e} de g\'{e}nie logiciel qui implique la transformation des diff\'{e}rents
aspects qui constituent une UI existante tels que les interactions (ou le dialogue entre l’utilisateur
et le syst\`{e}me) qu’il faut n\'{e}cessairement pr\'{e}server pour que l’utilisateur puisse toujours accomplir
les m\^{e}mes t\^{a}ches, les structures (organisations et types de donn\'{e}es des \'{e}l\'{e}ments graphiques), le positionnement et les styles des composants graphiques qui doivent \^{e}tre adapt\'{e}s pour \^{e}tre conformes
aux sp\'{e}cificit\'{e}s de la plateforme cible et toujours satisfaire les utilisateurs finaux dans les choix de
configuration qui ont put \^{e}tre fait sur la source
Au del\`{a} des probl\`{e}mes li\'{e}s \`{a} des diff\'{e}rences possibles entre les environnements d’ex\'{e}cution des
plateformes source et cible [TKB78] qui n\'{e}cessite le portage du code, les diff\'{e}rences des modalit\'{e}s
d’interactions entre la source et la cible impliquent n\'{e}cessairement la prise en compte des crit\`{e}res
usuels de conception des UI. La transformation de l’UI de d\'{e}part doit \^{e}tre guid\'{e}e non seulement par
les crit\`{e}res ergonomiques de conception mais aussi par les dispositifs d’interactions disponible sur la
plateforme cible qu’il peut \^{e}tre int\'{e}ressant d’utiliser.
Évidemment, en compl\'{e}ment de la volont\'{e} de rendre disponible des dispositifs d’interactions
nouveaux, peut-\^{e}tre plus intuitifs, les crit\`{e}res ergonomiques de conception constituent un ensemble de
principes \`{a} respecter pendant la mise en œuvre. Ils permettent de garantir l’utilisabilit\'{e} d’une UI. Par
exemple, l’utilisation d’une application pour desktop sur une table interactive sans aucune adaptation
pose des probl\`{e}mes d’utilisabilit\'{e}, car la simulation du clavier et de la souris n’est pas le meilleur
1. Les interfaces utilisateurs tangibles (TUI) permettent \`{a} une application de pouvoir interagir avec des objets physiques
directement manipulable par les utilisateurs [IU97]. Les objets tangibles sont les objets pouvant \^{e}tre manipul\'{e} par une
interface utilisateur tangible.
5
6
CHAPITRE 1. INTRODUCTION
mode d’interaction en terme d’utilisabilit\'{e}. Si l’application le permet, on peut m\^{e}me imaginer une
utilisation multi-utilisateurs avec une nouvelle UI \`{a} d\'{e}duire de l’UI d’origine.
1.2
Enjeux de la migration
Dans le cadre de la migration des applications pour desktops vers les tables interactives par
exemple, nous notons que les dialogues, la structure et le positionnement des \'{e}l\'{e}ments des UI graphiques doivent prendre en compte l’\'{e}ventualit\'{e} d’avoir plusieurs utilisateurs mais surtout la possibilit\'{e} d’utiliser des objets tangibles. Travailler sur ces deux exemples va nous permettre d’avoir une
r\'{e}elle \'{e}volutivit\'{e} des UI entre le terminal source et le terminal cible et permet de favoriser la collaboration dans un espace de travail co-localis\'{e} et multi-utilisateurs offert par la table interactive et donc
de l’utiliser pleinement.
Nous faisons l’hypoth\`{e}se que le cœur de l’application peut \^{e}tre migr\'{e} sans modification et qu’il ne
faut agir qu’au niveau de l’UI qu’il faut bien \'{e}videmment faire migrer. Cette migration, d’un terminal
source vers un terminal cible donn\'{e}e doit permettre de conserver dans la mesure du possible les
dialogues, la structure et les positionnements relatifs des \'{e}l\'{e}ments de l’interface, le respect du style
des \'{e}l\'{e}ments graphiques des UI de la source mais bien \'{e}videmment proposer une adaptation pour
utiliser au mieux, sans perdre les utilisateurs, les nouveaux moyens d’interaction mis \`{a} la disposition
des utilisateurs.
La transformation des dialogues d’une UI desktop vers une cible collaborative peut-elle garantir
\`{a} chaque utilisateur des interactions qui favorisent la collaboration ? De nombreux \'{e}l\'{e}ments sont \`{a}
prendre en compte. En effet, chaque dialogue ou message d’erreur ne doit pas perturber les activit\'{e}s
des autres utilisateurs ; par exemple les “feedbacks” ou les messages d’erreurs destin\'{e}s \`{a} un utilisateur
ne doivent pas bloquer un autre utilisateur si une r\'{e}ponse est attendue. Par ailleurs la transformation
des dialogues peut-elle assurer que chaque dialogue reste coh\'{e}rent avec l’intention et les actions des
utilisateurs ? Par exemple la modification d’une valeur par deux utilisateurs distincts peut les perturber.
Il faut en effet d\'{e}cider si le r\'{e}sultat est la somme des deux modifications, la moyenne, le min ou le
max. Évidemment le contexte et la nature de la valeur peuvent guider sur le choix \`{a} effectuer.
La transformation des dialogues sur les UI et les fonctionnalit\'{e}s de l’application source peuvent
provoquer de nombreux changements ? En effet, la transformation des dialogues des UI desktops
pour les tables interactives peuvent par exemple impliquer des modifications pour prendre en compte
la pr\'{e}sence de plusieurs utilisateurs ou d’objets tangibles.
La transformation de la structure et du positionnement des \'{e}l\'{e}ments d’une UI pour desktop vers
une table interactive consiste \`{a} assurer l’accessibilit\'{e} aux diff\'{e}rents utilisateurs des \'{e}l\'{e}ments graphiques pertinents 2. Par ailleurs la transformation des aspects structurels d’une UI doit elle aussi
pr\'{e}server au mieux l’utilisabilit\'{e} de l’UI. En effet les solutions de transformations propos\'{e}es doivent
produire des UI conformes aux sp\'{e}cificit\'{e}s de la plateforme cible et satisfaisantes pour les utilisateurs
finaux.
Pour terminer, la transformation du style a n\'{e}cessairement un impact sur la collaboration ou sur
l’utilisation des objets tangibles. G\'{e}n\'{e}ralement, chaque application adopte une charte graphique qu’il
faut soit respecter, soit r\'{e}nover.
Les questions soulev\'{e}es par les transformations des diff\'{e}rents aspects des UI lors d’une migration
sont usuellement trait\'{e}es de diff\'{e}rentes mani\`{e}res. En premier lieu, il est possible d’adopter une approche manuelle [WGM08]. Celle-ci permet d’avoir, au prix d’un travail minutieux, une UI conforme
aux attentes des utilisateurs finaux car les transformations sont flexibles et donc la qualit\'{e} de l’UI
produite d\'{e}pend directement des comp\'{e}tences de celui qui effectue la migration manuelle. Mais cette
2. Un menu par exemple
1.3. CONTRIBUTION DE LA THÈSE
7
approche est difficilement r\'{e}utilisable, \`{a} moins de la consigner dans un document, que ce soit pour
d’autres applications ou pour d’autres personnes car elle n\'{e}cessite une bonne connaissance des crit\`{e}res de transformation des diff\'{e}rents aspects et des technologies des UI des plateformes source et
cible.
Il est aussi possible d’adopter une approche automatique bas\'{e}e ou non sur des mod\`{e}les abstraits [Bes10, PZ10] qui permet une transformation des diff\'{e}rents aspects des UI. Bien que ces approches automatiques soient r\'{e}utilisables, elles sont n\'{e}cessairement moins flexibles car le concepteur
qui ne peut pas intervenir pendant la transformation, adapte l’UI produite dans un second temps.
Entre ces deux approches, il est aussi possible d’adopter une approche semi automatique[MR97]
qui pr\'{e}sente de nombreux inconv\'{e}nients car non seulement elle induit un travail suppl\'{e}mentaire pour
la personne en charge de la migration qui guide la transformation mais en permettant plus de flexibilit\'{e},
il est assez difficile de capitaliser sur le travail effectuer si la personne en charge de la migration
change. N\'{e}anmoins, cette approche permet une transformation interactive des diff\'{e}rents aspects de
l’UI. L’int\'{e}r\^{e}t d’une flexibilit\'{e} dans l’approche de migration d’UI permet d’avoir des UI migr\'{e}es
proches des attentes des utilisateurs finaux.
1.3
Contribution de la th\`{e}se
Cette th\`{e}se a pour premier objet d’\'{e}tudier la migration des UI entre plateformes ayant des dispositifs d’interactions diff\'{e}rents. Nous avons fait le choix de cibler en particulier la migration des
UI depuis une station poss\'{e}dant clavier et souris vers une table interactive. Nous souhaitons que
cette migration se fasse au moindre co\^{u}t pour les concepteurs ou les d\'{e}veloppeurs tout en prenant en
compte les sp\'{e}cificit\'{e}s de la cible et, bien \'{e}videmment, les crit\`{e}res ergonomiques usuels utilis\'{e}s pour
la conception des UI.
La solution que nous proposons repose sur un processus semi automatique de migration d’UI.
Celui-ci comporte plusieurs \'{e}tapes interactives pour que les les concepteurs ou les d\'{e}veloppeurs
puissent effectuer des choix conformes aux crit\`{e}res de conception qu’ils souhaitent privil\'{e}gi\'{e}s. Nous
avons fait le choix de l’interactivit\'{e}, bas\'{e} sur des choix simples pour garantir simultan\'{e}ment la r\'{e}utilisabilit\'{e}, minimiser les co\^{u}ts mais aussi pour accro\^{i}tre la flexibilit\'{e} et garantir ainsi l’UI la plus
pertinente.
Il est primordial de prendre en compte, lors de la migration des interfaces utilisateurs, les crit\`{e}res ergonomiques de conception. Nos travaux utilisent des concepts issus de plusieurs domaines de
recherche.
Dans le domaine de l’utilisabilit\'{e}, nous nous sommes int\'{e}ress\'{e}s aux travaux qui mod\'{e}lisent les
crit\`{e}res ergonomiques de conception pour les traduire en r\`{e}gles op\'{e}rationnelles et utilisables pendant
la conception. L’objectif est de pouvoir int\'{e}grer ces r\`{e}gles dans la plateforme de migration dans le but
de r\'{e}duire la charge de travail des personnes en charge de la migration.
Dans le domaine de l’ing\'{e}nierie des mod\`{e}les, nous nous sommes int\'{e}ress\'{e}s aux travaux permettant
de mod\'{e}liser une plateforme dans le but d’effectuer la migration \`{a} un niveau abstrait : de concept
\`{a} concept. L’objectif est de rendre notre travail r\'{e}utilisable si la source et la cible \'{e}voluent. Cette
approche nous permet d’abstraire non seulement les UI mais aussi les interactions. Nous proposons
ainsi un mod\`{e}le d’interactions bas\'{e}es sur les primitives d’interactions[KODPR12] pour d\'{e}crire les
actions atomiques qui constituent les dialogues entre l’utilisateur et le syst\`{e}me.
1.4
Plan du manuscrit
Notre manuscrit est structur\'{e} de la mani\`{e}re suivante :
8
CHAPITRE 1. INTRODUCTION
Le chapitre 2 pr\'{e}sente l’espace des probl\`{e}mes li\'{e}s \`{a} la transformations des diff\'{e}rents aspects d’une
UI. Il d\'{e}limite le p\'{e}rim\`{e}tre de la migration des UI d’un poste fixe vers une table interactive. Il fixe nos
objectifs.
Le chapitre 3 pr\'{e}sente une \'{e}tude du mod\`{e}le d’interactions des tables interactives dans le but
d’identifier les sp\'{e}cificit\'{e}s et les crit\`{e}res ergonomiques de conception \`{a} int\'{e}grer pour la migration
des UI vers cette cible.
Le chapitre 4 est un \'{e}tat de l’art des approches de migration des UI. Dans ce chapitre nous d\'{e}crivons les crit\`{e}res n\'{e}cessaires pour atteindre nos objectifs et nous \'{e}valuons les diff\'{e}rentes approches \`{a}
l’aide de ces crit\`{e}res. Ce chapitre se termine par une synth\`{e}se des diff\'{e}rentes approches pr\'{e}sent\'{e}es.
Le chapitre 5 propose un mod\`{e}le d’UI qui prend en compte deux aspects des UI : leurs interactions
et leur structure. Les objectifs de ce mod\`{e}le d’UI sont de d\'{e}crire les UI \`{a} migrer ind\'{e}pendamment des
plateformes mais aussi de d\'{e}crire des op\'{e}rateurs d’\'{e}quivalences entre les dispositifs d’interactions des
plateformes source et cible.
Le chapitre 6 pr\'{e}sente les m\'{e}canismes de transformation de la structure et des interactions de l’UI
source vers la cible. L’objectif de ce chapitre est de d\'{e}crire la prise en compte effective des crit\`{e}res
ergonomiques de conception sous forme de guidelines par les m\'{e}canismes de transformation.
Le chapitre 7 pr\'{e}sente le prototype que nous avons r\'{e}alis\'{e}. Il constitue une preuve de concept des
m\'{e}canismes propos\'{e}s. Plusieurs applications ont \'{e}t\'{e} migr\'{e}es afin de valider les \'{e}l\'{e}ments des diff\'{e}rents
mod\`{e}les, mettre en \'{e}vidence le respect des crit\`{e}res ergonomiques et surtout mettre en \'{e}vidence les
b\'{e}n\'{e}fices que l’on peut retirer d’une telle approche.
Le chapitre 8 est une conclusion. Elle rappelle les objectifs que nous souhaitions atteindre, met
en \'{e}vidence nos contributions et leurs apports. Elle donne aussi quelques \'{e}l\'{e}ments sur des travaux
compl\'{e}mentaires qui pourraient \^{e}tre men\'{e}s pour parfaire nos travaux.
CHAPITRE 2
Migration des applications vers les tables
interactives
Sommaire
2.1
Motivations . . . . . . . . . . . . . . . . . . . . . . . . . . . . . . . . . . . . . .
9
2.2
Sc\'{e}nario \textbackslash{}& Probl\`{e}mes . . . . . . . . . . . . . . . . . . . . . . . . . . . . . . . .
10
2.2.1
Cas d’une application desktop . . . . . . . . . . . . . . . . . . . . . . . .
10
2.2.2
Probl\`{e}mes li\'{e}s \`{a} la migration du dialogue . . . . . . . . . . . . . . . . . .
12
2.2.3
Probl\`{e}mes li\'{e}s \`{a} la migration de la structure et du positionnement . . . . .
13
2.2.4
Probl\`{e}mes li\'{e}s \`{a} la migration du style . . . . . . . . . . . . . . . . . . . .
15
2.2.5
Espace des probl\`{e}mes li\'{e}s la migration vers les tables interactives . . . . .
16
2.3
P\'{e}rim\`{e}tres de la migration des UI vers les tables interactives . . . . . . . . . .
17
2.3.1
Hypoth\`{e}se de travail . . . . . . . . . . . . . . . . . . . . . . . . . . . . .
17
2.3.2
Pistes de migration . . . . . . . . . . . . . . . . . . . . . . . . . . . . . .
18
2.4
Objectifs . . . . . . . . . . . . . . . . . . . . . . . . . . . . . . . . . . . . . . .
19
D
ANS le but de d\'{e}crire les diff\'{e}rentes \'{e}tapes et les probl\`{e}mes lier \`{a} la migration des interfaces
Hommes-Machines en g\'{e}n\'{e}ral, et vers une table interactive en particulier, nous avons pris le
parti d’utiliser un exemple “fil rouge”. Les motivations du choix de l’\'{e}tude de cette cible sont d\'{e}crites
dans la section 2.1, puis la section 2.2 d\'{e}crit pas \`{a} pas l’application initialement con\c{c}ue pour un
desktop que nous souhaitons mettre en œuvre une table surface. La section 2.4 pr\'{e}sente les objectifs
de cette th\`{e}se.
2.1
Motivations
Les tables interactives comme les desktops, les smartphones ou les tablettes peuvent \^{e}tre utilis\'{e}es
dans divers domaines d’activit\'{e}s [KLL+09]. Une table interactive peut servir comme plateforme de
jeux, pour la consultation de photos, de cartes, de dessins, pour noter les id\'{e}es lors d’une session
de brainstorming, etc. Les tables interactives permettent de mettre en œuvre un espace de travail
collaboratif.
N\'{e}anmoins, on remarque que les tables interactives ne se sont pas d\'{e}mocratis\'{e}es comme les tablettes ou les smartphones. Évidemment, le prix de vente en est une des raisons peut constituer un
frein \`{a} l’achat. Il est par exemple d’environ 6600AC pour la version 2.0 commercialis\'{e}es par Microsoft.
M\^{e}me si nous observons une nette baisse du prix par rapport \`{a} la premi\`{e}re version celui-ci constitue
indiscutablement un frein \`{a} l’achat. Par ailleurs, la faible diffusion de cette plateforme fait que les
applications pour les tables interactives sont loin d’\^{e}tre aussi nombreuses que pour les tablettes ou
les smartphones. Selon des donn\'{e}es provenant de Distimo 3, en mars 2012 nous avons not\'{e} plus de
350.000 applications d\'{e}velopp\'{e}es aux Etats-Unis pour des desktops, 150.000 applications pour les
3. Distimo : www.distimo.com
9
10
CHAPITRE 2. MIGRATION DES APPLICATIONS
smartphones et presque 100.000 sp\'{e}cifiquement pour les tablettes 4 (cf figure 2.1) et aucune pour les
tables interactives.
Pour accro\^{i}tre le nombre d’applications disponibles et toucher un \'{e}ventuel public, nous pensons
que la migration des applications existantes peut \^{e}tre une piste s\'{e}rieuse si la migration utilise effectivement les divers dispositifs d’interactions disponible sur la table. Dans cette th\`{e}se, nous \'{e}tudions
par cons\'{e}quent les divers probl\`{e}mes pos\'{e}s par le processus de migrations en esp\'{e}rant pouvoir en
compl\'{e}ment enrichir les applications d’une dimension collaborative et tangible.
Nous pr\'{e}sentons donc \`{a} la section 2.2 les probl\`{e}mes li\'{e}s \`{a} la migration d’une application desktop
vers une table interactive puis les objectifs de nos travaux. Les hypoth\`{e}ses et les diff\'{e}rentes approches
permettant de faire migrer des UI sont pr\'{e}sent\'{e}s dans la section 2.3. L’ensemble de ce chapitre nous
permet de d\'{e}crire le p\'{e}rim\`{e}tre de nos travaux.
FIGURE 2.1 – Applications disponibles sur les march\'{e}s de t\'{e}l\'{e}chargement par plateformes en mars
2012 aux États Unis, source Distimo
2.2
Probl\`{e}mes li\'{e}s \`{a} la migration des UI vers une table interactive
Cette section d\'{e}crit le fonctionnement d’une application permettant l’\'{e}laboration de bande dessin\'{e}e que nous souhaitons faire migrer vers une table interactive. Dans un second temps, nous pr\'{e}sentons les probl\`{e}mes li\'{e}s \`{a} cette migration et les diff\'{e}rents objectifs que nous souhaitons atteindre sont
d\'{e}crits \`{a} la fin de cette section.
2.2.1
Cas d’une application desktop
Consid\'{e}rons l’exemple suivant : soit une application con\c{c}ue pour un desktop qui permet l’\'{e}laboration des bandes dessin\'{e}es (BD) telle que Comics Book Application (CBA). Cette application est
4. En consid\'{e}rant Windows Phone 7 Marketplace, Google Play, Nokia Ovi Store, Apple Mac App Iphone, Apple App
Store - iPad, Apple Mac App Store, Amazon Appstore et BlackBerry App World
2.2. SCÉNARIO \textbackslash{}& PROBLÈMES
11
con\c{c}ue pour \^{e}tre utilis\'{e}e par un dessinateur de BD qui est assis devant son \'{e}cran en utilisant sa souris
et son clavier. La fen\^{e}tre principale de l’application CBA d\'{e}crite par la figure 2.2 permet d’identifier
trois zones majeures qui sont la zone des menus correspondant aux composants graphiques situ\'{e}s en
haut de la fen\^{e}tre principale, la zone d’espace de travail qui contient les cadres d’une page de bande
dessin\'{e}e et la zone de l\'{e}gende correspondant aux groupes, formulaires et listes \`{a} gauche de la fen\^{e}tre.
FIGURE 2.2 – Description de la fen\^{e}tre principale de l’application CBA
Nous consid\'{e}rons de mani\`{e}re globale que les applications \`{a} migrer doivent respecter une d\'{e}composition minimale permettant de s\'{e}parer ais\'{e}ment l’UI et le Noyau Fonctionnel (NF). Si cette
hypoth\`{e}se est v\'{e}rifi\'{e}e, les diff\'{e}rents aspects des UI des applications peuvent \^{e}tre d\'{e}crites de mani\`{e}re
g\'{e}n\'{e}rale comme suit :
– le dialogue entre l’utilisateur et le syst\`{e}me permet d’organiser les interactions et les comportements d’une UI graphique,
– la structure et le positionnement des \'{e}l\'{e}ments graphiques qui d\'{e}crivent l’organisation et les
types de donn\'{e}es d’une UI,
– le style pour d\'{e}crire les tailles, les couleurs et les polices d’une UI graphique.
Nous consid\'{e}rons ces trois aspects comme autant de dimensions qui permettent de construire
l’espace des probl\`{e}mes li\'{e} \`{a} la migration d’une UI pr\'{e}vue pour un desktop pour que l’application
s’ex\'{e}cute sur une table interactive. Pour chacune de ces dimensions nous \'{e}tudions les probl\`{e}mes soulev\'{e}s par l’utilisation des objets tangibles/physiques comme moyen d’interactions. Dans le cadre des
tables interactives, une interaction tangible consiste \`{a} utiliser des objets physiques (g\'{e}n\'{e}ralement
en contact avec l’\'{e}cran ou la surface d’affichage) comme des dispositifs de manipulation directes ou
comme des moyens d’acc\`{e}s \`{a} des fonctionnalit\'{e}s. Pour une application “conception de bandes dessin\'{e}es” (CBA) sur une table interactive, il est envisageable d’utiliser un stylo pour \'{e}crire dans les bulles
de dialogue d’une bande dessin\'{e}e. Dans ce cas le stylo est un objet tangible de manipulation directe.
Par ailleurs, il est \'{e}galement possible que les concepteurs des bandes dessin\'{e}es souhaitent utiliser un
cube comme objet tangible dans le but s\'{e}lectionner des images (repr\'{e}sentant des personnages par
exemple). Chaque face du cube sera associ\'{e}e \`{a} une cat\'{e}gorie de liste d’images ou d’objets graphiques
utilisables dans une bande dessin\'{e}e.
Nous \'{e}tudions aussi les probl\`{e}mes soulev\'{e}s par la transformation d’une UI mono-utilisateur en UI
multi-utilisateurs co-localis\'{e}es. Dans le cadre des tables interactives, les applications peuvent naturellement \^{e}tre utilis\'{e}es par plusieurs personnes ce qui est plus difficilement imaginable sur un smartphone. Il est alors n\'{e}cessaire de re-concevoir l’UI de l’application qui est alors destin\'{e}e \`{a} plusieurs
12
CHAPITRE 2. MIGRATION DES APPLICATIONS
personnes (UI multi-utilisateurs) mais qui partagent le m\^{e}me \'{e}cran (utilisateur co-localis\'{e}). L’application “conception de bandes dessin\'{e}es” que l’on souhaite faire migrer sur table interactive peut \^{e}tre
utilis\'{e}e simultan\'{e}ment par plusieurs cr\'{e}ateurs de bandes dessin\'{e}es qui ainsi collaborent pour concevoir ensemble une nouvelle BD.
Dans les sections 2.2.2, 2.2.3 et 2.2.4, nous illustrons ces probl\`{e}mes pour les diff\'{e}rents aspects de
l’application CBA d\'{e}crite par la figure 2.2 ci-dessus.
2.2.2
Probl\`{e}mes li\'{e}s \`{a} la migration du dialogue
Le dialogue est une dimension des UI de d\'{e}part qui varie fortement dans le cadre de la migration
vers les tables interactives. En effet le nombre d’utilisateurs et les modalit\'{e}s d’interactions changent
entre la source et la cible. Dans ce cadre, la transformation du dialogue de d\'{e}part soul\`{e}ve des questions par rapport \`{a} l’utilisation des objets tangibles, \`{a} l’\'{e}quivalence des interactions et au nombre
d’utilisateurs. Cette section expose ces questions.
2.2.2.1
Utilisation des objets tangibles comme moyen d’interactions
Dans le cadre de l’application CBA, les dessinateurs de BD peuvent utiliser diff\'{e}rents objets
(cube, disque, etc.) pour afficher un menu de mani\`{e}re contextuelle dans le but d’effectuer une t\^{a}che de
changement de couleur ou de police d’un texte. De mani\`{e}re g\'{e}n\'{e}rale, il est indispensable d’identifier
dans l’UI d\'{e}part les composants graphiques ou les fonctionnalit\'{e}s que les utilisateurs utiliseront par
des interactions tangibles.
Les objets physiques peuvent aussi servir d’instruments de manipulations directes ou pour entrer
des donn\'{e}es. L’exemple d’un stylo pour s\'{e}lectionner des options ou pour \'{e}crire dans le cadre de
l’application CBA est envisageable. Dans ce cas, la migration des dialogues impacte \`{a} la fois l’UI et
le NF. En effet, les interactions tangibles peuvent aussi n\'{e}cessiter l’ajout dans le NF des traitements
inexistants dans l’application de d\'{e}part et la transformation de l’UI. Par exemple l’\'{e}criture \`{a} main
lev\'{e}e avec un stylo \`{a} main lev\'{e}e peut n\'{e}cessiter l’ajout d’un module de reconnaissance d’\'{e}criture si la
plateforme cible n’offre pas cette fonctionnalit\'{e}.
L’utilisation des objets tangibles comme moyen d’interaction pour les UI de la cible soul\`{e}ve un
sous-ensemble de la dimension des probl\`{e}mes li\'{e}s \`{a} l’\'{e}quivalence des interactions entre la source et
la cible.
2.2.2.2
Équivalences des interactions entre la source et de la cible pendant la migration
Pour \'{e}tablir des \'{e}quivalences entre les plateformes source et cible, il est indispensable d’\'{e}tudier
les instruments qui permettent aux utilisateurs de dialoguer avec les UI sur ces plateformes. En effet
les diff\'{e}rents instruments d’interactions offrent diff\'{e}rentes modalit\'{e}s aux utilisateurs des applications.
Dans le cas de la migration de l’application CBA du desktop vers une table interactive, il est indispensable de transposer par les interactions du clavier et de la souris sur les tables par des interactions
tactiles ou par un objet tangible.
La transposition des dialogues de la source sur une table interactive impose des transformations
des interactions initiaux pr\'{e}sentes entre l’UI desktop et l’utilisateur et entre l’UI et le noyau fonctionnel. Dans le but d’assurer la coh\'{e}rence des dialogues apr\`{e}s la migration, il est alors indispensable
de s’assurer que les nouveaux dialogues permettent toujours une utilisation de l’application. Pour ce
faire, nous pensons que les \'{e}quivalences d’interactions doivent :
2.2. SCÉNARIO \textbackslash{}& PROBLÈMES
13
– Garantir que l’ensemble des fonctionnalit\'{e}s de l’application desktop restent accessible sur les
tables interactives apr\`{e}s la migration. En effet, une migration ne doit pas alt\'{e}rer les fonctionnalit\'{e}s de l’application initiale sauf si la personne en charge de la migration le souhaite.
– Assurer que les dialogues modifiant les donn\'{e}es d’une application prennent en compte les interactions de plusieurs utilisateurs.
2.2.2.3
Transformation du dialogue de la source pour une plateforme multi-utilisateurs et colocalis\'{e}e
Dans le cadre de la migration, ce probl\`{e}me consiste aussi \`{a} transformer une application monoutilisateur en une application multi-utilisateurs pour les tables interactives au cas o\`{u} le nombre d’utilisateurs est sup\'{e}rieur \`{a} un. La migration dans ce cas impacte \`{a} la fois l’UI et le NF.
En ce qui concerne les transformations au niveau de l’UI, comment \'{e}viter d’avoir un dialogue
g\^{e}nant pour les autres utilisateurs ? Dans le cadre de l’application CBA par exemple les feedbacks,
les alertes ou les messages d’erreurs provoqu\'{e}s par un utilisateur ne doivent pas \^{e}tre bloquants pour
d’autres. Par ailleurs, comment d\'{e}crire le couplage entre les utilisateurs et les dialogues ? En effet,
une t\^{a}che 5 ou une activit\'{e} est exprim\'{e}e par ensemble de dialogues entre l’utilisateur et le syst\`{e}me
pour une application desktop. Sur une table interactive, une t\^{a}che peut \^{e}tre effectu\'{e}e par plusieurs
utilisateurs de mani\`{e}re concurrente ou en collaboration. Dans ce cas, comment transformer une UI
pour prendre en compte les t\^{a}ches concurrentes ou les t\^{a}ches effectu\'{e}es en collaboration ? Comment
surtout garantir que le noyau fonctionnel g\`{e}re de mani\`{e}re coh\'{e}rente la multiplicit\'{e} des interfaces et
les diff\'{e}rents utilisateurs.
En ce qui concerne les transformations au niveau du NF, la migration de l’application CBA par
exemple implique que les t\^{a}ches de sauvegardes de fichiers ou de modifications des donn\'{e}es puissent
\^{e}tre faites par les utilisateurs en pr\'{e}servant la coh\'{e}rence des donn\'{e}es de l’application. Comment transformer le NF des applications mono-utilisateur pour prendre les dialogues multi-utilisateurs ?
G\'{e}n\'{e}ralement les sp\'{e}cifications conceptuelles des applications \`{a} migrer ne sont pas disponibles. Il
est indispensable de retrouver la description des t\^{a}ches d’une application \`{a} partir de son code source
par exemple pour transformer les dialogues. Comment retrouver et repr\'{e}senter la description des dialogues d’une application afin de les transformer ? Dans notre cadre de migration, si l’extraction des
t\^{a}ches est difficile ou impossible \`{a} partir du code source de l’application \`{a} migrer, est-il possible de
transformer les dialogues sans description des t\^{a}ches ? Quels sont les limites d’une transformation des
dialogues non bas\'{e}e sur les t\^{a}ches ?
2.2.3
Probl\`{e}mes li\'{e}s \`{a} la migration de la structure et du positionnement
Nous pensons que la migration des UI vers les tables interactives soul\`{e}ve des probl\`{e}mes li\'{e}s \`{a}
la transformation de sa structure et du positionnement de ses \'{e}l\'{e}ments. En effet, la structure et le
positionnement sont des caract\'{e}ristiques des UI qui se d\'{e}clinent de diff\'{e}rentes mani\`{e}res suivant les
plateformes et les utilisateurs. La transformation de la structure et du positionnement dans le cadre
des tables interactives a pour objectif de favoriser l’accessibilit\'{e} des \'{e}l\'{e}ments graphiques et de fa\c{c}on
globale \`{a} l’utilisabilit\'{e} des UI migr\'{e}es.
La transformation de la structure et du positionnement des UI de d\'{e}part soul\`{e}ve des questions
sur les \'{e}quivalences entre les \'{e}l\'{e}ments graphiques qui composent l’UI source et ceux propos\'{e}s par la
cible, le regroupement et le positionnement des \'{e}l\'{e}ments graphiques pour favoriser la collaboration.
5. Nous consid\'{e}rons les t\^{a}ches et les activit\'{e}s impliquant les utilisateurs dans cas
14
CHAPITRE 2. MIGRATION DES APPLICATIONS
2.2.3.1
Équivalences entre les \'{e}l\'{e}ments graphiques
Dans le cadre de l’application CBA par exemple, les menus ou les listes d’\'{e}l\'{e}ments peuvent
\^{e}tre pr\'{e}sent\'{e}s sur les tables interactives dans le but d’avoir des interactions directes et intuitives. Les
\'{e}l\'{e}ments graphiques propos\'{e}s par les tables interactives permettent d’atteindre ce but. Par exemple,
une liste contenant des items textuels peut \^{e}tre remplac\'{e}e par une liste d’\'{e}l\'{e}ments d’images associ\'{e}es
aux items sur la table interactive. L’ensemble des \'{e}l\'{e}ments de la liste textuelles doit \^{e}tre transpos\'{e} sur
la table interactive dans une liste \'{e}quivalente.
De mani\`{e}re globale, les \'{e}quivalences entre les \'{e}l\'{e}ments graphiques de la source et de la cible
doivent pr\'{e}server les donn\'{e}es qu’ils contiennent. La multitude d’\'{e}l\'{e}ments graphiques et des plateforme implique de mettre en place une solution \'{e}quivalente entre les \'{e}l\'{e}ments graphiques automatisable.
2.2.3.2
Favoriser la collaboration par un regroupement et un positionnement ad\'{e}quat des \'{e}l\'{e}ments graphiques
La fen\^{e}tre principale de l’application CBA est con\c{c}ue en respectant la m\'{e}taphore du bureau qui
permet de d\'{e}crire les applications pour les ordinateurs personnels [App95]. Les menus sont plac\'{e}s en
haut \`{a} des positions fixes, les outils ou les l\'{e}gendes sont toujours plac\'{e}s \`{a} gauche ou \`{a} droite et la zone
de travail au centre. Cette structuration en zones permet aux utilisateurs des desktops de d\'{e}velopper
des r\'{e}flexes quelque soit l’application dans une approche similaire.
 
Canvas1
Ressouces
 
FIGURE 2.3 – Migration de la structure d’une UI pour desktop sur une table interactive
Sur une table interactive, cette structuration n’est pas recommand\'{e}e car l’accessibilit\'{e} des composants graphiques d\'{e}pend directement de la position de l’utilisateur autour de la table : le haut pour un
utilisateur peut \^{e}tre le bas d’un autre. La table interactive ne poss\`{e}de pas d’axe bas-haut imm\'{e}diat et
il est donc n\'{e}cessaire de repenser la disposition de diff\'{e}rentes fen\^{e}tre. Par exemple, pour l’application
CBA, les diff\'{e}rentes zones de la structure de d\'{e}part doivent pouvoir \^{e}tre conserv\'{e}es mais chaque zone
n’est plus associ\'{e}e \`{a} un espace g\'{e}ographique sp\'{e}cifique de l’\'{e}cran. Sur la droite de la figure 2.3, les
diff\'{e}rentes zones de l’UI CBA de d\'{e}part n’ont plus de position fixe sur une table interactive. On peut
m\^{e}me imaginer dupliquer certaines zones pour qu’elle deviennent accessibles aux utilisateurs.
Les transformations de la structure et du positionnement de l’UI a pour objectif de favoriser
l’accessibilit\'{e}. Pour ce faire, il est important de d\'{e}terminer les groupes d’\'{e}l\'{e}ments graphiques \`{a}
transformer.
2.2. SCÉNARIO \textbackslash{}& PROBLÈMES
15
Dans le but de concevoir des UI utilisables, l’ing\'{e}nierie des IHM pr\'{e}conise de prendre en compte
des crit\`{e}res ergonomiques (ou des heuristiques) pendant la conception des UI [Sca86, NM90]. Nous
consid\'{e}rons dans cette th\`{e}se que le respect des crit\`{e}res ergonomiques de conception pendant la migration permet de favoriser l’utilisabilit\'{e} des UI obtenues. Dans ce cadre, il reste \`{a} d\'{e}finir comment
on peut prendre en compte ces crit\`{e}res ergonomiques pendant la migration des UI. Il est donc important de les identifier et de les int\'{e}grer aux processus de migration pour pouvoir ult\'{e}rieurement
affiner les propositions d’UI et \'{e}mettre des recommandations utilisables par la personne en charge de
la migration.
2.2.3.3
Remarque
– La migration de la structure et du positionnement sont des transformations uniquement au niveau des UI des applications \`{a} migrer.
2.2.4
Probl\`{e}mes li\'{e}s \`{a} la migration du style
Dans notre cadre, la migration du style consiste \`{a} modifier l’aspect visuel des diff\'{e}rents \'{e}l\'{e}ments
graphiques de l’UI d’arriv\'{e}e. Nous avons identifi\'{e} trois caract\'{e}ristiques des \'{e}l\'{e}ments graphiques pour
d\'{e}crire la modification du style.
– La taille des \'{e}l\'{e}ments graphiques est importante sur une table interactive car les dimensions
des composants graphiques peuvent favoriser ou non un travail commun. Selon ce que l’on
souhaite : v\'{e}ritable travail en commun ou chacun doit \^{e}tre capable de visualiser imm\'{e}diatement
l’activit\'{e} d’un autre utilisateur ou au contraire, un travail plus individuel ou l’activit\'{e} de l’un
ne doit pas perturber l’activit\'{e} de l’autre ; il sera n\'{e}cessaire d’ajuster la taille des composants
graphiques. En effet des composants graphiques trop larges peuvent g\^{e}ner les autres utilisateurs
et au contraire, des composants trop petit peuvent rendre difficile le travail commun.
– La couleur et la police des textes sont des propri\'{e}t\'{e}s subjectives car elles peuvent ne pas \^{e}tre
accept\'{e}es par tout le monde. Cependant dans le cadre de la migration nous pensons que leur
choix doivent \^{e}tre coh\'{e}rents entre les \'{e}l\'{e}ments graphiques d’une UI. La coh\'{e}rence du style
consiste \`{a} s’assurer que les \'{e}l\'{e}ments graphiques de m\^{e}me type doivent avoir des couleurs ou
des polices “compatibles” par exemple (respect d’une charte graphique).
Les tables interactives ont une surface d’affichage plus grande qu’un desktop ou une tablette. La
taille des \'{e}l\'{e}ments d’une UI doit n\'{e}cessairement \^{e}tre adapt\'{e}e en fonction de cette surface mais aussi
en fonction du nombre d’utilisateurs pr\'{e}vus afin de faciliter la participation de tous.
2.2.4.1
Coh\'{e}rence globale du style des UI migr\'{e}es
Il est important que les couleurs, les polices, les ic\^{o}nes, les images ou la taille des \'{e}l\'{e}ments graphiques soient d\'{e}finies suivant des recommandations pour assurer une coh\'{e}rence globale. L’homog\'{e}n\'{e}it\'{e} de caract\'{e}ristiques 6 pour chaque type d’\'{e}l\'{e}ments graphiques constitue la coh\'{e}rence globale
\`{a} un style. Il ne s’agit pas d’un probl\`{e}me sp\'{e}cifique aux tables interactives mais la taille de l’\'{e}cran
disponible pose d’une mani\`{e}re nouvelle le respect des chartes graphique et des recommandations de
style.
6. Les couleurs, les polices, les ic\^{o}nes, les images ou la taille des \'{e}l\'{e}ments graphiques
16
CHAPITRE 2. MIGRATION DES APPLICATIONS
2.2.5
Espace des probl\`{e}mes li\'{e}s la migration vers les tables interactives
Les probl\`{e}mes li\'{e}s \`{a} la migration que nous avons \'{e}tudi\'{e}s \`{a} la section 2.2 peuvent \^{e}tre repr\'{e}sent\'{e}s
\`{a} l’aide d’un espace \`{a} trois dimensions d\'{e}crit par la figure 2.4.
D2:Structure 
\textbackslash{}& Positionnement
D3: Style
D1: Dialogues
D11: Interactions 
tangibles
D12: Equivalence des 
interactions
D13: Multi-utilisateurs 
\textbackslash{}& co-localis\'{e}s
D21: Equivalences 
Structurelles
D22: Regroupement 
D32: Polices
D31: Taille
D23:Positionnement
D32: Couleurs
FIGURE 2.4 – Espace probl\`{e}mes de la migration des UI vers les tables interactives
La dimension D1 correspond aux probl\`{e}mes de migration des dialogues. Elle consiste \`{a} transformer l’UI et porter le NF des applications. Il est indispensable d’avoir des processus g\'{e}n\'{e}riques
pour prendre en compte plusieurs applications et r\'{e}duire les co\^{u}ts de la migration. Elle regroupe trois
sous-dimensions.
– La dimension D11 correspondant aux probl\`{e}mes li\'{e}s \`{a} la transformation des dialogues par ajout
des interactions tangibles aux applications de d\'{e}part.
– La dimension D12 correspond aux probl\`{e}mes d’\'{e}quivalence des interactions des UI migr\'{e}es.
– La dimension D13 correspondant aux probl\`{e}mes li\'{e}s \`{a} la transformation des dialogues monoutilisateur en multi-utilisateurs co-localis\'{e}s.
Ensuite la dimension D2 correspond aux probl\`{e}mes de la transformation de la structure et du
positionnement des \'{e}l\'{e}ments graphiques. Elle consiste \`{a} transformer l’UI de l’application source.
Dans ce cas, il est important d’avoir un processus qui prend en compte les crit\`{e}res ergonomiques de
conception afin de garantir l’utilisabilit\'{e}. La g\'{e}n\'{e}ricit\'{e} du processus de transformation est importante.
Elle regroupe deux sous-dimensions.
– La dimension D21 correspondant aux probl\`{e}mes li\'{e}s \`{a} l’\'{e}quivalence structurelle des \'{e}l\'{e}ments
graphiques de la source et de la cible.
– La dimension D22 correspondant aux probl\`{e}mes de regroupement des \'{e}l\'{e}ments graphiques
des UI migr\'{e}es.
– La dimension D23 correspondant aux probl\`{e}mes de positionnement des \'{e}l\'{e}ments graphiques
des UI migr\'{e}es.
2.3. PÉRIMÈTRES DE LA MIGRATION DES UI VERS LES TABLES INTERACTIVES
17
Enfin la dimension D3 correspond aux probl\`{e}mes li\'{e}s \`{a} la transformation du style de l’UI d\'{e}part
pour favoriser leur accessibilit\'{e}. Dans le but de r\'{e}duire la charge de travail des concepteurs et d’assurer
l’homog\'{e}n\'{e}it\'{e}, il est important d’utiliser des styles g\'{e}n\'{e}riques. Cependant la flexibilit\'{e} de la migration
du style permet aussi de personnaliser les UI cibles. Elle regroupe deux sous-ensembles de probl\`{e}mes :
– La dimension D31 correspondant aux probl\`{e}mes li\'{e}s \`{a} la taille des \'{e}l\'{e}ments graphiques sur un
espace de travail partag\'{e}.
– La dimension D32 correspondant aux probl\`{e}mes li\'{e}s \`{a} la police des textes des UI migr\'{e}es.
– La dimension D33 correspondant aux probl\`{e}mes li\'{e}s \`{a} la couleurs des \'{e}l\'{e}ments graphiques des
UI migr\'{e}es.
2.3
P\'{e}rim\`{e}tres de la migration des UI vers les tables interactives
Dans notre contexte, les applications \`{a} migrer sont des syst\`{e}mes interactifs (SI) con\c{c}us en respectant un mod\`{e}le d’architecture et identifiant un noyau fonctionnel (NF) et des interfaces utilisateurs
(UI). Les mod\`{e}les d’architecture sont des patrons de conception logiciel, ils pr\'{e}conisent des strat\'{e}gies de r\'{e}partition des services qui se traduisent par un ensemble de constituants logiciels. Il existe
plusieurs mod\`{e}les d’architecture qui permettent la s\'{e}paration UI-NF.
– Le mod\`{e}le d’architecture de ARCH [Dev92] par exemple, permet une s\'{e}paration entre le NF,
le Contr\^{o}leur de Dialogue (CD) et la Pr\'{e}sentation (P). Elle permet aussi la migration des composants de l’UI (CD et P) sans modification des composants du NF.
– Le mod\`{e}le d’architecture PAC-Amodeus [Nig94] est bas\'{e} sur ARCH, il permet de d\'{e}finir plusieurs contr\^{o}leurs de dialogue pour une application gr\^{a}ce aux agents PAC [Nig94].
– Le mod\`{e}le d’architecture MVC [KP+88] permet la conception des SI r\'{e}utilisables. Il divise
les applications en trois types de composants : le mod\`{e}le (M), la vue (V) et le contr\^{o}leur (C).
Le mod\`{e}le est une repr\'{e}sentation du domaine d’une application, il peut contenir des donn\'{e}es,
des services, etc. et il fait partie du NF d’une application. La vue est la structure de l’UI d’une
application, elle est constitu\'{e}e des \'{e}l\'{e}ments d’une biblioth\`{e}que graphique. Le contr\^{o}leur est
une interface entre le mod\`{e}le, la vue et les dispositifs d’interactions en entr\'{e}e. Cette architecture
permet une migration de la vue sans modification du mod\`{e}le et du contr\^{o}leur.
Nous pr\'{e}sentons \`{a} la section 2.3.1 notre hypoth\`{e}se de travail. La section 2.3.2 pr\'{e}sente les pistes
de migration des UI et la section 2.4, les objectifs de notre travail.
2.3.1
Hypoth\`{e}se de travail
Les applications \`{a} migrer ont une d\'{e}composition fonctionnelle minimale qui comprend une interface utilisateur (UI) et un noyau fonctionnel (NF). Cette d\'{e}composition identifie clairement la partie
qu’il faut migrer (l’UI) et la partie qui sera r\'{e}utilis\'{e}e sur la plateforme cible (le NF). Dans cette th\`{e}se,
pour ne pas aborder simultan\'{e}ment tous les probl\`{e}mes de migration et parce que ceux-ci ont \'{e}t\'{e} abondamment \'{e}tudi\'{e}s par ailleurs, nous consid\'{e}rons que nous pouvons r\'{e}utiliser le NF de l’application
source sur la plateforme cible sans modifier son code. Il offre donc les m\^{e}mes fonctionnalit\'{e}s. Il est
\'{e}vident que si ce NF devrait lui-m\^{e}me \^{e}tre migrer, quel qu’en soit les raisons, cela doit \^{e}tre effectivement trait\'{e}, \'{e}ventuellement par les propositions d\'{e}crites dans [TKB78, TKR10] du NF par exemple.
Les deux hypoth\`{e}ses simplificatrices que nous faisons sont r\'{e}alistes. La plupart des applications
interactives adoptent aujourd’hui une architecture proche du mod\`{e}le MVC tr\`{e}s largement r\'{e}pandu
permettant d’identifier tr\`{e}s clairement la partie UI et la partie NF de l’application. Par ailleurs, la
table surface que nous avons utilis\'{e} - une des rares tables interactives disponibles - utilisent le m\^{e}me
syst\`{e}me d’exploitation que celui utilis\'{e} par la tr\`{e}s grande majorit\'{e} des desktops rendant inutile le
18
CHAPITRE 2. MIGRATION DES APPLICATIONS
portable du NF. Si celui-ci s’av\'{e}rait n\'{e}cessaire, il serait alors possible d’utiliser une solution \`{a} base de
machine virtuelle.
Il est \'{e}vident que notre approche conserve les dialogues existants entre le NF et l’UI et que nous
n’avons que tr\`{e}s peu abord\'{e} les questions relatives \`{a} l’\'{e}volution de ceux-ci lors de la migration.
2.3.2
Pistes de migration
A partir des hypoth\`{e}ses pr\'{e}c\'{e}demment expos\'{e}es, la migration des UI des applications existantes
vers une nouvelle plateforme tout en conservant le NF peut se faire de plusieurs mani\`{e}res :
– Une nouvelle conception de l’UI pour la plateforme cible sans tenir compte de l’UI source.
– Une d\'{e}duction de l’UI \`{a} partir des sp\'{e}cifications du NF de d\'{e}part.
– Une adaptation de l’UI de d\'{e}part par rapport \`{a} la plateforme cible.
2.3.2.1
Nouvelle conception de l’UI
Cette approche consiste \`{a} concevoir une nouvelle UI pour la plateforme cible sans n\'{e}cessairement
tenir compte de l’UI de d\'{e}part. Cette nouvelle conception permet de prendre en compte les sp\'{e}cificit\'{e}s
de la plateforme d’arriv\'{e}e et d’int\'{e}grer ses crit\`{e}res de conception des UI (chaque plateforme est “livr\'{e}e” avec un guide du bon usage). Les approches de conception des UI telles que DIANE+ [Bar88]
ou MUSE [LL09] et les approches de conception bas\'{e}es sur des mod\`{e}les [MPV11] peuvent \^{e}tre utilis\'{e}es pour concevoir la nouvelle UI. DIANE+ et MUSE permettent d’inclure dans la conception de
l’UI des crit\`{e}res ergonomiques[Van97].
L’avantage d’une nouvelle conception est de construire une UI sur mesure, prenant en compte
l’ensemble des nouveaux besoins li\'{e}s \`{a} la plateforme cible. Cependant le temps pris pour la conception
de cette nouvelle UI comprends n\'{e}cessairement le temps pris pour la transformation des dialogues de
l’UI source et les temps n\'{e}cessaire \`{a} la conception de la structure, du positionnement et du style de
la nouvelle UI. Ces diff\'{e}rents temps constituent des co\^{u}ts de mise en œuvre non n\'{e}gligeables. Cette
piste de solution, ne capitalisant pas sur le travail d\'{e}j\`{a} effectu\'{e}, demandant des comp\'{e}tences m\'{e}tiers
fortes et un temps non n\'{e}gligeable de mise en œuvre, ne nous permet pas de r\'{e}duire les co\^{u}ts de
mise en œuvre des applications pour les tables interactives.
2.3.2.2
D\'{e}duction de l’UI \`{a} partir du NF
La deuxi\`{e}me piste de migration pr\'{e}conise de s’appuyer sur le NF de l’application de d\'{e}part,
particuli\`{e}rement sur les liens entre ce NF et l’UI. Il est en effet possible de d\'{e}duire une UI en se
basant sur la repr\'{e}sentation abstraite des liens entre le NF et l’UI. Il existe des travaux qui pr\'{e}conisent
la g\'{e}n\'{e}ration des UI \`{a} partir des descriptions de services Web par exemple [KKM03]. Par ailleurs
l’approche ALIAS [JOF11] propose des mod\`{e}les pour une repr\'{e}sentation abstraite du NF et des liens
entre NF et UI qui peuvent \^{e}tre utilis\'{e}s pour d\'{e}duire les UI \`{a} partir des NF. La d\'{e}duction des UI en se
basant sur les NF implique la d\'{e}duction \`{a} partir du NF des \'{e}l\'{e}ments suivants :
– les composants graphiques qui composent l’UI,
– la structure et le positionnement (layout) de ses composants graphiques,
– les dialogues et les encha\^{i}nement des fen\^{e}tres et des activit\'{e}s de l’UI.
Les travaux sur la g\'{e}n\'{e}ration d’une UI \`{a} partir du NF [TC02] montrent qu’en plus du NF, il est
indispensable d’avoir une sp\'{e}cification des t\^{a}ches sous forme d’un contr\^{o}leur de dialogue et d’adaptateur du NF (dans une architecture ARCH). Le NF ne permet usuellement pas de d\'{e}duire la structure
et les diff\'{e}rents types de regroupements des composants graphiques (fen\^{e}tres, panels, etc.). Les UI
2.4. OBJECTIFS
19
obtenues par une d\'{e}duction \`{a} partir du NF sont utilisables \`{a} condition d’avoir une description des
t\^{a}ches pour permettre l’organisation des dialogues et de sa structure.
2.3.2.3
Adaptation de l’UI par recherche d’\'{e}quivalences
La troisi\`{e}me piste de migration est celle qui pr\'{e}conise de s’appuyer sur les sp\'{e}cifications de l’UI
\`{a} migrer et de les adapter \`{a} la plateforme d’arriv\'{e}e. Ces sp\'{e}cifications sont d\'{e}crites par des mod\`{e}les
abstraits. Les mod\`{e}les servent \`{a} d\'{e}crire les diff\'{e}rentes dimensions des UI tels que les dialogues d\'{e}crits
par l’UI, le placement ou le layout des \'{e}l\'{e}ments de l’UI mais aussi les styles de pr\'{e}sentations des
caract\'{e}ristiques visuelles (tailles et couleurs des textes, etc.). Les diff\'{e}rents mod\`{e}les abstraits des UI
\`{a} migrer peuvent \^{e}tre retrouv\'{e} par abstraction \`{a} partir du code source par exemple [BS08, BV02,
GPF10, TS99].
L’adaptation des UI \`{a} migrer pour les tables interactives peut se faire en transformant partiellement
ou en totalit\'{e} les trois dimensions identifi\'{e}es.
– L’adaptation des dialogues consiste \`{a} produire des dialogues \'{e}quivalents \`{a} partir de l’UI de
d\'{e}part et en prenant en compte les caract\'{e}ristiques de la cible : introduction d’interactions tangibles et pr\'{e}sence potentielle de plusieurs utilisateurs.
– L’adaptation de la structure consiste \`{a} produire pour la plateforme cible une structure \'{e}quivalente de l’UI de d\'{e}part en prenant en compte les probl\`{e}mes li\'{e}s au regroupement et au positionnement des composants graphiques sur les tables interactives pour lesquels l’axe haut-bas
d\'{e}pend de la position de l’utilisateur. Selon les cas, la structure et le positionnement de d\'{e}part
peuvent \^{e}tre ignor\'{e}s ou r\'{e}utilis\'{e}s partiellement pendant la migration. Dans le cas o\`{u} ils sont
ignor\'{e}s, il est indispensable de proposer des outils pour les r\'{e}introduire facilement.
– En ce qui concerne le style, il est indispensable de le r\'{e}introduire pour la cible, en respectant la
charte graphique et en favorisant l’accessibilit\'{e} des \'{e}l\'{e}ments graphiques. En effet, la taille, la
couleur ou la police des textes des tables interactives sont diff\'{e}rentes de celles des desktops. Le
style de la cible peut se baser sur un mod\`{e}le de style li\'{e} aux tables interactives.
L’introduction manuelle du positionnement et du style permet d’avoir une approche de migration
des UI plus de flexibilit\'{e} mais avec un co\^{u}t de mise en œuvre non n\'{e}gligeable. Au contraire, une
adaptation automatique du positionnement et du style pr\'{e}sente un co\^{u}t de migration faible mais au
prix d’une moindre flexibilit\'{e} et un faible respect des crit\`{e}res ergonomiques de la cible.
2.4
Objectifs de la migration des UI vers les tables interactives
Nous avons fait le choix dans cette th\`{e}se d’explorer la troisi\`{e}me piste de migration (cf. section 2.3)
et d’appliquer nos travaux aux tables interactives qui sont des plateformes innovantes provocants
un changement assez important dans la conception des UI. En effet, il ne s’agit pas uniquement de
trouver les composants similaires entre les toolkits de la plateforme source et de la plateforme cible
mais d’introduire au sein de l’UI de nouveaux dispositifs d’interactions par l’utilisation d’interfaces
tangibles et de prendre en compte la r\'{e}elle possibilit\'{e} qu’ont des utilisateurs de pouvoir enfin partager
de mani\`{e}re intuitive, pour peu que l’UI soit bien construite, une m\^{e}me application. La migration d’une
application existante sans avoir \`{a} reconcevoir l’UI est \`{a} m\^{e}me d’enrichir le catalogue d’application
disponible sur les tables interactives par un abaissement important du co\^{u}t de mise en œuvre de ces
applications.
Ces arguments nous ont amen\'{e}s \`{a} nous poser deux questions :
– d’une part, nous devons adapter les diff\'{e}rentes dimensions d’une UI fondamentalement con\c{c}ue
pour une personne en une UI collaborative et tangible,
20
CHAPITRE 2. MIGRATION DES APPLICATIONS
– et d’autre part, il est n\'{e}cessaire d’inclure dans le processus de migration la prise en compte
de crit\`{e}res ergonomiques de conception[Sca86, NM90], li\'{e}s \`{a} la plateforme cible pour r\'{e}duire
effectivement le co\^{u}t de la migration tout en concevant une UI r\'{e}ellement utilisable.
La r\'{e}ponse \`{a} ces deux questions doit nous permettre d’atteindre les objectifs que nous nous sommes
fix\'{e}s dans le cadre de cette th\`{e}se. Pour ce faire, la deuxi\`{e}me partie du document pr\'{e}sente les sp\'{e}cificit\'{e}s des tables interactives qui constituent notre domaine d’\'{e}tude pour identifier les crit\`{e}res de conception. Nous consid\'{e}rons que ces crit\`{e}res sont les crit\`{e}res ergonomiques de conception qui doivent \^{e}tre
affin\'{e}s en recommandations lors de la migration des UI vers les tables interactives.
Deuxi\`{e}me partie
Domaine d’\'{e}tude
21
CHAPITRE 3
Tables interactives et Migration des UI
Sommaire
3.1
Mod\`{e}le d’interactions d’une table interactive . . . . . . . . . . . . . . . . . . .
23
3.1.1
Les dispositifs mat\'{e}riels d’interactions d’une table interactive
. . . . . . .
24
3.1.2
Biblioth\`{e}ques graphiques des tables interactives . . . . . . . . . . . . . . .
27
3.1.3
Modalit\'{e}s d’interactions . . . . . . . . . . . . . . . . . . . . . . . . . . .
29
3.1.4
Mod\`{e}le d’interactions abstraites pour la migration des UI . . . . . . . . . .
31
3.2
Principes de conception des UI pour les tables interactives
. . . . . . . . . . .
33
3.2.1
Propri\'{e}t\'{e}s caract\'{e}ristiques des tables interactives . . . . . . . . . . . . . .
33
3.2.2
Crit\`{e}res ergonomiques de conception des UI
. . . . . . . . . . . . . . . .
35
3.2.3
Recommandations pour la migration des UI vers les tables interactives . . .
38
3.3
Synth\`{e}se
. . . . . . . . . . . . . . . . . . . . . . . . . . . . . . . . . . . . . . .
43
N
OUS pr\'{e}sentons dans ce chapitre les sp\'{e}cificit\'{e}s des tables interactives de mani\`{e}re g\'{e}n\'{e}rale afin
d’identifier les recommandations pour la migration des UI. Pour atteindre cet objectif, nous
\'{e}tudions dans la section 3.1 les instruments d’interactions ou de dialogue qui sont \`{a} la fois des dispositifs mat\'{e}riels et logiciels. Puis les mod\`{e}les permettant de d\'{e}crire les dialogues ind\'{e}pendamment
des instruments d’une plateforme sont \'{e}tudi\'{e}s. La section 3.2 pr\'{e}sente les principes de conception en
d\'{e}crivant les crit\`{e}res ergonomiques de conception et les recommandations pour la migration des UI
vers les tables interactives. La section 3.3 synth\'{e}tise les principaux travaux de cette \'{e}tude sur les tables
interactives dans le but de mettre en \'{e}vidence les principes qui doivent \^{e}tre respect\'{e}s pour migrer une
UI vers cette plateforme cible.
3.1
Mod\`{e}le d’interactions d’une table interactive
Les instruments d’interactions constituent l’ensemble des dispositifs mat\'{e}riels et logiciels d’une
table interactive qui permettent d’interagir avec un SI. Les dispositifs mat\'{e}riels d’interactions (cf.
figure 3.1) sont des moyens d’interactions en entr\'{e}e (actions) ou en sortie (r\'{e}actions et feedback).
Les biblioth\`{e}ques graphiques encapsulent les paradigmes de manipulation des diff\'{e}rents dispositifs
pr\'{e}sent et permettent de d\'{e}crire des interfaces utilisateurs.
Dans cette section nous nous appuierons sur diff\'{e}rentes tables interactives : Microsoft PixelSense,
TangiSense, DiamondTouch, etc. Cette diversit\'{e} nous permet de caract\'{e}riser les dispositifs d’interactions (cf section 3.1.1), les biblioth\`{e}ques graphiques (cf section 3.1.2) et les modalit\'{e}s d’interactions
des tables interactives (cf section 3.1.3).
Nous \'{e}tudions ces trois tables interactives dans le but de caract\'{e}riser les diff\'{e}rents types d’UI qu’il
est possible de mette en œuvre en fonction des instruments d’interactions. Ces trois tables interactives
d\'{e}crivent des UI tactiles, des UI tangibles et des UI collaboratives.
23
24
CHAPITRE 3. TABLES INTERACTIVES ET MIGRATION DES UI
 
 
 
FIGURE 3.1 – Mod\`{e}le d’interactions instrumentale d’une table interactive
3.1.1
Les dispositifs mat\'{e}riels d’interactions d’une table interactive
Dans ce paragraphe nous \'{e}tudions les instruments d’interactions de trois tables interactives repr\'{e}sentatives. Nous \'{e}tudions d’abord DiamondTouch [DL01] qui est l’une des premi\`{e}re table interactive
utilis\'{e}e dans un cadre non exp\'{e}rimental, elle fut industrialis\'{e}e par MERL [Mit], nous l’\'{e}tudions car
elle comporte une biblioth\`{e}que graphique DiamondSpin qui permet de d\'{e}crire des interactions multiutilisateurs, collaboratives et tactiles. La seconde table \'{e}tudi\'{e}e est metaDESK [UI97] qui n’accepte
des interactions que via des objets tangibles. La troisi\`{e}me table \'{e}tudi\'{e}e est la famille Microsoft PixelSense 1.0 et 2.0 [Mic11] qui poss\`{e}dent une biblioth\`{e}que graphique compl\`{e}te qui permet de construire
des interactions collaboratives, tangibles et tactiles.
L’\'{e}tude effectu\'{e}e dans cette section a pour objectifs de caract\'{e}riser les diff\'{e}rentes cat\'{e}gories de
dispositifs d’interactions.
3.1.1.1
DiamondTouch
La table interactive DiamondTouch est multi-touche et multi-utilisateurs. Elle poss\`{e}de une surface
tactile capacitive pour les interactions en entr\'{e}e et un vid\'{e}o projecteur pour les interactions en sortie
qui sont reli\'{e}s \`{a} un ordinateur. Sa surface tactile fonctionne gr\^{a}ce \`{a} un r\'{e}seau d’antennes qui transmet
des signaux tr\`{e}s faible. Lorsqu’un utilisateur touche la surface tactile, des signaux capacitifs sont
associ\'{e}s \`{a} un r\'{e}cepteur en partant du point de contact entre l’utilisateur et la surface et en passant dans
son corps. Les r\'{e}cepteurs sont associ\'{e}s aux diff\'{e}rents utilisateurs de la table par l’interm\'{e}diaire d’un
tapis plac\'{e} sur leur chaise. La figure 3.2 pr\'{e}sente les diff\'{e}rents composants de la table DiamondTouch.
Les r\'{e}cepteurs (receiver) sont des circuits \'{e}lectroniques qui d\'{e}terminent les coordonn\'{e}es du point de
la surface touch\'{e}e par un utilisateur. Le transmetteur (transmitter) est constitu\'{e} de plusieurs antennes
de la surface tactile, il est coupl\'{e} de mani\`{e}re capacitif aux r\'{e}cepteurs de chaque utilisateur.
Cette table est aussi multi-utilisateurs, car elle permet simultan\'{e}ment plusieurs touches de diff\'{e}rents utilisateurs. Chaque touche peut \^{e}tre associ\'{e}e \`{a} un utilisateur pr\'{e}cis gr\^{a}ce \`{a} son tapis. Par sa
capacit\'{e} \`{a} reconna\^{i}tre les diff\'{e}rents utilisateurs, cette table permet de d\'{e}crire des UI collaboratives
qui permettent de partager la m\^{e}me UI et donc d’acc\'{e}der simultan\'{e}ment \`{a} la m\^{e}me application. Il est
aussi possible, de cr\'{e}er pour chaque utilisateur son propre espace de travail personnel. L’utilisation de
la table DiamondTouch impose une r\'{e}partition g\'{e}ographique fixe des utilisateurs autour de la table.
Elle ne supporte pas les interactions tangibles car elle ne peut pas d\'{e}tecter des objets ou des tags.
Pour faire migrer l’UI de l’application CBA vers la table DiamondTouch impose de connaitre le
nombre d’utilisateurs de l’UI migr\'{e}e car chaque utilisateur doit \^{e}tre en contact avec un tapis (assis ou
debout). Sur la figure 3.2, les tapis sont plac\'{e}s sur les chaises des utilisateurs pour servir de r\'{e}cepteur,
par contre la surface d’affichage (ou la table) est un transmetteur.
3.1. MODÈLE D’INTERACTIONS D’UNE TABLE INTERACTIVE
25
FIGURE 3.2 – Table interactive DiamondTouch
3.1.1.2
metaDESK
C’est la premi\`{e}re prototype de table interactive con\c{c}ue en 1997 par Tangible Media Group du MIT
Media Lab pour comprendre le conception des interface utilisateur tangible ou Tangible User Interface (TUI). Les moyens d’interactions de la table metaDESK sont constitu\'{e}s d’objets tangibles, des
cam\'{e}ras infrarouges, d’une cam\'{e}ra vid\'{e}o et d’un vid\'{e}o projecteur. Ces objets permettent \`{a} chaque utilisateur de manipuler des objets virtuels associ\'{e}s aux objets physiques et par cons\'{e}quent de construire
des TUI. La reconnaissance des objets tangibles se fait par la cam\'{e}ra vid\'{e}o et les cameras infrarouges
en utilisant une m\'{e}thode de reconnaissance et de tracking des objets sur la table. La table metaDESK
est utilis\'{e}e dans le cadre d’une application de localisation g\'{e}ographique ; les objets physiques dans
ce cas repr\'{e}sentent des b\^{a}timents, des ponts, etc. Il existe aussi des tables interactives con\c{c}ues pour
d’autres usages ; les tables AudioPad [PRI02] et Xenakis [BCL+08] par exemple disposent des TUI
pour composer et \'{e}couter de la musique. Les objets tangibles dans ce cadre ont des formes proches
des boutons d’une console de mixage pour faciliter les interactions des utilisateurs.
Contrairement \`{a} la table DiamondTouch, la table metaDESK ne limite pas le nombre d’utilisateurs
et ceux-ci ne sont pas assign\'{e}s \`{a} un emplacement fixe et peuvent se d\'{e}placer autour de la table.
La table metaDESK ne permet pas aux diff\'{e}rents utilisateurs d’avoir un espace personnel comme
DiamondTouch.
Ulmer et Ishii
[IU97] d\'{e}finissent un TUI comme un syst\`{e}me interactif qui utilise des objets
physiques pour repr\'{e}senter et utiliser des informations digitales. Le mapping entre le monde physique
et digital peut se faire en repr\'{e}sentant les diff\'{e}rents composants graphiques d’une UI graphique \`{a}
l’aide d’objets concrets. Par exemple, les concepteur de la table metaDESK [IU97], associent (cf.
figure 3.3) une lentille \`{a} une fen\^{e}tre, un plateau \`{a} un menu, etc.
Dans le cadre de la migration d’une UI desktop vers la table metaDESK, toutes les interactions de
l’UI de d\'{e}part doivent \^{e}tre adapt\'{e}es ou \'{e}mul\'{e}es par l’interm\'{e}diaire d’objets tangibles. Par exemple,
le formulaire Ressources de l’application CBA doit \^{e}tre affich\'{e} et rempli uniquement par l’utilisation des objets tangibles. Il en est de m\^{e}me pour la s\'{e}lection de la taille ou de la police qui peuvent
\^{e}tre \'{e}mul\'{e}es avec des objets circulaires en les faisant tourner par exemple. Par ailleurs, l’\'{e}mulation
d’un clavier \`{a} l’aide d’un objet tangible pour la saisie des textes est possible mais n’est pas facile
26
CHAPITRE 3. TABLES INTERACTIVES ET MIGRATION DES UI
FIGURE 3.3 – Instanciation physique des \'{e}l\'{e}ments GUI dans TUI
\`{a} utiliser [USJ+08]. Toutes les applications ne se pr\^{e}tent pas \`{a} \^{e}tre utilis\'{e}e sur des tables interactives uniquement tangibles comme metaDESK ou TangiSense [KLL+09]. Elles sont en g\'{e}n\'{e}ralement
utilis\'{e}es pour des applications de r\'{e}alit\'{e} augment\'{e}e ou pour des applications de cartographie.
3.1.1.3
Microsoft PixelSense
La premi\`{e}re version de la table Microsoft PixelSense est constitu\'{e}e d’un vid\'{e}o projecteur plac\'{e}
sous la surface tactile pour les interactions en sortie, de cinq cam\'{e}ras infrarouges pour d\'{e}tecter les
formes des objets ou des doigts des utilisateurs, d’un clavier virtuel et des dispositifs sonores. La
deuxi\`{e}me version est \'{e}quip\'{e}e d’un \'{e}cran \`{a} cristaux liquides (LCD) \'{e}quip\'{e} de la technologie PixelSense qui permet une reconnaissance des interactions en entr\'{e}e sans usage de cam\'{e}ras. Cette technologie d\'{e}tecte les doigts, les tags ou les objets tangibles gr\^{a}ce \`{a} plusieurs dispositifs de r\'{e}tro \'{e}clairage
qui produisent une lumi\`{e}re infrarouge. Cette lumi\`{e}re est d’abord r\'{e}fl\'{e}chit sur les objets ou les doigts
pos\'{e}s sur la surface tactile. Ensuite, un capteur LCD est utilis\'{e} pour recr\'{e}er l’image de ce qui est
pos\'{e} sur la surface \`{a} partir de la lumi\`{e}re r\'{e}fl\'{e}chie. Enfin l’image est interpr\'{e}t\'{e}e par des techniques
d’analyse et le r\'{e}sultat est envoy\'{e} \`{a} une unit\'{e} centrale.
Les versions 1 et 2 des Microsoft PixelSense permettent des UI tactiles multi-contacts. Elle supporte jusqu’\`{a} 50 contacts simultan\'{e}ment et permet donc de d\'{e}crire des UI multi-utilisateurs dans la
limite des contacts support\'{e}s. Elle permet aussi de d\'{e}crire des UI tangibles par la reconnaissance des
tags et la forme des objets. L’utilisation de tags permet de ne pas limiter les objets \`{a} utiliser dans la
description des UI. Les tags sont des codes barres \`{a} deux dimensions qui sont facilement identifiables
par la table surface. La Surface offre deux types tags : Identity tags (tags avec une plage de valeurs
sur 128 bits) et Byte tags(tags avec une plage de valeurs sur 8 bits). Les tags peuvent \^{e}tre utilis\'{e}s pour
diff\'{e}rentes actions :
– reconna\^{i}tre des objets physiques ou les distinguer parmi plusieurs,
– d\'{e}clencher une commande ou une action, par exemple : afficher un menu ou une application
par exemple,
– pointer et orienter une application par un objet tagu\'{e}, par exemple dans le but de s\'{e}lectionner
un \'{e}l\'{e}ment.
Tableaux blancs interactifs
La version 2.0 des tables PixelSense peut \^{e}tre utilis\'{e}e comme un tableau blanc interactif. Les tableaux blancs interactifs[GMAD05] (TBI) constituent aussi des tables. Le
TBI Magic Table [Ber03] est constitu\'{e} d’une surface blanche, d’un vid\'{e}o projecteur et d’une cam\'{e}ra
reli\'{e} \`{a} une unit\'{e} centrale. La cam\'{e}ra est coupl\'{e}e \`{a} un composant logiciel capable d’interpr\'{e}ter les symboles de la surface blanche. Les interactions en entr\'{e}e sont possibles avec des jetons (qui sont aussi
3.1. MODÈLE D’INTERACTIONS D’UNE TABLE INTERACTIVE
27
des objets tangibles) pour cr\'{e}er, d\'{e}placer, redimensionner ou supprimer des \'{e}l\'{e}ments graphiques. Les
jetons sont associ\'{e}s \`{a} des gestes pr\'{e}d\'{e}finis qui seront d\'{e}tect\'{e}s par la cam\'{e}ra.
Magic Table est aussi utilis\'{e} pour num\'{e}riser facilement le contenu d’un tableau blanc. Dans le
cadre des UI graphiques, il peut \^{e}tre utilis\'{e} simultan\'{e}ment par plus d’une personne en manipulant des
jetons. Contrairement aux Microsoft PixelSense, le nombre d’objets tangibles diff\'{e}rents utilisables
par le TBI MagicTable est limit\'{e} au nombre de formes d\'{e}tectables par sa camera.
3.1.1.4
Constat
Les tables interactives ont de mani\`{e}re g\'{e}n\'{e}rale deux types de dispositifs physiques d’interactions
(cf tableau 3.1) : les dispositifs d’interactions en entr\'{e}e qui peuvent \^{e}tre tactiles et/ou tangibles et
les dispositifs d’interactions en sortie qui sont des surfaces d’affichage. Ces surfaces sont de tailles
variables et de diff\'{e}rentes formes (rectangulaire, circulaire, etc.). Elles peuvent \^{e}tre dispos\'{e}es de mani\`{e}re horizontale (pour DiamondTouch, metaDESK, Microsoft PixelSense 1.0 et 2.0) ou de mani\`{e}re
verticale (pour Microsoft PixelSense 2.0 ou les tableaux blancs interactifs).
Le nombre d’utilisateurs et leurs dispositions autour de la surface d’affichage sont des \'{e}l\'{e}ments
qui permettent de caract\'{e}riser les interactions des tables interactives. En effet ces deux caract\'{e}ristiques
impactent la conception des UI car une UI destin\'{e}e \`{a} plusieurs personnes doit permettre l’accessibilit\'{e}
des diff\'{e}rentes fonctionnalit\'{e}s \`{a} tous les utilisateurs. Elle doit aussi prendre en compte le partage
des \'{e}l\'{e}ments de l’UI (menus, visualisation des contenus, etc.). Il est important aussi de savoir si les
utilisateurs ont ou non un espace qui leur est propre.
Certaines tables interactives poss\`{e}dent des \'{e}quipements mat\'{e}riels leur permettant de concevoir
des UI multi-utilisateurs, co-localis\'{e}es, tactiles et tangibles. Si l’UI de l’application \`{a} migrer ne poss\`{e}de pas ces caract\'{e}ristiques, il est alors n\'{e}cessaire de r\'{e}soudre les diff\'{e}rents probl\`{e}mes soulev\'{e}s
par l’introduction des ces nouvelles contraintes r\'{e}sultant de l’utilisation d’une table interactive et en
particulier de d\'{e}cider du nombre d’utilisateurs qui pourront interagir avec l’application et de leurs
dispositions par rapport \`{a} la surface d’affichage.
Tables Interactives
Dispositifs d’interactions
d’entr\'{e}e
Dispositifs d’interactions de
sortie
Diamond Touch
Écran Capacitif
Vid\'{e}o Projecteur
metaDesk
Objet Tangible, Camera
Infrarouge, Cam\'{e}ra Vid\'{e}o
Vid\'{e}o Projecteur
Microsoft PixelSense V1\textbackslash{}&2
Technologie PixelSense (V2),
Cam\'{e}ra Infrarouge(V1), Tags,
Objets Tangibles
Écran LCD
TABLE 3.1 – Synth\`{e}se des dispositifs physiques d’interactions pour 3 tables interactives
3.1.2
Biblioth\`{e}ques graphiques des tables interactives
Les biblioth\`{e}ques graphiques sont des bo\^{i}tes \`{a} outils logiciels qui contiennent des \'{e}l\'{e}ments pour
construire des UI graphiques et d’utiliser l’ensemble des possibilit\'{e}s offertes par le terminal que ce
soit en entr\'{e}e (reconnaissance d’objets tangibles, tag, multi-touch, multi-utilisateurs) ou en sortie
(affichage). Nous pr\'{e}sentons donc dans cette section, les principaux \'{e}l\'{e}ments que l’on retrouve dans
les biblioth\`{e}ques graphiques associ\'{e}es aux tables \'{e}tudi\'{e}es.
28
CHAPITRE 3. TABLES INTERACTIVES ET MIGRATION DES UI
3.1.2.1
DiamondSpin
DiamondSpin [SVFR04] est une biblioth\`{e}que graphique pour table interactive qui offre des composants graphiques adapt\'{e}s \`{a} une utilisation de la table par plusieurs utilisateurs. Les composants graphiques sont bas\'{e}s sur JavaSwing et sont compatible avec OpenGL. Les composants graphiques de
cette biblioth\`{e}que permettent une manipulation directe des documents visuels en utilisant des doigts,
un stylet ou un clavier comme des moyens d’interactions. Cette biblioth\`{e}que graphique permet aussi
de cr\'{e}er et de g\'{e}rer des espaces priv\'{e}s associ\'{e}s \`{a} chaque utilisateur d’une application collaborative
et co-localis\'{e}e. Les composants graphiques DSContainer, DSPanel, DSWindow, DSFrame (cf. figure
3.4) permettent par exemple d’avoir un container qui regroupe un ensemble de composants graphiques
qu’un utilisateur peut d\'{e}placer en fonction de sa position.
FIGURE 3.4 – Instances des containers DiamondSpin
3.1.2.2
Surface SDK 1.0 et 2.0
[Mic12a] sont des API et des bo\^{i}tes \`{a} outils permettant de d\'{e}velopper des applications pour une
table interactive Microsoft (1.0 et 2.0). Ils font partie du Framework .Net et offrent des biblioth\`{e}ques
graphiques pour concevoir des UI WPF [Mic12b] et XNA [Mic12c]. Ces biblioth\`{e}ques graphiques
permettent la reconnaissance des formes d’objets, l’utilisation des tags, la gestion des 50 points de
contacts simultan\'{e}s, la d\'{e}tection de l’orientation des touches, etc. La construction d’une UI avec ce
SDK est bien plus ais\'{e} qu’avec la biblioth\`{e}que graphique DiamondSpin car elle \'{e}vite par exemple
au programmeur la gestion du d\'{e}placement des \'{e}l\'{e}ments ou la gestion de la rotation des composants
graphiques. En effet un ScatterView (Figure 3.5) d\'{e}finie le d\'{e}placement ou la rotation de mani\`{e}re
intrins\`{e}que.
3.1.2.3
Constat
La biblioth\`{e}que Surface SDK permet de d\'{e}crire des TUI \`{a} l’aide de tags. De mani\`{e}re g\'{e}n\'{e}rale,
les biblioth\`{e}ques graphiques des tables interactives offrent des composants graphiques qui ont des
interactions (rotation, redimensionnement, d\'{e}placement, tags, etc.) adapt\'{e}es pour la description des
UI pour les tables interactives. Elles permettent d’affiner les caract\'{e}ristiques des tables interactives en
3.1. MODÈLE D’INTERACTIONS D’UNE TABLE INTERACTIVE
29
FIGURE 3.5 – Exemple de ScatterView
pr\'{e}cisant si une table interactive supporte ou non des interactions tangibles, si elle a des composants
graphiques accessibles par tous les utilisateurs par exemple.
Les biblioth\`{e}ques graphiques servent aussi pour sp\'{e}cialiser une table interactive dans un domaine
applicatif pr\'{e}cis. Dans ce cas, les biblioth\`{e}ques graphiques offrent des composants graphiques sp\'{e}cifiques \`{a} des domaines d’application. Par exemple la table metaDesk est uniquement destin\'{e}e \`{a} la
cartographie, \`{a} ce titre elle dispose des composants graphiques propres pour afficher facilement un
plan. Cependant, DiamondSpin et les Surface SDK 1.0 \textbackslash{}& 2.0 sont des biblioth\`{e}ques graphiques ind\'{e}pendantes des domaines applicatifs. Dans le cadre de la migration des UI desktop ces derni\`{e}res offrent
des composants graphiques \'{e}quivalents \`{a} la source.
3.1.3
Modalit\'{e}s d’interactions
Nigay [Nig94] propose la d\'{e}finition d’une modalit\'{e} d’interaction comme un couple <d, l> constitu\'{e} d’un dispositif d’interactions et d’un langage d’interactions :
– d d\'{e}signe un dispositif physique. Par exemple, une souris, une cam\'{e}ra, un \'{e}cran, un hautparleur,
– l d\'{e}note un langage d’interaction constitu\'{e} d’un ensemble structur\'{e} de signes permettant de d\'{e}crire les fonctions de communication. Il peut reposer sur un langage pseudo-naturel, un graphe
ou une table de correspondance.
La migration des UI vers les tables interactives implique bien souvent un changement des dispositifs d’interactions et impose donc de d\'{e}finir des \'{e}quivalences entre les modalit\'{e}s d’interactions des
plateformes de d\'{e}part et celles des tables interactives. Dans cette section nous abordons la probl\'{e}matique li\'{e}e aux changements des modalit\'{e}s d’interactions des UI \`{a} migrer. En effet comment d\'{e}crire les
interactions de l’UI de d\'{e}part \`{a} l’aide des modalit\'{e}s d’interactions de la plateforme d’arriv\'{e}e ? Pour
r\'{e}pondre \`{a} cette question nous \'{e}tudions au paragraphe 3.1.1 les probl\`{e}mes li\'{e}es aux changements de
modalit\'{e}s d’interactions pendant la migration. Puis au paragraphe 3.1.2 nous pr\'{e}sentons quelques mod\`{e}les d’interactions abstraites qui permet de d\'{e}crire les interactions ind\'{e}pendamment des modalit\'{e}s
d’interactions.
3.1.3.1
Changement de modalit\'{e} d’interactions
Consid\'{e}rons comme plateforme de d\'{e}part un desktop compos\'{e} d’un \'{e}cran comme moyen d’interactions en sortie et d’un clavier et d’une souris comme moyen d’interactions en entr\'{e}e.
30
CHAPITRE 3. TABLES INTERACTIVES ET MIGRATION DES UI
– La modalit\'{e} d’interactions en sortie M1 du desktop est d\'{e}crite par l’\'{e}cran et par le langage
de description des UI 2D (L1), M1=<Ecran, L1>. Le langage L1 correspond aux composants
graphiques tels que labels, champ de texte, bo\^{i}te de dialogue, etc. d’une biblioth\`{e}que graphique.
– La modalit\'{e} d’interactions en entr\'{e}e M2 du desktop est d\'{e}crite par le clavier et par le langage
des commandes (L2), M2=<Clavier, L2>. Le langage L2 d\'{e}crit pour chaque action les t\^{a}ches
r\'{e}alisables \`{a} l’aide d’un clavier sur une UI graphique, ces actions sont par exemple : copier
avec Ctrl+C, couper avec Ctrl+X, coller avec Ctrl+V, saisir un texte avec les touches alphanum\'{e}riques, valider avec la touche entr\'{e}e, etc.
– La modalit\'{e} d’interactions en entr\'{e}e M3 du desktop est d\'{e}crite par la souris et par le langage
de manipulation directe d’une UI 2D (L3), M3=<Souris, L3>. Le langage L3 d\'{e}crit les actions
r\'{e}alisables \`{a} l’aide d’une souris sur une UI graphique, ces actions sont par exemple : cliquer
et valider pour s\'{e}lectionner un composant graphique, cliquer et d\'{e}placer pour s\'{e}lectionner un
texte, d\'{e}placer un \'{e}l\'{e}ment, redimensionner, etc.
Consid\'{e}rons maintenant que la table interactive cible est compos\'{e}e d’un \'{e}cran tactile, d’un clavier
virtuel et de la reconnaissance des objets physiques comme moyens d’interactions.
– La modalit\'{e} d’interactions en sortie M’1 de la table interactive est d\'{e}crite par l’\'{e}cran tactile
et par le langage de description des UI 2D (L’1), M’1= <Ecran Tactile, L’1>. Le langage L’1
correspond \`{a} la biblioth\`{e}que graphique de la table interactive qui contient les composants graphiques tels que labels, champ de texte, images, fen\^{e}tre, etc.
– La modalit\'{e} d’interactions en entr\'{e}e M’2 de la table interactive est d\'{e}crite par l’\'{e}cran tactile
et par le langage de manipulation tactile d’une UI 2D (L’2), M’2= < Ecran Tactile, L’2>. Le
langage L’2 d\'{e}crit les actions utilisateurs sur un \'{e}cran tactile. Ces actions sont par exemple :
toucher avec un doigt pour activer ou s\'{e}lectionner un \'{e}l\'{e}ment, toucher avec deux doigts et
d\'{e}placer pour agrandir, r\'{e}duire, tourner des composants graphiques, etc.
– La modalit\'{e} d’interactions en entr\'{e}e M’3 de la table interactive est d\'{e}crite par le clavier virtuel
et par le langage de commande (L’3), M’3= < Ecran Tactile, L’3>. Le langage L’3 d\'{e}crit les
actions utilisateurs r\'{e}alisables avec un clavier virtuel tel que saisir un texte.
– La modalit\'{e} d’interactions en entr\'{e}e M’4 de la table interactive est d\'{e}crite par l’\'{e}cran tactile et
par le langage de manipulation des objets tangibles (L’4), M’4= < Objets Tangibles, L’4>. Le
langage L’4 d\'{e}crit les actions utilisateurs sur un \'{e}cran tactile \`{a} l’aide des objets tangibles. Ces
actions sont par exemple : poser un objet pour afficher un menu ou un formulaire, d\'{e}placer un
objet physique pour d\'{e}placer l’objet virtuel associ\'{e}, tourner un objet physique pour s\'{e}lectionner
une fonctionnalit\'{e}, etc.
Les langages d’interactions L1 et L’1 permettent de d\'{e}crire les interactions en sortie des UI des
applications source et cible. Ces langages peuvent \^{e}tre mod\'{e}lis\'{e}s en se basant sur des bo\^{i}tes \`{a} outils
ind\'{e}pendantes des dispositifs d’interactions en sortie. Crease dans [Cre01] propose une bo\^{i}te \`{a} outils d\'{e}crivant des widgets multimodales et ind\'{e}pendantes des dispositifs d’interactions en sortie. Les
widgets de Crease [Cre01] restent cependant li\'{e}es aux dispositifs d’interactions en entr\'{e}e tels que le
clavier et la souris. Les langages de description des interactions en sortie peuvent \^{e}tre mod\'{e}lis\'{e}s ind\'{e}pendamment des dispositifs de sortie en prenant compte les diff\'{e}rentes modalit\'{e}s de sortie. Cependant
ces mod\`{e}les peuvent \^{e}tre li\'{e}s aux dispositifs d’interactions en entr\'{e}e.
Les langages d’interactions L2, L3, L’2, L’3 et L’4 quant \`{a} eux permettent de d\'{e}crire et d’interpr\'{e}ter les interactions en entr\'{e}e des utilisateurs des plateformes de d\'{e}part et d’arriv\'{e}es. Ces langages
d’interactions permettent d’associer \`{a} chaque action de l’utilisateur un comportement ayant un sens
dans l’UI de l’application. Les actions utilisateurs telles que cliquer, s\'{e}lectionner, Ctrl+C, poser un
objet, etc. d\'{e}pendent des dispositifs d’interactions tandis que les comportements de l’UI d\'{e}pendent du
type d’UI et de l’interpr\'{e}tation souhait\'{e}e par le concepteur.
3.1. MODÈLE D’INTERACTIONS D’UNE TABLE INTERACTIVE
31
Le changement des modalit\'{e}s d’interactions doit pr\'{e}server les actions de l’UI de d\'{e}part et
l’adapter aux dispositifs d’interactions de l’UI d’arriv\'{e}e.
3.1.3.2
Équivalences des modalit\'{e}s d’interactions
La migration de la plateforme desktop vers une table interactive peut \^{e}tre consid\'{e}r\'{e}e comme un
processus de changement de modalit\'{e}s d’interactions. En effet dans le but de r\'{e}utiliser l’UI d’une
application de d\'{e}part avec les dispositifs d’interactions de la plateforme d’arriv\'{e}e, il doit \^{e}tre possible d’\'{e}tablir des \'{e}quivalences entre les dispositifs d’interactions et les langages d’interactions des
plateformes de d\'{e}part et d’arriv\'{e}e.
Pour d\'{e}crire les interactions d’une UI de la plateforme de d\'{e}part \`{a} l’aide des dispositifs d’interactions de la plateforme d’arriv\'{e}e (table interactive), l’une des approches \`{a} envisager peut \^{e}tre la mise
en correspondance des dispositifs d’interactions en \'{e}tablissant des \'{e}quivalences entre les diff\'{e}rentes
modalit\'{e}s d’interactions des plateformes source et cible. Ce qui consiste par exemple \`{a} d\'{e}crire des
\'{e}quivalences d’abord entre les modalit\'{e}s d’interactions en sortie M1 et M’1, puis entre les modalit\'{e}s
d’interactions en entr\'{e}e M2, M3 et M’2, M’3 M’4.
L’\'{e}quivalence entre M1 et M’1 consiste \`{a} comparer les deux dispositifs assez proches, des \'{e}crans,
qui peuvent afficher des UI graphiques mais aussi les langages L1 et L’1 qui permettent la manipulation de composants graphiques appartenant \`{a} des biblioth\`{e}ques graphiques diff\'{e}rentes. L’\'{e}quivalence
entre les modalit\'{e}s d’interactions en entr\'{e}e M2, M3 et M’2, M’3, M’4 n’est possible que si l’on peut
comparer les langages L2, L3, L’2, L’3 et L’4.
3.1.3.3
Constat
Le changement de modalit\'{e}s d’interactions est possible en d\'{e}crivant des \'{e}quivalences entre les
diff\'{e}rentes modalit\'{e}s d’interactions d’une UI. Ce qui peut se faire en se basant sur un mod\`{e}le ind\'{e}pendant des dispositifs d’interactions.
Un mod\`{e}le d’interactions abstraites permet de d\'{e}crire des \'{e}quivalences entre deux modalit\'{e}s d’interactions en jouant un r\^{o}le de langage pivot. Dans notre cadre, un langage pivot permet d’\'{e}tablir des
\'{e}quivalences entre les modalit\'{e} d’interactions des plateformes source et cible.
3.1.4
Mod\`{e}le d’interactions abstraites pour la migration des UI
Dans cette section nous pr\'{e}sentons deux mod\`{e}les d’interactions abstraites et nous en d\'{e}duirons
quelques caract\'{e}ristiques d’un mod\`{e}le d’interactions abstraites pour la migration.
3.1.4.1
Mod\`{e}le de Vlist
Vlist et al.
[vdVNHF11] proposent les primitives d’interactions (interaction primitives) qui
constituent les plus petits \'{e}l\'{e}ments d’interactions identifiables ayant une relation significative avec
l’interaction elle-m\^{e}me. Selon cette approche, il est possible de d\'{e}crire les capacit\'{e}s d’interactions
des diff\'{e}rents dispositifs d’interactions \`{a} l’aide des primitives d’interactions, puis de d\'{e}crire ensuite,
\`{a} l’aide d’un mapping s\'{e}mantique, la relation entre les dispositifs et l’UI. Les diff\'{e}rentes primitives
d’interactions sont associ\'{e}es \`{a} chaque dispositif d’interactions. Elles sont choisies, identifi\'{e}es et associ\'{e}es aux dispositifs par le concepteur de l’UI. Par exemple les primitives d’interactions “Up” et
“Down” correspondent aux touches \'{e}tiquet\'{e}es ’+’ et ’-’, ces primitives d’interactions ont des types
de donn\'{e}es et des valeurs, elles sont mapp\'{e}es sur des composants graphiques tels que les contr\^{o}les
de volume pour permettre l’utilisation des composants avec ces touches. L’objectif de ce mapping
s\'{e}mantique est de faciliter son adaptation en fonction des contextes.
32
CHAPITRE 3. TABLES INTERACTIVES ET MIGRATION DES UI
Dans le cadre de la migration, la mise en correspondance des dispositifs d’interactions des plateformes de d\'{e}part et d’arriv\'{e}e semble \^{e}tre possible en utilisant une description s\'{e}mantique des capacit\'{e}s des dispositifs d’interactions d’une plateforme par les primitives d’interactions. Pour cela il est
n\'{e}cessaire de trouver une correspodance \`{a} l’ensemble des primitives d’interactions de la plateforme de
d\'{e}part. Par exemple, si “Up” et “Down” existent dans la plateforme de d\'{e}part, il faut les red\'{e}finir pour
un table interactive en inventant le geste \`{a} d\'{e}crire pour les activer. “Up” pourrait \^{e}tre repr\'{e}sent\'{e} par le
touch\'{e} simple et “Down” par le touch\'{e} double. Mais on peut aussi imaginer l’association toucher et
glisser, etc.
3.1.4.2
Mod\`{e}le de Gellersen
Gellersen [GH95] propose un mod\`{e}le d’interactions abstraites qui a pour objectif de d\'{e}crire les
interactions en entr\'{e}e et en sortie ind\'{e}pendamment des modalit\'{e}s d’interactions. Ce mod\`{e}le est aussi
ind\'{e}pendant d’un domaine ou d’une application. Il est pr\'{e}sent\'{e} \`{a} la figure 3.6 et il permet de d\'{e}crire
une hi\'{e}rarchie d’interactions en entr\'{e}e et en sortie. Les interactions en entr\'{e}e sont raffin\'{e}es en deux
cat\'{e}gories, les interactions d’entr\'{e}e de donn\'{e}es telles que Editor, Valuator (\'{e}diter du texte) et Option
(s\'{e}lectionner un \'{e}l\'{e}ment d’une liste) qui sont des sous-classes de Entry d’une part et les interactions
de Command 7 et de Signal 8 d’autre part. Les interactions en sortie sont aussi de deux types : les
messages (alertes, confirmation, etc.) et la vue qui permet l’affichage des donn\'{e}es et des composants
de l’UI.
 
 
FIGURE 3.6 – Mod\`{e}le d’interactions abstraites de Gellersen
Dans le cadre du changement de modalit\'{e} d’interactions entrain\'{e} par la migration, ce mod\`{e}le permet de d\'{e}crire les interactions des diff\'{e}rents composants graphiques de l’UI de mani\`{e}re ind\'{e}pendante
des modalit\'{e}s d’interactions. Par exemple une fen\^{e}tre d’authentification comportant des champs de
7. Les Command sont des interactions en entr\'{e}e de contr\^{o}le avec des param\`{e}tres (par exemple “copier texte”, “coller
texte”, etc.)
8. Les Signal sont des interactions en entr\'{e}e sans param\`{e}tres (exemple valider)
3.2. PRINCIPES DE CONCEPTION DES UI POUR LES TABLES INTERACTIVES
33
texte et des boutons, les champs de texte seront repr\'{e}sent\'{e}s par les Editor qui sont des interactions \`{a} la
fois en entr\'{e}e et en sortie et les boutons seront repr\'{e}sent\'{e}s par des Command dans un mod\`{e}le abstrait.
Ce mod\`{e}le identifie \`{a} la fois les interactions en entr\'{e}e et en sortie de haut niveau ind\'{e}pendamment
des modalit\'{e}s d’interactions. Dans le cadre de la migration de l’UI CBA par exemple, les interactions
telles que la rotation, le d\'{e}placement ou le redimensionnement qui sont li\'{e}s aux guidelines de la table
interactives doivent \^{e}tre interpr\'{e}t\'{e}s comme des commandes. Cette interpr\'{e}tation n’est pas g\'{e}n\'{e}rique
et d\'{e}pend du concepteur des interactions. En effet le d\'{e}placement peut \^{e}tre consid\'{e}r\'{e} comme une
rotation en modifiant l’angle par exemple. De mani\`{e}re g\'{e}n\'{e}rale, le mod\`{e}le de Gellersen est un mod\`{e}le d’interactions abstraites qui n\'{e}cessite la sp\'{e}cification de certaines interactions ind\'{e}pendantes des
modalit\'{e}s d’interactions pendant la migration.
3.1.4.3
Constats
Deux caract\'{e}ristiques sont essentielles pour les mod\`{e}les d’interactions abstraites :
– leurs ind\'{e}pendances par rapport \`{a} une modalit\'{e} d’interactions
– et leur capacit\'{e} \`{a} d\'{e}crire toutes les interactions du langage d’une modalit\'{e}.
Le changement des modalit\'{e}s d’interactions est une transformation de la dimension dialogue (cf
section 2.2.5) d’une UI de d\'{e}part en \'{e}tablissant des \'{e}quivalences entre les instruments des plateformes
source et cible. L’\'{e}quivalence des modalit\'{e}s d’interactions permet de garantir l’accessibilit\'{e} des dialogues de l’UI de d\'{e}part avec des instruments \'{e}quivalents (cf section 2.2.2.2). Cette transformation ne
prend pas en compte tous les probl\`{e}mes li\'{e}s \`{a} la transformation des dialogues.
Les probl\`{e}mes de la sous-dimensions des dialogues multi-utilisateurs et co-localis\'{e}s par exemple
ne peuvent pas \^{e}tre pris en compte par des \'{e}quivalences simples sans tenir compte du nombre d’utilisateurs et des types d’interactions. Le changement des modalit\'{e}s d’interactions ne prend pas non plus
en compte les autres dimensions (Structure et Positionnement, Style). La section 3.2 suivante identifie
les principes qui guident la transformation des dimensions.
3.2
Principes de conception des UI pour les tables interactives
La conception ou la transformation des diff\'{e}rentes dimensions des UI pour les tables interactives
est une activit\'{e} d’ing\'{e}nierie qui est guid\'{e}e par des crit\`{e}res ou des heuristiques afin d’assurer l’utilisabilit\'{e} du r\'{e}sultat final. Les crit\`{e}res ergonomiques de conception par exemple constituent des propri\'{e}t\'{e}s
g\'{e}n\'{e}riques (de haut niveau d’abstraction) applicables \`{a} toutes les plateformes. Pour mettre en place
un processus de migration r\'{e}utilisable, il est indispensable d’affiner les crit\`{e}res ergonomiques g\'{e}n\'{e}ralement d\'{e}crit en langage naturel en r\`{e}gles applicables pour transformer les diff\'{e}rentes dimensions
d’une UI.
L’affinement des crit\`{e}res abstrait se fait par rapport aux sp\'{e}cificit\'{e}s de l’environnement cibl\'{e}. Dans
le cadre des tables interactives, nous identifions \`{a} la section 3.2.1 les propri\'{e}t\'{e}s qui les caract\'{e}risent et
qui influencent la migration des UI. La section 3.2.2 pr\'{e}sente les crit\`{e}res ergonomiques de conception
et leur affinement par rapport aux propri\'{e}t\'{e}s caract\'{e}ristiques des tables interactives. La section 3.2.3
pr\'{e}sentent l’ensemble des recommandations pour la transformation des diff\'{e}rentes dimensions des UI
pendant la migration.
3.2.1
Propri\'{e}t\'{e}s caract\'{e}ristiques des tables interactives
Les \'{e}l\'{e}ments du mod\`{e}le d’interactions instrumentales des tables interactives tels que les dispositifs mat\'{e}riels d’interactions en entr\'{e}e et en sortie, les biblioth\`{e}ques graphiques ou les modalit\'{e}s
34
CHAPITRE 3. TABLES INTERACTIVES ET MIGRATION DES UI
d’interactions nous permettent d’identifier des propri\'{e}t\'{e}s, caract\'{e}ristiques des tables interactives, qui
impactent la migration des UI vers des tables interactives.
3.2.1.1
Dispositifs d’interactions en entr\'{e}e
Ces dispositifs influencent la conception des dialogues. Dans le cadre des tables interactives nous
identifions deux propri\'{e}t\'{e}s qui correspondent aux moyens d’interactions tangibles et tactiles.
Propri\'{e}t\'{e} 1 Tangibilit\'{e} des interactions
Les dispositifs d’interactions en entr\'{e}e permettent d’associer des objets physiques aux fonctionnalit\'{e}s ou aux composants graphiques. En effet, les objets physiques peuvent aussi \^{e}tre
utilis\'{e}s pour activer des fonctionnalit\'{e}s et pour afficher ou d\'{e}placer des objets virtuels d’une
UI.
Propri\'{e}t\'{e} 2 Tactibilit\'{e} des interactions
Les dispositifs d’interactions en entr\'{e}e des tables interactives permettent de d\'{e}crire des
manipulations directes des UI telles que la s\'{e}lection, l’\'{e}dition, le redimensionnement, le
d\'{e}placement, etc.
3.2.1.2
Dispositifs d’interactions en sortie
Ce sont les surfaces d’affichage des tables interactives, les propri\'{e}t\'{e}s caract\'{e}ristiques li\'{e}es \`{a} la
taille et \`{a} la disposition des tables interactives qui impactent la conception de la structure et du positionnement des \'{e}l\'{e}ments des UI.
Propri\'{e}t\'{e} 3 Taille de la surface d’affichage
Cette propri\'{e}t\'{e} permet d’adapter la taille des composants graphiques par rapport \`{a} la taille
de l’\'{e}cran pour faciliter l’utilisation de l’UI.
Propri\'{e}t\'{e} 4 Disposition de la surface d’affichage
La surface d’affichage peut \^{e}tre dispos\'{e}e de mani\`{e}re horizontale comme une table de travail
ou de mani\`{e}re verticale comme un tableau collaboratif. Ces dispositions influencent l’orientation et l’utilisation des composants graphiques d’une UI en particulier si des utilisateurs
peuvent \^{e}tre tout autour de la table.
3.2.1.3
Utilisateurs des tables interactives
Le nombre d’utilisateurs influencent la conception des dialogues et de la structure des UI. En effet
le nombre permet de savoir si les dialogues de l’UI favorisent la collaboration ou ils sont destin\'{e}s
\`{a} un utilisateurs. Le nombre permet aussi de savoir si la structure de l’UI est accessible \`{a} plusieurs
utilisateurs ou si elle est destin\'{e}e \`{a} un utilisateur.
3.2. PRINCIPES DE CONCEPTION DES UI POUR LES TABLES INTERACTIVES
35
Propri\'{e}t\'{e} 5 Nombre d’utilisateurs
Cette propri\'{e}t\'{e} permet la description des dialogues et de la structure des \'{e}l\'{e}ments graphiques pour faciliter la collaboration et l’accessibilit\'{e} de l’UI des tables interactives.
Propri\'{e}t\'{e} 6 R\'{e}partition des utilisateurs
Cette propri\'{e}t\'{e} permet de savoir comment d\'{e}crire la collaboration entre les utilisateurs autour de la table. La r\'{e}partition de l’espace de travail de chaque utilisateur peut se faire par
une division g\'{e}ographique de la surface ou permettre aux utilisateurs d’acc\'{e}der \`{a} toute la
surface.
3.2.1.4
Types d’UI des tables interactives
Les tables interactives pr\'{e}sent\'{e}es \`{a} la section 3.1 pr\'{e}sentent des UI de types collaboratifs et UI
tangibles. Elles sont class\'{e}es dans le tableau 3.2.
Tables Interactives
Utilisateurs
Instruments
d’interactions
Types d’UI
Diamond
Touch [DL01]
1 \`{a} 4
Écran non capacitif,
DiamondSpin
UI Tactile, UI
Collaborative (Id
Utilisateurs)
metaDesk [UI97]
1 ou plusieurs
Écran, Camera
Infrarouge, Objets
Tangibles
UI Tangible, UI
Collaborative
Microsoft PixelSense
[Mic09]
1 ou Plusieurs (<= de
points de contact)
Écran capacitif, Tags et
Formes Objets,
SurfaceSDK
UI Tactile, UI
Tangible, UI
Collaborative
TABLE 3.2 – Type d’UI pour les tables interactives
3.2.2
Crit\`{e}res ergonomiques de conception des UI
Les propri\'{e}t\'{e}s caract\'{e}ristiques identifi\'{e}es ci-dessus (Tangibilit\'{e} des interactions, Tactibilit\'{e} des
interactions, Taille de la surface d’affichage, Disposition de la surface d’affichage, Nombre d’utilisateurs et R\'{e}partition des utilisateurs) nous permettent d’affiner les principaux crit\`{e}res ergonomiques
g\'{e}n\'{e}riques dans le cadre des tables interactives. Dans l’objectif de proposer un ensemble de principes
pour la migration des UI, nous pr\'{e}sentons dans cette section une grille d’affinement des crit\`{e}res ergonomiques de conception. Nous nous basons sur les huit crit\`{e}res ergonomiques de conception propos\'{e}s
par Scapin [Sca86] :
1. La compatibilit\'{e} est un principe qui repose sur le fait que les transferts d’information seront
plus rapides et plus efficaces si le recodage d’informations est r\'{e}duit [Sca86]. Ce principe pr\'{e}conise d’avoir des d\'{e}nominations et des commandes compatibles avec le vocabulaire de l’utilisateur.
2. L’homog\'{e}n\'{e}it\'{e} est un principe qui repose sur le fait que la prise de d\'{e}cision, le choix des
solutions, le rappel, etc., peuvent se r\'{e}p\'{e}ter de fa\c{c}on d’autant plus satisfaisante que l’environnement est constant [Sca86]. Ce principe pr\'{e}conise d’avoir des s\'{e}quences de commandes
36
CHAPITRE 3. TABLES INTERACTIVES ET MIGRATION DES UI
identiques pour arriver au m\^{e}me r\'{e}sultat et s’assurer que les labels, prompts et autres cat\'{e}gories
d’informations sont toujours localis\'{e}s au m\^{e}me endroit quelque soit l’\'{e}cran.
3. La concision est un principe qui repose sur l’existence d’une limite en m\'{e}moire \`{a} court terme
de l’op\'{e}rateur humain. Il convient donc de r\'{e}duire la charge mn\'{e}sique de l’utilisateur [Sca86].
L’application de ce principe pr\'{e}conise par exemple de r\'{e}duire les informations longues et trop
nombreuses, de pr\'{e}f\'{e}rer des proc\'{e}dures courtes pour favoriser la tache de m\'{e}morisation de
l’utilisateur.
4. La flexibilit\'{e} est une exigence li\'{e}e \`{a} l’existence de variations au sein de la population des utilisateurs. Il est pr\'{e}f\'{e}rable que le logiciel comporte diff\'{e}rents niveaux et qu’il prenne en consid\'{e}ration l’acquisition de l’exp\'{e}rience des utilisateurs [Sca86]. Par exemple l’usage du ctrl-c,
ctrl-v et ctrl-x rel\`{e}ve du respect de ce principe. Un utilisateur novice utilisera les menus tandis
qu’un utilisateur aguerri privil\'{e}giera les raccourcis.
5. Le feedback et le guidage reposent sur l’influence de la connaissance du r\'{e}sultat sur la qualit\'{e}
de la performance. Les r\'{e}sultats des actions des utilisateurs doivent toujours \^{e}tre r\'{e}percut\'{e}s
de fa\c{c}on explicite [Sca86]. Il n’est pas par exemple pas possible de d\'{e}placer une souris sans
feedback.
6. La charge de travail est un principe qui pr\'{e}conise que la prise en compte de la charge informationnelle de l’utilisateur est essentielle dans la mesure o\`{u} la probabilit\'{e} d’erreur humaine
augmente dans les situations \`{a} charge \'{e}lev\'{e}e [Sca86]. Par exemple il convient de minimiser le
nombre d’op\'{e}rations \`{a} effectuer par l’utilisateur ainsi que les temps de traitement.
7. Le contr\^{o}le explicite signifie que m\^{e}me si c’est le logiciel qui a le contr\^{o}le, l’interface doit
appara\^{i}tre comme \'{e}tant sous le contr\^{o}le de l’utilisateur et surtout ex\'{e}cuter, en r\`{e}gle g\'{e}n\'{e}rale,
des op\'{e}rations uniquement \`{a} la suite d’actions explicites [Sca86].
8. La gestion des erreurs est un principe qui consiste \`{a} fournir aux utilisateur des moyens pour
corriger leurs erreurs [Sca86]. La clart\'{e} des messages est un \'{e}l\'{e}ment fondamental.
Pour affiner ces huit crit\`{e}res ergonomiques \`{a} l’aide des propri\'{e}t\'{e}s caract\'{e}ristiques des tables interactives, nous utilisons le tableau 3.3 ci-dessous. Ce tableau nous permet de savoir comment une
propri\'{e}t\'{e} caract\'{e}ristique des tables influence un crit\`{e}re ergonomique. Les influences possibles sont :
– F : une propri\'{e}t\'{e} favorise un crit\`{e}re ergonomique de conception,
– RA : une propri\'{e}t\'{e} risque d’alt\'{e}rer un crit\`{e}re ergonomique de conception,
– C : une propri\'{e}t\'{e} conserve un crit\`{e}re ergonomique de conception,
– N : une propri\'{e}t\'{e} est neutre par rapport \`{a} un crit\`{e}re ergonomique de conception.
3.2.2.1
Exemple de lecture du tableau d’affinement
La tangibilit\'{e} favorise la flexibilit\'{e} dans le cadre de la migration des UI vers les tables interactives
car elle ajoute une nouvelle modalit\'{e} d’interactions par rapport \`{a} l’UI de d\'{e}part.
Le nombre d’utilisateurs risque d’alt\'{e}rer le feedback et le guidage de l’UI cible apr\`{e}s la migration
si les affichages des feedbacks sont bloquants pour les autres utilisateurs dans un contexte multiutilisateurs.
La tactibilit\'{e} permet de conserver la flexibilit\'{e} car les interactions tactiles et la manipulation directe (souris) des UI de la source sont \'{e}quivalentes.
La taille de la surface d’affichage est une propri\'{e}t\'{e} neutre par rapport \`{a} la compatibilit\'{e} car elle
n’implique pas de modifications des \'{e}l\'{e}ments du vocabulaire 9 de l’utilisateur final de l’UI.
9. Commandes et les d\'{e}nominations des \'{e}l\'{e}ments de l’UI
3.2. PRINCIPES DE CONCEPTION DES UI POUR LES TABLES INTERACTIVES
37
Tangibilit\'{e}
Tactibilit\'{e}
Taille
Disposition
Nombre
R\'{e}partition
Compatibilit\'{e}
F
RA
N
N
N
N
Homog\'{e}n\'{e}it\'{e}
RA
RA
RA
N
RA
RA
Concision
F
F
RA
N
RA
F
Flexibilit\'{e}
F
C
F
F
RA
F
Feedback
\textbackslash{}&
Guidage
N
C
F
C
RA
RA
Charge de Travail
F
F
RA
N
N
N
Contr\^{o}le Explicite
F
F
F
RA
RA
F
Gestion des erreurs
C
C
C
N
RA
RA
TABLE 3.3 – Liens entre crit\`{e}res ergonomiques de Scapin et les propri\'{e}t\'{e}s caract\'{e}ristiques des tables
interactives.
3.2.2.2
Constats
La projection des crit\`{e}res ergonomiques de conception sur les propri\'{e}t\'{e}s caract\'{e}ristiques des tables
int\'{e}ractives telles que pr\'{e}sent\'{e} dans le tableau 3.3 a pour objectif d’\'{e}mettre des recommandations
suivant les crit\`{e}res ergonomiques de conception avec pour objectif de garantir l’utilisabilit\'{e} de l’UI,
les coh\'{e}rences des dialogues et du style.
Dans le cas o\`{u} une propri\'{e}t\'{e} favorise (F) un crit\`{e}re ergonomique, les recommandations doivent
\^{e}tre d\'{e}crites pour transformer l’UI source et elles sont respect\'{e}es sans conditions. Par exemple pour
savoir comment les interactions tangibles favorisent la r\'{e}duction de la charge de travail d’une UI
source, il est indispensable de d\'{e}crire comment afficher et cacher des \'{e}l\'{e}ments des UI \`{a} l’aide d’objets
tangibles.
Dans le cas o\`{u} une propri\'{e}t\'{e} risque d’alt\'{e}rer (RA) un crit\`{e}re ergonomique de conception, il est
indispensable que les recommandations associ\'{e}es identifient les contraintes pour \'{e}viter les risques. Par
exemple pour savoir comment la propri\'{e}t\'{e} Nombre d’utilisateurs peut alt\'{e}rer le Feedback et le Guidage, il est indispensable d’identifier les \'{e}l\'{e}ments bloquants de l’UI source et proposer une mani\`{e}re
de les migrer.
Dans les cas o\`{u} une propri\'{e}t\'{e} conserve (C) un crit\`{e}re ergonomique de conception alors les \'{e}quivalences entre les \'{e}l\'{e}ments de l’UI source et cible suffisent pour assurer la conformit\'{e} \`{a} ces crit\`{e}res ergonomiques de conception. Il n’est pas indispensable de d\'{e}crire des recommandations. Par exemple les
\'{e}quivalences entre les modalit\'{e}s d’interactions permettent de pr\'{e}server la flexibilit\'{e} d’une UI source
et les \'{e}quivalences entre les \'{e}l\'{e}ments graphiques permettent de pr\'{e}server la gestion des erreurs des UI
sources.
Dans le cas o\`{u} une propri\'{e}t\'{e} est neutre (N) par rapport \`{a} un crit\`{e}re de conception alors aucune
recommandation n’est d\'{e}crite. Par exemple la taille de la surface d’affichage n’affecte pas la compatibilit\'{e} car les \'{e}l\'{e}ments du vocabulaire ne sont pas modifi\'{e}s par cette propri\'{e}t\'{e}.
Les propri\'{e}t\'{e}s qui favorisent ou qui risquent d’alt\'{e}rer des crit\`{e}res ergonomiques de conception
nous permettent d’identifier des recommandations de haut niveau pour les diff\'{e}rentes dimensions des
UI. Nous proposons dans la section suivante de d\'{e}finir et d’identifier les recommandations pour la
transformation des diff\'{e}rentes dimensions des UI \`{a} migrer.
38
CHAPITRE 3. TABLES INTERACTIVES ET MIGRATION DES UI
3.2.3
Recommandations pour la migration des UI vers les tables interactives
La section 3.1.3 met en \'{e}vidence les diff\'{e}rences de modalit\'{e}s d’interactions entre un desktop
et une table interactive. Ces diff\'{e}rences impactent aussi la conception des UI pour ces deux plateformes. Besacier et al. montrent que la r\'{e}utilisation des applications desktop sur les tables interactives
en adaptant les \'{e}l\'{e}ments de l’UI aux m\'{e}taphores du papier par exemple facilite l’utilisation des UI
[BRNB07].
Par ailleurs, une r\'{e}utilisation d’une application desktop sur des tables interactives sans prise en
compte de ces sp\'{e}cificit\'{e}s pose deux probl\'{e}matiques majeures : la transformation de l’UI de d\'{e}part
en UI collaborative d’une part et la transformation d’une GUI en TUI d’autre part. Ces deux caract\'{e}ristiques font partis de l’ensemble des guidelines qui guident la conception des UI pour les tables
interactives.
“A Guideline is a source of inspiration or a set of standards that gives recommandations, tips and
directives” [Que01]
Les guidelines pour la conception des UI constituent un ensemble de recommandations pour les
concepteurs des UI qui indiquent comment d\'{e}crire les aspects tels que les interactions instrumentales,
le layout (ou le placement des \'{e}l\'{e}ments graphiques) et aussi les styles de pr\'{e}sentations (couleurs,
polices, tailles, etc.).
Les guidelines de haut niveau doivent \^{e}tre traduites en r\`{e}gles formelles utilisables pendant la migration [Van97]. La synth\`{e}se pr\'{e}sent\'{e}e par le tableau 3.3 d’affinement des crit\`{e}res ergonomiques de
conception dans le cadre des tables interactives nous permet d’avoir un lien entre les crit\`{e}res ergonomiques de conception et les guidelines.
Dans cette section nous caract\'{e}risons les guidelines pour la migration des UI vers les tables interactives en deux cat\'{e}gories : les guidelines pour les UI tangibles et les guidelines pour les UI collaboratives et co-localis\'{e}es.
3.2.3.1
Corpus des recommandations de conception des UI pour Microsoft PixelSense
Les recommandations de conception des UI pour une table interactive Microsoft PixelSense
sont d\'{e}crites dans le document User Experience Design Guideline [Mic11]. Elles sont le fruit des
exp\'{e}riences des d\'{e}veloppeurs des premi\`{e}res applications sur tables interactives. L’objectif de ces
guidelines est de faciliter la conception des interfaces utilisateurs pour qu’elles soient plus naturelles [Res13] et plus intuitives. Ces guidelines couvrent plusieurs aspects du processus de conception
des UI tels que la conception des interactions entre les UI et les utilisateurs finaux, les guides de styles
pour une coh\'{e}rence visuelle, les guides d’utilisation des textes, etc.
3.2.3.2
Guidelines pour UI collaborative
Cette cat\'{e}gorie regroupe les recommandations pour la conception d’une UI collaborative et colocalis\'{e}e pour une table interactive. Les guidelines li\'{e}es \`{a} la propri\'{e}t\'{e} 5 du nombre d’utilisateurs
risquent d’alt\'{e}rer l’homog\'{e}n\'{e}it\'{e}, la concision, la flexibilit\'{e}, le feedback, le contr\^{o}le explicite et la
gestion des erreurs (cf tableau 3.3). Par ailleurs les guidelines li\'{e}es \`{a} la propri\'{e}t\'{e} 6 de la r\'{e}partition
des utilisateurs risquent d’alt\'{e}rer l’homog\'{e}n\'{e}it\'{e}, le feedback et la gestion des erreurs et elles favorisent
aussi le contr\^{o}le explicite (cf tableau 3.3). Les guidelines li\'{e}es \`{a} la propri\'{e}t\'{e} 3 de la taille de la surface
d’affichage risquent d’alt\'{e}rer l’homog\'{e}n\'{e}it\'{e}, la concision et la charge de travail, elles favorisent la
flexibilit\'{e} et le feedback (cf tableau 3.3). Enfin la propri\'{e}t\'{e} 4 de la disposition de la surface d’affichage
risque d’alt\'{e}rer le contr\^{o}le explicite et elle favorise la flexibilit\'{e} (cf tableau 3.3).
Les guidelines, pour favoriser la collaboration des UI migr\'{e}es, peuvent \^{e}tre d\'{e}clin\'{e}es suivant les
trois dimensions d’une UI :
3.2. PRINCIPES DE CONCEPTION DES UI POUR LES TABLES INTERACTIVES
39
– Les guidelines issues de la dimension des Dialogues (D1) permettent de d\'{e}crire les dialogues
entre les utilisateurs.
– Les guidelines issues de la dimension de la Structure et du Positionnement (D2) permettent
d’avoir une UI avec un espace de travail qui favorise le partage.
– Les guidelines issues de la dimension du Style (D3) favorisent aussi le partage dans un espace
de travail gr\^{a}ce \`{a} des contraintes sur la taille des \'{e}l\'{e}ments graphiques.
Guideline 1 Couplage des dialogues d’UI avec les utilisateurs
Cette guideline pr\'{e}conise d’adapter les composants graphiques d’une UI au
nombre d’utilisateurs en \'{e}liminant toutes les activit\'{e}s bloquantes pour les autres
utilisateurs de l’UI (bo\^{i}tes de dialogues ou fen\^{e}tres bloquantes).
Guideline 2 Partage de l’espace de travail
Cette guideline pr\'{e}conise de prendre en compte le nombre d’utilisateurs de l’UI
pendant la migration. Le nombre d’utilisateurs est un facteur important pour le
choix du comportement, de la structure, du positionnement et de la taille des
\'{e}l\'{e}ments d’une UI sur les tables interactives. Pour permettre le partage et la collaboration, l’espace de travail doit :
– avoir des \'{e}l\'{e}ments avec des tailles variables par les utilisateurs,
– avoir une structure (types de donn\'{e}es) et un regroupement d’\'{e}l\'{e}ments des UI
adapt\'{e}es,
– avoir des \'{e}l\'{e}ments avec des comportements (rotation, d\'{e}placement) pour favoriser le partage. Dans ce cas la propri\'{e}t\'{e} 360˚favorise l’accessibilit\'{e} des \'{e}l\'{e}ments graphiques.
Illustration
La figure 3.7 illustre un cas d’utilisation des composants graphiques utilisables \`{a} 360˚.
En effet, la figure de droite (barr\'{e}e d’un trait oblique rouge) montre le cas d’un composant sans la
propri\'{e}t\'{e} 360˚, ce composant oblige des utilisateurs \`{a} avoir une position fixe autour de la table. Cette
situation n’est pas recommand\'{e}e par la guideline d’accessibilit\'{e} des composants graphiques.
La figure de gauche montre des composants graphiques qui impl\'{e}mentent la propri\'{e}t\'{e} 360˚.
FIGURE 3.7 – Illustration de la propri\'{e}t\'{e} 360˚des composants graphiques sur une table interactive
Exemples
En consid\'{e}rant que l’UI de l’application CBA (cf. figure 2.2) est migr\'{e}e vers une table
Surface pour quatre dessinateurs de BD qui peuvent librement s’installer autour de la table. Les ressources (images, bulles, etc.) utilis\'{e}es pour la conception des BD doivent \^{e}tre accessibles par les quatre
40
CHAPITRE 3. TABLES INTERACTIVES ET MIGRATION DES UI
utilisateurs. La guideline 2 pr\'{e}conisant le partage de l’espace de travail et la guideline 2 pr\'{e}conisant
l’utilisation des objets 360˚permettent de d\'{e}crire une UI sans layout avec des groupes d’\'{e}l\'{e}ments
utilisables \`{a} 360˚. La guideline 1 par exemple permettra de supprimer les diff\'{e}rents composants graphiques bloquants tels que les bo\^{i}tes de dialogues de l’UI de l’application CBA.
Synth\`{e}se
Les guidelines qui favorisent la collaboration pendant la migration peuvent \^{e}tre repr\'{e}sent\'{e}es par le diagramme de la figure 3.8. Les rectangles repr\'{e}sentent les guidelines \`{a} diff\'{e}rents niveaux
d’abstraction. Les liens entre les rectangles repr\'{e}sentent les dimensions et les sous-dimensions des
UI.
FIGURE 3.8 – Repr\'{e}sentation synth\'{e}tique des guidelines selon les trois dimensions des UI
3.2.3.3
Guidelines pour TUI
Cette cat\'{e}gorie regroupe les recommandations et les contraintes qui permettent de d\'{e}crire le comportement des \'{e}l\'{e}ments de l’UI qui sont des objets virtuels d’une part et l’association entre ces objets
virtuels et les objets tangibles d’autre part. L’association peut se faire par des tags qui permettent de
marquer les objets physiques ou en se basant sur la forme des objets physiques. Les guidelines de
cette cat\'{e}gorie sont inspir\'{e}es par la propri\'{e}t\'{e} de la tangibilit\'{e} des interactions. Cette propri\'{e}t\'{e} risque
d’alt\'{e}rer l’homog\'{e}n\'{e}it\'{e} et elle favorise la concision, la flexibilit\'{e}, la charge de travail et le contr\^{o}le
explicite (cf tableau 3.3) .
Les guidelines pour TUI se d\'{e}clinent suivant la dimension des Dialogues (D1) et la dimensions de
la Structure et du Positionnement (D2) des \'{e}l\'{e}ments graphiques. Pour ajouter les interactions tangibles
sur les UI sources, nous pr\'{e}conisons dans cette section deux guidelines :
– La guideline pour associer un objet tangible \`{a} une fonctionnalit\'{e} est issue de la sous-dimension
D11 des Interactions Tangibles.
– La guideline pour associer un objet tangible \`{a} un \'{e}l\'{e}ment (ou groupe d’\'{e}l\'{e}ments) virtuel est
issue de la sous-dimension D11 des Interactions Tangibles et la sous-dimension D21 li\'{e}e \`{a}
l’Accessibilit\'{e} des \'{e}l\'{e}ments graphiques. Les objets tangibles facilitent l’acc\`{e}s aux composants
graphiques.
3.2. PRINCIPES DE CONCEPTION DES UI POUR LES TABLES INTERACTIVES
41
Guideline 3 Objets tangibles des objets virtuels
Cette guideline pr\'{e}conise d’associer le comportement des \'{e}l\'{e}ments \`{a} un objet
tangible. Les \'{e}l\'{e}ments (ou objets virtuels) de l’UI \`{a} migrer doivent \^{e}tre associ\'{e}s
\`{a} des objets physiques dans le but d’afficher des menus, des formulaires ou des
fen\^{e}tres et aussi dans le but d’activer ou d’utiliser des fonctionnalit\'{e}s.
Guideline 4 Objets tangibles et fonctionnalit\'{e}s
Les objets tangibles peuvent servir \`{a} activer une fonctionnalit\'{e}. Cette guideline
pr\'{e}conise quelles fonctionnalit\'{e}s associer \`{a} un objet tangible.
– Les fonctionnalit\'{e}s \`{a} associer doivent \^{e}tre s\'{e}lectionnables et activables \`{a} partir
de l’UI de d\'{e}part (les \'{e}l\'{e}ments d’un menu par exemple).
– Les param\`{e}tres des fonctionnalit\'{e}s doivent \^{e}tre connus ou finis (un formulaire
ou une liste par exemple).
Exemple
En consid\'{e}rant l’UI de l’application CBA \`{a} la figure 2.2, le menu principal et le formulaire Ressources peuvent \^{e}tre associ\'{e}s \`{a} un objet physique pour les afficher facilement sur l’\'{e}cran.
Les r\`{e}gles formelles issues de la guideline 3 permettront d’identifier et de transformer les \'{e}l\'{e}ments
concrets de l’UI. Les tags permettent par exemple d’utiliser deux objets de la m\^{e}me forme avec des
couleurs diff\'{e}rentes et marqu\'{e}s des deux diff\'{e}rents tags pour afficher un menu ou un formulaire.
Synth\`{e}se
Les guidelines pour l’utilisation des objets tangibles peuvent \^{e}tre repr\'{e}sent\'{e}es par le diagramme de la figure 3.9. Les rectangles repr\'{e}sentent les guidelines \`{a} diff\'{e}rents niveaux d’abstraction.
Les liens entre les rectangles repr\'{e}sentent les dimensions et les sous-dimensions des UI.
3.2.3.4
Autres guidelines pour les UI sur tables interactives
Les deux cat\'{e}gories de guidelines pr\'{e}sent\'{e}es ci-dessus qui permettent de d\'{e}crire des UI tangibles
et des UI collaboratives (cf. section 3.2.3.2) ne constituent pas le corpus des guidelines applicables
pendant la migration d’une UI vers une table interactive. En effet les aspects de l’UI li\'{e}s au style des
composants graphiques, des textes de l’UI et les aspects li\'{e}s aux interactions tactiles ne sont pas pris
en compte par ces deux cat\'{e}gories de guidelines. Cette troisi\`{e}me cat\'{e}gorie regroupe les guidelines de
mise en œuvre des interactions tactiles et des styles de l’UI.
Guidelines pour les interactions tactiles
Les tables interactives permettent en g\'{e}n\'{e}ral de d\'{e}crire
des interactions tactiles. L’ensemble des actions tactiles de l’utilisateur est interpr\'{e}t\'{e} par le dispositif
d’interactions en entr\'{e}e. Dans le cadre des UI tactiles, il existe des actions tactiles standards pour des
interactions de redimensionnement, de d\'{e}placement ou de rotation. Par exemple les guidelines des interactions tactiles pour une table Surface [Mic11] identifient l’ensemble des gestes recommand\'{e}s aux
concepteurs pour la s\'{e}lection, l’activation, le d\'{e}placement, la rotation, le zoom, etc. Cette cat\'{e}gorie
de guidelines est li\'{e}e \`{a} la propri\'{e}t\'{e} 2 de tactibilit\'{e} des interactions d’une table interactive.
Dans le cadre de la migration de l’UI de l’application CBA vers les tables interactives, il est
indispensable de pouvoir d\'{e}crire des correspondances entre les interactions tactiles et les interactions
du clavier et de la souris de l’UI de d\'{e}part.
42
CHAPITRE 3. TABLES INTERACTIVES ET MIGRATION DES UI
FIGURE 3.9 – Repr\'{e}sentation synth\'{e}tique des guidelines suivant deux dimensions
Formulaire de l’application de d\'{e}part
Formulaire migr\'{e} en prenant en compte les
guidelines de styles
 
 
TABLE 3.4 – Migration de l’aspect visuel d’un formulaire
Guidelines pour le style
Cette sous-cat\'{e}gorie regroupe l’ensemble des recommandations pour la
personnalisation des aspects visuels d’une UI pour une table interactive. Ces guidelines peuvent \^{e}tre
utilis\'{e}es dans le cadre d’une migration faisant intervenir les concepteurs comme des exemples pour les
inspirer du choix de la forme, des couleurs, des ic\^{o}nes et aussi de la disposition des textes d’une UI.
Dans le cas du sc\'{e}nario de migration, le formulaire Ressources par exemple peut \^{e}tre migr\'{e} comme
l’indique le tableau 3.4.
3.3. SYNTHÈSE
43
3.3
Synth\`{e}se
Dans ce chapitre, nous avons \'{e}tudi\'{e} le mod\`{e}le d’interactions des tables interactives. Cette \'{e}tude
nous montre que la migration des UI vers des tables interactives provoque dans la majorit\'{e} des cas
un changement de modalit\'{e} d’interactions. Ce changement peut \^{e}tre effectu\'{e} en d\'{e}crivant un mod\`{e}le
d’interactions abstraites qui est ind\'{e}pendant d’une modalit\'{e} d’interactions. Ce mod\`{e}le doit aussi \^{e}tre
capable de d\'{e}crire toutes les interactions du langage li\'{e}s aux diff\'{e}rentes modalit\'{e}. L’objectif d’un
mod\`{e}le d’interactions abstraites est d’\'{e}tablir des \'{e}quivalences entre les instruments d’interactions des
plateformes source et cible.
Nous avons aussi identifi\'{e} dans ce chapitre diff\'{e}rentes r\`{e}gles de bons usage (’guidelines’) qui
vont guider la transformation des UI desktop en une UI tangible et collaborative. Les guidelines sont
d\'{e}duites \`{a} partir des propri\'{e}t\'{e}s caract\'{e}ristiques des tables interactives et des crit\`{e}res ergonomiques de
conception de haut niveau. Les guidelines sont constitu\'{e}es des recommandations et/ou des contraintes
pour transformer les diff\'{e}rentes dimensions des UI de la source. Dans le chapitre, nous avons remarqu\'{e} que les guidelines qui favorisent la collaboration permettent de transformer les trois dimensions
(Dialogues, Structure et Positionnement, Style). Par ailleurs, les guidelines pour les TUI concernent
la dimension des Dialogues et la dimension de la Structure et du Positionnement.
Les processus de transformation des diff\'{e}rentes dimensions sont d\'{e}crits par des approches de
migration des UI. Nous pr\'{e}sentons dans le chapitre suivant les diff\'{e}rentes approches de migrations
dans le but d’identifier les m\'{e}canismes de migration des UI ad\'{e}quats pour nos objectifs.
44
CHAPITRE 3. TABLES INTERACTIVES ET MIGRATION DES UI
CHAPITRE 4
Approches de migration des UI
Sommaire
4.1
Introduction . . . . . . . . . . . . . . . . . . . . . . . . . . . . . . . . . . . . .
45
4.2
Approches sp\'{e}cifiques de migration des UI
. . . . . . . . . . . . . . . . . . . .
46
4.2.1
Migration manuelle . . . . . . . . . . . . . . . . . . . . . . . . . . . . . .
46
4.2.2
Portage des applications existantes sur des tables interactives . . . . . . . .
48
4.3
Approches de migration bas\'{e}es sur les mod\`{e}les de l’UI . . . . . . . . . . . . . .
51
4.3.1
Approches de migration automatiques des UI . . . . . . . . . . . . . . . .
52
4.3.2
Approche semi automatique de migration des UI . . . . . . . . . . . . . .
57
4.4
Synth\`{e}se et objectifs . . . . . . . . . . . . . . . . . . . . . . . . . . . . . . . . .
62
4.4.1
Synth\`{e}se
. . . . . . . . . . . . . . . . . . . . . . . . . . . . . . . . . . .
62
4.4.2
Objectifs
. . . . . . . . . . . . . . . . . . . . . . . . . . . . . . . . . . .
64
4.1
Introduction
La migration des applications est une activit\'{e} de d\'{e}placement et d’adaptation d’un logiciel d’un
environnement source vers un environnement cible. Elle est plus globale que le portage d’applications [Moo95] car elle ne se limite pas qu’au changement de langages de programmation ou au changement des syst\`{e}mes d’exploitation. La migration englobe les probl\'{e}matiques de r\'{e} engineering, de
reverse engineering, de forward engineering et de portage d’applications. McClure et al. [McC92]
d\'{e}finissent le r\'{e} engineering comme une am\'{e}lioration d’un syst\`{e}me existant pour le rendre conforme
aux standards. Le reverse engineering consiste \`{a} analyser un syst\`{e}me existant pour d\'{e}crire la repr\'{e}sentation d’origine de mani\`{e}re plus abstraite [McC92]. Le forward engineering est une concr\'{e}tisation
de la repr\'{e}sentation abstraite d’un syst\`{e}me dans une impl\'{e}mentation. Une approche de migration des
UI vers une table interactive implique un r\'{e} engineering de l’UI source en impliquant d’abord une
phase de reverse engineering de l’UI source suivit d’une phase de forward engineering.
Dans notre contexte, nous avons fait le choix d’imposer aux applications \`{a} migrer une d\'{e}composition fonctionnelle minimale qui permet de distinguer une UI et un NF. Cette d\'{e}composition facilite
la migration de l’application en permettant la r\'{e}utilisation du NF, si les plateformes sont compatibles.
Nos travaux se concentrent sur la migration de l’UI pour laquelle les diff\'{e}rentes approches possible
se basent sur un ensemble de m\'{e}canismes qui visent dans un premier temps \`{a} construire des \'{e}quivalences entre les instruments d’interactions des plateformes source et cible. Il est ensuite possible
d’adapter la structure, le positionnement et le style de l’UI source par rapport \`{a} celle cible. Dans un
dernier temps, il s’agit de respecter les guidelines de la cible. Cet enchainement d’actions permettant
la migration des UI peut \^{e}tre manuel, automatis\'{e} ou semi automatique.
Nous \'{e}tudions dans ce chapitre diff\'{e}rentes approches de migration des UI dans l’objectif d’\'{e}valuer :
– La flexibilit\'{e} et l’automaticit\'{e} de chaque approche afin d’\'{e}valuer le degr\'{e} d’intervention du
concepteur pendant le processus de migration.
45
46
CHAPITRE 4. APPROCHES DE MIGRATION DES UI
– La prise en compte des guidelines pour assurer le respect des crit\`{e}res ergonomiques de conception des UI.
– La r\'{e}utilisabilit\'{e} des m\'{e}canismes d’\'{e}quivalences et d’adaptation dans le but de r\'{e}duire les
co\^{u}ts de la migration.
La section 4.2 pr\'{e}sente par cons\'{e}quent les approches de migration sp\'{e}cifiques \`{a} des applications
ou \`{a} des bo\^{i}tes \`{a} outils. La section 4.3 pr\'{e}sente les approches de migration des UI bas\'{e}es sur des mod\`{e}les d’UI. La section 4.3 fait une synth\`{e}se des diff\'{e}rentes approches de migration des UI et pr\'{e}sente
les caract\'{e}ristiques de la solution que nous avons adopt\'{e}e.
4.2
Approches sp\'{e}cifiques de migration des UI
Dans cette section nous \'{e}tudions les processus de migrations sp\'{e}cifiques \`{a} une application ou \`{a}
une bo\^{i}te \`{a} outils. Nous pr\'{e}sentons d’abord un processus qui permet de migrer manuellement une
application vers une table interactive. Ce processus identifie des guidelines pour les UI collaboratives
et s’appuie sur une d\'{e}marche structur\'{e}e. Ensuite nous pr\'{e}sentons une famille de m\'{e}canismes de migration des UI des applications existantes sur les tables interactives sans re-conception de l’UI de
d\'{e}part.
4.2.1
Migration de l’application AgilePlanner vers une table interactive
AgilePlanner est une application de planification et de gestion de projet suivant la m\'{e}thode Agile ;
sur une table interactive, elle est utilis\'{e}e pour faire du brainstorming dans le cadre de la gestion de
plannings de projets. Nous \'{e}tudions cette approche dans le cadre de la migration de cette application
vers une table interactive [WGM08] afin d’identifier les diff\'{e}rents m\'{e}canismes manuels mis en œuvre.
Le processus mis en place repose sur quatre phases :
– La premi\`{e}re phase consiste \`{a} analyser l’UI de l’application \`{a} migrer. Elle permet d’identifier
les diff\'{e}rentes zones : menu, l\'{e}gendes, zones d’interaction, espace de travail, etc.
– La deuxi\`{e}me phase consiste \`{a} \'{e}valuer l’UI de l’application \`{a} migrer. Le but de cette \'{e}valuation
est d’identifier les diff\'{e}rences majeures entre la plateforme source et cible pour en d\'{e}duire des
recommandations qui seront des guidelines pour le concepteur. Il en est ressorti les 7 guidelines
suivantes :
– Avoir des composants de l’UI de l’application d\'{e}pla\c{c}ables et utilisables en 360°.
– Utiliser la reconnaissance gestuelle pour les interactions utilisateurs et \'{e}viter les menus traditionnels.
– Utiliser l’\'{e}criture \`{a} main lev\'{e}e au lieu du clavier pour la saisie des textes.
– Prendre en compte les interactions concurrentes pendant la conception de l’UI.
– Pouvoir utiliser des composants graphiques de taille assez grande pour faciliter les interactions tactiles.
– Éviter l’utilisation des bo\^{i}tes de dialogues pop-up.
– Permettre \`{a} l’UI de l’application de s’adapter aux diff\'{e}rentes tailles des tables interactives.
Ces guidelines sont conformes aux guidelines pour les UI collaboratives et aux guidelines pour
les interactions tactiles que nous avons pr\'{e}sent\'{e}es au chapitre pr\'{e}c\'{e}dent (cf. section 3.2.3).
– La troisi\`{e}me phase consiste \`{a} appliquer les guidelines de la phase pr\'{e}c\'{e}dente pour concevoir la
nouvelle UI de l’application AgilePlanner. Certaines des guidelines telles que la rotation et le
d\'{e}placement des composants graphiques, l’\'{e}criture \`{a} main lev\'{e}e, les reconnaissances gestuelle
et vocale peuvent \^{e}tre fournies par l’environnement logiciel de la table interactive. Évidemment,
cela d\'{e}pend de la table et des biblioth\`{e}ques graphiques disponibles.
4.2. APPROCHES SPÉCIFIQUES DE MIGRATION DES UI
47
– La derni\`{e}re phase du processus consiste \`{a} \'{e}valuer l’UI produite en demandant \`{a} des utilisateurs finaux d’\'{e}valuer l’utilisabilit\'{e} de chaque fonctionnalit\'{e} de l’application. Les diff\'{e}rentes
remarques \'{e}mises permettent de modifier l’UI migr\'{e}e. Une fois ces diff\'{e}rentes modifications
effectu\'{e}es, l’UI a \'{e}t\'{e} de nouveau \'{e}valu\'{e}e \`{a} nouveau pour mesurer la satisfaction des testeurs.
Par exemple, les r\'{e}sultats de cette r\'{e}\'{e}valuation ont montr\'{e} que l’utilisabilit\'{e} de l’\'{e}criture \`{a} main
lev\'{e}e sur la table interactive peut \^{e}tre am\'{e}lior\'{e}e.
La migration manuelle d’une UI est un processus qui comporte plusieurs phases. Dans cette section nous avons pr\'{e}sent\'{e} une approche qui nous permet d’identifier quatre phases importantes :
1. Analyser l’UI de d\'{e}part pour identifier la structure des \'{e}l\'{e}ments qui la compose et ses fonctionnalit\'{e}s.
2. Identifier les guidelines qui permettront la migration : cette identification se fait en se basant sur
les diff\'{e}rences entre la plateforme de d\'{e}part et celle d’arriv\'{e}e.
3. Appliquer les guidelines pour migrer l’UI de d\'{e}part.
4. Évaluer et am\'{e}liorer le r\'{e}sultat obtenu.
4.2.1.1
R\'{e}utilisation du processus pour d’autres applications
Les \'{e}tapes du processus de migration de l’application AgilePlanner vers une table interactive et
la mise en œuvre des guidelines identifi\'{e}es constituent des “savoir faire” r\'{e}utilisables pour d’autres
applications.
Cependant en ce qui concerne l’application AgilePlanner, les m\'{e}canismes d’\'{e}quivalences entre les
instruments d’interactions et l’adaptation du layout de la structure ne peuvent pas \^{e}tre r\'{e}utilis\'{e}s pour
d’autres applications car ils sont sp\'{e}cifiques \`{a} cette application. En effet, l’analyse et l’identification
des diff\'{e}rentes zones de l’UI ne s’appuient pas sur des mod\`{e}les ind\'{e}pendants de la plateforme.
La mise en œuvre des guidelines identifi\'{e}es par ce processus est bas\'{e}e sur l’intuition. Pour compl\'{e}ter l’\'{e}tude, nous avons consid\'{e}r\'{e} l’UI de deux autres applications que nous souhaiterions faire
migrer. L’UI de l’application CBA et celle d’une application agenda pr\'{e}sent\'{e}e par la figure 4.1. Nous
souhaitons effectuer cette migration en appliquant les guidelines identifi\'{e}es \`{a} la section 4.2.1. La migration de l’application agenda ne pose pas trop de probl\`{e}mes et par exemple, les divers composants
graphiques repr\'{e}sentant la barre de menu, la barre de recherche, la liste des cat\'{e}gories et le tableau
peuvent tous \^{e}tre d\'{e}pla\c{c}ables sur la table. On peut aussi consid\'{e}rer la liste des cat\'{e}gories et le tableau
de contacts comme tout indissociable. En ce qui concerne l’application CBA, le choix des groupes
utilisables \`{a} 360° peut se faire comme le propose la figure 2.3. Nous remarquons que l’utilisation
des guidelines pour deux applications diff\'{e}rentes repose sur l’interpr\'{e}tation faite par les personnes
en charge de la migration et laisse une certaine latitude car ces guidelines pr\'{e}cisent formellement
assez peu de choses. Par exemple, plusieurs possibilit\'{e} sont offertes pour regrouper les composants
graphiques \`{a} d\'{e}placer.
Le processus de migration d\'{e}crit par cette approche est tr\`{e}s flexible pour deux raisons, d’abord il
est essentiellement manuel et ensuite le processus int\`{e}gre une phase d’\'{e}valuation et d’am\'{e}lioration du
r\'{e}sultat. L’intervention humaine est omnipr\'{e}sente et permet d’am\'{e}liorer le r\'{e}sultat en continue.
Par contre, le fait que l’approche se base uniquement sur des m\'{e}canismes manuels rend plus lourd
le travail des personnes en charge de la migration. Les m\'{e}canismes d’\'{e}quivalence des instruments
d’interactions doivent \^{e}tre d\'{e}crit manuellement pour chaque application \`{a} migrer. Les m\'{e}canismes de
prise en compte des guidelines sont bas\'{e}s uniquement sur les connaissances des d\'{e}veloppeurs. Cette
approche n’offre que peu de m\'{e}canismes permettant de r\'{e}duire la charge de travail des d\'{e}veloppeurs
et par cons\'{e}quent le co\^{u}t de migration reste presque \'{e}quivalent \`{a} une nouvelle conception de l’UI.
48
CHAPITRE 4. APPROCHES DE MIGRATION DES UI
 
 
Barre de menu
Barre de recherche
Tableau de contacts
Liste des cat\'{e}gorie
 
FIGURE 4.1 – Application de consultation des contacts
4.2.1.2
Synth\`{e}se des approches manuelles
La figure 4.2 met en \'{e}vidence que la mise en \'{e}quivalence des composants graphiques ou des
dispositifs d’interactions des plateformes source et cible est manuelle tout comme les m\'{e}canismes de
prise en compte des guidelines. La qualit\'{e} de l’UI g\'{e}n\'{e}r\'{e}e est uniquement bas\'{e}e sur les connaissances
et les comp\'{e}tences des personnes en charge de la migration. Ces m\'{e}canismes sont donc difficilement
r\'{e}utilisables car non formalis\'{e}s. Par contre, cette approche est relativement flexible et permet d’avoir
une UI conforme aux crit\`{e}res ergonomiques pour peu que les concepteurs utilisent pleinement les
diff\'{e}rentes possibilit\'{e}s qui leur sont offertes et tirent partie des diff\'{e}rentes \'{e}valuations.
FIGURE 4.2 – Synth\`{e}se de la migration manuelle des UI
4.2.2
Portage des applications existantes sur des tables interactives
Le portage d’une UI vers sur une nouvelle plateforme consiste \`{a} adapter les instruments d’interactions de la plateforme cible pour l’UI source ainsi que le positionnement ou le style de l’UI source sans
4.2. APPROCHES SPÉCIFIQUES DE MIGRATION DES UI
49
adapter la structure. Dans cette section nous pr\'{e}sentons une approche de portage des UI propos\'{e}e par
Besacier [Bes10] qui a pour but de r\'{e}utiliser des applications d\'{e}velopp\'{e}es pour desktop sur une table
interactive sans reconception de l’UI. L’objectif premier d’adapter les diff\'{e}rents moyens d’interactions (clavier et souris) des desktops pour les rendre op\'{e}rants sur une table interactive DiamondTouch
selon diff\'{e}rentes technologies.
Pour Besacier [Bes10], six technologies permettent la r\'{e}utilisation souhait\'{e}e :
La capture d’\'{e}cran
est une technologie qui enregistre une copie exacte de l’\'{e}cran de l’UI de d\'{e}part
\`{a} intervalle r\'{e}gulier. La copie est stock\'{e}e dans une zone m\'{e}moire o\`{u} elle peut \^{e}tre modifi\'{e}e pour \^{e}tre
adapt\'{e}e \`{a} l’environnement de la table interactive. Cette approche se situe dans un contexte de portage
des applications existantes. Elle ne permet pas de porter des interactions en entr\'{e}e car l’UI port\'{e}e est
constitu\'{e}e d’images captur\'{e}es et ces images ne contiennent pas de m\'{e}ta donn\'{e}es.
La carte graphique virtuelle
est une am\'{e}lioration de la technologie pr\'{e}c\'{e}dente bas\'{e}e sur la capture
d’\'{e}cran. L’approche utilise le serveur Metisse [RCPsC05] qui stocke les fen\^{e}tres de l’application
de d\'{e}part sous forme d’images. Le serveur Metisse joue un r\^{o}le de carte graphique virtuelle car il
redessine les fen\^{e}tres sur l’\'{e}cran de la table interactive \`{a} travers un compositeur. Il redirige aussi les
\'{e}v\'{e}nements en provenance de la souris ou du clavier vers l’application. Cette technologie permet de
porter les fen\^{e}tres des applications desktops en offrant la possibilit\'{e} aux utilisateurs de les manipuler
facilement. Cependant les composants graphiques des fen\^{e}tres port\'{e}es ne sont toujours pas interactifs
car ils restent des images de l’UI de d\'{e}part.
L’\'{e}mulation du clavier et de la souris
est une technologie qui am\'{e}liore les deux technologies
pr\'{e}sent\'{e}es ci-dessus. Elle permet de porter aussi les interactions en entr\'{e}e. La technique permet de
porter plusieurs applications \`{a} la fois sur une table interactive et pour chaque application est associ\'{e}e
un utilisateur avec son clavier et sa souris virtuels [HCT06]. L’interaction d’un utilisateur verrouille la
table aux utilisateurs de la m\^{e}me application. Cette technologie permet d’associer \`{a} chaque utilisateur
son application mais ne permet pas de rendre collaborative une application. Cette technologie n’est
pas sp\'{e}cifique \`{a} une applications comme celles bas\'{e}es sur la capture d’\'{e}cran ; elle est r\'{e}utilisable pour
toutes les applications desktops (utilisant un clavier et une souris) que l’on souhaiterait migrer sur une
table interactive. Par contre elle reste sp\'{e}cifique aux dispositifs d’interactions (clavier et souris).
Le langage de script
est une technologie qui permet de porter les UI d’applications existantes vers
une table interactive en r\'{e}solvant les limites li\'{e}es aux multi-touches et aux multi-utilisateurs de la
technologie pr\'{e}c\'{e}dente. Le langage de script est sp\'{e}cifique \`{a} une application et consiste \`{a} \'{e}crire dans
un sc\'{e}nario un ensemble de scripts qui font chacun appel \`{a} une fonctionnalit\'{e} de l’application \`{a} porter.
A chaque \'{e}v\`{e}nement g\'{e}n\'{e}r\'{e} par la table interactive est associ\'{e} un script. Par exemple pour l’affichage
d’un fichier sur la table est r\'{e}alis\'{e} par l’association d’un script d’ouverture et de lecture de fichiers \`{a}
l’\'{e}v\'{e}nement d\'{e}clench\'{e} par la pose d’une feuille de papier, par exemple munie d’un tag, sur la table
interactive. Cette technique est difficilement r\'{e}utilisable d’une application \`{a} l’autre car chaque script
reste tr\`{e}s d\'{e}pendant de l’application pour laquelle il a \'{e}t\'{e} con\c{c}u. Cette technologie est utilis\'{e}e par
Besacier pour afficher les donn\'{e}es d’un tableur Excel sur une table interactive DiamondTouch tout
en prenant compte les guidelines li\'{e}es \`{a} la structure d’une UI \`{a} travers la m\'{e}taphore du papier. Cette
m\'{e}taphore stipule que les \'{e}l\'{e}ments graphiques d’une UI peuvent \^{e}tre utilis\'{e}s comme une feuille de
papier gr\^{a}ce aux interactions de d\'{e}placement, de rotation, de redimensionnement, de changement
d’\'{e}chelle et de pliage [BRNB07]. La technologie du langage de script n\'{e}cessite, dans le cas d’un
tableur Excel, l’impl\'{e}mentation d’un contr\^{o}leur et d’une vue pour afficher les donn\'{e}es d’un fichier
50
CHAPITRE 4. APPROCHES DE MIGRATION DES UI
Excel (ou du mod\`{e}le). La prise en compte des guidelines d\'{e}pend fortement des connaissances et de
l’intuition du concepteur en charge du portage. Par ailleurs le co\^{u}t de sa mise en œuvre reste \'{e}lev\'{e} par
rapport \`{a} sa r\'{e}utilisabilit\'{e}. En effet pour une UI \`{a} porter il faut d\'{e}crire des sc\'{e}narios d’utilisation afin
de construire les scripts.
Les API d’accessibilit\'{e} num\'{e}rique
permettent d’obtenir des informations s\'{e}mantiques \`{a} propos
des UI graphiques en cours d’ex\'{e}cution. Elles sont par exemple utilis\'{e}es pour offrir des interfaces
alternatives \`{a} des groupes d’utilisateurs sp\'{e}cifiques (d\'{e}ficients visuels). Ces informations constituent
des instances de mod\`{e}les de structure d’une UI donn\'{e}e comprenant les composants graphiques, les
donn\'{e}es, les positions, etc. Ce mod\`{e}le permet de d\'{e}crire une technique de migration qui segmente
les images obtenues par capture d’\'{e}cran de l’UI source. Les parties de l’UI de d\'{e}part peuvent donc
\^{e}tre dupliqu\'{e}es ou juxtapos\'{e}es en se basant sur le mod\`{e}le de structure et les images segment\'{e}es. Le
mod\`{e}le de structure et les images segment\'{e}es sont les \'{e}l\'{e}ments de l’UI source qui permettent de
g\'{e}n\'{e}rer l’UI cible. Compar\'{e}e aux technologies pr\'{e}c\'{e}dentes, cette technologie est moins r\'{e}utilisable
que la capture d’\'{e}cran, la carte graphique virtuelle et l’\'{e}mulation des dispositifs d’interactions. Elle
se base sur des technologies sp\'{e}cifiques \`{a} l’environnement Windows et toutes les applications ne
respectent pas forc\'{e}ment cette approche. Le co\^{u}t de sa mise en œuvre reste \'{e}lev\'{e} par rapport \`{a} sa
r\'{e}utilisabilit\'{e}. En effet, pour une UI \`{a} porter il faut \'{e}crire les algorithmes de segmentation de chaque
image captur\'{e}e et les m\'{e}canismes d’analyse des donn\'{e}es s\'{e}mantiques de l’UI.
La r\'{e}\'{e}criture d’une boite \`{a} outils d’interface homme machine
consiste \`{a} utiliser la bo\^{i}te \`{a} outils
de la table interactive \`{a} la place de celle du desktop. Le remplacement se fait en interceptant tous les
appels de fonction de l’application vers les objets de la bo\^{i}te \`{a} outils d’IHM de d\'{e}part et en les redirigeant vers les \'{e}l\'{e}ments de la bo\^{i}te \`{a} outils de la table interactive. On construit donc, \`{a} la mani\`{e}re d’une
machine virtuelle un \'{e}mulateur de la boite \`{a} outils source \`{a} l’aide de la boite \`{a} outils cible. L’avantage
de cette approche est d’utiliser des composants graphiques respectant les guidelines de la table interactive sans une nouvelle conception de l’application de d\'{e}part puisque celle-ci est inchang\'{e}e. Cette
solution peut \^{e}tre impl\'{e}ment\'{e}e en utilisant les wrappers [TKR10]. La mise place des m\'{e}canismes
d’\'{e}quivalences entre les instruments d’interactions de cette technologie n\'{e}cessite une bonne connaissance des biblioth\`{e}ques graphiques et des dispositifs d’interactions des plateformes source et cible.
L’utilisation de cette technologie n\'{e}cessite un temps d’apprentissage pour des personnes non expertes
des instruments d’interactions des plateformes source et cible. Cette approche ne permet pas le respect de l’ensemble des guidelines car la structure et le layout de l’UI par exemple ne sont pas modifi\'{e}s
pendant le portage. Elle permet de choisir les composants graphiques qui ont des interactions ou des
styles conformes \`{a} la table interactive.
4.2.2.1
Synth\`{e}se des approches de portage des UI
Le tableau 4.1 pr\'{e}sente une synth\`{e}se de ces diff\'{e}rentes approches en prenant en compte les crit\`{e}res d’\'{e}quivalences entre les \'{e}l\'{e}ments de la plateforme, les mod\`{e}les utilis\'{e}s et la prise en compte des
guidelines. Nous remarquons que les \'{e}quivalences entre les dispositifs d’interactions sont toujours
d\'{e}crites manuellement pendant la mise en œuvre de la solution. Ces approches utilisent une mod\'{e}lisation des interactions et seulement deux approches mod\'{e}lisent la structure de l’UI. La prise en compte
des guidelines d\'{e}pend de l’utilisation ou non de la bo\^{i}te \`{a} outils cible. En effet le langage de script, les
API d’accessibilit\'{e} et la r\'{e}\'{e}criture d’une bo\^{i}te \`{a} outils permettent d’utiliser les \'{e}l\'{e}ments de la bo\^{i}te \`{a}
outils cible. Cependant dans le cas de la r\'{e}\'{e}criture d’une bo\^{i}te \`{a} outils, la prise en compte est partielle
car la structure et le layout de l’UI de d\'{e}part ne sont pas modifi\'{e}s pendant le portage.
4.3. APPROCHES DE MIGRATION BASÉES SUR LES MODÈLES DE L’UI
51
Approches
Équivalences
Mod\'{e}lisations
Guidelines
Capture d’\'{e}cran
Aucune \'{e}quivalence
Aucune mod\'{e}lisation
Aucune prise en
compte
Carte graphique
virtuelle
Aucune \'{e}quivalence
Aucune mod\'{e}lisation
Aucune prise en
compte
Émulation du clavier
et de la souris
Équivalences entre les
dispositifs
d’interactions en entr\'{e}e
des plateformes source
et cible d\'{e}finies
manuellement
Aucune mod\'{e}lisation
Aucune prise en
compte des guidelines
Langage de script
Équivalences
manuelles des
dispositifs
d’interactions
Mod\'{e}lisation des
donn\'{e}es, approche non
r\'{e}utilisable
Prise en compte
partielle des guidelines
API d’accessibilit\'{e}
num\'{e}rique
Manuelles
Mod\'{e}lisation de la
structure
Prise en compte
partielle des guidelines
R\'{e}\'{e}criture d’une
boite \`{a} outils
Manuelles
Aucune
Prise en compte
partielle des guidelines
TABLE 4.1 – Synth\`{e}se des technologies de portage des UI sur tables interactives
Les m\'{e}canismes d’\'{e}quivalences des approches de portage des UI (cf. figure 4.3) ne sont pas
flexibles car les \'{e}quivalences entre les instruments d’interactions ou avec les wrappers en ce qui
concerne les biblioth\`{e}ques graphiques sont d\'{e}finis avant le portage de l’application. Aucune modification n’est possible pendant le processus de migration.
Le portage des UI ne prend pas en compte toutes les guidelines de la cible, les r\'{e}sultats obtenus
ne sont pas conformes aux crit\`{e}res ergonomiques des tables interactives. En effet, dans le cas de la
migration vers les tables interactives, le portage des UI pr\'{e}sente un inconv\'{e}nient car il ne modifie pas
la structure et le positionnement des \'{e}l\'{e}ments graphiques de l’UI source. Or la prise en compte des
guidelines pour les UI collaboratives impose l’adaptation de la structure et du positionnement de l’UI
source.
Contrairement \`{a} la migration manuelle, le portage des UI est une approche non sp\'{e}cifique \`{a} une
application, mais sp\'{e}cifique \`{a} des instruments d’interactions. Le portage des UI permet de migrer les
applications entre plateformes d\'{e}termin\'{e}es. Les m\'{e}canismes d’\'{e}quivalences d\'{e}finis pour une migration sont r\'{e}utilisables pour les autres applications mais pas pour d’autres instruments d’interactions.
Le portage des UI r\'{e}duit indiscutablement le co\^{u}t de la migration par rapport \`{a} une approche manuelle.
4.3
Approches de migration bas\'{e}es sur les mod\`{e}les de l’UI
Cette section pr\'{e}sente une famille d’approches de migration bas\'{e}es sur la d\'{e}finition de mod\`{e}les
permettant d’\^{e}tre ind\'{e}pendants des applications et des plateformes. En effet, m\^{e}me si les solutions de
portage des UI pr\'{e}sent\'{e}es \`{a} la section 4.2.2 utilisent des mod\`{e}les de donn\'{e}es et de structure de l’UI
de d\'{e}part pour g\'{e}n\'{e}rer l’UI cible. Ces mod\`{e}les comportent des informations s\'{e}mantiques qui incluent
pour chaque composant son nom, son r\^{o}le (bouton menu, case \`{a} cocher, etc.), sa position sur l’\'{e}cran
et dans la hi\'{e}rarchie en terme de contenant et de contenu. Ces informations sont exploit\'{e}es par des
m\'{e}canismes non r\'{e}utilisables et sp\'{e}cifiques aux applications car ces m\'{e}canismes sont constitu\'{e}s de
scripts sp\'{e}cifiques \`{a} une fonctionnalit\'{e} d’une UI. Ils ne sont dans pas facilement r\'{e}utilisables. Les
52
CHAPITRE 4. APPROCHES DE MIGRATION DES UI
FIGURE 4.3 – Synth\`{e}se des approches de portage des UI sur tables interactives
m\'{e}canismes de migration des UI dans notre contexte sont constitu\'{e}s des processus de transformation
des UI pour la plateforme cible.
Dans cette section nous \'{e}tudions les approches de migration qui se basent sur des mod\`{e}les qui
d\'{e}crivent la totalit\'{e} des caract\'{e}ristiques de l’UI. Nous d\'{e}crivons en premier les approches automatiques de migration bas\'{e}es sur des mod\`{e}les des UI, puis une approche semi automatique, bas\'{e}e sur
des mod\`{e}les de structure et d’interactions permettant la description des m\'{e}canismes de migration des
UI.
4.3.1
Approches de migration automatiques des UI
Dans cette section nous pr\'{e}sentons les approches de migration r\'{e}utilisables pour diff\'{e}rentes plateformes qui s’appuient sur des mod\`{e}les de l’UI et qui font appel \`{a} des m\'{e}canismes automatiques.
Ces solutions de migration des UI sont g\'{e}n\'{e}riques et sont utilis\'{e}es par des services de migration
des UI dans des contextes ubiquitaires comportant plusieurs types de plateformes. Les services de
migration des UI sont charg\'{e}s d’adapter une UI pendant son ex\'{e}cution. C’est une migration \`{a} la vol\'{e}e
ou entre deux sessions d’une UI d’une application o\`{u} l’\'{e}tat courant de l’UI est sauvegard\'{e} par des
mod\`{e}les abstraits de l’UI et adapt\'{e} \`{a} autre plateforme afin de permettre \`{a} un utilisateur de continuer
une t\^{a}che sans interruption [CCB+02]. Nous \'{e}tudions dans cette section les mod\`{e}les de l’UI et les
m\'{e}canismes adopt\'{e}s par ces approches.
4.3.1.1
Mod\`{e}les de l’UI
Les mod\`{e}les abstraits permettent de d\'{e}crire les UI comportent en g\'{e}n\'{e}ral plusieurs niveaux d’abstraction dans l’objectif de d\'{e}crire les diff\'{e}rents aspects des UI. Le framework de r\'{e}f\'{e}rence CAMELEON (CRF) [CCB+02] en propose quatre : le niveau des t\^{a}ches et concepts, le niveau interface
abstraite (AUI), le niveau interface concr\`{e}te (CUI) et le niveau interface finale (FUI).
Mod\`{e}le de t\^{a}ches et concepts
Ce mod\`{e}le exprime les t\^{a}ches et les concepts des UI dans un contexte
pr\'{e}cis tel que d\'{e}fini par le concepteur. Ce mod\`{e}le exprime les interactions de l’UI avec les utilisateurs
4.3. APPROCHES DE MIGRATION BASÉES SUR LES MODÈLES DE L’UI
53
ou le syst\`{e}me, le comportement et les diff\'{e}rents \'{e}tats des composants graphiques. Dans le cadre de
la migration des UI \`{a} la vol\'{e}e, le mod\`{e}le de t\^{a}ches permet de conserver l’\'{e}tat d’une UI en cours
d’ex\'{e}cution par exemple. Le mod\`{e}le de t\^{a}ches permet d’exprimer les interactions entre les utilisateurs
et l’UI ind\'{e}pendamment des modalit\'{e}s d’interactions.
Mod\`{e}le AUI
Ce mod\`{e}le permet de repr\'{e}senter les \'{e}l\'{e}ments d’une UI ind\'{e}pendamment des modalit\'{e}s d’interactions en exprimant leurs fonctionnalit\'{e}s essentielles. Ce mod\`{e}le concr\'{e}tise les \'{e}l\'{e}ments
du domaine utilis\'{e} par le mod\`{e}le de t\^{a}ches, par exemple une t\^{a}che de saisie de donn\'{e}es correspond
\`{a} un \'{e}l\'{e}ment de type Input dans le mod\`{e}le AUI (cf. figure 4.4). Les \'{e}l\'{e}ments du mod\`{e}le AUI permettent en g\'{e}n\'{e}ral d’exprimer les diff\'{e}rents types d’interactions tels que l’entr\'{e}e de donn\'{e}es (Input),
l’affichage d’informations (Output) et l’activation d’une commande (Command). Ce mod\`{e}le exprime
aussi la structure d’une UI \`{a} travers les regroupements (Container) ou les types des donn\'{e}es.
Le mod\`{e}le d’AUI exprime les interactions de haut niveau entre UI et NF et ind\'{e}pendamment des
dispositifs d’interactions. Ce mod\`{e}le n’exprime pas les interactions sur les propri\'{e}t\'{e}s visuelles des
composants graphiques (redimensionnement, d\'{e}placement, rotation, etc.) car il est ind\'{e}pendant de
toute modalit\'{e}, il peut \^{e}tre concr\'{e}tis\'{e} en UI graphique, vocale ou multimodale, etc.
Mod\`{e}le CUI
Ce mod\`{e}le permet de d\'{e}crire une repr\'{e}sentation concr\`{e}te de l’UI suivant une modalit\'{e}.
Par exemple, dans le cadre des UI graphiques qui nous concernent, des composants graphiques sont
choisis afin de raffiner les \'{e}l\'{e}ments du mod\`{e}le AUI. Ce mod\`{e}le exprime la structure, le positionnement
(layout) et m\^{e}me le style des \'{e}l\'{e}ments graphiques. Les composants graphiques de ce mod\`{e}le sont
exprim\'{e}s dans un langage ind\'{e}pendant des biblioth\`{e}ques graphiques pour garantir leur r\'{e}utilisabilit\'{e}.
Dans le cas d’une migration vers une table interactive, le mod\`{e}le de CUI exprime des UI de
modalit\'{e} graphique, interactive et tangible. On peut consid\'{e}rer la migration des UI desktop vers une
table interactive comme un changement de biblioth\`{e}que graphique qui n\'{e}cessite une adaptation des
interactions, de la structure, du layout et du style de l’UI de d\'{e}part.
Plusieurs impl\'{e}mentations du framework CAMELEON ont \'{e}t\'{e} mises en œuvre. Le langage XML
MARIA (cf. Paternò et al. [PSS09]) et le langage USIXML (User Interface eXtensible Markup Language) [VLM+04] d\'{e}crivent les UI aux niveaux des mod\`{e}les AUI et CUI. Le mod\`{e}le de t\^{a}ches et
concepts du CRF est impl\'{e}ment\'{e} par CTT (Concur Task Trees) [FCLD12], Il permet d’exprimer de
mani\`{e}re structur\'{e}e et hi\'{e}rarchis\'{e}e les t\^{a}ches d’une UI. La structure hi\'{e}rarchique de CTT permet aussi
d’exprimer les relations entre les sous t\^{a}ches gr\^{a}ce \`{a} des op\'{e}rateurs temporels.
Les diff\'{e}rents mod\`{e}les du CRF permettent de concevoir des UI multi plateformes, dans une approche de conception top down qui consiste \`{a} partir du mod\`{e}le le plus abstrait (mod\`{e}le de t\^{a}ches et
concepts), le concepteur g\'{e}n\`{e}re les UI par raffinements successifs des mod\`{e}les du CRF. Dans le cadre
de la migration des UI, ces mod\`{e}les sont utilis\'{e}s soit par des m\'{e}canismes de g\'{e}n\'{e}ration des UI finales
pour la nouvelle plateforme \`{a} partir des mod\`{e}les abstraits [TS99], soit par reverse engineering de l’UI
existante [PZ10].
4.3.1.2
M\'{e}canismes de migration des UI top down
Les m\'{e}canismes de migration des UI top down sont bas\'{e}s sur le raffinement des mod\`{e}les abstraits du CRF. Les concepteurs g\'{e}n\`{e}rent dans un premier temps le mod\`{e}le CUI sp\'{e}cifique \`{a} chaque
plateforme \`{a} partir des mod\`{e}les AUI et des t\^{a}ches. Cette approche permet de migrer les \'{e}tats des UI
\`{a} l’ex\'{e}cution, MigriXML [MVL06] par exemple d\'{e}finit gr\^{a}ce \`{a} ces m\'{e}canismes un environnement
virtuel comprenant diff\'{e}rentes plateformes (desktops avec diff\'{e}rentes tailles d’\'{e}cran ou smartphones).
MigriXML permet par exemple la migration d’une UI \`{a} l’ex\'{e}cution d’un desktop vers un smartphone
54
CHAPITRE 4. APPROCHES DE MIGRATION DES UI
 
 
 
FIGURE 4.4 – Exemple de transformation USIXML
si certaines conditions sont remplies. En effet l’UI de l’application doit \^{e}tre d\'{e}crit \`{a} l’aide d’USIXML
et les mod\`{e}les de CUI pour les desktops et les smartphones aussi doivent \^{e}tre d\'{e}crits avant la migration de l’UI. MigriXML permet alors de migrer un contexte d’ex\'{e}cution d’une UI du desktop vers un
smartphone et vice versa.
La g\'{e}n\'{e}ration des mod\`{e}les les moins abstraits dans cette approche doit \^{e}tre effectu\'{e}e d\`{e}s la phase
de conception. L’approche permet d’\'{e}tendre une application \`{a} d’autres plateformes en r\'{e}utilisant les
mod\`{e}les de t\^{a}ches et d’AUI dans le cas d’une nouvelle plateforme ayant des modalit\'{e}s diff\'{e}rentes de la
plateforme de d\'{e}part. Pour la migration d’une application existante, il est indispensable de construire
les diff\'{e}rents mod\`{e}les du CRF associ\'{e}s. Cette contrainte exclue toutes les UI des applications qui ne
sont pas con\c{c}ues suivant le CRF sauf \`{a} les construire par reverse engineering.
4.3.1.3
M\'{e}canismes de migration des UI bottom up
Ces m\'{e}canismes permettent la migration des UI par abstraction des UI existantes pour obtenir les
mod\`{e}les du CRF. Paternò et al. [PZ10] proposent par exemple, un service d’adaptation d’une page
web \`{a} un t\'{e}l\'{e}phone portable (cf. figure 4.5) bas\'{e} sur le reverse engineering et l’adaptation d’une UI
existante. Ce service abstrait les pages HTML dans les mod\`{e}les CUI et AUI du langage MARIA XML
et transforme les mod\`{e}les obtenus en fonction des principes suivants :
– adapter les tableaux pour qu’ils tiennent sur la surface d’\'{e}cran disponible en d\'{e}coupant les
tableaux trop larges ou en rempla\c{c}ant une donn\'{e}e d’un cellule trop large par un lien hypertexte,
– transformer tous les textes trop longs en liens hypertexte comme pour les donn\'{e}es des cellules
d’un tableau,
– transformer les images en r\'{e}duisant leur taille,
– convertir les listes, par exemple en menu drop down,
– adapter la taille et la disposition des composants graphiques, par exemple en les redimensionnant et en adoptant un layout vertical.
Le concepteur d\'{e}finit des m\'{e}canismes de transformations [LVMB05] pour appliquer ces principes
sur les mod\`{e}les de CUI et d’AUI. La migration ne change pas de biblioth\`{e}que graphique car l’UI reste
toujours \'{e}crite en HTML. Cependant le layout, la taille des composants graphiques, les donn\'{e}es et le
type de certains composants graphiques doivent \^{e}tre transform\'{e}s et remplac\'{e}s pendant la phase de migration. Les transformations des instances des mod\`{e}les AUI et CUI de l’UI de d\'{e}part sont horizontales
4.3. APPROCHES DE MIGRATION BASÉES SUR LES MODÈLES DE L’UI
55
car elle g\'{e}n\`{e}rent d’autres instances de ces mod\`{e}les pour la plateforme d’arriv\'{e}e.
L’avantage de cette approche est que seuls les mod\`{e}les indispensables \`{a} la migration sont abstraits
\`{a} partir de l’UI de d\'{e}part. Par exemple, si l’adaptation concerne le layout, le mod\`{e}le de t\^{a}che n’est pas
indispensable. Dans le cadre de la migration des UI vers les tables interactives qui implique un changement de dispositifs d’interactions, il est indispensable d’exprimer les interactions de mani\`{e}re abstraite
dans les diff\'{e}rents les mod\`{e}les T\^{a}ches et AUI du CRF. Cependant, comme les mod\`{e}les de t\^{a}ches et
AUI du CRF n’expriment pas toutes les interactions n\'{e}cessaire \`{a} la migration (cf. section 4.3.1.1), il
est donc n\'{e}cessaire que les applications respectent la s\'{e}paration entre UI et NF pour que toutes les
interactions puissent \^{e}tre d\'{e}couvertes. MVC ou ARCH sont des architectures respectant la s\'{e}paration
souhait\'{e}e.
 
 
FIGURE 4.5 – Service de migration des UI
4.3.1.4
Migration des UI desktop vers les tables interactives par une approche automatique
Dans cette section nous \'{e}tudions les adaptations n\'{e}cessaires dans le cas d’une r\'{e}utilisation du
service de migration des UI (propos\'{e}e par Paternò \`{a} la figure 4.5) pour la migration des UI vers les
tables interactives. Nous souhaitons montrer par cette \'{e}tudes les points forts et les difficult\'{e}s li\'{e}s \`{a}
l’adaptation de cette solution pour notre objectif. En consid\'{e}rant toujours notre exemple fil rouge :
l’application desktop de d\'{e}part (cf. section 2.2) respecte une architecture MVC et que la vue d\'{e}crite
\`{a} l’aide de la biblioth\`{e}que graphique Java Swing. La table interactive est la plateforme d’arriv\'{e}e et
elle permet de d\'{e}crire des UI en utilisant XAML. Une adaptation du service de migration des UI
n\'{e}cessite :
– Une table d’\'{e}quivalences (cf. tableau 4.2) entre les \'{e}l\'{e}ments du langage MARIA XML, XAML
et l’API JavaSwing pour abstraire et concr\'{e}tiser l’UI de d\'{e}part. La table d’\'{e}quivalences permet
de d\'{e}crire les correspondances entre les biblioth\`{e}ques graphiques des plateformes sources et
cibles. Ces \'{e}quivalences sont sym\'{e}triques car elles peuvent aussi \^{e}tre utilis\'{e}es pour la migration
des tables interactives vers les desktops.
– Une formalisation des guidelines des tables interactives en r\`{e}gles de migration pour d\'{e}crire des
UI qui respectent les crit\`{e}res ergonomiques de la cible. Nous remarquons que les guidelines formalis\'{e}es en r\`{e}gles ne sont pas sym\'{e}triques comme les tables d’\'{e}quivalences car la formalisation
des guidelines est sp\'{e}cifique \`{a} la plateforme cible.
56
CHAPITRE 4. APPROCHES DE MIGRATION DES UI
MARIA XML
Java Swing
XAML Surface
Activator
JButton
Button
Single choice
JCombox
ListBox
Text Edit
JTextField
TextBox
Object
Image
Image
Grouping
JPanel
ScatterViewItem, Grid
TABLE 4.2 – Table d’\'{e}quivalences
Les correspondances sont statiques et doivent \^{e}tre \'{e}tablies pour l’ensemble des composants graphiques d’une biblioth\`{e}que graphique. Évidemment, si une application n’utilise pas une partie de la
biblioth\`{e}que source, il est inutile de rechercher les correspondances. La construction de cette table est
donc \'{e}volutive et permet de capitaliser sur les migrations d\'{e}j\`{a} effectu\'{e}es. Les \'{e}quivalences statiques
sont \'{e}tablies en se basant sur les types des composants graphiques. Cependant le processus de migration se basant sur les instances des composants graphiques, chaque instance peut impl\'{e}menter ou
non les interactions d\'{e}crites par le type. Par exemple une liste d’\'{e}l\'{e}ments peut impl\'{e}menter ou non
le glisser d\'{e}poser, une instance qui ne l’impl\'{e}mente pas peut dans certains cas \^{e}tre transform\'{e}e en un
menu.
Par ailleurs, Silva et al. [SC12] proposent plusieurs crit\`{e}res pour la correspondance entre biblioth\`{e}ques graphiques bas\'{e}es sur les langages XML. Le premier crit\`{e}re est le comportement des composants graphiques car il caract\'{e}rise les actions utilisateurs ind\'{e}pendamment de la repr\'{e}sentation du
composant graphique. Les autres crit\`{e}res utilisables dans le cadre de la migration sont le style des UI
et les balises des \'{e}l\'{e}ments graphiques. Par exemple, en consid\'{e}rant la guideline de partage de l’espace de travail (Guideline 2 Partage de l’espace de travail), elle peut \^{e}tre traduite dans les termes
du langage MARIA XML, pour obtenir par exemple, \`{a} la r\`{e}gle suivante : tous les Groupings sont
transform\'{e}s en ScatterViewItem pour \^{e}tre conformes \`{a} la guideline d’utilisation 360° de l’UI.
4.3.1.5
Synth\`{e}se des approches automatiques de migration des UI
Cette section pr\'{e}sente des m\'{e}canismes de migration des UI bas\'{e}s sur les mod\`{e}les de CRF. Le
mod\`{e}le de CUI graphique permet de d\'{e}crire la structure hi\'{e}rarchique, les donn\'{e}es et le positionnement
d’une UI. Dans un m\'{e}canisme de migration des UI vers la table interactive, le mod\`{e}le de CUI est
transform\'{e} pour rendre l’UI de d\'{e}part conforme aux guidelines des UI cibles qui sont tangibles et
collaboratives.
Le mod\`{e}le d’AUI d\'{e}crit \`{a} la fois le regroupement des \'{e}l\'{e}ments abstraits de l’UI et les interactions
(en entr\'{e}e ou en sortie) entre l’utilisateur et le syst\`{e}me (NF). Les mod\`{e}les de t\^{a}ches et de concepts
d\'{e}crivent les activit\'{e}s de l’UI et les comportements de l’UI de fa\c{c}on globale dans un langage de haut
niveau.
Les interactions d\'{e}crites par les mod\`{e}les de t\^{a}ches et AUI expriment les interactions sur une UI
ind\'{e}pendamment des dispositifs d’interactions, mais elles n’expriment pas l’ensemble des comportements des composants graphiques n\'{e}cessaires pour \'{e}tablir des \'{e}quivalences.
Les approches automatiques de migrations sont plus r\'{e}utilisables que le portage des UI car elles
s’appuient sur des mod\`{e}les qui d\'{e}crivent tous les aspects d’une UI. Sur la figure 4.6, nous \'{e}valuons
la r\'{e}utilisabilit\'{e} par la valeur maximale car les m\'{e}canismes d’\'{e}quivalences et d’adaptations d\'{e}crits
sur les mod\`{e}les de l’UI sont ind\'{e}pendantes des plateformes et des applications. Cependant les m\'{e}canismes d’\'{e}quivalences des instruments d’interactions et les m\'{e}canismes d’adaptation (du layout
4.3. APPROCHES DE MIGRATION BASÉES SUR LES MODÈLES DE L’UI
57
par exemple [PZ10]) sont sp\'{e}cifiques aux plateformes source et cible, leur r\'{e}utilisation pour d’autres
plateformes implique une charge de travail pour la mise en place de nouveaux m\'{e}canismes.
Les services de migration des UI se basent sur un processus automatique, ce qui r\'{e}duit leur flexibilit\'{e}. En effet les concepteurs n’interviennent pas pendant le processus pour l’ajuster ou le modifier
par exemple. Sur la figure 4.6, nous \'{e}valuons la flexibilit\'{e} de cette approche par une valeur non nulle
mais inf\'{e}rieure \`{a} celle de l’approche manuelle. Le peu de flexibilit\'{e} des m\'{e}canismes d’adaptation des
services de migration des UI implique une charge de travail suppl\'{e}mentaire pendant la migration. En
effet si la transformation de la structure ou du layout d’une UI n’est pas conforme aux attentes des utilisateurs par exemple, le concepteur doit modifier les r\`{e}gles de transformation des mod\`{e}les abstraits.
Cette modification se fait avant une migration des UI car le processus est automatique.
Les UI migr\'{e}es par les services de migration des UI respectent les crit\`{e}res ergonomiques de la
plateforme cible plus que celles obtenues par portage des UI par exemple. Cependant le peu de flexibilit\'{e} de cette approche par rapport aux approches manuelles fait que la prise en compte des guidelines
d\'{e}pend des r\`{e}gles de transformation des mod\`{e}les d\'{e}finis avant la migration. Sur la figure 4.6, nous
\'{e}valuons le respect de crit\`{e}res ergonomiques par une valeur inf\'{e}rieure \`{a} la valeur maximale car toutes
les guidelines ne sont pas toujours formalis\'{e}es.
FIGURE 4.6 – Synth\`{e}se des approches automatiques de migration des UI
4.3.2
Approche semi automatique de migration des UI
Les processus de migration automatiques des UI sont limit\'{e}s en terme de flexibilit\'{e} pour les personnes en charge de la migration. Dans cette section nous \'{e}tudions un processus de migration semi
automatique. MORPH [MR97] est par exemple une solution de migration d’une UI textuelle vers une
UI graphique en se basant sur des mod\`{e}les abstraits de l’UI et un support pour les transformations vers
de nouvelles impl\'{e}mentations graphiques. Le processus de migration avec MORPH implique une reconception de l’UI textuelle en UI graphique car les UI graphiques de type WIMP [vD97] supportent
les dispositifs d’interactions de manipulations directes comme une souris. Cette nouvelle conception
est aussi un changement de modalit\'{e}s d’interactions et l’utilisation d’une bo\^{i}te \`{a} outils graphique.
Nous \'{e}tudions cette approche de migration car comme la migration des UI vers des tables interactives,
58
CHAPITRE 4. APPROCHES DE MIGRATION DES UI
les modalit\'{e}s d’interactions des plateformes d’arriv\'{e}e sont diff\'{e}rentes des plateformes de d\'{e}part avec
de nouveaux dispositifs d’interactions.
4.3.2.1
Processus de migration semi automatique
Il est repr\'{e}sent\'{e} par l’ensemble des m\'{e}canismes qui permettent l’extraction du mod\`{e}le de l’UI, sa
transformation et enfin sa g\'{e}n\'{e}ration pour la plateforme cible. Les mod\`{e}les abstraits sont utilis\'{e}s pour
la repr\'{e}sentation de l’UI. Par exemple, MORPH d\'{e}crit un processus de migration en trois \'{e}tapes : la
d\'{e}tection, la repr\'{e}sentation et la transformation (cf. figure 4.7).
1. La d\'{e}tection est une activit\'{e} de reverse engineering qui consiste \`{a} analyser le code source de
l’application \`{a} migrer dans le but d’identifier les mod\`{e}les de structure et d’interactions.
2. La g\'{e}n\'{e}ration est l’op\'{e}ration inverse de la d\'{e}tection qui consiste \`{a} produire le code source de
l’application migr\'{e}e \`{a} partir de mod\`{e}les abstraits transform\'{e}s.
3. La repr\'{e}sentation de l’UI de d\'{e}part est un ensemble de mod\`{e}les abstraits issu de la phase de
d\'{e}tection. La transformation consiste \`{a} modifier et enrichir les aspects visuels de l’UI \`{a} travers
les mod\`{e}les, \`{a} ajouter et compl\'{e}ter les composants graphiques ou les fonctionnalit\'{e}s de l’UI de
d\'{e}part, \`{a} modifier et restructurer les types de donn\'{e}es ou la structure hi\'{e}rarchique de l’UI des
mod\`{e}les abstraits de l’UI source pour \^{e}tre utilisables dans l’environnement cible.
 
 
FIGURE 4.7 – Processus de migration avec MORPH
4.3.2.2
Les mod\`{e}les abstraits des UI
Ce sont des \'{e}l\'{e}ments cl\'{e}s du processus MORPH car ils repr\'{e}sentent les diff\'{e}rents aspects (layout,
activit\'{e}s, etc.) de l’UI \`{a} migrer. Les mod\`{e}les sont d\'{e}crits suivant deux niveaux d’abstractions :
– Le niveau des t\^{a}ches d’interactions qui regroupe quatre interactions de base d’une UI
[FvDFH90] qui sont :
– la s\'{e}lection dans une liste,
– la quantification ou la saisie d’une donn\'{e}e num\'{e}rique,
– l’indication de la position 10 d’un \'{e}l\'{e}ment sur l’\'{e}cran
10. Sur une UI textuelle la position est exprim\'{e}e en terme de ligne et de colonne pour positionner le curseur du clavier
par exemple
4.3. APPROCHES DE MIGRATION BASÉES SUR LES MODÈLES DE L’UI
59
– et la saisie d’une donn\'{e}e textuelle.
Ces t\^{a}ches d’interactions sont identifi\'{e}es pour la migration des UI textuelle vers une UI graphique. Ce niveau d\'{e}crit les activit\'{e}s utilisateurs sur une UI textuelle dans le but de l’adapter \`{a}
une UI graphique. Il correspond \`{a} une partie du mod\`{e}le des t\^{a}ches du CRF (cf. section 4.3.1.1).
– Le niveau d’objet d’interactions abstraites qui repr\'{e}sente les \'{e}l\'{e}ments abstraits ind\'{e}pendants
d’une biblioth\`{e}que graphique (tel que Button, List, Menu, etc.), ce mod\`{e}le est sp\'{e}cifique aux UI
graphiques. Ce niveau permet de d\'{e}crire la structure de l’UI ind\'{e}pendamment de la biblioth\`{e}que
graphique, il correspond au niveau CUI du CRF (cf. section 4.3.1.1).
Dans le cadre de la migration, les t\^{a}ches d’interactions permettent de d\'{e}crire les interactions des
objets d’interactions abstraits ind\'{e}pendamment du dispositif d’interactions de d\'{e}part (clavier).
Les t\^{a}ches d’interactions constituent un mod\`{e}le d’interactions abstraites pour la migration d’une
UI textuelle vers une UI graphique.
Les objets d’interactions abstraits sont raffin\'{e}s \`{a} partir des attributs des diff\'{e}rentes t\^{a}ches d’interactions et en se basant sur le langage de repr\'{e}sentation des connaissances CLASSIC [RBB+95].
Extraction des mod\`{e}les abstraits
Les t\^{a}ches d’interactions et les objets d’interactions abstraites
sont identifi\'{e}s \`{a} partir des UI textuelles en se basant sur l’architecture de l’application \`{a} migrer et
sur des r\`{e}gles d’identifications [Moo96, RFJ08]. Pour identifier les t\^{a}ches d’interactions, Moore et
al. [Moo96] d\'{e}crivent les r\`{e}gles d’identifications sous la forme :
Si Condition est vraie; Alors identi fier une Tˆache d′interactions
o\`{u} Condition correspond \`{a} une situation d’une instance des UI qui repr\'{e}sente une T\^{a}che d’interactions du mod\`{e}le abstrait. Cette forme de r\`{e}gle n\'{e}cessite que l’ensemble des situations soit identifi\'{e}
par un mapping entre la bo\^{i}te \`{a} outils de l’UI source et les interactions abstraites.
Il existe d’autres approches pour extraire un mod\`{e}le \`{a} partir d’une UI existante, par exemple
Ratiu et al. [RFJ08] proposent une ontologie pour analyser et comprendre des API sp\'{e}cifiques \`{a} un
domaine comme les biblioth\`{e}ques graphiques. Le principe de l’approche consiste d’abord \`{a} construire
une ontologie capable de repr\'{e}senter les concepts d’une biblioth\`{e}que graphique que nous souhaitons
utiliser pour la migration ; dans notre cas ce sont la structure et les interactions abstraites d’une UI.
Toujours dans notre cas l’ontologie d\'{e}crirait les liens de contenance entre les composants graphiques,
les donn\'{e}es et les t\^{a}ches d’interactions en entr\'{e}e (s\'{e}lection, quantification, position, \'{e}dition).
Ensuite, \`{a} partir d’une UI d\'{e}crite \`{a} l’aide des instances des \'{e}l\'{e}ments d’une biblioth\`{e}que graphique, nous pouvons extraire les objets d’interactions abstraites gr\^{a}ce \`{a} l’algorithme d’extraction
d\'{e}crit dans [RFJ08] se basant sur l’ontologie propos\'{e}e. Les objets d’interactions abstraites de l’UI
source sont utilis\'{e}s comme pivot pour d\'{e}crire l’UI cible.
Cette approche d’extraction pr\'{e}sente un avantage par rapport aux r\`{e}gles d’identification de
[Moo96] car elle ne s’appuie pas sur des situations d\'{e}crivant des cas possibles dans une UI mais
elle se base sur les types des \'{e}l\'{e}ments utilis\'{e}s pour d\'{e}crire l’UI et elle est plus g\'{e}n\'{e}rique. Cependant
le choix de l’ontologie doit \^{e}tre exhaustive pour une API et pour les concepts pr\'{e}sents.
M\'{e}canismes de transformations et de g\'{e}n\'{e}ration de l’UI
Une fois l’UI source d\'{e}crite par les
t\^{a}ches d’interactions et ensuite raffin\'{e}e en objets d’interactions abstraites, la repr\'{e}sentation abstraite
de l’UI de d\'{e}part est transform\'{e}e en modifiant manuellement le mod\`{e}le de l’UI.
La transformation consiste d’une part \`{a} remplacer des objets d’interactions de l’UI source. Par
exemple, une t\^{a}che de s\'{e}lections dans une liste qui est raffin\'{e}e en une liste \`{a} choix unique peut \^{e}tre
remplac\'{e}e par un menu si le nombre d’\'{e}l\'{e}ments est inf\'{e}rieur \`{a} N (N \'{e}tant une constante \`{a} d\'{e}finir selon
la longueur de menu souhait\'{e}).
60
CHAPITRE 4. APPROCHES DE MIGRATION DES UI
D’autre part cette phase fait intervenir un utilisateur humain pour d\'{e}finir la position, la taille ou
pour modifier l’UI. Le concepteur intervient manuellement sur les mod\`{e}les de l’UI apr\`{e}s la transformation dans le but de placer les \'{e}l\'{e}ments de l’UI. En effet, l’UI textuelle de d\'{e}part n’est pas structur\'{e}e
comme une UI graphique et les mod\`{e}les de t\^{a}ches d’interactions et d’objets abstraits d’interactions ne
permettent pas, par exemple, de d\'{e}crire le layout.
La s\'{e}lection des \'{e}l\'{e}ments de la biblioth\`{e}que graphique cible se fait en recherchant dans une biblioth\`{e}que graphique d\'{e}crite par une ontologie des \'{e}l\'{e}ments qui correspondent le plus \`{a} un objet abstrait
du mod\`{e}le restructur\'{e}.
Le mapping entre les objets abstraits et les \'{e}l\'{e}ments d’une biblioth\`{e}que graphique est fait dynamiquement en se basant sur les ontologies. En effet les correspondances entre les objets abstraits et
les \'{e}l\'{e}ments de la biblioth\`{e}que graphique ne sont pas d\'{e}finies manuellement et de mani\`{e}re exhaustive
\`{a} la conception de la solution, mais elles sont \'{e}tablies en se basant sur les attributs de chaque \'{e}l\'{e}ment.
Le m\'{e}canisme d’\'{e}quivalences entre les biblioth\`{e}ques graphiques se base sur les propri\'{e}t\'{e}s de chaque
composant graphique d\'{e}crit par une ontologie. Cependant ce m\'{e}canisme ne prend pas en compte les
propri\'{e}t\'{e}s de l’instance des composants mais celles d\'{e}crites par le type.
Les m\'{e}canismes de transformations et g\'{e}n\'{e}rations utilis\'{e}s par MORPH sont bas\'{e}s aussi sur des
ontologies. Dans le cadre de la transformation de l’UI de d\'{e}part par exemple, il est possible d’introduire des guidelines dans la base de connaissances, en pr\'{e}f\'{e}rant remplacer une liste en menu si elle
contient moins de 10 \'{e}l\'{e}ments par exemple.
La transformation utilis\'{e}e par MORPH permet d’inclure les principes de conception des UI pour
la plateforme cible dans la base de connaissances [MR97].
4.3.2.3
Comment adapter MORPH pour la migration des UI vers les tables interactives ?
Dans l’objectif d’\'{e}tudier les points forts et les limites d’une approche de migration semi automatique comme MORPH, nous l’adaptons dans notre cadre.
La r\'{e}ponse \`{a} la question pos\'{e}e passe par l’ajout de la biblioth\`{e}que graphique de la table interactive
\`{a} la base de connaissances pour permettre l’extraction et la g\'{e}n\'{e}ration de l’UI finale. Ensuite, il faut
d\'{e}crire les mod\`{e}les abstraits. Par exemple, nous utilisons les mod\`{e}les abstraits de l’approche MORPH
pour repr\'{e}senter les interactions des \'{e}l\'{e}ments graphiques d’un art\'{e}fact de l’UI CBA (cf. figure 3.4) :
– le menu principal a une t\^{a}che d’interactions de type SELECTION-OBJET avec les r\^{o}les (action
= Procedural-Action, number-of-states = (2..10), variability = fixed, grouping= not-grouped)
– la liste d\'{e}roulante (ComboBox) a une t\^{a}che d’interactions de type SELECTION-OBJET avec
les r\^{o}les (action= Procedural-Action, number-of-states = (2..10), variability = fixed, grouping =
not-grouped)
– la liste d’images a une t\^{a}che d’interactions de type SELECTION-OBJET avec les r\^{o}les (action
= Procedural-Action, number-of-states = (2..10), variability = fixed, grouping = not-grouped)
Enfin, il faut mettre en place des r\`{e}gles de transformations des mod\`{e}les abstraits en prenant en
compte les guidelines. Par exemple, la prise en compte de la guideline 2 de partage de l’espace de
travail qui pr\'{e}conise de remplacer les composants graphiques qui en contiennent d’autres avec ceux
qui sont d\'{e}pla\c{c}ables. Cependant, nous remarquons que le mod\`{e}le de t\^{a}ches d’interactions pr\'{e}sent\'{e}
ci-dessus ne permet pas de d\'{e}crire les comportements des composants graphiques d’une biblioth\`{e}que
mais les activit\'{e}s possibles des utilisateurs. Par exemple ce mod\`{e}le permet de caract\'{e}riser les fonctionnalit\'{e}s d’un menu, mais ne permet pas de pr\'{e}ciser si un menu est d\'{e}pla\c{c}able. Dans le cadre de la
migration vers les tables interactives, les comportements des \'{e}l\'{e}ments graphiques sont indispensables
pour prendre en compte les guidelines pour les UI collaboratives. Les mod\`{e}les abstraits utilis\'{e}s dans
ce cadre doivent \^{e}tre capables de d\'{e}crire et de s\'{e}lectionner les composants graphiques conformes \`{a}
une guideline.
4.3. APPROCHES DE MIGRATION BASÉES SUR LES MODÈLES DE L’UI
61
4.3.2.4
Synth\`{e}se du processus semi automatique de migration des UI
Cette approche propose une solution de migration r\'{e}utilisable qui prend en compte une phase
de reverse engineering (d\'{e}tection), une phase de re engineering (transformation) et une phase de
forward engineering (g\'{e}n\'{e}ration). Le tableau 4.3 pr\'{e}sente l’approche MORPH suivant les caract\'{e}ristiques d’\'{e}quivalences des \'{e}l\'{e}ments des plateformes, de mod\'{e}lisation et de prise en compte des
guidelines.
Équivalences
Mod\'{e}lisations
Prise en compte des
guidelines
MORPH
Dynamique des
biblioth\`{e}ques
graphiques bas\'{e}es sur
des mod\`{e}les de
connaissances
Mod\'{e}lisation des
interactions abstraites,
Mod\'{e}lisation de la
structure
Prise en compte par les
r\`{e}gles de
transformations des
mod\`{e}les abstraits
TABLE 4.3 – R\'{e}capitulatif de MORPH
L’approche est bas\'{e}e sur des mod\`{e}les abstraits qui d\'{e}crivent la structure et les interactions (et les
comportements) des UI \`{a} migrer. Le layout et le style des UI ne sont pas mod\'{e}lis\'{e}s par cette approche
et leur migration est effectu\'{e}e manuellement par un utilisateur humain. Cependant les mod\`{e}les abstraits propos\'{e}s par MORPH ne sont pas adapt\'{e}s pour \^{e}tre utilis\'{e}s sur une plateforme comme une table
interactive. Le mod\`{e}le de t\^{a}ches d’interactions par exemple ne prend pas en compte les mouvements
(rotation, d\'{e}placement, etc) des composants graphiques des tables interactives. Le mod\`{e}le de t\^{a}ches
d’interactions permet d’\'{e}tablir des \'{e}quivalences entre dispositifs d’interactions en se basant sur des
t\^{a}ches de haut niveau d’abstraction et ind\'{e}pendantes des plateformes. Cependant ce mod\`{e}le n’est pas
adapt\'{e} \`{a} la migration des UI vers les tables interactives.
Les \'{e}quivalences entre les biblioth\`{e}ques graphiques sont bas\'{e}es sur des mod\`{e}les de connaissances
et peuvent \^{e}tre \'{e}tablies dynamiquement. Ces \'{e}quivalences sont uniquement bas\'{e}es sur les types des
composants graphiques.
Les \'{e}quivalences dynamiques entre les biblioth\`{e}ques graphiques permettent de s\'{e}lectionner les
composants graphiques appropri\'{e}s pour chaque cas en fonction de leurs caract\'{e}ristiques et de leur
utilisation.
Cette approche permet aussi la prise en compte des guidelines \`{a} travers les transformations (remplacement des objets abstraits ou s\'{e}lection des composants graphiques sp\'{e}cifiques). Cependant l’identification des guidelines et leur prise en compte pendant la transformation doivent \^{e}tre d\'{e}crit par le
concepteur de la solution de migration sur une plateforme comme les tables interactives.
L’\'{e}valuation de cette approche suivant l’axe de la r\'{e}utilisabilit\'{e} des diff\'{e}rents m\'{e}canismes r\'{e}v\`{e}le
qu’elle est moins r\'{e}utilisable que les approches bas\'{e}es sur des processus automatiques car tous les
aspects de l’UI ne sont pas d\'{e}crits par des mod\`{e}les abstraits. Cette approche pr\'{e}sente un co\^{u}t de mise
en œuvre important. Dans le cadre de la migration des UI vers les tables interactives, il est d’abord
indispensable de d\'{e}crire un mod\`{e}le d’interactions abstraites capables de prendre en compte toutes les
interactions. Ensuite il est n\'{e}cessaire de d\'{e}crire les transformations automatiques de la structure. Sur
la figure 4.8, nous \'{e}valuons la r\'{e}utilisabilit\'{e} de cette approche par une valeur inf\'{e}rieure \`{a} celle de
l’approche automatique et sup\'{e}rieure au portage des UI.
La flexibilit\'{e} de la migration semi automatique est sup\'{e}rieure \`{a} celle des approches automatiques
car le processus fait intervenir un humain pour adapter les aspects qui ne sont pas pris en compte par
les mod\`{e}les abstraits. Sur la figure 4.8, nous \'{e}valuons la flexibilit\'{e} de cette approche par une valeur
sup\'{e}rieure \`{a} celle de l’approche automatique et inf\'{e}rieure \`{a} l’approche manuelle.
62
CHAPITRE 4. APPROCHES DE MIGRATION DES UI
La prise en compte des guidelines dans cette approche est faite par les processus automatiques
et en se basant sur les connaissances des personnes en charge de la migration pour les adaptations
manuelles. L’UI produite permet un respect des crit\`{e}res ergonomiques \'{e}gal \`{a} l’approche manuelle.
FIGURE 4.8 – Synth\`{e}se de l’approche semi automatique
4.4
Synth\`{e}se et objectifs
Dans cette section nous faisons une synth\`{e}se des approches de migration des UI d\'{e}crites dans ce
chapitre et nous sp\'{e}cifions nos objectifs en nous basant sur les r\'{e}sultats de cette synth\`{e}se.
4.4.1
Synth\`{e}se
Nous avons pr\'{e}sent\'{e} \`{a} la figure 4.9 le r\'{e}capitulatif des approches de migration des UI \'{e}tudi\'{e}es suivant les crit\`{e}res d’\'{e}valuation d\'{e}crits \`{a} la section 4.1. Ces approches de migration des UI nous montrent
que les solutions de migration peuvent \^{e}tre sp\'{e}cifiques \`{a} une application ou \`{a} une biblioth\`{e}que graphique. Ces solutions peuvent aussi \^{e}tre r\'{e}utilisables en se basant sur une mod\'{e}lisation des diff\'{e}rents
aspects d’une UI.
Cette \'{e}tude nous a permis de raffiner les m\'{e}canismes d’\'{e}quivalences. Nous avons identifi\'{e} deux
types d’\'{e}quivalences entre les \'{e}l\'{e}ments des plateformes :
– les \'{e}quivalences statiques entre les biblioth\`{e}ques graphiques qui sont d\'{e}finies par les concepteurs \`{a} l’aide d’une table d’\'{e}quivalences par exemple,
– et les \'{e}quivalences dynamiques qui sont \'{e}tablies en se basant sur les caract\'{e}ristiques des \'{e}l\'{e}ments \`{a} comparer en utilisant les inf\'{e}rences d’un mod\`{e}le de connaissances par exemple.
En ce qui concerne les m\'{e}canismes de prise en compte des guidelines, nous avons constat\'{e} que
le concepteur peut se baser sur ses connaissances dans une approche manuelle pour d\'{e}crire les m\'{e}canismes de transformations et d’\'{e}quivalences. Pour des approches qui mod\'{e}lisent l’UI \`{a} migrer, les
guidelines \`{a} consid\'{e}rer sont traduites en r\`{e}gles de transformation des diff\'{e}rents aspects. Les guidelines permettent aussi d’\'{e}tablir des \'{e}quivalences entre les biblioth\`{e}ques graphiques.
4.4. SYNTHÈSE ET OBJECTIFS
63
FIGURE 4.9 – Synth\`{e}se des approches de migration des UI
Les m\'{e}canismes d’adaptation bas\'{e}s sur des mod\`{e}les de l’UI utilis\'{e}s par les diff\'{e}rentes approches
pr\'{e}sent\'{e}es nous permettent d’affirmer que :
– les interactions des UI peuvent \^{e}tre mod\'{e}lis\'{e}es partiellement par les mod\`{e}les de t\^{a}ches et
AUI [PMM97] ainsi que les t\^{a}ches d’interactions [FvDFH90, KSM99a].
– la structure des UI comprend \`{a} la fois les donn\'{e}es de l’UI et les relations hi\'{e}rarchiques entre
les diff\'{e}rents composants de l’UI. Cet aspect de l’UI peut \^{e}tre mod\'{e}lis\'{e} par des m\'{e}ta donn\'{e}es [BRNB07] ou des mod\`{e}les de l’UI [VLM+04, PSS09] pour d\'{e}crire les composants graphiques et les donn\'{e}es de l’UI.
– le positionnement des \'{e}l\'{e}ments de l’UI graphiques peut aussi \^{e}tre d\'{e}crit dans un mod\`{e}le ind\'{e}pendant d’une plateforme. Les langages de description des UI (UIDL) tels que USIXML,
MARIA, UIML permettent de d\'{e}crire le layout.
– le style des UI peut \^{e}tre mod\'{e}lis\'{e} au travers d’UIDL comme UIML 11.
4.4.1.1
Les points forts des approches pr\'{e}sent\'{e}es
La figure 4.9 pr\'{e}sente la synth\`{e}se des approches \'{e}tudi\'{e}es suivant les axes de la r\'{e}utilisabilit\'{e}, la
flexibilit\'{e} des m\'{e}canismes et le respect des crit\`{e}res ergonomiques \`{a} travers la prise en compte des
guidelines. Nous identifions les points forts suivants :
– Les solutions de migration des UI [WB03, MR97] flexibles et qui font intervenir les utilisateurs humains permettent d’avoir des UI qui respectent des crit\`{e}res ergonomiques et qui sont
proche des attentes des utilisateurs finaux. En effet ces approches permettent une reconception
manuelle pendant la migration et les prises en compte des guidelines sont plus fines pour les
utilisateurs ayant une bonne connaissance des plateformes cibles. Parmi les approches r\'{e}utilisables, les processus semi automatiques sont les plus flexibles.
– Les solutions de migration bas\'{e}es sur les mod\`{e}les permettent de d\'{e}crire des m\'{e}canismes de
transformations et d’\'{e}quivalences r\'{e}utilisables pour des applications respectant une architecture
11. UIML : An Appliance-Independent XML User Interface Language
64
CHAPITRE 4. APPROCHES DE MIGRATION DES UI
d\'{e}finie. L’utilisation des mod\`{e}les de l’UI permet d’accro\^{i}tre la r\'{e}utilisabilit\'{e} d’une approche et
permet aussi la r\'{e}duction du temps de la migration des applications une fois que les mod\`{e}les et
les m\'{e}canismes de transformation sont mis en place.
4.4.1.2
Les limites
La r\'{e}utilisation des approches bas\'{e}es sur des mod\`{e}les de l’UI dans le cadre des tables interactives
nous permet d’identifier les limites ci-dessous :
– Le portage des UI n’est pas une approche de migration des UI adapt\'{e}es pour la migration
vers les tables interactives car elle ne permet pas la prise en compte des guidelines des UI
collaboratives. Cette approche peut par contre \^{e}tre adopt\'{e}e pour la migration entre plateformes
proches, par exemple pour faire migrer des UI Microsoft PixelSense 1.0 vers la version 2.0
gr\^{a}ce aux techniques de types wrappers.
– Les approches de migration manuelles pr\'{e}sentent des co\^{u}ts de mise en œuvre comparables \`{a}
ceux d’une nouvelle conception. L’automatisation des m\'{e}canismes d’\'{e}quivalences et d’adaptations de l’UI source permet une r\'{e}duction de ces co\^{u}ts.
– Les services de migration des UI sont des solutions avec des processus ferm\'{e}s et tr\`{e}s peu
flexibles. La flexibilit\'{e} d’une approche permet d’avoir des r\'{e}sultats conformes aux attentes des
utilisateurs finaux.
– L’approche semi automatique est flexible mais les m\'{e}canismes d’adaptation ou de prise en
compte des guidelines n’aident pas les concepteurs ou d\'{e}veloppeurs pendant la migration. En
effet les m\'{e}canismes d’\'{e}quivalences des composants graphiques peuvent proposer les \'{e}quivalences en pr\'{e}cisant ceux qui sont conformes aux guidelines de la cible par exemple.
4.4.2
Objectifs
Les principaux objectifs de notre travail de th\`{e}se est de r\'{e}duire le co\^{u}t de la migration et de
prendre en compte les guidelines de plateforme d’arriv\'{e}e lors de la migration des UI existantes vers
les tables interactives. Nous faisons l’hypoth\`{e}se que les applications \`{a} migrer s\'{e}pare le NF et l’UI.
Pour atteindre ces objectifs, nous choisissons d’adapter l’UI source par un processus de migration
semi automatique en trois \'{e}tapes : d’abord extraire l’UI source par reverse engineering en l’abstrayant
dans des mod\`{e}les abstraits. Ensuite, transformer les mod\`{e}les abstraits tout en faisant intervenir le
concepteur pour personnaliser les aspects non mod\'{e}lis\'{e}s. Et pour terminer, nous g\'{e}n\'{e}rons l’UI finale
\`{a} partir des mod\`{e}les abstraits transform\'{e}s.
En consid\'{e}rant cet objectif principal, l’\'{e}tude des diff\'{e}rentes approches de migration des UI de
fa\c{c}on globale, mais aussi plus particuli\`{e}rement ceux traitant de la migration vers les tables interactives,
nous montre des mod\`{e}les d’interactions qui ne facilitent pas les \'{e}quivalences entre les plateformes.
Pour atteindre notre objectif, nos contributions sont :
– de proposer un mod\`{e}le d’interactions qui permet d’accro\^{i}tre la flexibilit\'{e} des m\'{e}canismes de
changement de modalit\'{e} d’interactions et la pr\'{e}servation des interactions de l’UI de d\'{e}part tout
en prenant en compte celles de la plateforme d’arriv\'{e}e,
– de d\'{e}crire des m\'{e}canismes d’adaptations r\'{e}utilisables et flexibles en prenant en compte les
guidelines de la plateforme cible.
Dans la partie suivante, nous pr\'{e}sentons le cœur de notre approche. En nous servant de cette \'{e}tude,
nous proposons une approche bas\'{e}e sur un mod\`{e}le d’interactions abstraites qui permet d’\'{e}tablir des
\'{e}quivalences dynamiques entre les plateformes en prenant en compte les interactions des instances.
Nous souhaitons inclure aussi la prise en compte des guidelines pour les tables interactives.
Troisi\`{e}me partie
M\'{e}thodes
65
CHAPITRE 5
Mod\`{e}les pour la migration de UI vers les
tables interactives
Sommaire
5.1
Introduction . . . . . . . . . . . . . . . . . . . . . . . . . . . . . . . . . . . . .
67
5.2
Primitives d’interactions
. . . . . . . . . . . . . . . . . . . . . . . . . . . . . .
68
5.2.1
Les primitives d’interactions en entr\'{e}e . . . . . . . . . . . . . . . . . . . .
68
5.2.2
Les primitives d’interactions en sortie . . . . . . . . . . . . . . . . . . . .
72
5.3
Mod\`{e}les de composants graphiques
. . . . . . . . . . . . . . . . . . . . . . . .
73
5.3.1
Un mod\`{e}le de types de composants graphiques . . . . . . . . . . . . . . .
75
5.3.2
Un mod\`{e}le d’instance d’une UI
. . . . . . . . . . . . . . . . . . . . . . .
81
5.3.3
Synth\`{e}se des mod\`{e}les abstraits . . . . . . . . . . . . . . . . . . . . . . . .
90
5.4
Op\'{e}rateurs d’\'{e}quivalences
. . . . . . . . . . . . . . . . . . . . . . . . . . . . .
91
5.4.1
Op\'{e}rateurs d’\'{e}quivalences . . . . . . . . . . . . . . . . . . . . . . . . . .
91
5.4.2
Op\'{e}rateurs d’\'{e}quivalences \textbackslash{}& types de donn\'{e}es . . . . . . . . . . . . . . .
95
5.5
Synth\`{e}se
. . . . . . . . . . . . . . . . . . . . . . . . . . . . . . . . . . . . . . .
98
5.1
Introduction
Les approches de migration des UI g\'{e}n\'{e}riques et r\'{e}utilisables que nous avons \'{e}tudi\'{e}es dans le
chapitre pr\'{e}c\'{e}dent se basent sur des mod\`{e}les abstraits qui d\'{e}crivent diff\'{e}rents aspects 12 des UI ind\'{e}pendamment des plateformes et des applications.
La migration des UI desktops vers les tables interactives implique n\'{e}cessairement une prise en
compte des nouveaux dispositifs d’interactions disponibles dans cet environnement. Afin d’effectuer cette migration de mani\`{e}re coh\'{e}rente il nous semble essentiel de mettre en place un m\'{e}canisme
d’\'{e}quivalences entre les instruments d’interactions disponibles sur les plateformes source et cible.
Ce m\'{e}canisme d’\'{e}quivalences prend en compte les principes de conception n\'{e}cessaires \`{a} la mise en
œuvre des UI collaboratives et tangibles. Pour ce faire, nous nous basons sur l’id\'{e}e de d\'{e}crire les
correspondances entre les instruments d’interactions de la source et de la cible \`{a} l’aide d’un mod\`{e}le
d’interactions abstraites (cf section 3.1.4). Ce mod\`{e}le d’interactions abstraites permet de caract\'{e}riser
les instruments d’interactions ind\'{e}pendamment des plateformes et des langages.
Ce chapitre propose donc \`{a} la section 5.2 un mod\`{e}le d’interactions abstraites qui d\'{e}crit les interactions atomiques possibles sur les composants graphiques ind\'{e}pendamment des instruments d’interactions et celles qui sont n\'{e}cessaires pour la migration des UI vers les tables interactives. Nous
utiliserons ult\'{e}rieurement ce mod\`{e}le pour d\'{e}crire des \'{e}quivalences entre composants graphiques. La
section 5.3 pr\'{e}sente les caract\'{e}ristiques des composants graphiques en fonction du mod\`{e}le d’interactions abstraites propos\'{e}. La section 5.4 propose un ensemble d’op\'{e}rateurs d’\'{e}quivalences entre les
composants graphiques en se basant sur le mod\`{e}le d’interactions abstraites propos\'{e}.
12. Interactions [GH95, KSM99b], Structure, Positionnement [PSS09, VLM+04] et Style
67
68
CHAPITRE 5. MODÉLISATION DES UI
5.2
Un mod\`{e}le d’interactions abstraites
Les mod\`{e}les de syst\`{e}mes interactifs [Nig94, RSP11, Weg97] distinguent deux types d’interactions
entre les utilisateurs et une UI [W3C03] : les interactions en entr\'{e}e et les interactions en sortie. Les
interactions en entr\'{e}e permettent de fournir des donn\'{e}es ou d’invoquer des fonctionnalit\'{e}s de l’application gr\^{a}ce \`{a} une UI et \`{a} des dispositifs d’entr\'{e}e (souris, clavier, microphone, camera de reconnaissance gestuelle, \'{e}cran tactile, acc\'{e}l\'{e}rom\`{e}tre, etc.). Les interactions en sortie permettent d’effectuer des
rendus des donn\'{e}es de l’application \`{a} travers des dispositifs de sortie (tel un \'{e}cran, des haut-parleurs,
etc.). Les interactions (en entr\'{e}e ou en sortie) se d\'{e}clinent en plusieurs modalit\'{e}s d’interactions en
fonction des langages et des dispositifs d’entr\'{e}e ou des dispositifs de sortie utilis\'{e}s. Chaque modalit\'{e}
d’interactions est utilis\'{e}e comme un canal de communication entre un utilisateur et une application
pour transmettre ou acqu\'{e}rir des informations [Nig94].
Les composants graphiques sont caract\'{e}ris\'{e}s par les donn\'{e}es qu’ils contiennent, leurs propri\'{e}t\'{e}s
graphiques et leurs comportements. Ils appartiennent \`{a} des biblioth\`{e}ques graphiques qui permettent
de d\'{e}crire des UI graphiques en 2D, des images de synth\`{e}se en 3D, des jeux vid\'{e}o, des graphiques
de donn\'{e}es [BCW+06, Lon10, SWND03], etc. Ainsi, les interactions en entr\'{e}e modifient l’\'{e}tat d’un
composant graphique soit \`{a} travers ses propri\'{e}t\'{e}s (donn\'{e}es contenues, taille, position et repr\'{e}sentation), soit par son comportement (tels les appels des fonctionnalit\'{e}s du NF et les interactions sur
d’autres composants graphiques d’une UI, etc.). Les interactions en sortie quant \`{a} elles permettent
d’afficher des donn\'{e}es de l’application. Les donn\'{e}es des composants graphiques sont soit des donn\'{e}es graphiques si elles d\'{e}crivent le composant en question, soit des donn\'{e}es de l’application si
elles proviennent du NF ou des utilisateurs.
Une interaction en entr\'{e}e ou en sortie sur les composants graphiques d’une UI peut \^{e}tre d\'{e}compos\'{e}e en une succession d’actions \'{e}l\'{e}mentaires que nous appelons primitives d’interactions. Il est possible d’obtenir une d\'{e}composition atomique. Cette section pr\'{e}sente une caract\'{e}risation des primitives
d’interactions en fonction des deux types d’interactions (entr\'{e}e et sortie) ainsi qu’une mod\'{e}lisation
des primitives d’interactions qui sont utiles pour d\'{e}crire des \'{e}quivalences entre les plateformes source
et cible.
5.2.1
Les primitives d’interactions en entr\'{e}e
Elles d\'{e}crivent de mani\`{e}re abstraite toutes les actions atomiques des utilisateurs qui permettent de
s\'{e}lectionner des composants graphiques et d’en modifier une propri\'{e}t\'{e}. Les primitives d’interactions
en entr\'{e}e permettent de caract\'{e}riser aussi les interactions qui font appel aux fonctionnalit\'{e}s du NF.
Par exemple, l’action de cliquer \`{a} l’aide d’une souris sur un bouton qui fait appel \`{a} une fonctionnalit\'{e}
du NF se d\'{e}compose en plusieurs primitives d’interactions en entr\'{e}e.
Nous avons identifi\'{e} comme des interactions en entr\'{e}e sur les contenus : l’\'{e}dition (Data Edition),
la s\'{e}lection (Data Selection) et le d\'{e}placement (Data Move In et Data Move Out). Ces primitives
sont ind\'{e}pendantes des dispositifs d’interactions en entr\'{e}e et peuvent \^{e}tre effectu\'{e}es soit avec des
dispositifs de manipulation directe (souris, \'{e}cran tactile, reconnaissance gestuelle, etc.), soit avec des
dispositifs de manipulation indirecte (clavier physique ou virtuel, reconnaissance vocale, etc.) ou soit
par la combinaison de plusieurs dispositifs (clavier et souris).
Primitive d’interactions 1 : Data Edition
Elle consiste \`{a} modifier les donn\'{e}es des composants graphiques qui ont des
donn\'{e}es de l’application.
5.2. PRIMITIVES D’INTERACTIONS
69
La modification du contenu d’un composant graphique peut se faire avec un clavier, une reconnaissance vocale, gestuelle ou de forme en fonction des plateformes. Dans le cadre des tables interactives,
l’\'{e}dition se fait avec un clavier virtuel ou \`{a} main lev\'{e}e sur un \'{e}cran tactile. En ce qui concerne les
desktops, l’\'{e}dition d’un contenu se fait avec un clavier, une synth\`{e}se vocale ou avec un clavier virtuel
et une souris.
Primitive d’interactions 2 : Data Selection
Elle permet d’acc\'{e}der aux donn\'{e}es de l’application d’un composant graphique.
Les composants graphiques peuvent contenir plusieurs donn\'{e}es de m\^{e}me type ; dans le cas des listes
par exemple, cette primitive d’interactions permet la s\'{e}lection d’un contenu pr\'{e}cis. Elle se fait \`{a}
l’aide d’un clavier, d’une souris, d’un \'{e}cran tactile, d’un dispositif de reconnaissance gestuelle, etc.
Dans le cadre d’une table interactive, la s\'{e}lection de contenus se fait de mani\`{e}re directe sur l’\'{e}cran
tactile ou aussi avec un objet tangible. Par exemple, la s\'{e}lection dans une liste d’images avec un objet
tangible consiste \`{a} poser l’objet sur l’\'{e}l\'{e}ment en question.
Primitive d’interactions 3 : Data Move Out
Elle permet d’exporter la (les) donn\'{e}e(s) de l’application d’un composant
graphique vers un autre composant graphique comportant des donn\'{e}es de m\^{e}me
type.
Les composants graphiques peuvent s’\'{e}changer des contenus. Le drag et le drop permettent de
d\'{e}placer des contenus \`{a} l’aide d’une souris ou d’un geste tactile. Cette primitive d’interactions
correspond \`{a} la premi\`{e}re action d’un drag \textbackslash{}& drop. L’interaction peut \^{e}tre r\'{e}alis\'{e}e sur une table
interactive en utilisant l’\'{e}cran tactile ou \^{e}tre \'{e}mul\'{e}e avec un objet tangible. Avec un objet tangible
comme un stylo par exemple, cette primitive d’interactions repr\'{e}sente l’initiation de l’action drag \textbackslash{}&
drop apr\`{e}s la primitive de Data Selection du contenu \`{a} d\'{e}placer. La primitive d’interactions Data
Move Out exprime la possibilit\'{e} d’exporter la(les) donn\'{e}e(s) d’un composant graphique vers un autre.
Primitive d’interactions 4 : Data Move In
Elle permet de recevoir la (les) donn\'{e}e(s) de l’application d’un autre composant graphique de type compatible.
Cette primitive d’interactions correspond \`{a} la seconde \'{e}tape du drag \textbackslash{}& drop, l’interaction peut \^{e}tre
r\'{e}alis\'{e}e sur une table interactive en utilisant l’\'{e}cran tactile ou \^{e}tre \'{e}mul\'{e}e avec un objet tangible.
Avec un objet tangible comme un stylo par exemple, cette primitive d’interactions repr\'{e}sente la fin de
l’action drag \textbackslash{}& drop apr\`{e}s la primitive de Data Move Out. La primitive d’interactions Data Move
In exprime la possibilit\'{e} d’importer des donn\'{e}es vers un composant graphique.
Les quatre primitives d’interactions en entr\'{e}e sur le(s) contenu(s) des composants graphiques d\'{e}crivent de mani\`{e}re exhaustive l’ensemble des actions possibles qu’un utilisateur peut faire sur un
\'{e}l\'{e}ment graphique contenant des donn\'{e}es de l’application. En effet, ces primitives d’interactions permettent :
70
CHAPITRE 5. MODÉLISATION DES UI
– la cr\'{e}ation (ou “Create”, en \'{e}ditant les donn\'{e}es de l’application par la primitive d’interactions
Data Edition),
– l’acc\`{e}s en lecture (ou “Read”, en s\'{e}lectionnant le(s) contenu(s) par la primitive d’interactions
Data Selection ou Data Move Out),
– la modification (ou “Update”, en \'{e}ditant le(s) contenu(s) par les primitives d’interactions Data
Edition ou Data Move In) et
– la suppression des donn\'{e}es (ou “Delete”, des contenus avec les primitives d’interactions Data
Move Out et Data Edition).
Ces primitives d’interactions couvrent les op\'{e}rations Create, Read, Update et Delete (CRUD) d\'{e}finies
dans [D. 06] sur les donn\'{e}es des composants graphiques. Les primitives d’interactions en entr\'{e}e sur
les contenus constituent un ensemble complet d’interactions atomiques et permet de d\'{e}crire toutes les
actions utilisateurs sur les donn\'{e}es d’applications des composants graphiques.
Par ailleurs, les actions utilisateurs sur une UI peuvent aussi modifier les propri\'{e}t\'{e}s graphiques
lors d’un d\'{e}placement (Widget Move), ou par une rotation sans changement d’emplacement (Widget
Rotation), ou un redimensionnement (Widget Resize). Un composant graphique peut \^{e}tre s\'{e}lectionn\'{e}
de mani\`{e}re directe (Widget Selection) ou indirecte en naviguant \`{a} travers un container (Navigation).
Les primitives d’interactions sur les propri\'{e}t\'{e}s graphiques sont ind\'{e}pendantes des dispositifs d’interactions en entr\'{e}e. Elles peuvent \^{e}tre effectu\'{e}es avec des dispositifs de manipulation directe (souris,
\'{e}cran tactile, reconnaissance gestuelle, etc.), avec des dispositifs de manipulation indirecte (claviers
physiques ou virtuels, reconnaissance vocale, etc.) ou en combinant l’utilisation de plusieurs dispositifs (clavier et souris).
Primitive d’interactions 5 : Widget Move
Elle permet d’exprimer le changement de la position d’un composant graphique.
La position des composants graphiques peut \^{e}tre chang\'{e}e par un utilisateur \`{a} l’aide d’un dispositif de
manipulation directe ou des raccourcis d’un clavier. Cette primitive d’interactions est d\'{e}finie pour les
composants graphiques d\'{e}pla\c{c}ables \`{a} l’aide d’un dispositif d’interactions.
Primitive d’interactions 6 : Widget Rotation
Elle permet d’exprimer le changement d’orientation d’un composant graphique.
L’orientation des composants graphiques est une caract\'{e}ristique importante dans le cadre des UI collaboratives et co localis\'{e}es. En effet l’orientation des composants graphiques permet l’utilisation d’une
UI \`{a} partir de n’importe quelle position autour de la table. Cette opportunit\'{e} favorise l’utilisation d’une
UI par plusieurs personnes autour de la table. La rotation de composants graphiques n’est pas une interaction en entr\'{e}e indispensable pour les composants graphiques de la plateforme desktop car l’UI
est destin\'{e} \`{a} une utilisation dans une position fix\'{e}e. Dans le cadre de la migration vers des plateformes
multi-utilisateurs cette primitive d’interactions est indispensable.
Primitive d’interactions 7 : Widget Resize
Elle permet de modifier les dimensions d’un composant graphique.
Les dimensions d’un composant graphique peuvent \^{e}tre red\'{e}finies par un utilisateur d’une UI. Cette
action peut \^{e}tre effectu\'{e}e avec un clavier ou une souris sur un desktop et \`{a} l’aide d’un geste tactile
5.2. PRIMITIVES D’INTERACTIONS
71
sur une table interactive. Le redimensionnement des composants graphiques en contenant d’autres
implique de redimensionner leurs composants fils tout en pr\'{e}servant le crit\`{e}re ergonomique de guidage
(cf 3.2.2). La pr\'{e}servation du guidage implique d’avoir les m\^{e}mes composants graphiques apr\`{e}s le
redimensionnement. Les tailles des composants graphiques redimensionables peuvent \^{e}tre encadr\'{e}es
pour pr\'{e}server le guidage.
Primitive d’interactions 8 : Widget Selection
Elle permet d’exprimer la s\'{e}lection imm\'{e}diate de composants graphiques
avec un dispositif de manipulation directe.
Cette primitive d’interactions d\'{e}crit une s\'{e}lection directe d’un composant graphique par un utilisateur
\`{a} l’aide d’un dispositif de manipulation directe. Cette action peut \^{e}tre effectu\'{e}e avec un clavier ou une
souris sur un desktop et \`{a} l’aide d’un geste tactile sur une table interactive. C’est l’une des deux interactions atomiques pr\'{e}alables \`{a} toute action utilisateur sur les donn\'{e}es d’application ou les propri\'{e}t\'{e}s
graphiques d’un composant graphique.
Primitive d’interactions 9 : Navigation
Elle permet d’exprimer la s\'{e}lection d’un composant graphique de mani\`{e}re
s\'{e}quentielle.
Cette primitive d’interactions d\'{e}crit une s\'{e}lection s\'{e}quentielle des composants graphiques d’une UI
par un utilisateur \`{a} l’aide de dispositifs de manipulation non directe. C’est l’une des deux interactions
atomiques pr\'{e}alables \`{a} toutes les actions des utilisateurs sur les donn\'{e}es d’application ou les propri\'{e}t\'{e}s
graphiques d’un composant graphique.
Les primitives d’interactions ci-dessus 13 caract\'{e}risent les actions utilisateurs sur les propri\'{e}t\'{e}s
graphiques (position, orientation, taille et repr\'{e}sentation) d’un composant graphique. Elles permettent
de d\'{e}crire l’ensemble des actions utilisateurs sur l’orientation, la taille et la position des composants
graphiques.
Primitive d’interactions 10 : Activation
Elle permet de d\'{e}crire les interactions des composants graphiques avec les
autres composants graphiques ou le NF d’une application.
L’activation repr\'{e}sente le lien entre la structure d’une UI et les fonctionnalit\'{e}s. Concr\`{e}tement, dans le
cadre d’une architecture MVC, cette primitive permet de repr\'{e}senter un appel de m\'{e}thode du contr\^{o}leur par un \'{e}l\'{e}ment de la vue.
5.2.1.1
Exemple
Les primitives d’interactions en entr\'{e}e caract\'{e}risent les actions possibles des utilisateurs sur le(s)
contenu(s), la taille, la position, l’orientation des composants graphiques d’une UI. Elles sont des
interactions abstraites ind\'{e}pendantes des instruments d’interactions, elles sont d\'{e}riv\'{e}es en interactions
concr\`{e}tes sur les composants graphiques. En consid\'{e}rant l’art\'{e}fact d’une UI de l’application CBA (cf
chapitre 3), la liste d’images permet aux dessinateurs de BD de s\'{e}lectionner des images pr\'{e}d\'{e}finies
et de les ajouter \`{a} la BD en cours de r\'{e}alisation. Les \'{e}l\'{e}ments de cette liste peuvent \^{e}tre ajout\'{e}s au
canevas de dessin par une interaction de glisser-d\'{e}poser (drag \textbackslash{}& drop). Les primitives d’interactions
en entr\'{e}e de la liste d’images sont :
13. Widget Move, Widget Rotation, Widget Resize, Widget Selection et Navigation
72
CHAPITRE 5. MODÉLISATION DES UI
– Navigation et Widget Selection pour acc\'{e}der \`{a} chaque \'{e}l\'{e}ment de la liste
– Data Display pour afficher chaque image de la liste
– Data Selection, Data Move Out et Activation pour s\'{e}lectionner et pour effectuer le drag d’une
image de la liste.
Les primitives d’interactions en entr\'{e}e du canevas sont :
– Navigation et Widget Selection pour acc\'{e}der \`{a} chaque objet du canevas
– Data Selection, Data Move In et Activation pour effectuer le drop d’un objet sur le canevas.
– Data Selection, Data Edition et Activation pour s\'{e}lectionner et modifier les propri\'{e}t\'{e}s d’un
objet du canevas. Par exemple pour le redimensionnement d’une image, il faut s\'{e}lectionner
l’objet repr\'{e}sentant l’image, modifier la taille et faire appel au NF pour valider la modification.
Le tableau 5.1 pr\'{e}sente les primitives d’interactions en entr\'{e}e de la liste d’images, de la liste d\'{e}roulante
et du canevas.
Nous remarquons que les primitives d’interactions en entr\'{e}e permettent de caract\'{e}riser les actions
des utilisateurs sur un contenant (liste ou canevas) et sur les \'{e}l\'{e}ments contenus (images, objets des
canevas). La structure des composants graphiques au niveau abstrait doit tenir compte \`{a} la fois des
contenants et des contenus.
FIGURE 5.1 – Un art\'{e}fact d’une UI
Composants
graphiques
Primitives d’interactions en entr\'{e}e
Liste d’images
Widget Selection, Navigation, Data Display, Data Selection, Data Move Out,
Activation
Canevas
Widget Selection, Navigation, Data Selection, Data Edition, Data Move In,
Activation
TABLE 5.1 – Exemples de primitives d’interactions en entr\'{e}e
5.2.2
Les primitives d’interactions en sortie
Le mod\`{e}le d’interactions instrumentales pr\'{e}sent\'{e} \`{a} la section 3.1 caract\'{e}rise les interactions en
sortie en deux types d’interactions atomiques. D’une part, les r\'{e}ponses correspondent aux interactions du NF sur les donn\'{e}es d’applications ou les propri\'{e}t\'{e}s graphiques des composants graphiques.
5.3. MODÈLES DE COMPOSANTS GRAPHIQUES
73
D’autre part, les feedbacks et les r\'{e}actions des \'{e}l\'{e}ments d’une UI correspondent aux comportements
des composants graphiques. Les feedbacks et les r\'{e}actions sont les effets visibles d’une interaction,
par exemple l’apparition d’un pop up.
Dans cette section nous caract\'{e}risons les r\'{e}ponses sur le(s) contenu(s) et les propri\'{e}t\'{e}s graphiques
d’un composant graphique. Nous avons identifi\'{e} deux primitives d’interactions en sortie : la primitive
d’interactions Data Display et la primitive d’interactions Widget Display.
Primitive d’interactions 11 : Data Display
Elle permet d’exprimer l’affichage des donn\'{e}es d’application par un composant graphique.
Les donn\'{e}es d’application sont modifi\'{e}es par les actions utilisateurs ou par le NF. Cette primitive
exprime la modification et l’affichage des donn\'{e}es d’application par une r\'{e}ponse du NF.
Primitive d’interactions 12 : Widget Display
Elle exprime les r\'{e}ponses du NF ou d’autres composants graphiques sur l’aspect visuel d’un composant graphique.
Cette primitive d’interactions exprime les r\'{e}ponses du NF ou d’autres composants graphiques qui
modifient les propri\'{e}t\'{e}s graphiques d’un composant graphique. Par exemple l’affichage d’une bo\^{i}te
de dialogue ou d’une autre fen\^{e}tre apr\`{e}s une action utilisateur est une r\'{e}ponse, la bo\^{i}te de dialogue ou
la fen\^{e}tre concern\'{e}e par cette r\'{e}ponse doit d\'{e}crire la primitive d’interactions Widget Display.
FIGURE 5.2 – Bo\^{i}te de dialogue
5.2.2.1
Exemple
Les primitives d’interactions en sortie caract\'{e}risent les r\'{e}ponses du NF ou d’autres composants
graphiques.
Consid\'{e}rons par exemple la bo\^{i}te de dialogue de la figure 5.2 de l’application CBA : elle est
affich\'{e}e apr\`{e}s une demande de fermeture d’un document. La bo\^{i}te de dialogue aura la primitive d’interactions Widget Display telle que d\'{e}crite dans le tableau 5.2. Le message est affich\'{e} dans un label
qui obtient le nom du fichier \`{a} partir du NF, ce label aura donc la primitive d’interactions Data Display.
5.3
Mod\`{e}les de composants graphiques
Cette section a pour objectif d’expliquer comment les primitives d’interactions (PI) sont utilis\'{e}es
pour caract\'{e}riser les composants graphiques dans le but de d\'{e}crire des m\'{e}canismes d’\'{e}quivalences
et d’adaptation pour la migration des UI. Pour que ces m\'{e}canismes soient r\'{e}utilisables et flexibles
74
CHAPITRE 5. MODÉLISATION DES UI
Composants
graphiques
Primitives d’interactions en sortie
Label Message
Data Display
Bo\^{i}te de dialogue
Widget Display
TABLE 5.2 – Exemple de primitives d’interactions en sortie
nous allons caract\'{e}riser aussi les composants graphiques ind\'{e}pendamment des plateformes \`{a} l’aide de
mod\`{e}les abstraits.
Les composants graphiques appartiennent \`{a} des biblioth\`{e}ques graphiques (ou bo\^{i}tes \`{a} outils) et
sont utilis\'{e}s sous la forme d’instances pour d\'{e}crire des UI. Les composants graphiques d’instance
n’impl\'{e}mentent pas syst\'{e}matiquement toutes les interactions possibles de leur type. Nous consid\'{e}rons
la figure 5.3 qui illustre le lien entre des instances et des types de composants graphiques. Par exemple
un composant graphique de type ListBox d\'{e}finit les PI Data Move Out et Data Move In car il
supporte l’interaction glisser-d\'{e}poser (drag \textbackslash{}& drop) de mani\`{e}re intrins\`{e}que. Cependant de mani\`{e}re
effective, il est possible de l’instancier sans impl\'{e}menter ces PI dans le cas d’une liste qui ne permet
pas d’effectuer un glisser-d\'{e}poser de son contenu.
FIGURE 5.3 – Types et instances de composants graphiques
Les PI pr\'{e}sent\'{e}es ci-dessus (cf. section 5.2) sont utilis\'{e}es pour d\'{e}crire l’ensemble des interactions
possibles sur les composants graphiques. Elles font partie du mod\`{e}le abstrait des UI. Les aspects
structurels de l’UI tels que les types et les cardinalit\'{e}s de donn\'{e}es d’application d’une part, les liens
de contenance et les groupements des \'{e}l\'{e}ments graphiques d’autre part, sont d\'{e}crits dans un mod\`{e}le
de structure. Nous caract\'{e}risons dans cette section deux cat\'{e}gories de PI en fonction des types et
des instances des composants graphiques. Une UI finale est d\'{e}crite par les instances des composants
graphiques et les types de composants graphiques sont contenus dans des biblioth\`{e}ques graphiques.
Chaque type de composants graphiques d\'{e}finit des PI intrins\`{e}ques. Par ailleurs, les \'{e}l\'{e}ments d’une UI
finale sont mod\'{e}lis\'{e}s \`{a} l’aide d’un mod\`{e}le de structure qui d\'{e}crit les types de donn\'{e}es, leur cardinalit\'{e}
et le regroupement des \'{e}l\'{e}ments graphiques dont chaque instance impl\'{e}mente des PI effectives.
Dans le cadre de la migration des UI, cette diff\'{e}rentiation entre PI intrins\`{e}ques et PI effectives
5.3. MODÈLES DE COMPOSANTS GRAPHIQUES
75
permet de ne conserver pendant la migration que les interactions d’un composant graphique r\'{e}ellement
utilis\'{e}es par l’application.
Les PI effectives peuvent \^{e}tre identifi\'{e}es pour chaque biblioth\`{e}que graphique suivant deux approches. La premi\`{e}re approche consiste \`{a} d\'{e}crire manuellement les PI effectives de chaque composant
graphique en se basant sur les d\'{e}finitions des PI. Cette approche implique une identification exhaustive des PI effectives de chaque composant graphique d’une biblioth\`{e}que graphique. La deuxi\`{e}me
approche consiste d’abord \`{a} \'{e}tiqueter les \'{e}v\'{e}nements et les propri\'{e}t\'{e}s 14 des composants graphiques,
puis \`{a} d\'{e}crire des r\`{e}gles bas\'{e}es sur les \'{e}tiquettes de chaque composant graphique.
Nous avons opt\'{e} pour une identification des PI effectives bas\'{e}es sur l’\'{e}tiquetage en identifiant le
comportement de chaque \'{e}v\'{e}nement par rapport aux contenus ou aux propri\'{e}t\'{e}s graphiques des composants graphiques. En effet, cette seconde approche peut \^{e}tre reproduite en d\'{e}crivant les correspondances de chaque \'{e}tiquette dans une biblioth\`{e}que graphique et en appliquant les r\`{e}gles d’identifications. Le nombre d’\'{e}tiquettes caract\'{e}risant les \'{e}v\'{e}nements et les propri\'{e}t\'{e}s est inf\'{e}rieur aux nombres
de composants graphiques car il existe des \'{e}v\'{e}nements communs \`{a} plusieurs composants graphiques.
Par exemple l’\'{e}v\'{e}nement TextChanded est d\'{e}fini pour tout composant ayant une propri\'{e}t\'{e} Text en
XAML. Nous pr\'{e}sentons \`{a} la section 5.3.1.3 les \'{e}tiquettes permettant de caract\'{e}riser les \'{e}v\'{e}nements
des composants graphiques. L’\'{e}tiquetage est une t\^{a}che manuelle qui n\'{e}cessite une connaissance de la
syntaxe des \'{e}v\'{e}nements des composants graphiques des diff\'{e}rentes biblioth\`{e}ques graphiques. Elle est
effectu\'{e}e de mani\`{e}re unique pour les biblioth\`{e}ques graphiques des plateformes source et cible.
Les PI effectives sont ensuite identifi\'{e}es \`{a} partir des UI \`{a} l’aide des r\`{e}gles d’identifications sp\'{e}cifiques \`{a} chaque biblioth\`{e}que graphique en fonction des \'{e}v\'{e}nements effectivement impl\'{e}ment\'{e}s. Nous
proposons \`{a} la section 5.3.2.6 les r\`{e}gles d’identification des PI effectives qui sont utilis\'{e}es par notre
solution de migration des UI vers les tables interactives. Ces r\`{e}gles sont d\'{e}crites et valid\'{e}es pour la
biblioth\`{e}que graphique XAML.
La suite de cette section pr\'{e}sente d’abord un mod\`{e}le de types de composants graphiques ainsi que
les r\`{e}gles qui permet d’identifier les PI intrins\`{e}ques \`{a} chaque composant graphique. Nous pr\'{e}sentons
ensuite un mod\`{e}le de composants graphiques d’instances pour exprimer la structure d’une UI ainsi
que les r\`{e}gles d’identifications des PI effectives.
5.3.1
Un mod\`{e}le de types de composants graphiques
Les composants graphiques appartiennent \`{a} des biblioth\`{e}ques graphiques sp\'{e}cifiques (JavaSwing,
DiamondSpin, etc). Ce sont des \'{e}l\'{e}ments qui d\'{e}finissent des interactions, des comportements et des
donn\'{e}es sp\'{e}cifiques [Cre01]. Nous appelons dans cette th\`{e}se ”Widget”, les types de composants graphiques au niveau mod\`{e}le qui repr\'{e}sentent les \'{e}l\'{e}ments d’une UI d’une biblioth\`{e}que graphique.
De mani\`{e}re transversale et ind\'{e}pendamment des biblioth\`{e}ques graphiques, les Widgets permettent
de caract\'{e}riser par les diff\'{e}rents aspects d’une UI : les interactions entre les utilisateurs et l’application,
la structure (type et cardinalit\'{e} des donn\'{e}es), le positionnement dans une UI et le style (couleur, police,
taille, etc.). Nous proposons un mod\`{e}le de types de composants graphiques qui d\'{e}crit uniquement la
structure des \'{e}l\'{e}ments d’une biblioth\`{e}que graphique et les comportements possibles des composants
graphiques \`{a} la suite d’interactions (des utilisateurs, du NF ou d’autres composants graphiques). En
effet, dans le cadre d’un processus semi automatique (cf. section 4.3.2) de migration des UI, il n’est
pas indispensable de d\'{e}crire au niveau du mod\`{e}le tous les aspects de l’UI \`{a} migrer car les autres
aspects sont ajout\'{e}s manuellement.
L’objectif du mod\`{e}le de types de composants graphiques d\'{e}crit \`{a} la figure 5.4 est de repr\'{e}senter
l’ensemble des \'{e}l\'{e}ments d’une biblioth\`{e}que graphique avec ses caract\'{e}ristiques structurelles (type
14. Un \'{e}v\'{e}nement d\'{e}crit un comportement d’un composant graphique apr\`{e}s une interaction
76
CHAPITRE 5. MODÉLISATION DES UI
FIGURE 5.4 – Mod\`{e}le de types de composants graphiques
et cardinalit\'{e} de donn\'{e}es) et ses primitives d’interactions intrins\`{e}ques pour d\'{e}crire un m\'{e}canisme
d’\'{e}quivalences r\'{e}utilisables. La r\'{e}utilisabilit\'{e} des \'{e}quivalences est assur\'{e}e par un mod\`{e}le de structure
et les PI ind\'{e}pendantes des dispositifs d’interactions et des biblioth\`{e}ques graphiques.
Le mod\`{e}le de la figure 5.4 ci-dessus d\'{e}crit les classes Widget, PrimitiveInteractions et Event.
5.3.1.1
Widget
La classe Widget est identifi\'{e}e par son nom (name) et le nombre d’\'{e}l\'{e}ments qu’elle peut contenir
(cardinality). Le champ name permet de l’identifier de mani\`{e}re unique dans une biblioth\`{e}que graphique. La cardinalit\'{e} permet quant \`{a} elle de pr\'{e}ciser le nombre de donn\'{e}es ou de Widget maximal
qu’elle peut contenir.
L’attribut contentType pr\'{e}cise le type de donn\'{e}es d’un Widget. Les donn\'{e}es d’un Widget sont
des donn\'{e}es de l’application et/ou des donn\'{e}es graphiques. Par exemple : en XAML, les donn\'{e}es
d’un TextBox sont des donn\'{e}es de l’application pouvant \^{e}tre modifi\'{e}es par l’utilisateur tandis que les
donn\'{e}es d’un label sont des donn\'{e}es graphiques car elles ne sont pas destin\'{e}es \`{a} \^{e}tre modifi\'{e}es par
l’utilisateur.
Une biblioth\`{e}que graphique est constitu\'{e}e d’un ensemble d’instances de la classe Widget qui repr\'{e}sente les types de composants graphiques.
5.3. MODÈLES DE COMPOSANTS GRAPHIQUES
77
5.3.1.2
PrimitiveInteractions
La classe PrimitiveInteractions du mod\`{e}le de types de composants graphiques d\'{e}crit les PI intrins\`{e}ques d’un Widget \`{a} travers le r\^{o}le IntrinsicPI. Elle est caract\'{e}ris\'{e}e par l’attribut name qui exprime
le nom d’une PI.
Les PI intrins\`{e}ques d’un Widget sont exprim\'{e}es par des instances de la classe PrimitiveInteractions avec des noms diff\'{e}rents.
5.3.1.3
Event
La classe Event caract\'{e}rise le comportement des propri\'{e}t\'{e}s structurelles des Widgets par rapport
aux interactions sur les propri\'{e}t\'{e}s graphiques et sur les contenus. Ces comportements peuvent \^{e}tre des
appels de m\'{e}thodes, des modifications ou s\'{e}lections des propri\'{e}t\'{e}s graphiques des composants ou des
donn\'{e}es d’application. L’attribut name pr\'{e}cise le nom de l’\'{e}v\'{e}nement du composant graphique dans
sa biblioth\`{e}que graphique.
Étiquettes des Events
L’attribut inputDeviceType caract\'{e}rise les types des dispositifs d’interactions
en entr\'{e}e qui peuvent provoquer un comportement d’un Widget. Le type DirectManipulationType correspond aux dispositifs de manipulations directes tels que les souris, les \'{e}crans tactiles, les objets
tangibles, etc. Le type SequentialManipulationType correspond aux dispositifs de manipulations s\'{e}quentielles tels que les claviers physique et virtuel. Un comportement peut \^{e}tre d\'{e}clench\'{e} par plusieurs
dispositifs d’interactions en entr\'{e}e. Dans le cas o\`{u} le comportement est d\'{e}clench\'{e} par une interaction
en sortie, le type de dispositif d’interactions en entr\'{e}e est Null.
L’attribut eventType pr\'{e}cise le type de comportement d’un \'{e}v\'{e}nement par une \'{e}tiquette. Ces \'{e}tiquettes repr\'{e}sentent les diff\'{e}rents comportements des Widgets :
– La modification des propri\'{e}t\'{e}s si eventType = Change.
– La s\'{e}lection de donn\'{e}es d’application ou d’un composant graphique si eventType = Select.
– L’appel de m\'{e}thodes si eventType = Call.
Dans le but sp\'{e}cifier si une \'{e}tiquette concerne le contenu ou les propri\'{e}t\'{e}s graphiques, l’attribut
propertyType pr\'{e}cise sur quel type de propri\'{e}t\'{e} porte un comportement de l’\'{e}v\'{e}nement concern\'{e} de
mani\`{e}re suivante :
– si propertyType = Content alors le comportement porte sur les donn\'{e}es d’application ou les
donn\'{e}es graphiques,
– si propertyType = Size alors le comportement porte sur une propri\'{e}t\'{e} pr\'{e}cisant la taille de
l’\'{e}l\'{e}ment graphique,
– si propertyType = Position alors le comportement porte sur une propri\'{e}t\'{e} pr\'{e}cisant la position
de l’\'{e}l\'{e}ment graphique,
– si propertyType = Orientation alors le comportement porte sur une propri\'{e}t\'{e} pr\'{e}cisant l’orientation de l’\'{e}l\'{e}ment graphique.
Un \'{e}v\'{e}nement peut avoir plusieurs types. Par exemple une liste d\'{e}roulante permet de s\'{e}lectionner
un item de la liste et aussi de d\'{e}clencher l’appel d’une m\'{e}thode. Cette liste aura un attribut Event avec
\`{a} la fois eventType = Call et eventType = Select.
Les comportements des Widgets sont impl\'{e}ment\'{e}s de diff\'{e}rentes mani\`{e}res par les biblioth\`{e}ques
graphiques. Par exemple :
– En JavaSwing, les comportements de type Change sont d\'{e}finis pour un composant graphique
si la propri\'{e}t\'{e} isEditable est d\'{e}finie pour un composant graphique. Cette propri\'{e}t\'{e} permet de
dire qu’un composant graphique d\'{e}finit de mani\`{e}re intrins\`{e}que un \'{e}v\'{e}nement pour changer son
contenu.
78
CHAPITRE 5. MODÉLISATION DES UI
– En XAML, les comportements de type Select sont d\'{e}finis si la m\'{e}thode SelectedItem est d\'{e}finie
pour un composant graphique. Cette m\'{e}thode permet la s\'{e}lection d’un contenu.
– En javaSwing, les comportements de type Call sont d\'{e}finis si la m\'{e}thode addActionListener
est d\'{e}finie pour un composant graphique. Elle permet de dire qu’un composant graphique peut
d\'{e}crire un handler.
5.3.1.4
Remarques sur le mod\`{e}le de types de composants graphiques
Pour chaque biblioth\`{e}que graphique, nous proposons de construire une table qui d\'{e}crit les correspondances entre chaque type d’\'{e}v\'{e}nements et leur repr\'{e}sentation sp\'{e}cifique. Cette table permet
d’identifier les instances des Widgets. Des exemples de cette table sont d\'{e}crits de mani\`{e}re exhaustive
pour les biblioth\`{e}ques graphiques XAML et XAML Surface dans l’annexe A.1.
Dans le mod\`{e}le que nous proposons, nous avons identifi\'{e} un ensemble de types de donn\'{e}es d’application ou donn\'{e}es graphiques des composants graphiques. Les donn\'{e}es prises en compte par ce
mod\`{e}le de composant graphique sont de type bool\'{e}en (Boolean), entier (Integer), cha\^{i}ne de caract\`{e}res (String), image (Image), son ou vid\'{e}o (MediaElement), classes ind\'{e}pendantes de la biblioth\`{e}que
graphique (Object) ou composants graphiques (Widget). L’identification des types est faite \`{a} l’aide
d’une seconde table de correspondances entre les types sp\'{e}cifiques aux biblioth\`{e}ques graphiques et
les types d\'{e}crits dans notre mod\`{e}le.
Les m\'{e}canismes qui permettent de d\'{e}crire les instances du mod\`{e}le de types de composants graphiques sont d\'{e}crits \`{a} l’annexe A.1 pour les biblioth\`{e}ques graphiques XAML et XAML Surface. Il
existe autant d’instances de ce mod\`{e}le que de biblioth\`{e}ques graphiques.
5.3.1.5
Identification des primitives d’interactions intrins\`{e}ques
En nous basant sur le mod\`{e}le de types de composants graphiques propos\'{e} \`{a} la figure 5.4, nous caract\'{e}risons les diff\'{e}rentes PI en nous basant aussi sur leur d\'{e}finition et sur les types de comportements
des Widgets apr\`{e}s une interaction en entr\'{e}e (action) ou en sortie (r\'{e}ponse).
Dans cette section, nous pr\'{e}sentons les r\`{e}gles d’identification des PI intrins\`{e}ques en fonction
du mod\`{e}le de types de composant pr\'{e}sent\'{e} \`{a} la figure 5.4. Les r\`{e}gles pr\'{e}cisent les caract\'{e}ristiques
correspondant \`{a} chaque PI intrins\`{e}que.
R\`{e}gle 1 : Widget Selection \textbackslash{}& Navigation
Les PI Widget Selection et Navigation sont d\'{e}finies de mani\`{e}re intrins\`{e}que pour tous les Widgets
qui sont accessibles \`{a} l’aide d’un dispositif d’interactions directes (Widget Selection) ou indirectes
(Navigation).
∀w ∈\textbackslash{}{Widget\textbackslash{}}, ∃pi ∈w.InstrinsicPI,
∧
∃event ∈w.events,
DirectManipulationType ∈event.inputDeviceType ⇒pi.name =′ WidgetSelection′
∧
SequentialManipulationType ∈event.inputDeviceType ⇒pi.name =′ Navigation′
R\`{e}gle 2 : Widget Resize
5.3. MODÈLES DE COMPOSANTS GRAPHIQUES
79
La PI Widget Resize est d\'{e}finie de mani\`{e}re intrins\`{e}que pour tous les Widgets qui ont un comportement permettant de changer sa taille.
∀w ∈\textbackslash{}{Widget\textbackslash{}}, ∃pi ∈w.InstrinsicPI,




∃event ∈Widget.events,
event.eventType = Change
V
event.propertyType = Size



⇒pi.name =′ WidgetResize′
R\`{e}gle 3 : Widget Move
La PI Widget Move est d\'{e}finie de mani\`{e}re intrins\`{e}que pour tous les Widgets qui ont un comportement permettant de changer sa position.
∀w ∈\textbackslash{}{Widget\textbackslash{}}, ∃pi ∈w.InstrinsicPI,




∃event ∈Widget.events,
event.eventType = Change
V
event.propertyType = Position



⇒pi.name =′ WidgetMove′
R\`{e}gle 4 : Widget Rotation
La PI Widget Rotation est d\'{e}finie de mani\`{e}re intrins\`{e}que pour tous les Widgets qui ont un comportement permettant de changer son orientation.
∀w ∈\textbackslash{}{Widget\textbackslash{}}, ∃pi ∈w.InstrinsicPI,




∃event ∈Widget.events,
event.eventType = Change
V
event.propertyType = Orientation



⇒pi.name =′ WidgetRotation′
R\`{e}gle 5 : Widget Display
La PI Widget Display est d\'{e}finie de mani\`{e}re intrins\`{e}que pour tous les Widgets, car tout composant
graphique peut \^{e}tre visible sauf indication contraire du concepteur pour une instance.
∀w ∈\textbackslash{}{Widget\textbackslash{}}, ∃pi ∈w.InstrinsicPI,
pi.name ∈\textbackslash{}{′WidgetDisplay′\textbackslash{}}
R\`{e}gle 6 : Data Edition
La PI Data Edition est d\'{e}finie de mani\`{e}re intrins\`{e}que pour tous les Widgets qui ont un comporte80
CHAPITRE 5. MODÉLISATION DES UI
ment qui permet de changer des propri\'{e}t\'{e}s de contenu de type entier ou cha\^{i}ne de caract\`{e}res.
∀w ∈\textbackslash{}{Widget\textbackslash{}}, ∃pi ∈w.InstrinsicPI,








w.contentType ∈\textbackslash{}{Integer,String\textbackslash{}}
∃event ∈Widget.events,
V
event.eventType = Change
V
event.propertyType = Content








⇒pi.name =′ DataEdition′
R\`{e}gle 7 : Data Selection
La PI Data Selection est d\'{e}finie de mani\`{e}re intrins\`{e}que pour tous les Widgets qui ont un comportement qui permet la s\'{e}lection de son contenu.
∀w ∈\textbackslash{}{Widget\textbackslash{}}, ∃pi ∈w.InstrinsicPI,








w.contentType /∈\textbackslash{}{Null, Widget\textbackslash{}}
∃event ∈Widget.events,
V
event.eventType = Select
V
event.propertyType = Content








⇒pi.name =′ DataSelection′
R\`{e}gle 8 : Data Move In et Data Move Out
Les PI Data Move In et Data Move Out sont d\'{e}finies de mani\`{e}re intrins\`{e}que pour tous les Widgets
qui ont un comportement qui permet la s\'{e}lection et le changement des contenus.
∀w ∈\textbackslash{}{Widget\textbackslash{}}, ∃pi ∈w.InstrinsicPI,










w.contentType /∈\textbackslash{}{Null\textbackslash{}}
∃event ∈Widget.events,
V
event.eventType = Change
V
event.eventType = Select V
event.propertyType = Content










⇒pi.name ∈\textbackslash{}{′DataMoveIn′, ′DataMoveOut′\textbackslash{}}
R\`{e}gle 9 : Data Display
La PI Data Display est d\'{e}finie de mani\`{e}re intrins\`{e}que pour tous les Widgets qui n’ont pas de
contenu nul.
∀w ∈\textbackslash{}{Widget\textbackslash{}}, ∃pi ∈w.InstrinsicPI,
w.contentProperty /∈\textbackslash{}{Null\textbackslash{}} ⇒pi.name ∈\textbackslash{}{′DataDisplay′\textbackslash{}}
5.3. MODÈLES DE COMPOSANTS GRAPHIQUES
81
La PI est d\'{e}finie de mani\`{e}re intrins\`{e}que pour tous les Widgets qui ont des \'{e}v\'{e}nements de type
Call.
R\`{e}gle 10 : Activation
∀w ∈\textbackslash{}{Widget\textbackslash{}}, ∃pi ∈w.InstrinsicPI,

event ∈w.events
event.eventType = Call

⇒pi.name ∈\textbackslash{}{′Activation′\textbackslash{}}
Remarque
Nous avons valid\'{e} les r\`{e}gles d’identification des PI propos\'{e}es dans cette section pour
les biblioth\`{e}ques graphiques XAML et XAML Surface (cf. l’annexe A.3).
5.3.1.6
Synth\`{e}se
Le mod\`{e}le de types de composants graphiques d\'{e}crit \`{a} la section 5.3.1 permet de repr\'{e}senter les
composants graphiques d’une biblioth\`{e}que graphique et les primitives d’interactions intrins\`{e}ques associ\'{e}es. La repr\'{e}sentation des \'{e}l\'{e}ments d’une biblioth\`{e}que graphique dans ce mod\`{e}le se fait \`{a} l’aide
de deux tables de correspondances. La premi\`{e}re d\'{e}crit les correspondances entre les types de propri\'{e}t\'{e}s et leur repr\'{e}sentation dans les biblioth\`{e}ques graphiques (cf annexe 8.2.2). La seconde d\'{e}crit
les correspondances entre les types d’\'{e}v\'{e}nements et leur repr\'{e}sentation dans les biblioth\`{e}ques graphiques.
Les primitives d’interactions intrins\`{e}ques sont identifi\'{e}es \`{a} l’aide des r\`{e}gles se basant sur le mod\`{e}le des composants graphiques types. Par ailleurs, l’identification des PI intrins\`{e}ques des Widgets
d’une biblioth\`{e}que de mani\`{e}re exhaustive \`{a} l’aide d’une table de correspondances est possible. Cependant cette approche manuelle d\'{e}pend de l’interpr\'{e}tation des PI par les personnes en charge de la mise
en place et elle n\'{e}cessite le parcours de tous les Widgets d’une biblioth\`{e}que graphique. L’approche que
nous proposons, bas\'{e}e sur les r\`{e}gles d’identification, utilise aussi des tables de correspondances mais
elle est facilement r\'{e}utilisable car les tailles des tables de correspondances sont inf\'{e}rieures au nombre
de Widgets d’une biblioth\`{e}que graphique. Dans le cadre de notre approche de migration des UI vers
les tables interactives, nous consid\'{e}rons que ces tables de correspondances sont d\'{e}crites pour chaque
biblioth\`{e}que graphique pendant la mise en œuvre de la solution et fournies avec notre processus de
migration des UI.
5.3.2
Un mod\`{e}le d’instance d’une UI
La structure des instances des UI finales \`{a} migrer peut \^{e}tre d\'{e}crite par diff\'{e}rents formats (XML,
Archives Java, etc.). Cependant tous ces formats peuvent \^{e}tre repr\'{e}sent\'{e}s par un arbre dont la racine
est une fen\^{e}tre d’une UI et les feuilles, les diff\'{e}rents composants graphiques \'{e}l\'{e}mentaires.
Comme pour le mod\`{e}le de types de composants (cf. section 5.3.1), nous proposons un mod\`{e}le
d’instance d’une UI qui d\'{e}crit sa structure et ses interactions effectives. En effet nous avons pour objectif de d\'{e}crire un processus semi automatique de migration des UI qui prend en compte l’adaptation
automatique des aspects structurels et les \'{e}quivalences des instruments d’interactions. Dans cette optique, nous consid\'{e}rons qu’un mod\`{e}le de CUI capable de d\'{e}crire les liens de contenance, la cardinalit\'{e}
et les types de donn\'{e}es d’une UI peut \^{e}tre utilis\'{e}. Le but d’un mod\`{e}le de structure pour le processus
de migration est de pr\'{e}server les donn\'{e}es graphiques et d’adapter la structure de l’UI de d\'{e}part.
82
CHAPITRE 5. MODÉLISATION DES UI
FIGURE 5.5 – Mod\`{e}le de structure d’instance d’une UI
Le mod\`{e}le de structure d’une UI de la figure 5.5 que nous utilisons dans notre contexte comporte
des Interactors. Dans cette th\`{e}se nous appelons ”Interactor” les instances abstraites de composants
graphiques, c’est une repr\'{e}sentation au niveau mod\`{e}le des instances de Widget.
5.3.2.1
Interactor
Un Interactor est caract\'{e}ris\'{e} par un identifiant unique pour une UI et un nom qui correspond
\`{a} celui du composant graphique de l’UI \`{a} migrer. Un Interactor pr\'{e}serve les valeurs des donn\'{e}es
structurelles de l’UI \`{a} migrer par la classe Content, par exemple les \'{e}tiquettes des labels, des boutons,
des images, etc. Un interacteur est reli\'{e} \`{a} la classe Widget par une relation de type. Nous distinguons
deux types d’instances de Widget :
– Container repr\'{e}sentant l’ensemble des composants graphiques pouvant contenir d’autres instances de Widget, les Containers sont des nœuds de la structure arborescente.
– UIComponent repr\'{e}sentant l’ensemble des composants graphiques ne pouvant pas en contenir
d’autres. Ils sont identifi\'{e}s \`{a} partir des Widgets dont les contenus ne sont pas d’autres Widgets.
Les UIComponents sont les feuilles de la structure arborescente.
5.3.2.2
AUIStructure
La classe caract\'{e}rise l’ensemble des nœuds racines de la structure arborescente de l’UI. La structure des UI \`{a} migrer est abstraite \`{a} l’aide des \'{e}l\'{e}ments d\'{e}crits par le mod\`{e}le propos\'{e} \`{a} la figure 5.5.
5.3. MODÈLES DE COMPOSANTS GRAPHIQUES
83
5.3.2.3
ImplementedEvent
La classe caract\'{e}rise les comportements r\'{e}ellement impl\'{e}ment\'{e}s par les Interactors au niveau
mod\`{e}le. Elle est caract\'{e}ris\'{e}e par les attributs name, type et property.
– L’attribut name repr\'{e}sente le nom d’un \'{e}v\'{e}nement impl\'{e}ment\'{e}.
– L’attribut type repr\'{e}sente le comportement effectivement impl\'{e}ment\'{e}.
– L’attribut property repr\'{e}sente la propri\'{e}t\'{e} affect\'{e}e par le comportement impl\'{e}ment\'{e}. Les valeurs
de cette attribut sont Size, Content, Position, Orientation, WidgetStructure et Visible 15.
– L’attribut inputDeviceType correspond \`{a} l’attribut de la classe Event du mod\`{e}le des Widgets (cf
figure 5.4).
Étiquettes des \'{e}v\'{e}nements impl\'{e}ment\'{e}s
Les comportements effectivement impl\'{e}ment\'{e}s des Interactors sont identifi\'{e}s \`{a} partir de l’UI finale des applications \`{a} migrer.
Les comportements de type Change sont impl\'{e}ment\'{e}s pour une instance de la biblioth\`{e}que JavaSwing par exemple si la propri\'{e}t\'{e} isEditable = True, ce qui permet de dire qu’un composant graphique
d\'{e}finit de mani\`{e}re effective un comportement pour changer son contenu.
Les comportements de type Select sont impl\'{e}ment\'{e}s pour une instance de XAML par exemple si
la m\'{e}thode SelectedItem est d\'{e}finie et la propri\'{e}t\'{e} isEnabled = True.
Les comportements de type Call sont impl\'{e}ment\'{e}s pour une instance de la biblioth\`{e}que JavaSwing
par exemple si un ActionListener est d\'{e}fini pour cette instance.
Les comportements de type Delete sont impl\'{e}ment\'{e}s pour une instance ou pour les donn\'{e}es d’application d’une instance. Dans une biblioth\`{e}que graphique comme XAML par exemple ce comportement est d\'{e}finie de mani\`{e}re effective pour un composant qui impl\'{e}mente un Handler \`{a} l’\'{e}v\'{e}nement
DragEnter.
Les comportements de type Update sont impl\'{e}ment\'{e}s pour une instance. Ils ont des donn\'{e}es d’application modifi\'{e}es par le NF ou par d’autres composants graphiques. Dans la biblioth\`{e}que JavaSwing
par exemple si une propri\'{e}t\'{e} de contenu est modifi\'{e}e dans le code d’un listener alors ce comportement
est impl\'{e}ment\'{e} par l’instance du composant graphique.
Pour chaque biblioth\`{e}que graphique une table qui d\'{e}crit les correspondances entre les types d’\'{e}v\'{e}nements et leurs repr\'{e}sentations sp\'{e}cifiques, permet d’instancier les Interactors et leurs types d’\'{e}v\'{e}nements. Nous avons d\'{e}crit de mani\`{e}re exhaustive cette table pour les biblioth\`{e}ques graphiques XAML
et XAML Surface \`{a} l’annexe A.2.
5.3.2.4
Content
La classe Content du mod\`{e}le de structure pr\'{e}sent\'{e} \`{a} la figure 5.5 permet de pr\'{e}server les donn\'{e}es
graphiques et les donn\'{e}es d’application. Le processus de migration dans notre cas permet un changement de la modalit\'{e} d’interactions mais les donn\'{e}es d’une UI doivent \^{e}tre pr\'{e}serv\'{e}es pour conserver
la s\'{e}mantique de l’UI source. Pour ce faire, les types de donn\'{e}es (dataType) d’une UI source, leur
cardinalit\'{e} (cardinality) et leur(s) valeur(s) (value) sont d\'{e}crits dans notre mod\`{e}le pour permettre leur
pr\'{e}servation et \'{e}ventuellement leur adaptation \`{a} la nouvelle plateforme.
5.3.2.5
Types de Container
Les containers peuvent \^{e}tre cat\'{e}goris\'{e}s en fonction des interacteurs qu’ils contiennent. Nous caract\'{e}risons les types de container : Homog\`{e}ne, H\'{e}t\'{e}rog\`{e}ne, R\'{e}cursif et Racine. Le tableau 5.3 d\'{e}crit
15. Cette propri\'{e}t\'{e} permet d’impl\'{e}menter le comportement de changement de visibilit\'{e} d’une instance de Widget
84
CHAPITRE 5. MODÉLISATION DES UI
les types de container en fonction des types de composants graphiques qu’ils contiennent. En parcourant la structure arborescente du mod\`{e}le d’instance d’une UI, les types de container sont identifi\'{e}s
de mani\`{e}re r\'{e}cursive en identifiant d’abord les types des interacteurs fils. Un container peut \^{e}tre de
plusieurs types \`{a} la fois.
Ces quatre types de containers caract\'{e}risent les diff\'{e}rents regroupements de composants graphiques d’une UI \`{a} migrer. Nous avons remarqu\'{e} \`{a} la section 3.2.3 qui identifie les guidelines pour
la migration des UI vers les tables interactives que la plupart des transformations des UI de d\'{e}part en
UI collaboratives et tactiles porte sur les regroupements des composants graphiques. Les cat\'{e}gories
identifi\'{e}es dans cette section constituent des \'{e}l\'{e}ments basiques pour appliquer les guidelines d’une UI
collaboratives et tangibles pendant l’adaptation de la structure des UI sources.
Container Racine
C’est l’\'{e}l\'{e}ment racine d’une structure d’une UI. Il peut contenir des containers
et des interacteurs simples. En consid\'{e}rant la figure 5.6, le container bleu est une illustration d’un
Container de type Racine.
Container R\'{e}cursif
Il ne contient que des Containers. Ce type permet d’identifier un regroupement
de plusieurs containers sur lequel on pourra appliquer des transformations. Un container de type R\'{e}cursif ne peut donc pas contenir des interacteurs de types UIComponent. En consid\'{e}rant la figure 5.6,
le container vert est une illustration d’un Container de type R\'{e}cursif.
Container Homog\`{e}ne
Il contient uniquement des UIComponents dont les donn\'{e}es sont de m\^{e}me
type et de m\^{e}me cardinalit\'{e}. Un container de ce type permet de d\'{e}finir un groupe de composants
graphiques de m\^{e}me type (tel qu’une liste, un tableau, un menu, etc.). Les containers jaune et orange
de la figure 5.6 sont une illustration d’un Container de type Homog\`{e}ne.
Container H\'{e}t\'{e}rog\`{e}ne
C’est un container pouvant contenir \`{a} la fois des UIComponents et des
Containers. Ce type permet d’identifier les groupes de composants graphiques contenant des donn\'{e}es de types diff\'{e}rents. Les containers de types H\'{e}t\'{e}rog\`{e}ne peuvent repr\'{e}senter des formulaires qui
seront adapt\'{e}s aux tables interactives pour les rendre facilement accessibles ou pour les afficher par
un objet tangible. En consid\'{e}rant la figure 5.6, le container rouge est une illustration d’un Container
de type H\'{e}t\'{e}rog\`{e}ne.
Contenu
Types de container
\textbackslash{}{UIComponent\textbackslash{}}
Homog\`{e}ne : donn\'{e}es du m\^{e}me type
H\'{e}t\'{e}rog\`{e}ne : donn\'{e}es de type diff\'{e}rent
\textbackslash{}{UIComponent\textbackslash{}}
∪
\textbackslash{}{Container\textbackslash{}}
H\'{e}t\'{e}rog\`{e}ne : si a un parent
\textbackslash{}{Container\textbackslash{}}
R\'{e}cursif : si contient des containers
TABLE 5.3 – Types de container
5.3.2.6
Identification des primitives d’interactions effectives
Les primitives d’interactions r\'{e}ellement utilis\'{e}es par chaque interacteur sont identifi\'{e}es \`{a} partir de
la structure de l’UI de d\'{e}part. Dans le but de pr\'{e}server l’assemblage des composants graphiques de
5.3. MODÈLES DE COMPOSANTS GRAPHIQUES
85
FIGURE 5.6 – Illustration des types de container
l’UI \`{a} migrer, nous extrayons au moyen d’interacteurs cette structure arborescente \`{a} l’aide du mod\`{e}le
de la figure 5.5. Ce mod\`{e}le d\'{e}crit la structure de l’UI comme un ensemble d’arbres dont les racines
sont des fen\^{e}tres, les nœuds des interacteurs et les arcs la relation de contenance entre les interacteurs.
Dans cette section, nous pr\'{e}sentons les r\`{e}gles d’identification des PI effectives en fonction du
mod\`{e}le de structure d’une UI pr\'{e}sent\'{e} \`{a} la figure 5.5. Les r\`{e}gles pr\'{e}cisent les caract\'{e}ristiques correspondant \`{a} chaque PI effective.
R\`{e}gle 11 : Widget Selection \textbackslash{}& Navigation
Les PI Widget Selection et Navigation sont d\'{e}finies de mani\`{e}re effective pour tous les interacteurs
qui impl\'{e}mentent les comportements de s\'{e}lection de leur structure.
86
CHAPITRE 5. MODÉLISATION DES UI
∀interactor ∈AUIStructure, ∃pi ∈interactor.E f fectivePI,
∧
∃iEvt ∈interactor.events, iEvt.type = Select
∧
iEvt.property = WidgetStructure
∧
pi.name =′ WidgetSelection′ ⇒




∃p ∈interactor.type.IntrinsicPI,
p.name =′ WidgetSelection′
∧
DirectManipulationType ∈iEvt.inputDevice




∧
pi.name =′ Navigation′ ⇒




∃p ∈interactor.type.IntrinsicPI,
p.name =′ Navigation′
∧
SequentialManipulationType ∈iEvt.inputDevice




R\`{e}gle 12 : Widget Resize
La PI Widget Resize est d\'{e}finie de mani\`{e}re effective pour les interacteurs qui impl\'{e}mentent les
comportements de changement de taille.
∀interactor ∈AUIStructure, ∃pi ∈interactor.E f fectivePI,
∧
∃iEvt ∈interactor.events, iEvt.type = Change
∧
iEvt.property = Size
∧
 ∃p ∈interactor.type.IntrinsicPI,
p.name =′ WidgetResize′

⇒pi.name =′ WidgetResize′
R\`{e}gle 13 : Widget Move
La PI Widget Move est d\'{e}finie de mani\`{e}re effective pour les interacteurs qui impl\'{e}mentent les
5.3. MODÈLES DE COMPOSANTS GRAPHIQUES
87
comportements de changement de position.
∀interactor ∈AUIStructure,
∃pi ∈interactor.E f fectivePI,
∧
∃iEvt ∈interactor.events,
iEvt.type = Change
∧
iEvt.property = Position
∧
 ∃p ∈interactor.type.IntrinsicPI,
p.name =′ WidgetMove′

⇒pi.name =′ WidgetMove′
R\`{e}gle 14 : Widget Rotation
La PI Widget Rotation est d\'{e}finie de mani\`{e}re effective pour les interacteurs qui impl\'{e}mentent les
comportements de changement d’orientation.
∀interactor ∈AUIStructure,
∃pi ∈interactor.E f fectivePI,
∧
∃iEvt ∈interactor.events,
iEvt.type = Change
∧
iEvt.property = Orientation
∧
 ∃p ∈interactor.type.IntrinsicPI,
p.name =′ WidgetRotation′

⇒pi.name =′ WidgetRotation′
R\`{e}gle 15 : Widget Display
La PI Widget Display est d\'{e}finie de mani\`{e}re effective pour tous les interacteurs qui impl\'{e}mentent
88
CHAPITRE 5. MODÉLISATION DES UI
les PI effectives Widget Selection ou Navigation et pour ceux qui sont visibles.
∀interactor ∈AUIStructure, ∃pi ∈interactor.E f fectivePI,
∧
∃iEvt ∈interactor.events, iEvt.type = Select
∧
iEvt.property = WidgetStructure
∨
iEvt.type = Change∧iEvt.property = Visible
∧
 ∃p ∈interactor.type.IntrinsicPI,
p.name =′ WidgetDisplay′

⇒pi.name =′ WidgetDisplay′
R\`{e}gle 16 : Data Edition
La PI Data Edition est d\'{e}finie de mani\`{e}re effective pour tous les interacteurs qui ont des contenus
et qui d\'{e}crivent un comportement de changement de contenus de mani\`{e}re effective.
∀interactor ∈AUIStructure, ∃pi ∈interactor.E f fectivePI,






∃iEvt ∈interactor.events, iEvt.type = Change
∧
iEvt.property = Content
∧
interactor.ContentData! = Null






⇒pi.name =′ DataEdition′
R\`{e}gle 17 : Data Selection
La PI Data Selection est d\'{e}finie de mani\`{e}re effective pour tous les interacteurs qui ont un contenu
et qui impl\'{e}mentent un comportement de s\'{e}lection de ce contenu de mani\`{e}re effective.
∀interactor ∈AUIStructure, ∃pi ∈interactor.E f fectivePI,






∃iEvt ∈interactor.events, iEvt.type = Select
∧
iEvt.property = Content
∧
interactor.ContentData! = Null






⇒pi.name =′ DataSelection′
R\`{e}gle 18 : Data Move In
La PI Data Move In est d\'{e}finie de mani\`{e}re effective pour tous les interacteurs qui impl\'{e}mentent
des comportements de s\'{e}lection et de modification des contenus.
5.3. MODÈLES DE COMPOSANTS GRAPHIQUES
89
∀interactor ∈AUIStructure,
∃pi ∈interactor.E f fectivePI,














∃iEvt ∈interactor.events,
iEvt.type = Select
∧
∃iEvt ∈interactor.events,
iEvt.type = Change
∧
iEvt.property = Content
∧
interactor.ContentData! = Null














⇒pi.name =′ DataMoveIn′
R\`{e}gle 19 : Data Move Out
Les PI Data Move Out est d\'{e}finie de mani\`{e}re effective pour tous les interacteurs qui impl\'{e}mentent
des comportements de s\'{e}lection et de suppression des contenus.
∀interactor ∈AUIStructure,
∃pi ∈interactor.E f fectivePI,














∃iEvt ∈interactor.events,
iEvt.type = Select
∧
∃iEvt ∈interactor.events,
iEvt.type = Delete
∧
iEvt.property = Content
∧
interactor.ContentData! = Null














⇒pi.name =′ DataMoveOut′
R\`{e}gle 20 : Data Display
La PI Data Display est d\'{e}finie de mani\`{e}re effective pour tous les interacteurs qui ont un contenu
ou dont la propri\'{e}t\'{e} de contenu est modifi\'{e}e par une m\'{e}thode.
∀interactor ∈AUIStructure, ∃pi ∈interactor.E f fectivePI,
 ∃iEvt ∈interactor.events,
iEvt.type = U pdate

⇒pi.name =′ DataDisplay′
R\`{e}gle 21 : Activation
La PI Activation est d\'{e}finie de mani\`{e}re effective pour tous les interacteurs qui font appel \`{a} une
m\'{e}thode du contr\^{o}leur (dans une architecture MVC).
∀interactor ∈AUIStructure, ∃pi ∈interactor.E f fectivePI,
 ∃iEvt ∈interactor.events,
iEvt.type = Call

⇒pi.name =′ Activation′
90
CHAPITRE 5. MODÉLISATION DES UI
Remarque
Nous avons valid\'{e} les r\`{e}gles d’identification propos\'{e}es dans cette section pour les UI
d\'{e}crites \`{a} l’aide des biblioth\`{e}ques graphiques XAML et XAML Surface (cf. l’annexe A.3.). Pour la
validation, nous avons d\'{e}crit le mod\`{e}le d’instance et le mod\`{e}le de types comme des m\'{e}ta mod\`{e}les
avec EMF [Fou13] et nous avons impl\'{e}ment\'{e} les diff\'{e}rentes r\`{e}gles en Java.
5.3.2.7
R\'{e}sum\'{e}
Une UI est exprim\'{e}e au niveau abstrait par un mod\`{e}le d\'{e}crivant les primitives d’interactions effectives, les liens de contenances et les donn\'{e}es des UI \`{a} migrer.
Le mod\`{e}le de structure d\'{e}crit dans cette section permet d’exprimer, sous la forme d’un arbre,
l’instance de l’UI ind\'{e}pendamment des biblioth\`{e}ques graphiques et des dispositifs d’interactions de
d\'{e}part pour la transformer par rapport aux guidelines des tables interactives.
5.3.3
Synth\`{e}se des mod\`{e}les abstraits
Dans cette section nous avons pr\'{e}sent\'{e} un mod\`{e}le de composants graphiques et un mod\`{e}le de
structure pour une UI. Le mod\`{e}le de composants graphiques d\'{e}crit de mani\`{e}re abstraite les primitives
d’interactions intrins\`{e}ques et la structure (les donn\'{e}es d’application et les donn\'{e}es graphiques et la
cardinalit\'{e}) des Widgets ind\'{e}pendants des biblioth\`{e}ques graphiques. En effet, le mod\`{e}le de la figure 5.4
d\'{e}crit le contenu, la cardinalit\'{e}, les comportements sur les propri\'{e}t\'{e}s (donn\'{e}es, taille, position, orientation, etc.) et les interactions d’un composant graphique ind\'{e}pendamment de sa repr\'{e}sentation par
une bo\^{i}te \`{a} outils donn\'{e}e.
FIGURE 5.7 – Mod\`{e}le abstrait et UI finale
Par ailleurs le mod\`{e}le de structure (cf. figure 5.5) repr\'{e}sente les liens de contenance et les donn\'{e}es
d’une UI \`{a} migrer. Ce mod\`{e}le est certes compos\'{e} d’instances de composants graphiques mais il ne
d\'{e}crit que les aspects structurels et les interactions d’une UI \`{a} migrer.
Les mod\`{e}les de structure et d’interactions propos\'{e}s dans ce chapitre permettent de d\'{e}crire les
types et les instances des UI. La figure 5.7 par exemple repr\'{e}sente le type et l’instance d’une ListBox.
5.4. OPÉRATEURS D’ÉQUIVALENCES
91
Dans le cadre le la migration, les PI effectives des instances permettent d’\'{e}tablir des \'{e}quivalences avec
les \'{e}l\'{e}ments de la plateforme cible.
Les sections 5.2 et 5.3 d\'{e}crivent les aspects (d’interactions et de structures) des composants graphiques que nous prenons en compte dans notre approche de migration des UI semi automatique.
Ces mod\`{e}les permettent de d\'{e}crire des m\'{e}canismes d’\'{e}quivalences des instruments d’interactions qui
passent par l’\'{e}quivalence des composants graphiques bas\'{e}e sur les PI. Dans la section suivante nous
pr\'{e}sentons les op\'{e}rateurs d’\'{e}quivalences des composants graphiques bas\'{e}s sur les primitives d’interactions dans le but de d\'{e}crire un m\'{e}canisme d’\'{e}quivalences des instruments d’interactions.
5.4
Op\'{e}rateurs d’\'{e}quivalences des composants graphiques bas\'{e}s sur
des PI
L’objectif des primitives d’interactions identifi\'{e}es \`{a} la section 5.2 est de caract\'{e}riser les composants graphiques ind\'{e}pendamment des dispositifs physiques et des biblioth\`{e}ques graphiques. À la
section 5.3, nous avons propos\'{e} une mod\'{e}lisation des composants graphiques (types et instances). Ces
mod\`{e}les nous permettent de caract\'{e}riser les composants graphiques en nous basant sur leur contenu
(type de donn\'{e}es et cardinalit\'{e}) et leurs primitives d’interactions (effectives et intrins\`{e}ques). L’utilisation des op\'{e}rateurs par les m\'{e}canismes d’\'{e}quivalences a pour objectif d’\'{e}tablir des correspondances
\`{a} la vol\'{e}e entre les composants graphiques de la source et de la cible. Ce type d’\'{e}quivalences accro\^{i}t
la flexibilit\'{e} du processus de migration des UI car les \'{e}quivalences ne sont pas d\'{e}finies de mani\`{e}re
statique dans une table d’\'{e}quivalences.
Les mod\`{e}les de type et d’instance des composants graphiques d\'{e}crits \`{a} la section 5.3 caract\'{e}risent
la structure et les interactions. Les op\'{e}rateurs que nous proposons dans cette section se basent sur les
PI (effectives et intrins\`{e}ques), les types et les cardinalit\'{e}s des donn\'{e}es.
La section 5.4.1 pr\'{e}sente les op\'{e}rateurs d’\'{e}quivalences bas\'{e}s uniquement sur les PI et la section 5.4.2 pr\'{e}sente ceux qui sont bas\'{e}s sur PI et les types de donn\'{e}es.
5.4.1
Op\'{e}rateurs d’\'{e}quivalences bas\'{e}s sur des PI
Le premier enjeu pendant la migration des UI est d’avoir les interactions n\'{e}cessaires pour l’utilisation de l’UI source sur la cible. Pour \'{e}valuer la s\'{e}lection d’un composant graphique de la cible
pendant le processus de migration, nous avons identifi\'{e} trois op\'{e}rateurs d’\'{e}quivalences qui se basent
sur les primitives d’interactions des composants graphiques \`{a} migrer.
Nous consid\'{e}rons que deux PI sont \'{e}gales si elles ont le m\^{e}me nom.
∀pi1, pi2 ∈\textbackslash{}{PirmitiveInteractions\textbackslash{}},
pi1 = pi2 ⇒pi1.name = pi2.name
5.4.1.1
Équivalence stricte : ≡
Cet op\'{e}rateur permet de retrouver sur la cible les composants graphiques qui ont les m\^{e}mes interactions que la source. Cet op\'{e}rateur s’assure aussi que les op\'{e}randes ont les m\^{e}mes types de donn\'{e}es
et les m\^{e}mes cardinalit\'{e}s.
∀int ∈instanceO f(Interactor),∀wid ∈instanceO f(Widget),
92
CHAPITRE 5. MODÉLISATION DES UI
int ≡wid ⇒



























int.contentData.cardinality = wid.cardinality
((int.contentData! = Null∧
int.contentData.dataType = wid.ContentType)
∨(int.contentData = Null ∧wid.contentType = Null))
∀pi1 ∈int.e f fectivePI,
∃pi2 ∈wid.IntrinsicPI,
(pi1 = pi2
∧
card(int.e f fectivePI) = card(wid.IntrinsicPI))
Exemple d’\'{e}quivalence stricte
Pour une UI source d\'{e}crite en JavaSwing qui instancie un JButton
que l’on souhaite migrer vers une table Microsoft PixelSense, button1 est une instance de JButton et
sbutton est une instance de SurfaceButton. Nous supposons que :
– int est une instance de Interactor qui repr\'{e}sente button1 au niveau mod\`{e}le
– wid est une instance de Widget qui repr\'{e}sente sbutton au niveau mod\`{e}le
int ⇔











id = 1
name = button1
instrinsicPI = \textbackslash{}{Navigation,WidgetSelection,Activation\textbackslash{}}
contentData = \textbackslash{}{dataType = String,cardinality = 1,
value = ‘Apply′, propertyName = ‘Text′\textbackslash{}}
wid ⇔



name = sbutton
e f fectivePI = \textbackslash{}{Navigation,WidgetSelection,Activation\textbackslash{}}
contentType = String,cardinality = 1
⇒int ≡wid
5.4.1.2
Équivalence large : ≦AdditionalPI
Cet op\'{e}rateur permet de retrouver les composants graphiques de la cible qui ont toutes les
primitives d’interactions du composant graphique de la source et qui proposent d’autres interactions additionnelles. Les primitives d’interactions suppl\'{e}mentaires 16 sont pr\'{e}cis\'{e}es dans l’ensemble
AdditionalPI. Cet op\'{e}rateur permet par exemple de choisir des composants graphiques qui permettent
de mieux respecter des guidelines (que l’\'{e}quivalence stricte) sur la cible tout en conservant l’ensemble
des interactions de la source.
∀int ∈instanceO f(Interactor),∀wid ∈instanceO f(Widget),
16. Ces PI sont pr\'{e}conis\'{e}es par les guidelines de migration des UI
5.4. OPÉRATEURS D’ÉQUIVALENCES
93
int ≦AdditionalPI wid ⇒











































int.contentData.cardinality = wid.cardinality
((int.contentData! = Null
∧
int.contentData.dataType = wid.ContentType)
∨
(int.contentData = Null ∧wid.contentType = Null))
∀pi1 ∈int.e f fectivePI,
∃pi2 ∈wid.IntrinsicPI,
pi1 = pi2
∧
card(int.e f fectivePI) < card(wid.IntrinsicPI)
∧
∃pi ∈wid.IntrinsicPI ∖int.e f fectivePI ∧pi ∈AdditionalPI
Exemple d’\'{e}quivalence large
Pour une UI source d\'{e}crite en JavaSwing, les instances des JFrame
sont migr\'{e}es vers une table DiamondTouch en y ajoutant l’interaction de rotation pour \^{e}tre conforme
\`{a} la guideline du partage de l’espace de travail, frame1 est une instance de JFrame et dsFrame est une
instance de DSFrame. Nous supposons que :
– int est une instance de Interactor qui repr\'{e}sente frame1 au niveau mod\`{e}le
– wid est une instance de Widget qui repr\'{e}sente dsFrame au niveau mod\`{e}le
– AdditionalPI = \textbackslash{}{WidgetRotation,WidgetResize\textbackslash{}}
int ⇔















id = 2
name = frame1
instrinsicPI = \textbackslash{}{Navigation,WidgetSelection,
WidgetDisplay,WidgetMove,
WidgetResize,Activation\textbackslash{}}
contentData = Null
wid ⇔











name = dsFrame
e f fectivePI = \textbackslash{}{Navigation,WidgetSelection,
WidgetDisplay,WidgetMove,
WidgetResize,Activation,
WidgetRotation\textbackslash{}}
⇒int ≦AdditionalPI wid
5.4.1.3
Équivalence faible : ≧EssentialPI
Cet op\'{e}rateur d’\'{e}quivalence recherche les composants graphiques de la cible qui ont un minimum
d’interactions essentielles pour permettre l’utilisation de l’UI sur la cible. Cet op\'{e}rateur permet une
migration en mode ”d\'{e}grad\'{e}” dans le cas o\`{u} la cible ne peut pas fournir certaines interactions. Les
primitives d’interactions essentielles \`{a} conserver sont identifi\'{e}es dans l’ensemble EssentialPI. Cet
ensemble de PI est d\'{e}finit en fonction de la plateforme d’arriv\'{e}e par les personnes en charge de la
migration. Dans le cadre des tables interactives, nous avons identifi\'{e} l’ensemble des PI essentielles en
prenant en compte les guidelines qui pr\'{e}servent l’homog\'{e}n\'{e}it\'{e} des interactions sources et cibles.
∀int ∈instanceO f(Interactor),∀wid ∈instanceO f(Widget),
94
CHAPITRE 5. MODÉLISATION DES UI
int ≦EssentialPI wid ⇒











































int.contentData.cardinality = wid.cardinality
(int.contentData! = Null
∧
int.contentData.dataType = wid.ContentType
∨
int.contentData = Null ∧wid.contentType = Null)
∀pi1 ∈int.e f fectivePI,
∃pi2 ∈wid.IntrinsicPI,
(pi1 = pi2
∧
card(int.e f fectivePI) > card(wid.IntrinsicPI)
∧
∃pi ∈int.e f fectivePI ∩wid.IntrinsicPI ∧pi ∈EssentialPI)
Exemple d’\'{e}quivalence faible
Nous illustrons cet op\'{e}rateur par un exemple de remplacement d’une
liste \'{e}ditable par une liste non \'{e}ditable. Si une UI source d\'{e}crite en XAML instancie un ComboBox
\'{e}ditable, nous pouvons la remplacer par une SurfaceListBox non \'{e}ditable sur une table Microsoft
PixelSense. Dans cet exemple, nous pensons que les PI Activation, Data Selection sont essentielles car
elles pr\'{e}servent la fonction de s\'{e}lection d’une liste m\^{e}me si elle n’est pas \'{e}ditable apr\`{e}s la migration.
La liste \'{e}ditable list1 est une instance de ComboBox et surfacelist est une instance de SurfaceListBox. Nous supposons que :
– int est une instance de Interactor qui repr\'{e}sente list1 au niveau mod\`{e}le
– wid est une instance de Widget qui repr\'{e}sente surfacelist au niveau mod\`{e}le
– EssentialPI = \textbackslash{}{Activation,DataSelection\textbackslash{}}
int ⇔























id = 4
name = list1
instrinsicPI = \textbackslash{}{Navigation,WidgetSelection,
WidgetDisplay, DataSelection,
DataMoveIn/Out,
DataEdition, Activation\textbackslash{}}
contentData = \textbackslash{}{dataType = String, cardinality = N,
value = ‘′, propertyName = ‘Content′\textbackslash{}}
wid ⇔



















name = dsFrame
e f fectivePI = \textbackslash{}{Navigation, WidgetSelection,
WidgetDisplay, DataSelection,
DataMoveIn/Out,
Activation\textbackslash{}}
contentType = String,
cardinality = N
⇒int ≧EssentialPI wid
5.4.1.4
Remarques
Les op\'{e}rateurs d’\'{e}quivalences d\'{e}crits par cette section 5.4.1 permettent d’\'{e}tablir des correspondances \`{a} la vol\'{e}e entre les composants graphiques des plateformes source et cible en s’appuyant sur
les PI, les types de donn\'{e}es et les cardinalit\'{e}s.
L’op\'{e}rateur d’\'{e}quivalence stricte s’applique dans le cas o\`{u} l’on souhaite avoir des composants
graphiques cibles identiques \`{a} la source. Cependant dans le cadre des UI pour les tables interactives,
5.4. OPÉRATEURS D’ÉQUIVALENCES
95
il est indispensable de pendre en compte des guidelines qui pr\'{e}conisent des composants graphiques
accessibles \`{a} plusieurs personnes par exemple.
L’op\'{e}rateur d’\'{e}quivalence large permet d’\'{e}tablir des correspondances entre les composants graphiques des plateformes source et cible mais en int\'{e}grant les interactions pr\'{e}conis\'{e}es par les guidelines de la cible. Par exemple pour un menu d’une UI desktop, cet op\'{e}rateur permet de retrouver un
composant graphique \'{e}quivalent et avec les interactions de d\'{e}placement pour rendre le menu accessible.
L’op\'{e}rateur d’\'{e}quivalence faible est l’inverse de l’\'{e}quivalence large et l’ensemble EssentialPI est
constitu\'{e}
– des PI li\'{e}es \`{a} la s\'{e}lection ou \`{a} la navigation (Widget Select ou Navigation)
– des PI li\'{e}es \`{a} la modification de contenus (Data Edition ou Data Selection)
– de la PI Activation.
– et des PI d’affichage de contenus Display Data
Ces PI sont fondamentales pour des interactions en entr\'{e}e et en sortie. Concernant une UI en mode
”d\'{e}grad\'{e}e” par exemple, les PI li\'{e}es \`{a} la taille, la position ou l’orientation ne sont pas essentielles.
M\^{e}me si les guidelines de la plateforme cible ne sont pas respect\'{e}es, l’UI produite conserve l’accessibilit\'{e} aux fonctionnalit\'{e}s de l’application \`{a} migrer.
Nous remarquons que ces trois op\'{e}rateurs ne prennent pas en compte les composants graphiques
\'{e}quivalents du point de vue des interactions mais avec des types de donn\'{e}es diff\'{e}rents. Pour la migration des UI vers les tables interactives, un menu contenant des sous menus textuels par exemple,
les donn\'{e}es graphiques textuelles peuvent \^{e}tre remplac\'{e}es par des ic\^{o}nes. Les op\'{e}rateurs d\'{e}crits \`{a}
section 5.4.2 r\'{e}pondent \`{a} cette limite.
5.4.2
Op\'{e}rateurs d’\'{e}quivalences bas\'{e}s sur des PI et des donn\'{e}es
Dans l’objectif de combler les limites des op\'{e}rateurs identifi\'{e}s \`{a} la section 5.4.1 nous d\'{e}finissons
dans cette section trois autres op\'{e}rateurs qui prennent en compte la diff\'{e}rence du type de donn\'{e}es entre
la source et la cible. Les trois op\'{e}rateurs d’\'{e}quivalences bas\'{e}s uniquement sur les PI sont modifi\'{e}s ici
pour tenir compte de la variation des types de donn\'{e}es.
5.4.2.1
Équivalence stricte des PI avec types de donn\'{e}es diff\'{e}rents : ∼=TargetType
Cet op\'{e}rateur permet de s\'{e}lectionner les composants graphiques qui ont les m\^{e}mes interactions
que celles de la source mais avec des types de donn\'{e}es diff\'{e}rents. L’ensemble TargetType exprime
les types de donn\'{e}es souhait\'{e}s pour un composant graphique sur la cible.
∀int ∈instanceO f(Interactor),∀wid ∈instanceO f(Widget),
int ∼=TargetType wid ⇒























int.contentData.cardinality = wid.cardinality
(int.contentData! = Null∧
wid.ContentType ∈TargetType)
∀pi1 ∈int.e f fectivePI,
∃pi2 ∈wid.IntrinsicPI,
pi1 = pi2
∧
card(int.e f fectivePI) = card(wid.IntrinsicPI)
Exemple
Dans cet exemple nous montrons que ListBox ∼=TargetType LibraryBar avec :
96
CHAPITRE 5. MODÉLISATION DES UI
– int est une instance de Interactor qui repr\'{e}sente list2 au niveau mod\`{e}le
– wid est une instance de Widget qui repr\'{e}sente library au niveau mod\`{e}le
– TargetType = \textbackslash{}{Image,Object,MediaElement\textbackslash{}}
int ⇔



















id = 5
name = list2
instrinsicPI = \textbackslash{}{Navigation, WidgetSelection,
WidgetDisplay, DataSelection, DataMoveIn/Out,
DataDisplay, Activation\textbackslash{}}
contentData = \textbackslash{}{dataType = String, cardinality = N,
value = ‘′, propertyName = ‘Content′\textbackslash{}}
wid ⇔















name = library
e f fectivePI = \textbackslash{}{Navigation, WidgetSelection,
WidgetDisplay, DataSelection, DataMoveIn/Out,
DataDisplay, Activation\textbackslash{}}
contentType = Object ∈TargetType,
cardinality = 1
⇒int ∼=TargetType wid
5.4.2.2
Équivalence large des PI avec types de donn\'{e}es diff\'{e}rents : ≲AdditionalPI,TargetType
Cet op\'{e}rateur permet de retrouver les composants graphiques offrant des interactions suppl\'{e}mentaires et comprenant un type de donn\'{e}es diff\'{e}rent de la source. C’est une am\'{e}lioration de l’op\'{e}rateur
≦AdditionalPI en prenant en compte la diff\'{e}rence entre les donn\'{e}es. L’ensemble TargetType exprime
les types de donn\'{e}es souhait\'{e}s pour un composant graphique sur la cible.
∀int ∈instanceO f(Interactor),∀wid ∈instanceO f(Widget),
int ≲AdditionalPI,TargetType wid ⇒



































int.contentData.cardinality = wid.cardinality
(int.contentData! = Null
∧
wid.ContentType ∈TargetType)
∀pi1 ∈int.e f fectivePI,
∃pi2 ∈wid.IntrinsicPI,
pi1 = pi2
∧
card(int.e f fectivePI) < card(wid.IntrinsicPI)
∧
∃pi ∈wid.IntrinsicPI ∖int.e f fectivePI ∧pi ∈AdditionalPI
Exemple
Dans cet exemple, nous montrons que Button ≲AdditionalPI,TargetType Image avec :
– int est une instance de Interactor qui repr\'{e}sente button2 au niveau mod\`{e}le
– wid est une instance de Widget qui repr\'{e}sente image au niveau mod\`{e}le
– TargetType = \textbackslash{}{Image,Object,MediaElement\textbackslash{}}
– AdditionalPI = \textbackslash{}{WidgetRotation,WidgetResize\textbackslash{}}
5.4. OPÉRATEURS D’ÉQUIVALENCES
97
int ⇔















id = 6
name = button2
instrinsicPI = \textbackslash{}{Navigation, WidgetSelection,
WidgetDisplay, Activation\textbackslash{}}
contentData = \textbackslash{}{dataType = String, cardinality = N,
value = ‘′, propertyName = ‘Content′\textbackslash{}}
wid ⇔















name = dsFrame
e f fectivePI = \textbackslash{}{Navigation, WidgetSelection,
WidgetDisplay, DataDisplay,
Activation\textbackslash{}}
contentType = Image ∈TargetType,
cardinality = 1
⇒int ≲AdditionalPI,TargetType wid
5.4.2.3
Équivalence faible des PI avec types de donn\'{e}es diff\'{e}rents : ≳EssentialPI,TargetType
Cet op\'{e}rateur permet d’\'{e}tablir des UI \'{e}quivalentes en mode ”d\'{e}grad\'{e}” et en prenant la diff\'{e}rence
de type de donn\'{e}es. C’est une am\'{e}lioration de l’op\'{e}rateur ≧EssentialPI en prenant compte de la diff\'{e}rence entre les donn\'{e}es. L’ensemble TargetType exprime les types de donn\'{e}es souhait\'{e}s pour un
composant graphique sur la cible.
∀int ∈instanceO f(Interactor),∀wid ∈instanceO f(Widget),
int ≳EssentialPI,TargetType wid ⇒



































int.contentData.cardinality = wid.cardinality
(int.contentData! = Null
∧
wid.ContentType ∈TargetType)
∀pi1 ∈int.e f fectivePI,
∃pi2 ∈wid.IntrinsicPI,
pi1 = pi2
∧
card(int.e f fectivePI) > card(wid.IntrinsicPI)
∧
∃pi ∈int.e f fectivePI ∩wid.IntrinsicPI ∧pi ∈EssentialPI
Exemple
Nous montrons que Image ≳EssentialPI,SourceType Sur faceButton avec :
– int est une instance de Interactor qui repr\'{e}sente image2 au niveau mod\`{e}le
– wid est une instance de Widget qui repr\'{e}sente button3 au niveau mod\`{e}le
– SourceType = \textbackslash{}{Image, Object, MediaElement\textbackslash{}}
– EssentialPI = \textbackslash{}{Activation, WidgetResize\textbackslash{}}
int ⇔



















id = 3
name = image2
instrinsicPI = \textbackslash{}{Navigation, WidgetSelection,
WidgetDisplay,DataDisplay,
Activation\textbackslash{}}
contentData = \textbackslash{}{dataType = Image ∈SourceType,cardinality = N,
value = ‘′, propertyName = ‘Content′\textbackslash{}}
98
CHAPITRE 5. MODÉLISATION DES UI
wid ⇔











name = button3
e f fectivePI = \textbackslash{}{Navigation, WidgetSelection,
WidgetDisplay, Activation\textbackslash{}}
contentType = String,
cardinality = 1
⇒int ≳EssentialPI,TargetType wid
Remarque
Les op\'{e}rateurs d’\'{e}quivalences bas\'{e}s sur les PI et prenant en compte des types de donn\'{e}es
diff\'{e}rents entre la source et la cible, permettent de s\'{e}lectionner les composants graphiques \'{e}quivalents
du point de vue des PI. Les changements des types de donn\'{e}es impliquent des conversions des donn\'{e}es
sources vers la cible. Les conversions peuvent \^{e}tre automatis\'{e}es s’il existe des relations entre les types
de donn\'{e}es sources et cibles (par exemple les entiers peuvent \^{e}tre transform\'{e}s en cha\^{i}ne de caract\`{e}res).
Dans d’autres cas, les changements de types n\'{e}cessitent une intervention humaine pour \'{e}tablir un
mapping entre les donn\'{e}es de la source avec celles de la cible (par exemple la conversion de cha\^{i}ne de
caract\`{e}res en image n\'{e}cessite le choix des images repr\'{e}sentant les diff\'{e}rentes cha\^{i}nes de caract\`{e}res).
5.5
Synth\`{e}se
Les PI pr\'{e}sent\'{e}es constituent un mod\`{e}le d’interactions abstraites. L’objectif de ce mod\`{e}le est
d’\'{e}tablir les \'{e}quivalences entre instruments d’interactions en prenant en compte les sp\'{e}cificit\'{e}s de la
cible 17. Les PI intrins\`{e}ques des Widgets, les PI effectives des Interactors et les \'{e}l\'{e}ments structurels
offrent des \'{e}l\'{e}ments de comparaison entre les composants graphiques de mani\`{e}re g\'{e}n\'{e}rique.
Nous avons propos\'{e} dans ce chapitre une repr\'{e}sentation abstraite des UI \`{a} migrer en prenant en
compte les interactions (\`{a} travers les PI) et la structure (\`{a} travers les mod\`{e}les de type et d’instance
d’une UI). En faisant ces choix, nous avions les objectifs suivant :
– Identifier dans un mod\`{e}le abstrait les interactions en entr\'{e}e et en sortie des UI graphiques en
prenant en compte les diff\'{e}rences entre les types et les instances des composants graphiques
qui sont importantes dans un processus de migration.
– Montrer qu’il est possible d’identifier les \'{e}l\'{e}ments concrets (menu, fen\^{e}tre, tableau, boutons,
etc.) d’une UI en se basant sur un mod\`{e}le de structure minimal et les PI.
Nous avons montr\'{e} aussi qu’en se basant sur les primitives d’interactions et la structure d’un
composant graphique, il est possible de d\'{e}crire des op\'{e}rateurs d’\'{e}quivalences. Cependant nous remarquons que ces op\'{e}rateurs peuvent produire plus d’un composant graphique \'{e}quivalent. La question du
choix du meilleur se pose pour les concepteurs en charge de la partie manuelle de la migration . Leurs
choix en g\'{e}n\'{e}ral sont guid\'{e}s par le soucis du respect des crit\`{e}res ergonomiques et la baisse du co\^{u}t
de la mise en œuvre. Le chapitre 6 suivant pr\'{e}sente les m\'{e}canismes de classement des composants
graphiques \'{e}quivalents en prenant en compte les guidelines de la cible.
Au-del\`{a} de ces questions de s\'{e}lection et de respect des crit\`{e}res ergonomiques, le mod\`{e}le d’une UI
que nous avons pr\'{e}sent\'{e} dans ce chapitre a pour objectif de d\'{e}crire des m\'{e}canismes d’adaptations des
interactions et de la structure de l’UI qui prennent en compte les guidelines de la cible. Le chapitre
r\'{e}pond aux questions soulev\'{e}es dans cette synth\`{e}se. Nous proposons aussi un processus de migration
qui se base sur ce mod\`{e}le d’interactions abstraites et sur un mod\`{e}le de structure minimal (contenu,
type de donn\'{e}es, cardinalit\'{e} de donn\'{e}es et lien de contenance) et qui prend en compte les guidelines.
17. Dans notre cas ces sp\'{e}cificit\'{e}s sont caract\'{e}ris\'{e}es par les guidelines pour la migration vers les tables interactives
CHAPITRE 6
M\'{e}canismes de migration des UI vers les
tables interactives
Sommaire
6.1
Introduction . . . . . . . . . . . . . . . . . . . . . . . . . . . . . . . . . . . . .
99
6.2
Prise en compte des guidelines . . . . . . . . . . . . . . . . . . . . . . . . . . . 100
6.2.1
Conception assist\'{e}e des UI . . . . . . . . . . . . . . . . . . . . . . . . . . 101
6.2.2
Interpr\'{e}tation des guidelines pour la migration
. . . . . . . . . . . . . . . 102
6.2.3
Utilisation effective des guidelines . . . . . . . . . . . . . . . . . . . . . . 104
6.3
Transformations du mod\`{e}le de l’UI source
. . . . . . . . . . . . . . . . . . . . 108
6.3.1
Transformation des groupes d’\'{e}l\'{e}ments graphiques . . . . . . . . . . . . . 108
6.3.2
Transformation d’un \'{e}l\'{e}ment . . . . . . . . . . . . . . . . . . . . . . . . . 119
6.4
Classement des \'{e}l\'{e}ments \'{e}quivalents . . . . . . . . . . . . . . . . . . . . . . . . 124
6.4.1
Conformit\'{e} des composants graphiques aux guidelines . . . . . . . . . . . 124
6.4.2
Charge de travail . . . . . . . . . . . . . . . . . . . . . . . . . . . . . . . 127
6.4.3
Algorithme de classement . . . . . . . . . . . . . . . . . . . . . . . . . . 129
6.5
Synth\`{e}se
. . . . . . . . . . . . . . . . . . . . . . . . . . . . . . . . . . . . . . . 133
6.1
Introduction
Le mod\`{e}le de l’UI pr\'{e}sent\'{e} au chapitre 5 permet de d\'{e}crire les UI \`{a} migrer ind\'{e}pendamment des
biblioth\`{e}ques graphiques et des dispositifs d’interactions des plateformes source et cible. Ce mod\`{e}le
d\'{e}crit particuli\`{e}rement la structure et les interactions abstraites des UI \`{a} migrer. Dans ce chapitre nous
pr\'{e}sentons le processus de migration des UI vers les tables interactives en d\'{e}crivant les m\'{e}canismes
de prise en compte des guidelines et les transformations de l’UI source en UI collaborative et tangible.
Ces m\'{e}canismes sont guid\'{e}s par un ensemble de guidelines pour la migration des UI vers les tables
interactives identifi\'{e}es \`{a} la section 3.2.3. Les guidelines sont des recommandations d\'{e}crites en langage
naturel qui doivent \^{e}tre interpr\'{e}t\'{e}es en r\`{e}gles de migrations des \'{e}l\'{e}ments concrets (tels que fen\^{e}tres,
pop-ups, menus, tableaux, formulaires, boutons, champs de texte, etc.) d’une UI.
Nous montrons dans ce chapitre qu’en utilisant notre mod\`{e}le d’UI, il est possible de transformer la
structure et les interactions des UI desktops pour les tables interactives et en respectant les guidelines
des UI collaboratives et tangibles.
Pour atteindre cet objectif, nous pr\'{e}sentons \`{a} la section 6.2 notre processus de migration des UI
ainsi que les interpr\'{e}tations des guidelines identifi\'{e}es au chapitre 3 pour d\'{e}crire les m\'{e}canismes de
transformation d’un mod\`{e}le d’UI source. La section 6.3 pr\'{e}sente les m\'{e}canismes de transformations
de la structure et les interactions des UI \`{a} migrer. La section 6.4 pr\'{e}sente le m\'{e}canisme de classement
des \'{e}l\'{e}ments graphiques \'{e}quivalents en prenant en compte les guidelines pour la migration. Le classement des \'{e}l\'{e}ments \'{e}quivalents en fonction des guidelines facilite le choix des composants graphiques
\`{a} utiliser pour d’\'{e}crire l’UI cible. Le chapitre se termine par une synth\`{e}se qui pr\'{e}sente les points forts
99
100
CHAPITRE 6. MÉCANISMES DE MIGRATION DES UI
et les limites de notre approche de transformation d’une UI source en UI collaborative et tangible pour
les tables interactives.
6.2
Prise en compte des guidelines dans le processus de migration des
UI vers les tables interactives
Nous proposons un processus semi automatique de migration des UI vers les tables interactives.
Nous optons pour une approche semi automatique \`{a} la section 4.3.2 afin de b\'{e}n\'{e}ficier de la flexibilit\'{e}, de la r\'{e}utilisabilit\'{e} et du respect des crit\`{e}res ergonomiques qu’elle offre. Notre processus se
d\'{e}compose en trois \'{e}tapes (cf. figure 6.1).
– La premi\`{e}re \'{e}tape consiste \`{a} abstraire de mani\`{e}re automatique la structure et les interactions de
l’UI source dans les mod\`{e}les pr\'{e}sent\'{e}s au chapitre 5. Nous consid\'{e}rons que chaque application
de d\'{e}part respecte les contraintes d’architecture pr\'{e}sent\'{e}es \`{a} la section 2.3.
– La deuxi\`{e}me \'{e}tape consiste \`{a} restructurer le mod\`{e}le abstrait de l’UI en s’appuyant sur les
guidelines et sur les op\'{e}rateurs d’\'{e}quivalences. La restructuration est automatique dans un premier temps puis manuelle pour permettre \`{a} l’utilisateur de personnaliser l’UI cible apr\`{e}s avoir
concr\'{e}tiser le mod\`{e}le restructur\'{e}. La restructuration automatique consiste \`{a} retrouver les interacteurs \'{e}quivalents pour la cible et \`{a} substituer les interacteurs de la source par ceux de la
cible dans le but de proposer une UI concr\`{e}te. La restructuration manuelle permet \`{a} la personne en charge de la migration de modifier l’UI propos\'{e}e automatiquement et de r\'{e}introduire
le positionnement et le style dans l’UI migr\'{e}e.
– La troisi\`{e}me \'{e}tape consiste \`{a} g\'{e}n\'{e}rer automatiquement l’application pour la table interactive
en consid\'{e}rant l’UI transform\'{e}e et le NF de l’application source. Cette \'{e}tape consiste d’abord \`{a}
g\'{e}n\'{e}rer l’UI finale pour la plateforme cible qui d\'{e}crit les interactions, la structure, le positionnement et le style. Ensuite cette \'{e}tape consiste \`{a} g\'{e}n\'{e}rer le lien entre l’UI et le NF en tenant
compte du mod\`{e}le d’architecture de l’application d\'{e}part.
FIGURE 6.1 – Processus de migration assist\'{e}e
La prise en compte des guidelines par les m\'{e}canismes automatiques et manuels de la restructuration a pour but de respecter les crit\`{e}res ergonomiques de conception et les sp\'{e}cificit\'{e}s des tables
interactives. Pendant la restructuration, les guidelines doivent \^{e}tre traduites en recommandations ou
en contraintes de migration des UI dans le but de favoriser la collaboration et l’utilisation des objets
tangibles.
Une approche de prise en compte des guidelines dans un processus de conception des UI est
propos\'{e}e par Vanderdonckt [Van97] dans le but de minimiser les co\^{u}ts de d\'{e}veloppement des UI.
Nous \'{e}tudions \`{a} la section 6.2.1 cette approche afin de d\'{e}crire les guidelines de migration des UI
utilisables pendant la restructuration et la concr\'{e}tisation. La section 6.2.2 est une interpr\'{e}tation des
6.2. PRISE EN COMPTE DES GUIDELINES
101
guidelines \`{a} partir des crit\`{e}res ergonomiques de conception et la section 6.2.3 pr\'{e}sente les m\'{e}canismes
d’utilisation effective des guidelines interpr\'{e}t\'{e}es.
6.2.1
Conception assist\'{e}e des UI
La conception assist\'{e}e comporte quatre \'{e}tapes pour transformer des t\^{a}ches interactives en UI
finale tout en prenant en compte les guidelines de la plateforme cible.
1. La premi\`{e}re \'{e}tape consiste \`{a} d\'{e}duire des facteurs d’utilit\'{e} et d’utilisabilit\'{e} que l’UI \`{a} produire doit satisfaire, \`{a} partir d’une description des t\^{a}ches interactives, des utilisateurs et de
la plateforme. L’utilit\'{e} concerne l’ad\'{e}quation qui doit exister entre les fonctions s\'{e}mantiques
pr\'{e}sent\'{e}es par une IHM et les actions n\'{e}cessaires pour l’utilisateur en vue d’accomplir sa
t\^{a}che [Van97]. L’utilisabilit\'{e} concerne l’ad\'{e}quation entre la mani\`{e}re dont une t\^{a}che interactive est accomplie par un utilisateur particulier dans un contexte donn\'{e} et le profil cognitif de
cet utilisateur [FPV95]. Les facteurs d’“utilit\'{e}” et d’“utilisabilit\'{e}” concr\'{e}tisent les facteurs critiques de succ\`{e}s suivant les deux perspectives d’utilit\'{e} et d’utilisabilit\'{e}. Dans le cadre de la
migration des UI pour les tables interactives, nous avons identifi\'{e} \`{a} la section 3.2 les propri\'{e}t\'{e}s
qui favorisent la mise en place des UI collaboratives et tangibles. Nous consid\'{e}rons que les
propri\'{e}t\'{e}s de tangibilit\'{e} et tactibilit\'{e} des interactions, de la taille et la disposition de la surface d’affichage et celles du nombre et la r\'{e}partition des utilisateurs constituent des facteurs
critiques qui influencent la migration des UI vers les tables interactives.
2. La deuxi\`{e}me \'{e}tape consiste \`{a} faire ressortir \`{a} partir des facteurs identifi\'{e}s pr\'{e}c\'{e}demment une
pond\'{e}ration des crit\`{e}res ergonomiques de conception 18 pour les t\^{a}ches interactives [Van97].
Dans le cadre de la migration des UI, cette \'{e}tape consiste \`{a} consid\'{e}rer le tableau 3.3 de raffinement des crit\`{e}res ergonomiques de conception par rapport aux propri\'{e}t\'{e}s caract\'{e}ristiques (qui
sont aussi consid\'{e}r\'{e}es comme facteur critiques dans notre contexte). Ce tableau nous a permis
\`{a} la section 3.2.3 d’identifier les guidelines pour la migration des UI vers les tables interactives.
3. La troisi\`{e}me \'{e}tape consiste \`{a} d\'{e}duire l’ensemble des guidelines non conflictuelles. Dans le cadre
de la conception des UI, les guidelines constituent des recommandations pour d\'{e}crire la structure, les interactions, le layout et le style de l’UI finale. Dans le cadre de la migration des UI, les
guidelines permettent de transformer automatiquement les interactions et la structure de l’UI
source. Nous avons identifi\'{e} \`{a} la section 3.2.3 les guidelines pour favoriser les UI collaboratives
et les UI tangibles.
4. La quatri\`{e}me \'{e}tape consiste \`{a} g\'{e}n\'{e}rer automatiquement les objets interactifs concrets de l’UI
en se basant sur un mod\`{e}le des UI et en s’appuyant aussi sur les guidelines identifi\'{e}es. Dans
le cadre de notre processus de migration des UI, cette phase consiste \`{a} transformer selon les
guidelines de migration le mod\`{e}le de l’UI source pr\'{e}alablement extrait.
Consid\'{e}rons l’hypoth\`{e}se de la section 2.3 qui pr\'{e}conise la non modification du NF de l’application
de l’UI source. Et consid\'{e}rons que notre mod\`{e}le des UI d\'{e}crit les interactions et la structure de l’UI
source (cf chapitre 5). Alors les deux premi\`{e}res \'{e}tapes du processus de Vanderdonckt qui consistent
\`{a} adapter les t\^{a}ches par rapport \`{a} la plateforme cible ne peuvent pas \^{e}tre consid\'{e}r\'{e}es dans le cadre de
la migration assist\'{e}e. En effet, la transformation des activit\'{e}s d’une UI mono-utilisateur pour prendre
en compte des t\^{a}ches collaboratives n\'{e}cessite une modification du NF.
Cependant les interactions et la structure de l’UI source peuvent \^{e}tre adapt\'{e}es pour favoriser la
collaboration et permettre l’utilisation des objets tangibles si les recommandations des guidelines sont
d\'{e}crites de mani\`{e}re non ambigu\"{e}.
18. Nous consid\'{e}rons les huit crit\`{e}res ergonomiques de conception propos\'{e}s par Scapin [Sca86] (cf section 3.2.2).
102
CHAPITRE 6. MÉCANISMES DE MIGRATION DES UI
6.2.2
Interpr\'{e}tation des guidelines pour la migration
Les guidelines peuvent \^{e}tre comprises et interpr\'{e}t\'{e}es de diff\'{e}rentes mani\`{e}res par les concepteurs
des UI car elles sont d\'{e}crites par un langage naturel. Dans le cadre de la migration des UI vers les
tables interactives et en consid\'{e}rant les deux groupes de guidelines identifi\'{e}s \`{a} la section 3.2.3, nous
interpr\'{e}tons les guidelines par des actions sur une repr\'{e}sentation abstraite des UI \`{a} migrer. Les interpr\'{e}tations propos\'{e}es dans cette section ont pour objectifs de mettre en place des r\`{e}gles de transformations r\'{e}utilisables pour diff\'{e}rents types d’UI.
6.2.2.1
Guidelines pour favoriser la collaboration
Dans le but de d\'{e}crire les guidelines identifi\'{e}es \`{a} la section 3.2.3 \`{a} partir des propri\'{e}t\'{e}s caract\'{e}ristiques et des crit\`{e}res ergonomiques de conception, nous pr\'{e}sentons dans cette section une interpr\'{e}tation des guidelines qui favorise la collaboration. Les propri\'{e}t\'{e}s caract\'{e}ristiques de taille, de
disposition de la surface d’affichage, du nombre et de la r\'{e}partition des utilisateurs permettent d’affiner
les guidelines qui favorisent la collaboration. Le tableau 3.3 d’affinement des crit\`{e}res ergonomiques
de conception nous permet d’identifier deux types de guidelines.
FIGURE 6.2 – Crit\`{e}res ergonomiques de conception et guidelines pour favoriser la collaboration
– Les guidelines d\'{e}duites des propri\'{e}t\'{e}s qui favorisent des crit\`{e}res ergonomiques ; ces guidelines
constituent des recommandations \`{a} appliquer sur les PI et la structure de l’UI.
6.2. PRISE EN COMPTE DES GUIDELINES
103
– Les guidelines d\'{e}duites \`{a} partir des propri\'{e}t\'{e}s qui risquent d’alt\'{e}rer des crit\`{e}res ergonomiques ; ces guidelines comportent des contraintes \`{a} respecter pendant la migration.
Le diagramme de la figure 6.2 est un arbre qui regroupe les crit\`{e}res ergonomiques de conception et les guidelines. Les nœuds oranges sont les crit\`{e}res ergonomiques qui risquent d’\^{e}tre alt\'{e}r\'{e}s si
les contraintes (repr\'{e}sent\'{e}es par les feuilles rectangulaires) ne sont pas respect\'{e}es pendant la migration. Les nœuds gris sont des crit\`{e}res ergonomiques de conception favoris\'{e}s si les recommandations
(repr\'{e}sent\'{e}es par les feuilles rectangulaires) sont respect\'{e}es.
6.2.2.2
Guidelines pour des UI tangibles
La propri\'{e}t\'{e} caract\'{e}ristique de tangibilit\'{e} permet d’affiner les guidelines pour l’utilisation des
objets tangibles. La figure 6.3 pr\'{e}sente les recommandations et les contraintes pour chaque crit\`{e}re
ergonomique.
FIGURE 6.3 – Crit\`{e}res ergonomiques de conception et guidelines pour des UI tangibles
6.2.2.3
Synth\`{e}se
Les recommandations et les contraintes permettront de transformer la structure et les interactions
du mod\`{e}le de l’UI source.
Les guidelines sont aussi \`{a} affiner sous la forme des contraintes sur les types de donn\'{e}es, la taille
des \'{e}l\'{e}ments graphiques, le comportement des groupes d’\'{e}l\'{e}ments graphiques, etc.
Les diagrammes des figures 6.2 et 6.3 sont des interpr\'{e}tations des guidelines \`{a} l’aide des crit\`{e}res
ergonomiques de conception qui constitueront des r\`{e}gles de transformation de l’UI source.
L’interpr\'{e}tation des guidelines est effectu\'{e}e avant l’impl\'{e}mentation de notre processus de migration des UI. Elle n\'{e}cessite un affinement des crit\`{e}res ergonomiques de conception \`{a} l’aide des
sp\'{e}cificit\'{e}s de la plateforme cible. L’interpr\'{e}tation dans notre cadre est une prise de position sur la
d\'{e}finition et l’effet des guidelines dans le but de produire des r\`{e}gles de transformation de la source.
104
CHAPITRE 6. MÉCANISMES DE MIGRATION DES UI
Dans la section suivante nous proposons un mod\`{e}le de r\`{e}gles pour les transformations du mod\`{e}le
des UI pour permettre une utilisation effective des guidelines par le processus de migration des UI
vers les tables interactives.
6.2.3
Utilisation effective des guidelines
Les guidelines sont utilis\'{e}es par les m\'{e}canismes de restructuration pendant les transformations
des mod\`{e}les abstraits [MVG06]. Notre processus de migration des UI vers les tables interactives
impl\'{e}mente deux types de transformation des mod\`{e}les.
– Les transformations exog\`{e}nes et verticales [MVG06] qui comprennent le m\'{e}canisme d’abstraction de l’UI finale dans le mod\`{e}le abstrait et le m\'{e}canisme de concr\'{e}tisation du mod\`{e}le abstrait
en UI finale \`{a} l’aide des \'{e}l\'{e}ments de la plateforme cible.
– Les transformations endog\`{e}nes et horizontales [MVG06] qui consistent \`{a} restructurer le mod\`{e}le abstrait de l’UI. Dans notre cas la restructuration de l’UI de d\'{e}part s’appuie sur les guidelines de la plateforme d’arriv\'{e}e afin d’avoir une UI conforme aux crit\`{e}res ergonomiques de la
cible.
Pour utiliser les guidelines, nous proposons de les traduire sous la forme des r\`{e}gles utilisables par
les m\'{e}canismes de restructuration et de concr\'{e}tisation (cf figure 6.4). Le mod\`{e}le d’UI \`{a} restructurer
se pr\'{e}sente sous la forme d’un arbre dont les nœuds et les feuilles correspondent \`{a} des Interactors. .
FIGURE 6.4 – Utilisations effectives des guidelines par les m\'{e}canismes de migration des UI
Nous avons opt\'{e} pour une strat\'{e}gie de restructuration bas\'{e}e sur la substitution des Interactors de la
source pour d\'{e}duire la cible. La substitution des Interactors conserve la structure arborescente de l’UI
source tout en utilisant des Interactors de la cible qui ont des interactions et une structure conforme
aux guidelines. Une des recommandations du diagramme de la figure 6.2 pr\'{e}conise de “d\'{e}placer et
faire tourner des groupes d’\'{e}l\'{e}ments graphiques”. En consid\'{e}rant que les groupes d’\'{e}l\'{e}ments correspondent aux Containers, nous pourront d\'{e}crire une r\`{e}gle de substitution qui remplace les Containers
de la source par ceux qui impl\'{e}mentent les PI de d\'{e}placement et de rotation.
Par ailleurs, les r\`{e}gles de concr\'{e}tisation ont pour objectifs de g\'{e}n\'{e}rer une UI concr\`{e}te qui sera
personnalis\'{e}e par la personne en charge de la migration. Les guidelines sont traduites dans les r\`{e}gles
de concr\'{e}tisation. Les recommandations et les contraintes issues de la propri\'{e}t\'{e} de tangibilit\'{e} sont
prises en compte pendant la concr\'{e}tisation. Les contraintes sur la taille des \'{e}l\'{e}ments graphiques sont
appliqu\'{e}es par les r\`{e}gles de concr\'{e}tisation.
Dans l’objectif d’assister les personnes en charge de la migration pendant la phase de personnalisation en pr\'{e}cisant la guideline et les crit\`{e}res ergonomiques de conception associ\'{e}s \`{a} chaque substitution
ou concr\'{e}tisation. Nous pr\'{e}sentons d’abord un mod\`{e}le des r\`{e}gles de substitution (cf section 6.2.3.1),
6.2. PRISE EN COMPTE DES GUIDELINES
105
ensuite un mod\`{e}le des r\`{e}gles de concr\'{e}tisation (cf section 6.2.3.2). Ces mod\`{e}les permettent de d\'{e}crire
des r\`{e}gles formelles issues des guidelines et des crit\`{e}res ergonomiques de conception.
6.2.3.1
Mod\`{e}le de r\`{e}gles de substitution
Les r\`{e}gles de substitution ont pour objectif d’indiquer pour un ou plusieurs interacteur(s) de l’UI
source les caract\'{e}ristiques des interacteurs \'{e}quivalents pour le mod\`{e}le cible. Ces r\`{e}gles de substitution sont constitu\'{e}es de deux membres : le membre de gauche (de la figure 6.5) correspond aux
caract\'{e}ristiques des interacteurs de l’UI source et le membre de droite (de la figure 6.5) correspond
aux caract\'{e}ristiques souhait\'{e}es pour les \'{e}l\'{e}ments \'{e}quivalents du mod\`{e}le cible. Les composants graphiques correspondant au membre de droite de la r\`{e}gle de substitution sont identifi\'{e}s \`{a} l’aide des
op\'{e}rateurs d’\'{e}quivalences.
Chaque r\`{e}gle de substitution est conforme \`{a} une recommandation ou une contrainte des guidelines
des figures 6.2 et 6.3 ci-dessus.
Le mod\`{e}le de r\`{e}gles de substitution sert \`{a} impl\'{e}menter un feedback du respect des crit\`{e}res ergonomiques de conception pendant les substitutions manuelles des interacteurs par les personnes en charge
de la migration. En effet ce mod\`{e}le permet de savoir si la substitution d’un Interactor est conforme \`{a}
une guideline.
Nous sp\'{e}cifions formellement les r\`{e}gles de substitution \`{a} l’aide du mod\`{e}le de la figure 6.5.
FIGURE 6.5 – Mod\`{e}le de r\`{e}gles de substitution
La classe SubstitutionRule
caract\'{e}rise une r\`{e}gle de substitution, elle est constitu\'{e}e d’un identifiant
unique pour chaque r\`{e}gle et des r\^{o}les :
– operators correspond aux op\'{e}rateurs d’\'{e}quivalences \`{a} appliquer pour d\'{e}terminer les Widgets
\'{e}quivalents de la plateforme cible
– leftMember correspond \`{a} la caract\'{e}ristique des interacteurs du mod\`{e}le source sur lesquels s’applique ces r\`{e}gles
– rightMember correspond \`{a} la caract\'{e}ristique souhait\'{e}e des interacteurs de la cible
L’attribut guideline pr\'{e}cise la guideline qui a permis de d\'{e}duire la r\`{e}gle de substitution conform\'{e}ment aux diagrammes de raffinement des figures 6.2 et 6.3 ci-dessus.
106
CHAPITRE 6. MÉCANISMES DE MIGRATION DES UI
L’attribut constraint pr\'{e}cise si une r\`{e}gle de substitution est issue d’une guideline qui comporte des
contraintes.
La classe InteractorType
caract\'{e}rise les membres d’une r\`{e}gle de substitution. Elle est une m\'{e}ta
classe de l’interface Interactor et correspond a l’ensemble des interacteurs qui se d\'{e}finie par les attributs suivants :
– type est le type d’interacteur de l’ensemble (UIComponent ou Container)
– dataType est le type de donn\'{e}es des interacteurs de l’ensemble
– effectivePI correspond aux PI effectives des interacteurs de l’ensemble
La classe EquivalenceOperator
caract\'{e}rise les op\'{e}rateurs d’\'{e}quivalences utilis\'{e}s, elle se d\'{e}finie
par les attributs suivants :
– symbol est l’un des symboles des six op\'{e}rateurs d\'{e}crits \`{a} la section 5.4. L’\'{e}num\'{e}ration SymbolType pr\'{e}sente les six op\'{e}rateurs d’\'{e}quivalences pr\'{e}sent\'{e}s \`{a} la section 5.4.
– piSet correspond \`{a} l’ensemble des PI utilis\'{e} comme param\`{e}tre par les op\'{e}rateurs d’\'{e}quivalences
– dataTypeSet correspond aux types de donn\'{e}es de l’ensemble TargetType des op\'{e}rateurs d’\'{e}quivalences
6.2.3.2
Mod\`{e}le de r\`{e}gle de concr\'{e}tisation
En ce qui concerne la concr\'{e}tisation du mod\`{e}le obtenu apr\`{e}s l’application des r\`{e}gles de substitution, nous proposons aussi un mod\`{e}le qui permet de d\'{e}crire de mani\`{e}re formelle les r\`{e}gles de
concr\'{e}tisation. La concr\'{e}tisation permet de proposer une instance du mod\`{e}le de l’UI en UI finale en
utilisant les composants graphiques et les dispositifs d’interactions de la plateforme cible. Les guidelines permettent de d\'{e}crire l’utilisation de ces instruments d’interactions. Le mod\`{e}le de la figure 6.6
permet de d\'{e}crire les r\`{e}gles de concr\'{e}tisation.
FIGURE 6.6 – Mod\`{e}le de r\`{e}gles de concr\'{e}tisation
La classe ConcretizationRule
caract\'{e}rise les r\`{e}gles de concr\'{e}tisation des interacteurs d’un mod\`{e}le
de l’UI. Chaque r\`{e}gle a un identifiant (id) unique et elle se d\'{e}finie aussi par les r\^{o}les :
– devices qui correspond aux dispositifs d’interactions utilisables pendant la concr\'{e}tisation
6.2. PRISE EN COMPTE DES GUIDELINES
107
– abstractMember qui correspond au type d’interacteur (InteractorType) sur le quel s’applique
cette r\`{e}gle de concr\'{e}tisation ; \`{a} chaque type d’interacteur correspond une seule r\`{e}gle de concr\'{e}tisation
– concreteMember qui correspond aux Widgets \'{e}quivalents de la plateforme cible
L’attribut guideline pr\'{e}cise la guideline qui a permis de d\'{e}duire la r\`{e}gle de concr\'{e}tisation conform\'{e}ment aux diagrammes d’affinement des figures 6.2 et 6.3 ci-dessus.
L’attribut constraint pr\'{e}cise si une r\`{e}gle de concr\'{e}tisation est issue d’une guideline qui comporte
des contraintes.
La classe InteractionDevice
caract\'{e}rise les dispositifs physiques d’interactions qui permettent de
sp\'{e}cifier les interactions utilisateurs dans le cadre d’une guideline. Elle est caract\'{e}ris\'{e}e par l’attribut
name du dispositif et l’attribut definition qui est une description textuelle du dispositif.
La classe InteractorType
correspond \`{a} celle d\'{e}crite par le mod\`{e}le de r\`{e}gles de substitution \`{a} la
figure 6.5.
La classe Widget
correspond \`{a} celle d\'{e}crite par le mod\`{e}le des types de composants graphiques \`{a} la
figure 5.4
6.2.3.3
Remarques
Si une r\`{e}gle de substitution issue d’une contrainte n’est pas respect\'{e}e pendant la personnalisation
d’une UI, il est indispensable d’alerter la personne en charge de la migration en pr\'{e}cisant la guideline
et les crit\`{e}res ergonomiques de conception associ\'{e}s.
Il en est de m\^{e}me si une r\`{e}gle de concr\'{e}tisation issue d’une contrainte est modifi\'{e}e.
Les guidelines de migration des UI dans notre contexte ont pour objectif de permettre la description des caract\'{e}ristiques des membres de droite au sein r\`{e}gles de substitution et concr\'{e}tisation. Ces
caract\'{e}ristiques d\'{e}pendent de l’interpr\'{e}tation de chaque guideline en fonction des \'{e}l\'{e}ments de notre
mod\`{e}le de l’UI (PI, Container, UIComponent, Type de donn\'{e}es, etc.).
Dans la section suivante nous d\'{e}crivons les diff\'{e}rentes r\`{e}gles de substitution et concr\'{e}tisation.
108
CHAPITRE 6. MÉCANISMES DE MIGRATION DES UI
6.3
Transformations du mod\`{e}le de l’UI source
Notre processus semi automatique propose aux personnes en charge de la migration, des UI destin\'{e}es aux tables interactives sans le positionnement des \'{e}l\'{e}ments graphiques et sans le style. L’UI
propos\'{e}e est ensuite personnalis\'{e}e en respectant les guidelines pour les UI collaboratives et pour les
UI tangibles. La proposition et la personnalisation des UI pour la cible sont des transformations du
mod\`{e}le de l’UI. La restructuration du mod\`{e}le source est une forme de transformation par la substitution et la concr\'{e}tisation des interacteurs \'{e}quivalents.
Dans cette section nous d\'{e}crivons les r\`{e}gles de substitution et de concr\'{e}tisation pour chaque type
d’interacteurs caract\'{e}ris\'{e} par la classe InteractorType. Les transformations concernent des ensembles
de composants graphiques caract\'{e}ris\'{e}s par les PI, les types de donn\'{e}es et les types d’interacteurs.
Par exemple, les composants graphiques de contr\^{o}le tels que les menus ou les groupes de boutons
sont caract\'{e}ris\'{e}s par la PI Activation. Les composants graphiques d’affichage de contenus (tels que
les tableaux, les listes, les carrousels, etc.) sont caract\'{e}ris\'{e}s par la PI Data Display. Les composants
graphiques de modifications de contenus tels que les champs de texte par exemple peuvent \^{e}tre caract\'{e}ris\'{e}s par la PI Data Edition.
Nous avons opt\'{e} pour deux strat\'{e}gies de substitution et de concr\'{e}tisation du mod\`{e}le source qui
prend en compte les deux types d’interacteurs. La premi\`{e}re consiste \`{a} consid\'{e}rer les groupes d’\'{e}l\'{e}ments (section 6.3.1) et la seconde consid\`{e}re les \'{e}l\'{e}ments graphiques de mani\`{e}re individuelle (section 6.3.2) pendant la restructuration. En effet, le mod\`{e}le d’instance que nous avons pr\'{e}sent\'{e} \`{a} la
section 5.3.2 d\'{e}crit deux types d’interacteurs : Container et UIComponent), ils repr\'{e}sentent respectivement les groupes et les \'{e}l\'{e}ments individuels. La substitution et la concr\'{e}tisation de ces deux types
d’\'{e}l\'{e}ments graphiques n’ont pas les m\^{e}mes objectifs. Dans notre cadre, la substitution et la concr\'{e}tisation des Containers permettent de les d\'{e}placer, de les roter ou de les utiliser avec des objets tangibles
tout en respectant les contraintes des guidelines. Par ailleurs, la substitution et la concr\'{e}tisation des
UIComponent ont pour but de conserver les PI de d\'{e}part et d’utiliser les types de donn\'{e}es adapt\'{e}s
pour les tables interactives.
6.3.1
Transformation des groupes d’\'{e}l\'{e}ments graphiques
La transformation des groupes d’\'{e}l\'{e}ments graphiques a pour objectifs de favoriser les activit\'{e}s collaboratives en rendant accessible chaque groupe \`{a} tous les utilisateurs. Elle permet aussi d’associer les
groupes d’\'{e}l\'{e}ments graphiques \`{a} des objets tangibles. Les r\`{e}gles de substitution et de concr\'{e}tisation
des groupes de composants graphiques ont pour objectif de favoriser la flexibilit\'{e} 19, le contr\^{o}le explicite 20, la charge de travail 21 sans alt\'{e}rer l’homog\'{e}n\'{e}it\'{e} 22 et le contr\^{o}le explicite 23 (cf. figures 6.2
et 6.3).
Pour les UI graphiques de mani\`{e}re g\'{e}n\'{e}rale, les groupes d’\'{e}l\'{e}ments graphiques ont une s\'{e}mantique pour l’utilisateur. En effet les panels, les formulaires, les menus ou les tableaux repr\'{e}sentent des
groupes d’\'{e}l\'{e}ments dans une UI graphique qui permettent d’afficher des donn\'{e}es ou d’acc\'{e}der \`{a} des
fonctionnalit\'{e}s. Dans notre mod\`{e}le des instances d’une UI pr\'{e}sent\'{e} \`{a} la section 5.3.2 nous proposons
de d\'{e}crire les groupes d’\'{e}l\'{e}ments d’une UI graphique en nous basant sur la nature du container. Nous
19. Le redimensionnement, le d\'{e}placement et la rotation de certains groupes facilitent l’accessibilit\'{e}
20. L’affichage des groupes d’\'{e}l\'{e}ments \`{a} l’aide des objets tangibles par exemple est un contr\^{o}le explicite
21. Afficher les groupes cach\'{e}s par des objets tangibles r\'{e}duit la charge de travail car l’utilisateur aura moins d’\'{e}l\'{e}ments
graphiques \`{a} l’\'{e}cran
22. Les containers sans taille limite ne permettent pas d’avoir une UI homog\`{e}ne par exemple
23. Le non respect de l’accessibilit\'{e} des menus ne permet pas une utilisation ais\'{e}e de l’UI par exemple
6.3. TRANSFORMATIONS DU MODÈLE DE L’UI SOURCE
109
avons identifi\'{e} les containers de types Racine 24, R\'{e}cursif 25, Homog\`{e}ne 26 et H\'{e}t\'{e}rog\`{e}ne 27.
La transformation des groupes d’\'{e}l\'{e}ments de l’UI source passe par la substitution et la concr\'{e}tisation de ces quatre types de containers en fonction des guidelines interpr\'{e}t\'{e}es \`{a} la section 6.2.2.
L’avantage de ces quatre types de containers est qu’ils permettent de repr\'{e}senter l’UI source en se
basant sur les interactions et les donn\'{e}es des instances de composants graphiques et non sur les types
des composants graphiques. En effet les diff\'{e}rences entre les types et les instances des composants
graphiques sont importantes pour les processus de migration car le choix des \'{e}l\'{e}ments graphiques
\'{e}quivalents en d\'{e}pend. Une instance de composant graphique peut impl\'{e}menter moins de PI effectives que son type, le processus de migration ne doit prendre en compte que les PI effectives pour
\'{e}tablir les \'{e}quivalences avec la cible. Cette contrainte a pour but de ne pr\'{e}server que les dialogues des
UI de d\'{e}part pendant la migration. Dans ce cas deux instances de composant d’un m\^{e}me type peuvent
avoir des \'{e}quivalences diff\'{e}rentes, si elles ont des PI effectives diff\'{e}rentes.
Pour transformer la structure des UI de d\'{e}part, nous consid\'{e}rons quatre groupes d’\'{e}l\'{e}ments graphiques en fonction de leurs utilisations dans une UI graphique. Les PI, les containers et les types
de donn\'{e}es nous permettent de les caract\'{e}riser. Le but de cette consid\'{e}ration est de permettre une
sp\'{e}cialisation des groupes en fonction des interactions.
1. Les groupes de contr\^{o}le (ControlGroup) repr\'{e}sentent un ensemble d’\'{e}l\'{e}ments graphiques qui
permettent d’activer des fonctionnalit\'{e}s d’une UI graphique (cf section 6.3.1.1).
2. Les groupes d’affichage de contenus (DisplayGroup) repr\'{e}sentent un ensemble de composants
graphiques qui permettent de visualiser des donn\'{e}es (cf section 6.3.1.2).
3. Les groupes de modification de contenus (UpdateGroup) repr\'{e}sentent un ensemble de composants graphiques qui permettent des modifications des donn\'{e}es d’application (cf section
6.3.1.3).
4. Les groupes mixtes (MixedGroup) repr\'{e}sentent un ensemble de composants graphiques qui
contient des groupes de contr\^{o}le, des groupes d’affichage et/ou des groupes de modification de
contenus (cf section 6.3.1.4).
La figure 6.7 illustre un exemple de ces diff\'{e}rents groupes en s’appuyant sur l’application Comic
Book d\'{e}crite \`{a} la section 2.2. Dans la suite de cette section nous pr\'{e}sentons les r\`{e}gles de substitution
de concr\'{e}tisation des diff\'{e}rents groupe d’\'{e}l\'{e}ments.
6.3.1.1
Transformation des groupes de contr\^{o}le
Les groupes de contr\^{o}le repr\'{e}sentent les composants graphiques tels que les menus, les groupes de
boutons, les listes de s\'{e}lections, etc. Ces groupes sont caract\'{e}ris\'{e}s par un container de type Homog\`{e}ne
dont chaque \'{e}l\'{e}ment impl\'{e}mente les PI Activation et Navigation et/ou Widget Selection (\'{e}ventuellement Display Data pour les boutons munis des ic\^{o}nes).
Ce groupe est caract\'{e}ris\'{e} par la classe InteractorType d\'{e}crit par le tableau 6.1
Les ControlGroups sont transform\'{e}s pour les tables interactives suivant la guideline “G22 : Accessibilit\'{e} des Menus” qui a pour objectif de rendre les groupes de contr\^{o}le accessibles. Cette accessibilit\'{e} se traduit par l’ajout des comportements de d\'{e}placement et de rotation. Si le concepteur le
souhaite il peut associer un groupe de contr\^{o}le \`{a} un objet tangible en appliquant la guideline “G3 :
Objets Tangibles et Objets virtuels”.
24. Ce sont les containers qui contiennent tous les \'{e}l\'{e}ments d’une UI
25. Ce sont les containers qui ne contiennent que les \'{e}l\'{e}ments de type container
26. Ce sont les containers qui contiennent les composants graphiques avec des donn\'{e}es du m\^{e}me type
27. Ces containers contiennent \`{a} la fois les Containers et des UIComponents avec des types de donn\'{e}es diff\'{e}rents
110
CHAPITRE 6. MÉCANISMES DE MIGRATION DES UI
FIGURE 6.7 – Exemple des groupes d’\'{e}l\'{e}ments graphiques \`{a} transformer
InteractorType
type
type ∈\textbackslash{}{Container\textbackslash{}}
type.containerType() = \textbackslash{}{Homogeneous\textbackslash{}}
∀i ∈type.contains,
i.e f fectivePI ⊆
 Activation,WidgetSelection,
Navigation, DisplayData

i.type.contentType ∈\textbackslash{}{String,Image\textbackslash{}}
dataType
type.type.contentType ∈\textbackslash{}{String,Image\textbackslash{}}
effectivePI
type.e f fectivePI ⊆\textbackslash{}{ WidgetSelection, Navigation\textbackslash{}}
TABLE 6.1 – Caract\'{e}ristiques des ControlGroups
R\`{e}gle de substitution
Cette r\`{e}gle est issue de la contrainte de la guideline “G22 : Accessibilit\'{e} des
Menus”. Le tableau 6.2 caract\'{e}rise la r\`{e}gle de substitution des groupes de contr\^{o}le. Cette r\`{e}gle pr\'{e}conise l’application de l’op\'{e}rateur d’\'{e}quivalence large ≦AdditionalPI en privil\'{e}giant les PI WidgetMove
et WidgetRotation pour le container. L’utilisation de l’op\'{e}rateur d’\'{e}quivalence stricte avec diff\'{e}rence
de donn\'{e}es ∼=TargetType doit privil\'{e}gier les donn\'{e}es de type Image pour les \'{e}l\'{e}ments des menus.
R\`{e}gle de concr\'{e}tisation
Cette r\`{e}gle est issue de la guideline “G3 : Objets Tangibles et Objets
virtuels”. Le tableau 6.3 caract\'{e}rise la concr\'{e}tisation d’un ControlGroup. L’utilisation des objets tangibles a pour but d’afficher le ControlGroup qui sera cach\'{e} \`{a} tout moment. Chaque utilisateur qui
souhaite acc\'{e}der \`{a} un groupe de contr\^{o}le doit avoir un objet capable de l’afficher.
Remarques
– Nous avons not\'{e} que l’affichage multiple d’un ControlGroup entra\^{i}ne des incoh\'{e}rences si deux
utilisateurs font appel \`{a} une fonctionnalit\'{e} qui modifie une donn\'{e}e. En effet le NF de l’application n’\'{e}tant pas modifi\'{e}, les appels de fonctionnalit\'{e}s se font de mani\`{e}re s\'{e}quentielle et les
6.3. TRANSFORMATIONS DU MODÈLE DE L’UI SOURCE
111
SubstitutionRule
id
ControlGroupSRule
guideline
G22
constraint
True
leftMember
InteractorType du tableau 6.1
rightMember
InteractorType du tableau 6.1
operators
≦AdditionalPI et ∼=TargetType
avec AdditionnalPI = \textbackslash{}{WidgetMove, WidgetRotation\textbackslash{}}
et TargetType = \textbackslash{}{Image\textbackslash{}}
TABLE 6.2 – R\`{e}gle de substitution des groupes de contr\^{o}le
ConcretizationRule
id
ControlGroupCRule
guideline
G3
constraint
True
abstractMember
InteractorType du tableau 6.11
concreteMember
w ∈\textbackslash{}{Widget\textbackslash{}}
devices
Écran Tactile, Objet Tangible unique pour chaque ControlGroup
TABLE 6.3 – R\`{e}gle de concr\'{e}tisation des groupes de contr\^{o}le
modifications (changement de couleur ou de taille par exemple) se feront sur la m\^{e}me donn\'{e}e
pour les deux utilisateurs. Cependant les appels vers le NF qui permettent d’afficher des composants graphiques (une nouvelle fen\^{e}tre par exemple), cr\'{e}ent des instances diff\'{e}rentes de ces
composants graphiques pour chaque utilisateur.
– Ces observations nous ont pouss\'{e}s \`{a} opter pour l’affichage unique d’un ControlGroup dans le
but d’\'{e}viter les incoh\'{e}rences. Dans le cas o\`{u} deux utilisateurs ont des objets tangibles pour le
m\^{e}me ControlGroup, un seul ControlGroup sera affich\'{e} s’ils les posent sur la surface d’affichage au m\^{e}me moment. Dans les autres cas, le dernier objet pos\'{e} sera consid\'{e}r\'{e}.
– La transformation des ControlGroups que nous proposons permet de r\'{e}soudre les probl\`{e}mes
d’accessibilit\'{e} des \'{e}l\'{e}ments graphiques et des menus en particulier. Par exemple les menus File
et Edit seront accessibles par deux utilisateurs autour de la table au m\^{e}me moment sans qu’ils
se g\^{e}nent.
– Cependant cette transformation ne permet pas \`{a} deux utilisateurs d’acc\'{e}der \`{a} un menu au m\^{e}me
moment. Par exemple, le menu File ne peut \^{e}tre affich\'{e} plusieurs fois au m\^{e}me moment. Cette
restriction permet de garder des dialogues coh\'{e}rents et qui ne perturbent pas les utilisateurs. Par
exemple si deux utilisateurs souhaitent ouvrir deux fichiers (repr\'{e}sentant des bandes dessin\'{e}es)
diff\'{e}rents, cette transformation ne le permet pas car l’application ne peut afficher qu’un seul
document au m\^{e}me moment.
Exemple de ControlGroup
Nous consid\'{e}rons la figure 6.8 qui est un art\'{e}fact de l’application CBA
pr\'{e}sent\'{e}e au chapitre 3, repr\'{e}sentant les menus de cette application. En effet, les menus File, Edit, etc.
de la figure 6.8 repr\'{e}sentent des groupes de contr\^{o}le conformes au type d’interacteurs d\'{e}crit par le
tableau 6.1. Chaque sous-menu impl\'{e}mente les PI Activation, Navigation et Widget Selection car ils
sont accessibles par un clavier ou une souris et activent des fonctionnalit\'{e}s. Les donn\'{e}es de ces menus
et sous-menus sont des cha\^{i}nes de caract\`{e}res.
112
CHAPITRE 6. MÉCANISMES DE MIGRATION DES UI
FIGURE 6.8 – Exemple de ControlGroup
L’exemple pr\'{e}sent\'{e} est migr\'{e} sur une table interactive en appliquant d’abord la r\`{e}gle de substitution d\'{e}crite par le tableau 6.2 ensuite la r\`{e}gle de concr\'{e}tisation d\'{e}crite par le tableau 6.3. Les menus
File, Edit, View et Help auront les PI Widget Move et Widget Rotation et ils seront affich\'{e}s chacun par
des Tags associ\'{e}s \`{a} des objets dans le cadre d’une table Micorsoft PixelSense. La figure 6.9 pr\'{e}sente
une illustration du menu File pour table interactive.
FIGURE 6.9 – Exemple de ControlGroup sur table interactive Micorsoft PixelSense
6.3.1.2
Transformation des groupes d’affichage de contenus
Les groupes d’affichage de contenus repr\'{e}sentent les composants graphiques tels que les listes,
les tableaux, les groupes de donn\'{e}es, etc. Ces groupes sont caract\'{e}ris\'{e}s par les containers de type
Homog\`{e}ne ou H\'{e}t\'{e}rog\`{e}ne. Chaque \'{e}l\'{e}ment de DisplayGroup est un UIComponent impl\'{e}mentant les
PI Data Display, Widget Selection et Navigation. Les \'{e}l\'{e}ments DisplayGroup n’impl\'{e}mentent pas les
PI Activation, Data Edition, Data Move In/Out et Data Selection. Ce groupe est caract\'{e}ris\'{e} par la
classe InteractorType d\'{e}crit par le tableau 6.4.
La transformation des groupes d’affichage de contenus pour les tables interactives se fait en fonction des types de donn\'{e}es et dans le but de partager l’espace de travail en permettant \`{a} plusieurs
personnes de consulter les contenus affich\'{e}s.
– Si les contenus des \'{e}l\'{e}ments de la cible sont des Images, MediaElements ou Objects
– La guideline “G11 : Taille des \'{e}l\'{e}ments graphiques” permet d’ajouter la PI WidgetResize \`{a}
chaque \'{e}l\'{e}ment de DisplayGroup.
– La guideline “G12 : Structure et Positionnement des \'{e}l\'{e}ments graphiques” permet le changement de type de donn\'{e}es des contenus si possible en favorisant les types Images par exemple.
6.3. TRANSFORMATIONS DU MODÈLE DE L’UI SOURCE
113
InteractorType
type
type ∈\textbackslash{}{Container\textbackslash{}}
type.containerType() ∈\textbackslash{}{Homogeneous,Heterogeneous\textbackslash{}}
∀i ∈type.contains, i ∈\textbackslash{}{UIComponent\textbackslash{}}

Activation, DataEdition,
DataMoveIn/Out, DataSelection

⊈i.e f fectivePI
DataDisplay ∈i.e f fectivePI
i.type.contentType ∈\textbackslash{}{Image, MediaElement, Object\textbackslash{}}
∨
i.type.contentType ∈\textbackslash{}{Integer, String, Boolean\textbackslash{}}
dataType
type.type.contentType ∈DataType
effectivePI
type.e f fectivePI ̸= /0
TABLE 6.4 – Caract\'{e}ristiques des groupes d’affichage de contenus
Cependant les recommandations li\'{e}es aux positionnements des \'{e}l\'{e}ments par la guideline
G12 28 ne peuvent pas \^{e}tre appliqu\'{e}es dans ce cas car les \'{e}l\'{e}ments de ce groupe sont d\'{e}pla\c{c}ables conform\'{e}ment \`{a} la guideline G13 29.
– La guideline “G13 : Comportement des \'{e}l\'{e}ments graphiques” permet d’ajouter les PI WidgetResize et WidgetMove \`{a} chaque \'{e}l\'{e}ments de DisplayGroup
– Si les contenus des \'{e}l\'{e}ments de la cible sont de type Boolean, Integer ou String
– La guideline “G11 : Taille des \'{e}l\'{e}ments graphiques” permet d’ajouter la PI WidgetResize au
container et les \'{e}l\'{e}ments sont migr\'{e}s sans changement. Dans le cas d’un tableau contenant
des donn\'{e}es num\'{e}riques, il est pr\'{e}f\'{e}rable de conserver la structure du tableau pour faciliter
sa lecture et sa compr\'{e}hension.
– La guideline “G13 : Comportement des \'{e}l\'{e}ments graphiques” permet d’ajouter les PI WidgetResize, WidgetMove au container et les \'{e}l\'{e}ments sont migr\'{e}s sans changement.
R\`{e}gle de substitution des DisplayGroups
La substitution de ce groupe d\'{e}pend des types de donn\'{e}es de ses \'{e}l\'{e}ments. Si tous les \'{e}l\'{e}ments sont des cha\^{i}nes de caract\`{e}res ou des donn\'{e}es num\'{e}riques
provenant du NF, il est difficile de transformer ces donn\'{e}es en images par exemple. En effet les donn\'{e}es provenant du NF peuvent constituer un ensemble infini et le concepteur ne peut pas pr\'{e}voir
les images repr\'{e}sentant chaque donn\'{e}e du NF. Le tableau 6.5 caract\'{e}rise la r\`{e}gle de substitution des
groupes d’affichage des contenus dans le cas o\`{u} les donn\'{e}es sont de type Boolean, Integer, String et
proviennent du NF. Cette r\`{e}gle pr\'{e}conise l’application de l’op\'{e}rateur d’\'{e}quivalence large ≦AdditionalPI
pour le container en ajoutant les PI WidgetMove, WidgetRotation et WidgetResize. L’op\'{e}rateur d’\'{e}quivalences stricte pour les \'{e}l\'{e}ments du groupe container.
Dans le cas o\`{u} les donn\'{e}es graphiques peuvent \^{e}tre traduites vers les types Image, MediaElement
ou Object, la r\`{e}gle d\'{e}crite par le tableau 6.5 est utilis\'{e}e en appliquant l’op\'{e}rateur d’\'{e}quivalence large
≦AdditionalPI pour les \'{e}l\'{e}ments du groupe en ajoutant les PI WidgetMove, WidgetRotation et WidgetResize. L’op\'{e}rateur d’\'{e}quivalence stricte est appliqu\'{e} pour le container du groupe.
La r\`{e}gle de substitution des DisplayGroups est issue des guidelines “G11 : Taille des \'{e}l\'{e}ments
graphiques”, “G13 : Comportement des \'{e}l\'{e}ments graphiques” et de la contrainte de la guideline
“G12 : Structure et Positionnement des \'{e}l\'{e}ments graphiques”.
28. Structure et Positionnement des \'{e}l\'{e}ments graphiques
29. Comportement des \'{e}l\'{e}ments graphiques d’une UI
114
CHAPITRE 6. MÉCANISMES DE MIGRATION DES UI
SubstitutionRule
id
DisplayGroupSRule1
guideline
G11, G12, G13
constraint
True
leftMember
InteractorType du tableau 6.4)
rightMember
InteractorType du tableau 6.4)
operators
≦AdditionalPI et ≡
avec AdditionnalPI =
 WidgetMove, WidgetRotation,
WidgetResize

TABLE 6.5 – R\`{e}gle de substitution des groupes DisplayGroups
R\`{e}gle de concr\'{e}tisation des DisplayGroup
– Les groupes contenant des donn\'{e}es num\'{e}riques ou textuelles sont concr\'{e}tis\'{e}s en tableaux avec
les comportements WidgetMove, WidgetResize et WidgetRotation.
– Les groupes contenant des donn\'{e}es de type Image, MediaElement ou Object sont concr\'{e}tis\'{e}s
en impl\'{e}mentant pour chaque \'{e}l\'{e}ment les comportements WidgetMove, WidgetResize et WidgetRotation.
Remarques
– Nous pr\'{e}conisons dans le cas o\`{u} une UI migr\'{e}e contient plusieurs groupes d’affichage de d\'{e}limiter pour chaque groupe une zone ou de cacher certains groupes (par les groupes contenant les
donn\'{e}es num\'{e}riques) et l’afficher avec des objets tangibles en appliquant la guideline Objets
Tangibles et Objets virtuels (G3). Dans ce cas plusieurs instances d’un m\^{e}me DisplayGroup
peuvent \^{e}tre affich\'{e}es pour consultation.
– La transformation des DisplayGroups concerne la dimension Structure et Positionnement des
\'{e}l\'{e}ments d’une UI. Elle traite les questions d’accessibilit\'{e} des \'{e}l\'{e}ments graphiques pour pr\'{e}senter des donn\'{e}es pour tous les utilisateurs. La contrainte li\'{e}e \`{a} la taille des \'{e}l\'{e}ments graphiques
r\'{e}sout l’un des probl\`{e}mes li\'{e}s \`{a} la dimension Style de l’espace des probl\`{e}mes.
Exemple d’un DisplayGroup
Nous consid\'{e}rons la figure 6.10 qui repr\'{e}sente une liste d’images.
Chaque \'{e}l\'{e}ment de la liste d’images impl\'{e}mente les PI Navigation, Widget Selection et Data Display
car ils sont accessibles par un clavier ou une souris et activent des fonctionnalit\'{e}s.
L’exemple pr\'{e}sent\'{e} est migr\'{e} sur une table interactive en appliquant d’abord la r\`{e}gle de substitution d\'{e}crite par le tableau 6.5 ensuite chaque \'{e}l\'{e}ment de la liste sera un ScatterViewItem pour \^{e}tre
accessible \`{a} tous les utilisateurs dans le cadre d’une table Micorsoft PixelSense. La figure 6.11 illustre
un exemple de DisplayGroup pour une table Micorsoft PixelSense
FIGURE 6.10 – Exemple de DisplayGroup
6.3. TRANSFORMATIONS DU MODÈLE DE L’UI SOURCE
115
FIGURE 6.11 – Exemple de DisplayGroup migr\'{e} sur une table interactive
6.3.1.3
Transformation des groupes de modification de contenus
Les groupes de modification de contenus repr\'{e}sentent les composants graphiques tels que les
formulaires, des canevas, des tableaux \'{e}ditables, etc. Ces groupes sont caract\'{e}ris\'{e}s par des containers
de type H\'{e}t\'{e}rog\`{e}ne ou Homog\`{e}ne. Un UpdateGroup contient des interacteurs impl\'{e}mentant les PI
Activation, Data Edition, Data Selection ou Data Move In/Out. Nous caract\'{e}risons ces groupes par la
classe InteractorType d\'{e}crit par le tableau 6.6.
InteractorType
type
type ∈\textbackslash{}{Container\textbackslash{}}
type.containerType() = \textbackslash{}{Heterogeneous,Homogeneous\textbackslash{}}
∃i ∈type.contains, i ∈\textbackslash{}{UIComponent\textbackslash{}}

Activation, DataEdition,
DataMoveIn/Out, DataSelection

⊆i.e f fectivePI
i.type.contentType ∈DataType
dataType
type.type.contentType ∈DataType
effectivePI
type.e f fectivePI ̸= /0
TABLE 6.6 – Caract\'{e}ristiques des groupes de modification de contenus
La transformation des UpdateGroups pour les tables interactives se fait selon la guideline “G21 :
Activit\'{e}s Bloquantes” qui recommande d’avoir des activit\'{e}s de modification de contenus qui n’emp\^{e}chent pas les diff\'{e}rents utilisateurs d’acc\'{e}der \`{a} d’autres fonctionnalit\'{e}s de l’application. Dans le
cadre d’une UI desktop la saisie de deux champs de texte distincts n’est pas possible, les tables interactives offrent n\'{e}anmoins la possibilit\'{e} de le faire. Nous consid\'{e}rons que la guideline “G21 : Activit\'{e}s
Bloquantes” permet d’avoir plusieurs UpdateGroups distincts. Cependant il n’est pas possible d’avoir
plusieurs instances d’un UpdateGroup dans modifier le NF pour prendre en compte la modification
d’une donn\'{e}e du NF par deux utilisateurs distincts.
La guideline “G13 : Comportement des \'{e}l\'{e}ments graphiques” permet d’ajouter les PI WidgetResize et WidgetMove au container repr\'{e}sentant le groupe.
La guideline “G12 : Structure et Positionnement des \'{e}l\'{e}ments graphiques” est appliqu\'{e}e pour les
\'{e}l\'{e}ments du container pour permettre le changement de type de donn\'{e}es des contenus si possible en
favorisant le type Images par exemple. Les recommandations li\'{e}es aux positionnements des \'{e}l\'{e}ments
par la guideline G12 30 sont appliqu\'{e}es dans ce cas pour les \'{e}l\'{e}ments.
R\`{e}gle de substitution d’un UpdateGroup
Les \'{e}l\'{e}ments \'{e}quivalents \`{a} ce groupe sont obtenus en
appliquant l’op\'{e}rateur d’\'{e}quivalence large ≦AdditionalPI pour le container en y ajoutant les PI Widget30. Structure et Positionnement des \'{e}l\'{e}ments graphiques
116
CHAPITRE 6. MÉCANISMES DE MIGRATION DES UI
Move, WidgetRotation. Le tableau 6.7 caract\'{e}rise la r\`{e}gle de substitution. Les \'{e}l\'{e}ments du container
sont substitu\'{e}s en appliquant l’op\'{e}rateur d’\'{e}quivalence stricte avec une diff\'{e}rence dans les donn\'{e}es
afin de privil\'{e}gier des donn\'{e}es de type Image ou Object.
SubstitutionRule
id
DisplayGroupSRule1
guideline
G11, G12, G13
constraint
True
leftMember
InteractorType du tableau 6.4)
rightMember
InteractorType du tableau 6.4)
operators
≦AdditionalPI et ∼=TargetType
avec AdditionnalPI =
 WidgetMove,WidgetRotation,
WidgetResize

et TargetType = \textbackslash{}{Image\textbackslash{}}
TABLE 6.7 – R\`{e}gle de substitution des UpdateGroups
R\`{e}gle de concr\'{e}tisation des UpdateGroups
La concr\'{e}tisation des groupes de modification de
contenu se fait comme les groupes de contr\^{o}le. En effet les objets tangibles peuvent \^{e}tre associ\'{e}s
\`{a} des UpdateGroups pour r\'{e}duire le nombre de groupes visibles dans l’espace de travail.
Nous remarquons que la duplication d’un groupe de modifications de contenus dans le but de
permettre \`{a} plusieurs personnes de modifier des contenus diff\'{e}rents est possible si le NF est adapt\'{e}.
Cependant nous pr\'{e}conisons de cacher les UpdateGroups si une UI en compte plusieurs et de les
afficher \`{a} l’aide d’un objet tangible.
Remarques
– La transformation des UpdateGroups a pour objectif de rendre accessibles ses groupes \`{a} tous
les utilisateurs. Ces groupes peuvent \^{e}tre affich\'{e}s par les utilisateurs \`{a} l’aide des objets tangibles
dans le but modifier des donn\'{e}es. Le type et la forme des objets tangibles sont choisis par les
personnes en charge de la migration. Le processus de migration garantit que chaque objet est
associ\'{e} \`{a} une fonctionnalit\'{e} ou \`{a} un groupe d’\'{e}l\'{e}ments graphiques.
– Pour \'{e}viter les probl\`{e}mes li\'{e}s \`{a} la coh\'{e}rence de la modification des donn\'{e}es, un UpdateGroup
ne peut \^{e}tre affich\'{e} plus d’une fois au m\^{e}me moment.
– Les UpdateGroups contenant des interacteurs avec des donn\'{e}es de types Boolean, Integer ou
String repr\'{e}sentent des panels de configurations par exemple. Ils sont migr\'{e}s sur les tables
interactives pour un dialogue mono-utilisateur.
Exemple d’un UpdateGroup
Nous consid\'{e}rons la figure 6.12 qui repr\'{e}sente un formulaire de changement de taille, de police et de contenu d’une cha\^{i}ne de caract\`{e}res. Le container repr\'{e}sentant le
groupe est concr\'{e}tis\'{e} en ScaterViewItem pour \^{e}tre accessible \`{a} tous les utilisateurs dans le cadre
d’une table Micorsoft PixelSense. Ce formulaire est associ\'{e} \`{a} un Tag pour l’afficher \`{a} l’aide d’un
objet dans le cadre d’une table Micorsoft PixelSense. La figure 6.13 est un UpdateGroup migr\'{e}.
6.3. TRANSFORMATIONS DU MODÈLE DE L’UI SOURCE
117
FIGURE 6.12 – Illustration d’un UpdateGroup
FIGURE 6.13 – Exemple d’un UpdateGroup sur une table interactive
6.3.1.4
Transformation des groupes mixtes
Ils regroupent les ControlGroups, des DisplayGroups ou des UpdateGroups. Ces groupes repr\'{e}sentent des fen\^{e}tres ou des containers d’une UI graphique. Les groupes mixtes sont caract\'{e}ris\'{e}s par
des containers de type R\'{e}cursif ou Racine. Chaque \'{e}l\'{e}ment d’un groupe mixte est un container de
type Homog\`{e}ne, H\'{e}t\'{e}rog\`{e}ne ou R\'{e}cursif. Cette cat\'{e}gorie de regroupement a pour objectif de consid\'{e}rer une UI graphique dans sa globalit\'{e} pour la prise en compte des guidelines li\'{e}es au partage de
l’espace de travail. Nous caract\'{e}risons ce groupe par le tableau 6.8.
InteractorType
type
type ∈\textbackslash{}{Container\textbackslash{}}
\textbackslash{}{Recursive,Root\textbackslash{}} ∈type.containerType()
∃i ∈type.contains, i ∈\textbackslash{}{Container\textbackslash{}}



Homogeneous,
Heterogeneous,
Recursive


⊆i.containerType()
dataType
type.type.contentType ∈DataType
effectivePI
type.e f fectivePI ̸= /0
TABLE 6.8 – Caract\'{e}ristiques des groupes mixtes
R\`{e}gle de substitution de MixedGroup
La transformation des groupes mixtes pour les tables interactives d\'{e}pend du type d’\'{e}l\'{e}ments qu’ils contiennent.
Si la fen\^{e}tre principale d’une application est un MixedGroup, alors dans l’objectif de rendre accessibles les diff\'{e}rents sous-groupes nous recommandons d’appliquer les r\`{e}gles de transformations
des diff\'{e}rents groupes identifi\'{e}es ci-dessus.
118
CHAPITRE 6. MÉCANISMES DE MIGRATION DES UI
Si un des sous-groupe d’un MixedGroup est aussi mixte alors le container de type Recursif correspondant \`{a} ce sous-groupe doit avoir une position fix\'{e}e pour garder dans un cadre les groupes qu’il
contient.
Exemple d’un MixedGroup
En consid\'{e}rant l’illustration des diff\'{e}rents groupes \`{a} la figure 6.7, la
fen\^{e}tre principale est un groupe mixte et le container des canevas est aussi un groupe mixte.
La transformation de la fen\^{e}tre principale permettra d’avoir des groupes de contr\^{o}le d\'{e}la\c{c}ables
et utilisables avec des objets tangibles. Les DisplayGroups et UpdateGroups de l’application seront
aussi transform\'{e}s en suivant leurs r\`{e}gles substitution et de concr\'{e}tisation d\'{e}crites ci-dessus. Le groupe
mixte contenant des canevas ne sera pas d\'{e}pla\c{c}able ou redimensionnable pour rassembler les canevas
dans un m\^{e}me groupe visible \`{a} l’\'{e}cran. Le cadre vert de la figure 6.14 ci-dessous illustre le groupe
mixte contenant les canevas. Cependant, il est possible que les dimensions d’un groupe mixte soient
\'{e}quivalentes \`{a} la surface d’affichage pour que les \'{e}l\'{e}ments contenus soient accessibles partout sur
l’\'{e}cran.
FIGURE 6.14 – Repr\'{e}sentation de l’UI CBA pour tables interactives
Remarque
– La transformation des MixedGroups a pour objectif d’accro\^{i}tre l’accessibilit\'{e} des \'{e}l\'{e}ments de
l’UI de d\'{e}part. Cette transformation permet aussi d’obtenir des \'{e}l\'{e}ments graphiques de tailles
variables. Elle concerne la migration de la dimension Structure et positionnement des UI de
d\'{e}part.
6.3.1.5
M\'{e}canismes de substitution et de concr\'{e}tisation des groupes
Les r\`{e}gles de substitution des groupes 31 favorisent une utilisation collaborative de l’UI de d\'{e}part
en permettant le d\'{e}placement des groupes d’\'{e}l\'{e}ments graphiques. L’identification des groupes est
effectu\'{e}e en se basant sur la structure et le type de PI de leurs \'{e}l\'{e}ments.
L’application de ces r\`{e}gles de substitution est faite de mani\`{e}re automatique ou interactive. Cependant les r\`{e}gles de concr\'{e}tisation sont appliqu\'{e}es par un m\'{e}canisme automatique.
La substitution automatique des groupes
permet de proposer une premi\`{e}re version de l’UI pour
la personnaliser. Le m\'{e}canisme en charge parcourt tous les interacteurs du mod\`{e}le de l’UI source et
applique les r\`{e}gles de substitution \`{a} chaque interacteur de ce mod\`{e}le afin de d\'{e}duire le mod\`{e}le cible.
Nous remarquons que plusieurs r\`{e}gles peuvent \^{e}tre appliqu\'{e}es sur les containers comportant plusieurs
types (par exemple Racine et R\'{e}cursive, Racine et Homog\`{e}ne, etc.)
31. ControlGroup, DisplayGroup, UpdateGroup, MixedGroup
6.3. TRANSFORMATIONS DU MODÈLE DE L’UI SOURCE
119
La substitution interactive des containers
a pour but de permettre \`{a} la personne en charge de la
migration des UI de les personnaliser en se basant sur l’ensemble d’\'{e}quivalences de chaque container.
En effet les op\'{e}rateurs d’\'{e}quivalences utilis\'{e}s par les r\`{e}gles de substitution permettent d’avoir plusieurs \'{e}quivalences pour un interacteur du mod\`{e}le source. Ces interacteurs \'{e}quivalents sont utilis\'{e}s
pour la personnalisation de l’UI source. Pour chaque interacteur s\'{e}lectionn\'{e} par l’utilisateur, le m\'{e}canisme lui propose la liste des interacteurs \'{e}quivalents. Par ailleurs, pour chaque r\`{e}gle de substitution
\`{a} appliquer, le m\'{e}canisme est en mesure d’identifier les guidelines correspondantes, ces guidelines
permettront aux utilisateurs moins exp\'{e}riment\'{e}s d’appr\'{e}hender facilement la migration des UI sur les
tables interactives.
La concr\'{e}tisation des groupes
permet d’avoir un rendu graphique des UI apr\`{e}s les substitutions. Ce
m\'{e}canisme cr\'{e}e des instances de chaque Widget des diff\'{e}rentes r\`{e}gles de concr\'{e}tisation en appliquant
les contraintes associ\'{e}es (par exemple les contraintes li\'{e}es \`{a} la taille des containers)
6.3.1.6
Synth\`{e}se des transformations des groupes
– L’interpr\'{e}tation des guidelines et leurs utilisations par les m\'{e}canismes de transformation se
font \`{a} travers les interacteurs et les caract\'{e}ristiques des composants graphiques (PI et type de
donn\'{e}es).
– Le m\'{e}canisme de substitution op\`{e}re de mani\`{e}re automatique ou de mani\`{e}re interactive. La
substitution interactive a pour objectif d’assister les personnes en charge de la migration des
UI. Concernant le choix des Widgets \'{e}quivalents, un m\'{e}canisme de classement de ces derniers
en prenant compte des guidelines est indispensable.
– L’ordre de transformation des groupes du mod\`{e}le de l’UI d\'{e}pend de la strat\'{e}gie du parcours
de l’arbre repr\'{e}sentant sa structure. Nous avons remarqu\'{e} que pour consid\'{e}rer les plus petits
groupes d’\'{e}l\'{e}ments contenus dans l’arbre, il faut un parcours en profondeur pour identifier les
types de groupes et appliquer les r\`{e}gles associ\'{e}es. Un parcours en largeur de l’arbre repr\'{e}sentant
le mod\`{e}le de structure de l’UI \`{a} transformer n’est pas possible car pour identifier le type d’un
groupe il faut parcourir tous ses \'{e}l\'{e}ments fils. La section 6.3.2 pr\'{e}sente les r\`{e}gles de substitution
et de concr\'{e}tisation des \'{e}l\'{e}ments simples d’un groupe.
– Le tableau 6.9 ci-dessous pr\'{e}sente les sous probl\`{e}mes trait\'{e}s par les transformations des groupes
d\'{e}crites dans cette section.
6.3.2
Transformation d’un \'{e}l\'{e}ment
Les m\'{e}canismes de transformation de la section 6.3.1 ne pr\'{e}cisent pas comment substituer les \'{e}l\'{e}ments fils qui ne sont pas des groupes. En effet une r\`{e}gle de substitution d’un UpdateGroup ne pr\'{e}cise
pas comment substituer les composants graphiques simples (tels que les labels, boutons, champs de
texte) d’un container. Cette section d\'{e}crit les r\`{e}gles de substitution et de concr\'{e}tisation d’un composant graphique simple sans alt\'{e}rer l’homog\'{e}n\'{e}it\'{e} 32, la concision 33 et favoriser le contr\^{o}le explicite 34
(cf figures 6.2 et 6.3).
Les interacteurs de type UIComponent constituent l’ensemble des \'{e}l\'{e}ments concern\'{e}s par cette
substitution. Nous cat\'{e}gorisons ces interacteurs en fonction de leurs PI effectives.
32. Les contraintes issues de ce crit\`{e}re ergonomique de conception pr\'{e}conisent d’avoir une taille minimale pour les
donn\'{e}es graphiques et cette taille doit \^{e}tre respect\'{e}e pour toutes les donn\'{e}es de l’UI
33. L’utilisation des donn\'{e}es de type Image ou vid\'{e}o favorise la compr\'{e}hension des interactions et r\'{e}duit la charge de
travail
34. L’activation d’une fonctionnalit\'{e} \`{a} l’aide d’un objet tangible par exemple est une action explicite
120
CHAPITRE 6. MÉCANISMES DE MIGRATION DES UI
Transformation des
groupes
Dialogues
Structure et
Positionnement
Style
ControlGroup
Mono-utilisateur,
Interactions Tangibles
(D11)
Regroupement (D22)
et Accessibilit\'{e} des
Menus
DisplayGroup
Multi-utilisateurs
(D13)
Regroupement (D22)
et Positionnement
(D23)
Taille variable (D31)
UpdateGroup
Mono-utilisateur,
Interactions Tangibles
(D11)
Regroupement (D22)
et
Positionnement(D23)
MixedGroup
Regroupement (D22)
et Positionnement
(D23)
TABLE 6.9 – Synth\`{e}se des probl\`{e}mes trait\'{e}s par les transformations des groupes
– Les interacteurs qui n’impl\'{e}mentent que les PI en sortie 35 et les PI en entr\'{e}e 36 peuvent \^{e}tre
class\'{e}s dans la cat\'{e}gorie des interacteurs des donn\'{e}es en sortie car ils sont transform\'{e}s par
les contraintes sur les types de donn\'{e}es 37 et sur la taille 38.
– Les interacteurs qui impl\'{e}mentent les PI en entr\'{e}e sur les donn\'{e}es 39 appartiennent \`{a} la cat\'{e}gorie des interacteurs des donn\'{e}es en entr\'{e}e car ils repr\'{e}sentent des \'{e}l\'{e}ments \'{e}ditables ou
modifiables par les utilisateurs. Ils sont transform\'{e}s suivant les contraintes issues de la guideline
“G21 : Activit\'{e}s Bloquantes”.
– Les interacteurs d’activation impl\'{e}mentent la PI Activation en plus des PI sur les Widgets 40
car ces interacteurs repr\'{e}sentent les \'{e}l\'{e}ments d’activation de fonctionnalit\'{e}, ils sont transform\'{e}s
en suivant la guideline “G12 : Structure et Positionnement des \'{e}l\'{e}ments graphiques” pour
les transformer en image interactive. La guideline “G4 : Objets tangibles et fonctionnalit\'{e}s”
permet d’associer ces interacteurs \`{a} un objet tangible dans le cas o\`{u} la guideline G12 n’est
pas appliqu\'{e}e. Le choix des interacteurs de cette transformation est laiss\'{e} \`{a} la personne en
charge de la migration, car il est le seul capable de juger de l’utilit\'{e} de la transformation pour
un interacteur d’activation. Cependant le processus l’assiste en lui proposant les interacteurs
d’activation d’une UI.
6.3.2.1
Transformation des interacteurs de donn\'{e}es en sortie
Ce sont des interacteurs de type UIComponent appartenant \`{a} une UI et qui permettent d’afficher
des donn\'{e}es d’application. Ils n’appartiennent pas aux containers de type Homog\`{e}ne, car ce cas est
trait\'{e} par la transformation des containers de ce type. Ils sont caract\'{e}ris\'{e}s par l’InteractorType du
tableau 6.10.
Ces interacteurs peuvent contenir des donn\'{e}es graphiques ou des donn\'{e}es d’application. Dans le
cas o\`{u} ils contiennent des donn\'{e}es graphiques, la transformation consiste d’abord \`{a} lui substituer un
35. Widget Display et Data Display
36. Widget Selection et Navigation
37. “G12 : Structure et Positionnement des \'{e}l\'{e}ments graphiques”
38. “G11 : Taille des \'{e}l\'{e}ments graphiques”
39. Data Edition, Data Selection, Data Move In/Out
40. cf. note36
6.3. TRANSFORMATIONS DU MODÈLE DE L’UI SOURCE
121
InteractorType
type
type ∈\textbackslash{}{UIComponent\textbackslash{}}
dataType
type.type.contentType! = null
effectivePI
type.e f fectivePI ⊂







WidgetDisplay,
DataDisplay,
WidgetSelection,
Navigation







TABLE 6.10 – Caract\'{e}ristiques des interacteurs de donn\'{e}es en sortie
interacteur pouvant contenir des donn\'{e}es graphiques de type Images en utilisant l’op\'{e}rateur d’\'{e}quivalence stricte avec diff\'{e}rences de donn\'{e}es. Ensuite la concr\'{e}tisation et la personnalisation permettront
de le positionner manuellement et de pr\'{e}ciser une taille.
Dans le cas o\`{u} ces interacteurs contiennent des donn\'{e}es d’application, ils sont substitu\'{e}s par des
\'{e}l\'{e}ments strictement \'{e}quivalents en appliquant l’op\'{e}rateur d’\'{e}quivalence stricte.
La r\`{e}gle de substitution des interacteurs des donn\'{e}es en sortie
est d\'{e}crite en consid\'{e}rant la
contrainte issue de la guideline “G12 : Structure et Positionnement des \'{e}l\'{e}ments graphiques” qui
pr\'{e}conise de favoriser les donn\'{e}es de type image ou vid\'{e}o. Nous appliquons cette contrainte pour les
donn\'{e}es graphiques. En ce qui concerne les donn\'{e}es de l’application, il est indispensable soit de modifier le NF soit d’identifier de mani\`{e}re exhaustive l’ensemble des valeurs d’une donn\'{e}e et d’\'{e}tablir
des correspondances avec des images par exemple.
La R\`{e}gle de concr\'{e}tisation des donn\'{e}es en sortie
est d\'{e}crite en consid\'{e}rant la contrainte issue de
la guideline “G11 : Taille des \'{e}l\'{e}ments graphiques” 41.
Remarques
– La substitution de ces interacteurs se fait d’abord par un m\'{e}canisme automatique ensuite manuellement si les concepteurs souhaitent modifier les transformations propos\'{e}es.
– La concr\'{e}tisation ne permet pas le placement de ces composants graphiques. Les concepteurs
doivent positionner l’ensemble des interacteurs de ce type apr\`{e}s la concr\'{e}tisation. Cependant, il
est possible de proposer un positionnement par d\'{e}faut, mais il n’est pas possible de garantir une
coh\'{e}rence des positions propos\'{e}es car notre mod\`{e}le ne conserve pas les positions ou les liens
de proximit\'{e} entre les composants graphiques de l’UI source.
– Cette transformation concerne les probl\`{e}mes li\'{e}s aux sous-dimensions Taille et Regroupement
et de Positionnement des \'{e}l\'{e}ments graphiques.
6.3.2.2
Transformation des interacteurs des donn\'{e}es en entr\'{e}e
Ce sont des interacteurs de type UIComponent qui permettent de modifier ou d’\'{e}diter des donn\'{e}es d’une UI graphique. Ils appartiennent \`{a} tous les types de containers. Ce type d’interacteur est
caract\'{e}ris\'{e} par l’InteractorType du tableau 6.11
Ces interacteurs sont aussi caract\'{e}ris\'{e}s par les PI de modification des donn\'{e}es Data Edition, Data
Selection et Activation, Data Move Out et Activation et Data Move In.
41. Cette contrainte pr\'{e}conise une taille maximale et une taille minimale pour les donn\'{e}es et les composants graphiques
(cf figure 6.2)
122
CHAPITRE 6. MÉCANISMES DE MIGRATION DES UI
InteractorType
type
type ∈\textbackslash{}{UIComponent\textbackslash{}}
dataType
type.type.contentType! = null
effectivePI







DataEdition,
DataSelection,Activation,
DataMoveOut,
DataMoveIn







⊂type.e f fectivePI
TABLE 6.11 – Caract\'{e}ristiques des interacteurs de donn\'{e}es en entr\'{e}e et en sortie
La r\`{e}gle de substitution des interacteurs des donn\'{e}es en entr\'{e}e
est d\'{e}crite en consid\'{e}rant les
contraintes issues de la guideline “G21 : Activit\'{e} Bloquantes”. Les interactions en entr\'{e}es des UI
desktops sont des activit\'{e}s bloquantes. Sur une table interactive elles doivent permettre aux autres
utilisateurs d’acc\'{e}der \`{a} d’autres fonctionnalit\'{e}s.
Ces interacteurs sont substitu\'{e}s par les m\^{e}mes types d’interacteurs avec des types de donn\'{e}es
\'{e}quivalents en appliquant les op\'{e}rateurs d’\'{e}quivalence stricte.
La r\`{e}gle de concr\'{e}tisation des interacteurs de donn\'{e}es en entr\'{e}e
est issue de la contrainte de la
guideline “G11 : Taille des \'{e}l\'{e}ments graphiques”. La taille des Widgets doit \^{e}tre d\'{e}finie en prenant
compte de la taille de son container. Son positionnement est effectu\'{e} manuellement par la personne
en charge de la migration.
Remarques
– Les containers des interacteurs des donn\'{e}es en entr\'{e}e ne sont pas accessibles par plusieurs utilisateurs au m\^{e}me moment. Cette restriction \'{e}vite les probl\`{e}mes li\'{e}s \`{a} la coh\'{e}rence des donn\'{e}es
modifi\'{e}es.
– La transformation de ces interacteurs a pour objectif de conserver les dialogues (interactions
en entr\'{e}e) de l’UI de d\'{e}part. Cependant cette transformation ne permet pas les dialogues multiutilisateurs.
6.3.2.3
Transformation des interacteurs d’activation
Ce sont des interacteurs de type UIComponent qui permettent de d\'{e}clencher des fonctionnalit\'{e}s.
Ils n’appartiennent qu’aux containers de type Homog\`{e}ne car ce cas est trait\'{e} par la r\`{e}gle de substitution
des menus. Ce type d’interacteur est caract\'{e}ris\'{e} par l’InteractorType du tableau 6.12.
InteractorType
type
type ∈\textbackslash{}{UIComponent\textbackslash{}}
dataType
type.type.contentType! = null
effectivePI
type.e f fectivePI ⊆\textbackslash{}{Navigation, WidgetSelection, Activation\textbackslash{}}
TABLE 6.12 – Caract\'{e}ristiques des interacteurs d’activation
La r\`{e}gle de substitution des interacteurs d’activation
est issue de la contrainte qui favorise les
donn\'{e}es de type image ou vid\'{e}o 42. Ils sont substitu\'{e}s par les interacteurs de m\^{e}me type avec une dif42. “G12 : Structure et Positionnement des \'{e}l\'{e}ments graphiques”
6.3. TRANSFORMATIONS DU MODÈLE DE L’UI SOURCE
123
f\'{e}rence de type de donn\'{e}es. La r\`{e}gle de substitution des \'{e}l\'{e}ments correspondant au type d’interacteur
∼=TargetType, avec TargetType = \textbackslash{}{Image\textbackslash{}} est appliqu\'{e}e.
La r\`{e}gle de concr\'{e}tisation des interacteurs d’activation
est issue de la contrainte qui favorise
l’image et la vid\'{e}o et de la recommandation qui pr\'{e}conise l’utilisation des objets tangibles pour activer
les fonctionnalit\'{e}s. Le concepteur pr\'{e}cise les interacteurs d’activation \`{a} associer aux objets tangibles.
Remarques
– La transformation des interacteurs d’activation a pour objectif l’accessibilit\'{e} des fonctionnalit\'{e}s
par des dialogues homog\`{e}nes et par un contr\^{o}le explicite \`{a} l’aide d’objets physiques.
– Ces interacteurs ne permettent pas des dialogues multi-utilisateurs car un seul utilisateur a acc\`{e}s
\`{a} un interacteur au m\^{e}me moment.
6.3.2.4
Synth\`{e}se des transformations des interacteurs
– Les r\`{e}gles de substitution des interacteurs de type UIComponent pr\'{e}servent les PI et adaptent
les types de donn\'{e}es de la source.
– Le positionnement des \'{e}l\'{e}ments substitu\'{e}s se fait au niveau de l’UI finale par la personne en
charge de la migration.
– L’association d’un objet tangible \`{a} une fonctionnalit\'{e} se fait au niveau de l’UI finale par la
personne en charge de la migration. Le processus de migration des UI semi automatique propose
une liste d’interacteurs appartenant au type d\'{e}crit par le tableau 6.12.
– Les transformations des interacteurs propos\'{e}es dans cette section traitent les probl\`{e}mes d’accessibilit\'{e} des interacteurs et des fonctionnalit\'{e}s par des transformations automatiques et manuelles. Le tableau 6.13 pr\'{e}sente de mani\`{e}re synth\'{e}tique les probl\`{e}mes abord\'{e}s selon les dimensions des UI.
Transformation des
groupes
Dialogues
Structure et
Positionnement
Style
Interacteur des
donn\'{e}es en entr\'{e}e
Mono-utilisateur
Positionnement
Manuel (D23)
Taille d\'{e}finie
Manuellement (D31)
Interacteur des
donn\'{e}es en sortie
Mono-utilisateur
Accessibilit\'{e},
Positionnement
Manuel (D23)
Taille variable, Taille
d\'{e}finie Manuellement
(D31)
Interacteur
d’activation
Mono-utilisateur,
Interactions Tangibles
Accessibilit\'{e} des
fonctionnalit\'{e}s,
Positionnement
Manuel (D23)
Taille d\'{e}finie
Manuellement (D31)
TABLE 6.13 – Synth\`{e}se des probl\`{e}mes trait\'{e}s par les transformations des interacteurs
124
CHAPITRE 6. MÉCANISMES DE MIGRATION DES UI
6.4
Classement des \'{e}l\'{e}ments \'{e}quivalents
Nous proposons dans ce manuscrit une approche de migration des UI semi automatique qui assiste les concepteurs en proposant des options de migration des UI conform\'{e}ment aux guidelines
de la cible. Cependant le choix de certaines options peut impliquer des t\^{a}ches suppl\'{e}mentaires non
automatisables 43 pour les concepteurs.
Les strat\'{e}gies de transformation de la structure de l’UI source telles que pr\'{e}sent\'{e}es \`{a} la section 6.3
utilisent les m\'{e}canismes d’\'{e}quivalences des composants graphiques. Nos op\'{e}rateurs d’\'{e}quivalences
permettent de retrouver plusieurs \'{e}l\'{e}ments \'{e}quivalents pour un composant graphique de l’UI source.
Dans cette section, nous proposons de classer les composants graphiques \'{e}quivalents dans le but de
faciliter leur choix pour le concepteur en prenant en compte les guidelines de la cible. Par ailleurs le
choix d’un composant graphique \'{e}quivalent est guid\'{e} par l’objectif de r\'{e}duire le co\^{u}t de sa mise en
œuvre pour le concepteur.
Nous pr\'{e}sentons dans cette section deux crit\`{e}res pour classer les \'{e}l\'{e}ments \'{e}quivalents : la conformit\'{e} des composants graphiques aux guidelines et la charge de travail engendr\'{e}e par le choix d’un
\'{e}l\'{e}ment \'{e}quivalent. Le but de ces crit\`{e}res est d’accro\^{i}tre le respect des guidelines tout en r\'{e}duisant les
co\^{u}ts des interventions humaines pendant le processus de migration.
6.4.1
Conformit\'{e} des composants graphiques aux guidelines
Nous caract\'{e}risons les composants graphiques \`{a} la section 5.3 par les PI (intrins\`{e}ques et effectives)
et les \'{e}l\'{e}ments structurels (types et cardinalit\'{e}s des donn\'{e}es). Ces caract\'{e}ristiques nous permettent
d’\'{e}valuer la conformit\'{e} de chaque composant graphique aux guidelines de la cible.
La section 6.4.1.1 pr\'{e}sente une m\'{e}thode d’\'{e}valuation de la conformit\'{e} des composants graphiques
aux guidelines qui favorisent la collaboration \`{a} l’aide des PI. La section 6.4.1.2 pr\'{e}sente une \'{e}valuation de la conformit\'{e} aux guidelines qui favorisent la collaboration \`{a} l’aide des caract\'{e}ristiques
structurelles des composants graphiques. La section 6.4.1.3 pr\'{e}sente la m\'{e}thode d’\'{e}valuation de la
conformit\'{e} d’un composant graphique aux guidelines en prenant en compte les PI et les caract\'{e}ristiques structurelles.
6.4.1.1
Évaluation de la conformit\'{e} aux guidelines en fonction des PI
Les PI d\'{e}crivent les interactions atomiques des composants graphiques ind\'{e}pendamment des instruments d’interactions. Nous identifions trois cat\'{e}gories de PI ; celles qui sont conformes \`{a} des guidelines, celles qui ne favorisent pas la conformit\'{e} \`{a} des guidelines et celles qui sont neutres par rapport
aux guidelines.
PI favorisant le respect des guidelines :
Les PI Widget Move, Widget Rotation et Widget Resize
favorisent la mise en œuvre d’un espace de travail partag\'{e} si elles sont appliqu\'{e}es \`{a} des groupes
d’\'{e}l\'{e}ments graphiques (cf section 6.3).
La guideline “G11 : Taille des \'{e}l\'{e}ments graphique” est interpr\'{e}t\'{e}e \`{a} la figure 6.2, elle recommande des tailles variables pour les groupes d’affichage de contenus. Dans ce cas, les composants
graphiques impl\'{e}mentant la PI Widget Resize sont conformes \`{a} G11 car ils peuvent \^{e}tre redimensionn\'{e}s.
43. La transformation des donn\'{e}es textuelles en ic\^{o}nes ou images n\'{e}cessite d’associer \`{a} chaque texte une repr\'{e}sentation
graphique ayant la m\^{e}me s\'{e}mantique
6.4. CLASSEMENT DES ÉLÉMENTS ÉQUIVALENTS
125
La guideline “G13 : Comportement des \'{e}l\'{e}ments graphiques” recommande des \'{e}l\'{e}ments graphiques d\'{e}pla\c{c}ables pour favoriser leur accessibilit\'{e}. Les groupes des \'{e}l\'{e}ments graphiques impl\'{e}mentant les PI Widget Move et Widget Rotation sont conformes \`{a} la G13.
PI ne favorisant pas le respect des guidelines :
Les interactions d’\'{e}dition de contenu (Data Edition) limitent une utilisation collaborative car elles d\'{e}crivent des activit\'{e}s bloquantes pour d’autres
utilisateurs. La guideline “G21 : Activit\'{e}s Bloquantes” recommande l’utilisation des interactions non
bloquantes pour les autres utilisateurs.
PI neutres par rapport aux guidelines :
Les PI Widget Selection, Navigation, Data Selection,
Data Move In/Out, Data Display et Activation sont neutres par rapport aux guidelines des UI collaboratives et tangibles pr\'{e}sent\'{e}es \`{a} la section 6.2.
Comment \'{e}valuer la conformit\'{e} aux guidelines \`{a} l’aide des PI ?
Pour \'{e}valuer la conformit\'{e} des
composants ayant des PI aux guidelines des UI collaboratives, nous avons pond\'{e}r\'{e} toutes les PI suivant
les trois cat\'{e}gories \'{e}voqu\'{e}es pr\'{e}c\'{e}demment.
– Les PI favorisant les guidelines ont une pond\'{e}ration sup\'{e}rieure \`{a} z\'{e}ro
– Les PI neutres ont une pond\'{e}ration nulle
– Les PI ne favorisant pas les guidelines ont une pond\'{e}ration inf\'{e}rieure \`{a} z\'{e}ro
Le poids de chaque PI d\'{e}termine son importance pour une migration. Ceci dans le but de rendre
flexible le m\'{e}canisme d’\'{e}quivalences pour les personnes en charge de la migration. Les poids des PI
sont d\'{e}termin\'{e}s par la fonction
InteractionsWeight : Widget →\textbackslash{}{Interger\textbackslash{}}
et en utilisant le tableau 6.14 (ci-dessous) des poids des PI par rapport aux guidelines. Dans ce tableau,
les poids Pf > 0 expriment les PI qui favorisent des guidelines. Les poids Pn = 0 expriment les PI qui
sont neutres et les poids Pa < 0 ne favorisent pas les guidelines. La fonction InteractionsWeight
calcule la somme des poids des PI d’un Widget.
∀w ∈\textbackslash{}{Widget\textbackslash{}}∧∀pi ∈w.e f fectivePI
InteractionsWeight(w) =∑Ppi, Ppi ∈\textbackslash{}{Pa,Pj,Pn\textbackslash{}}
6.4.1.2
Évaluation de la conformit\'{e} aux guidelines en fonction des types de donn\'{e}es
Les types de donn\'{e}es des composants graphiques permettent comme les PI d’\'{e}valuer leur conformit\'{e} aux guidelines. Nous constatons que les contenus de certains types permettent de d\'{e}crire des UI
faciles \`{a} appr\'{e}hender et qui favorisent la collaboration dans le cadre des tables interactives.
Types de donn\'{e}es conseill\'{e}s pour les tables interactives
Les donn\'{e}es graphiques de type Image,
MediaElement ou Object sont adapt\'{e}es pour les tables interactives. En effet la guideline “G12 :
Structure et Positionnement des \'{e}l\'{e}ments graphiques” conseille aussi l’usage de ces types de donn\'{e}es
car ils illustrent mieux les interactions ou les informations d’une UI.
126
CHAPITRE 6. MÉCANISMES DE MIGRATION DES UI
Primitives
d’interactions
Contraintes sur le
poids
Guidelines
Widget Resize,
Widget Move, Widget
Rotation
Pf >0
G11, G13
Widget Selection,
Navigation, Data
Selection, Data Move
In/Out, Data Display,
Activation
Pn =0
/
Data Edition
Pa <0
G21
TABLE 6.14 – Poids des PI par rapport aux guidelines
Types de donn\'{e}es non conseill\'{e}s pour les tables interactives
Les contenus de type String, Boolean ou Integer ne sont pas adapt\'{e}s en tant que donn\'{e}es graphiques pour les tables interactives. En
effet ces types de donn\'{e}es n’illustrent pas les interactions mais obligent les utilisateurs \`{a} lire les donn\'{e}es d’un menu par exemple.
Comment \'{e}valuer la conformit\'{e} \`{a} l’aide des types de donn\'{e}es ?
Pour \'{e}valuer la conformit\'{e} des
composants \`{a} la guideline pr\'{e}conisant les types de donn\'{e}es, nous avons pond\'{e}r\'{e} tous les types de
donn\'{e}es suivant les deux cat\'{e}gories cit\'{e}es pr\'{e}c\'{e}demment.
– Les types de donn\'{e}es conseill\'{e}s ont une pond\'{e}ration sup\'{e}rieure \`{a} z\'{e}ro
– Les types de donn\'{e}es non conseill\'{e}s ont une pond\'{e}ration inf\'{e}rieure \`{a} z\'{e}ro.
Les poids des types des donn\'{e}es d’un composant graphique sont d\'{e}termin\'{e}s par la fonction
DataTypeWeight : Widget →\textbackslash{}{Integer\textbackslash{}}
∀w ∈\textbackslash{}{Widget\textbackslash{}},
DataTypeWeight(w) = Pi, Pi ∈\textbackslash{}{Pbl,Pin,Pst,Pim,Pme,Pob,Pat\textbackslash{}}
et en utilisant le tableau 6.15. Dans ce tableau, les donn\'{e}es de types Image, MediaElement et Object
ont des poids positifs car elles sont adapt\'{e}es aux les tables interactives. Les donn\'{e}es de types Boolean
et Integer ont aussi des pond\'{e}rations positives car elles sont d\'{e}crites pour composants graphiques qui
permettent de d\'{e}finir des valeurs discr\`{e}tes et continues des contenus qui peuvent \^{e}tre s\'{e}lectionn\'{e}s. La
s\'{e}lection est une interaction qui favorise la collaboration contrairement \`{a} l’\'{e}dition de contenus.
6.4.1.3
Évaluation de la conformit\'{e} aux guidelines
Elle est calcul\'{e}e \`{a} l’aide des deux crit\`{e}res pr\'{e}sent\'{e}s pr\'{e}c\'{e}demment. La fonction
GuidelinesRate : Widget →Integer
6.4. CLASSEMENT DES ÉLÉMENTS ÉQUIVALENTS
127
Data Type
Contraintes sur le
poids
Guidelines
Image
Pim >0
G12
MediaElement
Pme >0
G12
Object
Pob >0
G12
Boolean
Pbl <0
G12
Integer
Pin <0
G12
String
Pst <0
G12
TABLE 6.15 – Poids des types de donn\'{e}es par rapport aux guidelines
permet de faire la somme des fonctions repr\'{e}sentant les \'{e}valuations selon les PI et les types de donn\'{e}es
pour prendre en compte la conformit\'{e} aux guidelines :
∀w ∈\textbackslash{}{Widget\textbackslash{}},
GuidelinesRate(w) = InteractionsWeight(w)+DataTypeWeight(w)
6.4.1.4
Remarques
La fonction GuidelinesRate permet d’\'{e}valuer la conformit\'{e} des composants graphiques aux guidelines “G11 : Taille des \'{e}l\'{e}ments graphiques” , “G12 : Structure et Positionnement des \'{e}l\'{e}ments
graphiques”, “G13 : Comportement des \'{e}l\'{e}ments graphiques” et “G21 : Activit\'{e}s Bloquantes”. Cependant la guideline “G22 : Accessibilit\'{e} des Menus” ne peut pas \^{e}tre v\'{e}rifi\'{e}e car elle est sp\'{e}cifique
\`{a} un type de composant graphique (menus) dont les comportements sont d\'{e}crits par la guideline G13.
Les fonctions InteractionsWeight et DataTypeWeight se basent sur le tableau 6.14 et le tableau 6.15 dont les valeurs des poids de chaque PI ou de chaque type de donn\'{e}es peuvent \^{e}tre pr\'{e}cis\'{e}es
par la personne en charge de la migration.
La conformit\'{e} de certains types de donn\'{e}es entraine une charge de travail pour les personnes en
charge de la migration car elle n\'{e}cessite la mise en place de nouvelles ressources. Dans la section
suivante nous \'{e}valuons cette charge de travail pour le classement des composants graphiques \'{e}quivalents.
6.4.2
Charge de travail
Le choix d’un composant graphique conforme aux principes des guidelines peut entrainer un ajout
d’un connecteur pour l’adaptation de type de donn\'{e}es, un ajout de nouvelles donn\'{e}es pour l’interface
utilisateur, un ajout de codes suppl\'{e}mentaires pour la prise en compte des nouvelles primitives d’interactions ou encore l’utilisation des objets tangibles.
L’ajout d’un connecteur [MMP00] consiste \`{a} g\'{e}n\'{e}rer un code qui permet de faire un changement
de type entre celui de la source et de la cible, cette t\^{a}che est automatisable pour les types de donn\'{e}es Boolean, Integer et String et n’entraine pas une charge de travail pour le programmeur. Les
connecteurs sont fournis dans ce cas.
Cependant, si les donn\'{e}es de la cible sont de type Image, MediaElement ou Object alors un
ajout de donn\'{e}es suppl\'{e}mentaires est indispensable. En effet dans ce cas la substitution n’est pas
automatisable car l’utilisateur doit trouver ou cr\'{e}er ces donn\'{e}es et ensuite faire le ‘mapping’ avec les
types de donn\'{e}es de d\'{e}part. Par exemple le remplacement d’un menu classique dont les \'{e}tiquettes
128
CHAPITRE 6. MÉCANISMES DE MIGRATION DES UI
sont des cha\^{i}nes de caract\`{e}res avec un menu dont les \'{e}tiquettes sont des ic\^{o}nes (Images) n\'{e}cessite la
recherche ou la cr\'{e}ation de ces ic\^{o}nes.
Par ailleurs, dans le cas o\`{u} les composants graphiques \'{e}quivalents ont des PI suppl\'{e}mentaires par
rapport \`{a} la source, l’impl\'{e}mentation de ces PI est une t\^{a}che automatisable pour une biblioth\`{e}que
graphique car les codes \`{a} ajouter correspondent \`{a} des PI intrins\`{e}ques \`{a} impl\'{e}menter. Par exemple
l’ajout de la PI Widget Rotation implique par exemple d’avoir la propri\'{e}t\'{e} canRotate=True dans le
cas de la biblioth\`{e}que graphique Microsoft Surface.
L’association d’un objet virtuel ou d’une fonctionnalit\'{e} avec un objet tangible est une t\^{a}che automatisable pour un ensemble de types de composants graphiques. Cependant, il est n\'{e}cessaire de
limiter le nombre d’objets tangibles \`{a} utiliser dans une UI pour faciliter la prise en main de l’UI. Le
choix des fonctionnalit\'{e}s ou des objets virtuels doit \^{e}tre fait par le concepteur pendant la personnalisation. Le processus propose l’ensemble des objets virtuels et des fonctionnalit\'{e}s utilisables avec des
objets tangibles conform\'{e}ment aux r\`{e}gles de concr\'{e}tisation pr\'{e}sent\'{e}es \`{a} la section 6.3 . Cette t\^{a}che
implique un co\^{u}t pour les personnes en charge de la migration.
6.4.2.1
Évaluation de la charge de travail
Les crit\`{e}res permettant d’\'{e}valuer la charge de travail engendr\'{e}e par le choix d’un Widget sont :
– la diff\'{e}rence entre le type de donn\'{e}es de l’instance de la cible et celui de la source 44 et
– le choix des fonctionnalit\'{e}s et des objets virtuels.
À chacun de ces crit\`{e}res, le programmeur associe un co\^{u}t en se basant sur ses comp\'{e}tences pour
r\'{e}aliser une t\^{a}che. L’estimation de ce co\^{u}t peut se faire en se basant sur la technique Pomodoro [GV08]
qui est utilis\'{e}e en Extreme Programming [Bec00] pour permettre aux programmeurs d’optimiser leur
temps de programmation en \'{e}valuant au mieux le temps des diff\'{e}rentes t\^{a}ches \`{a} effectuer.
Au niveau du mod\`{e}le, les composants graphiques de la source sont des instances de type Interactor
et les composants graphiques \'{e}quivalents sont des Widgets. La fonction
WorkLoad : Interactor ×Widget →\textbackslash{}{Integer\textbackslash{}}
calcule la charge de travail engendr\'{e}e par le choix d’un Widget \'{e}quivalent \`{a} l’Interactor de la source.
C’est une somme des co\^{u}ts de chaque crit\`{e}re v\'{e}rifi\'{e} par le Widget \`{a} choisir. Les crit\`{e}res sont :
1. La n\'{e}cessit\'{e} des donn\'{e}es suppl\'{e}mentaires \`{a} la suite d’un changement de type. Ce crit\`{e}re est
v\'{e}rifi\'{e} si l’Interactor et le Widget \`{a} choisir n’ont pas le m\^{e}me type de donn\'{e}es et si le type de
donn\'{e}es du Widget est Image, MediaElement ou Object. Le co\^{u}t de ce crit\`{e}re est donn\'{e} par la
fonction
DataTypeCost : Interactor ×Widget →Integer
et \`{a} l’aide du tableau ci-dessous qui permet aussi au programmeur de pr\'{e}ciser avant la migration
le co\^{u}t de chaque op\'{e}ration.
2. La concr\'{e}tisation des primitives d’interactions intrins\`{e}ques. Ce crit\`{e}re est calcul\'{e} par la fonction
InteractionCost : Interactor ×Widget →Integer
en utilisant le tableau 6.16.
Dans le tableau 6.16, le co\^{u}t de chaque t\^{a}che manuelle est estim\'{e} par le programmeur avant de
commencer une migration. De nouvelles op\'{e}rations manuelles peuvent \^{e}tre d\'{e}finies, cette liste n’est
pas exhaustive.
44. Dans le cas des types Image, MediaElement ou Object
6.4. CLASSEMENT DES ÉLÉMENTS ÉQUIVALENTS
129
Op\'{e}rations Manuelles
Co\^{u}t estim\'{e} par le
programmeur
Nouvelles ressources de type
Image
C1 > 0
Nouvelles ressources de type
Media Element
C1 > 0
Nouvelles ressources de type
Object
C2 > 0
S\'{e}lection des fonctionnalit\'{e}s
C3 > 0
S\'{e}lection des objets virtuels
C4 > 0
TABLE 6.16 – Co\^{u}t des interventions manuelles
La charge de travail est calcul\'{e}e par la fonction :
WorkLoad (interactor, widget) = ∑Ci,
Ci ∈DataTypeCost (interactor,widget) ∪InteractionCost (interactor,widget)
6.4.3
Algorithme de classement
L’algorithme de classement d\'{e}termine le rang de chaque Widget de la classe d’\'{e}quivalence d’un
Interactor en faisant la diff\'{e}rence entre la conformit\'{e} aux guidelines et la charge de travail. Cette
diff\'{e}rence d\'{e}termine la valeur de l’objectif d’un Widget appartenant \`{a} une classe d’\'{e}quivalence. On
suppose que l’ensemble EquivalentWidgetinteractor contient les Widgets \'{e}quivalents \`{a} un Interactor.
∀i ∈EquivalentGroup∪EquivalentElement,
∀w ∈EquivalentWidgeti,
Le ‘meilleur’ Widget est obtenu en classant par ordre d\'{e}croissant la diff\'{e}rence GuidelinesRate(w)−
WorkLoad(i,w) pour les Widgets \'{e}quivalents
BestEquivalentWidget (i, EquivalentWidgeti)
(6.1)
=
max


[
∀w∈EquivalentWidgeti(i)
(GuidelinesRate(w)−WorkLoad(i,w))


6.4.3.1
Exemple de classement
La figure 6.15 repr\'{e}sente le mod\`{e}le de structure de l’UI CBA. Consid\'{e}rons quelques interacteurs
de l’exemple d\'{e}crit \`{a} la figure 6.7 ; le mod\`{e}le d’interacteurs de la figure 6.15 et les Widgets de leurs
classes d’\'{e}quivalences sont pr\'{e}sent\'{e}s par le tableau 6.17.
Consid\'{e}rons que le programmeur souhaite utiliser les Widgets les plus conformes aux principes
des guidelines et pour ce faire fixe la valeur du poids des primitives d’interactions \`{a} Pf = 2 et les co\^{u}ts
\`{a} Ci = 1, avec 0 < i < 6
Le Widget SurfaceListBox de la classe d’\'{e}quivalence de l’Interactor id=32 a pour type de donn\'{e}es
String.
130
CHAPITRE 6. MÉCANISMES DE MIGRATION DES UI
FIGURE 6.15 – Mod\`{e}le UI CBA
GuidelineRate(Sur faceListBox) =


InteractionWeight (Sur faceListBox) = 0
+
DataTypeWeight (Sur faceListBox) = 0

c f6.1
=


0
+
0


= 0
Workload (listBox,Sur faceListBox) = 0+0+0+0+0,
6.4. CLASSEMENT DES ÉLÉMENTS ÉQUIVALENTS
131
id, name
≡/ ∼=TargetType
≦AdditionalPI
/ ≲AdditionalPI,TargetType
≧EssentialPI
/ ≳EssentialPI,TargetType
id=2, name=menu
- / ScatterView / - / id=21, name=menuFile
- / ElementMenu,
SurfaceMenu / - / id=32, name=listImage
- / SurfaceListBox /
LibraryContainer,
LibraryBar
- / TABLE 6.17 – Widgets \'{e}quivalents
avec Ci = 0, 0 < i < 6 car la source et la cible ont les m\^{e}mes types de donn\'{e}es.
GuidelinesRate(Sur faceListBox) −Workload (listBox,Sur faceListBox) = 0−0
= 0
Les Widgets LibraryContainer et LibraryBar ont pour type de donn\'{e}es Object.
GuidelineRate(LibraryContainer) =


InteractionWeight (LibraryContainer)
+
DataTypeWeight (LibraryContainer)


=



0
+
2
= 2
(6.2)
Workload (listImage,LibraryContainer) = 0+0+1+0+0,
avec C2 = 1, et Ci = 0, 0 < i < 6 et i ̸= 2 car la source est de type String et la cible de type Object et
le co\^{u}t d’une nouvelle ressource est fix\'{e} \`{a} 1.
GuidelinesRate(LibraryContainer) −Workload (listImage,LibraryContainer) = 2−1
= 1
132
CHAPITRE 6. MÉCANISMES DE MIGRATION DES UI
Le Widget ScatterView a en plus les PI Widget Rotation et Widget Move.
GuidelineRate(ScatterView) =


InteractionWeight (ScatterView)
+
DataTypeWeight (ScatterView)


=


4, car InteractionWeight (ScatterView) = Pf +Pf = 2+2
+
0, car pas de type de donn´ees


= 4
(6.3)
Workload (mainMenu,ScatterView) = 0+0+0+0+0,
avec Ci = 0, 0 < i < 6 car aucune intervention manuelle n’est n\'{e}cessaire pour instancier ce container.
GuidelinesRate(ScatterView) −Workload (menu,ScatterView) = 4−2
= 2
Le Widget Grid n’a pas de PI en plus ou des types de donn\'{e}es diff\'{e}rents
GuidelineRate(Grid) =


InteractionWeight (Grid)
+
DataTypeWeight (Grid)


=


0
+
0, car pas de type de donn´ees


= 0
(6.4)
GuidelinesRate(Grid)−Workload (menu,Grid) = 0−0
= 0
id, name
Widget Equivalent
GuidelinesRate(w)-WorkLoad(i,w)
Rang
id=32, name=listImage
SurfaceListBox
0
3
LibraryContainer
1
1
LibraryBar
1
1
id=2 Name=menu
Grid
0
0
ScatterView
2
1
Remarques
– L’algorithme de classement propose le(s) meilleur(s) Widget(s) en tenant compte de la conformit\'{e} aux guidelines et de la charge de travail.
– Les \'{e}l\'{e}ments propos\'{e}s garantissent des comportements, des interactions et une structure qui
facilitent la collaboration entre les utilisateurs.
6.5. SYNTHÈSE
133
– Dans le cas o\`{u} l’algorithme de classement propose plusieurs \'{e}l\'{e}ments \'{e}quivalents pour un \'{e}l\'{e}ment de la source, le processus s\'{e}lectionne d’abord au hasard un \'{e}l\'{e}ment \'{e}quivalent car nous
sommes s\^{u}rs qu’ils respectent tous les crit\`{e}res fix\'{e}s. Ensuite l’utilisateur peut modifier l’\'{e}l\'{e}ment s\'{e}lectionn\'{e} automatiquement en se basant sur le style du composant graphique.
6.5
Synth\`{e}se
Dans ce chapitre nous avons propos\'{e} une interpr\'{e}tation des guidelines en r\`{e}gles de substitution
et de concr\'{e}tisation pour restructurer le mod\`{e}le de l’UI de d\'{e}part. Ces interpr\'{e}tations illustr\'{e}es par
les diagrammes 6.2 et 6.3 permettent de d\'{e}crire de mani\`{e}re formelle les guidelines qui favorisent la
collaboration et l’utilisation des objets tangibles.
Nous avons aussi pr\'{e}sent\'{e} les m\'{e}canismes de restructuration du mod\`{e}le de l’UI source en se
basant sur une strat\'{e}gie de substitution d’\'{e}l\'{e}ments graphiques guid\'{e}e par les guidelines. Le choix des
’meilleurs’ composants graphiques est fait gr\^{a}ce \`{a} l’algorithme de classement qui prend en compte
les guidelines et la charge de travail.
Les transformations des groupes ou des interacteurs simples ont pour but d’accro\^{i}tre l’accessibilit\'{e}
des groupes d’\'{e}l\'{e}ments et des fonctionnalit\'{e}s. Les dialogues transform\'{e}s pendant la migration restent
mono-utilisateur pour les interactions en entr\'{e}e et les activations. Toutefois les interactions en sortie
sont des dialogues multi-utilisateurs apr\`{e}s la migration. Les transformations des styles sont effectu\'{e}es
manuellement en proposant des dimensions des \'{e}l\'{e}ments graphiques.
Le mod\`{e}le r\'{e}sultant apr\`{e}s les substitutions est concr\'{e}tis\'{e} en utilisant les \'{e}l\'{e}ments de biblioth\`{e}que
graphique cible. La personnalisation de l’UI propos\'{e}e permet de d\'{e}finir d’abord un layout pour les
\'{e}l\'{e}ments des groupes de modification de contenu (UpdateGroup) et pour les groupes mixtes (MixedGroup). Ensuite la personnalisation permet l’association des objets tangibles aux fonctionnalit\'{e}s et
aux objets virtuels. Enfin cette phase permet aussi de modifier le style de l’UI en fonction des besoins
de l’utilisateur.
Les m\'{e}canismes de transformations propos\'{e}s dans ce chapitre ont pour objectifs de d\'{e}crire un
processus de migration semi automatique. L’une des \'{e}tapes manuelles de ce processus concerne le
positionnement d’\'{e}l\'{e}ments graphiques. En effet les personnes en charge de la migration doivent positionner les \'{e}l\'{e}ments fils des containers.
Le processus de migration des UI vers les tables interactives que nous proposons se d\'{e}compose
en trois \'{e}tapes :
1. L’abstraction de la structure et des interactions de l’UI source, cette \'{e}tape est automatisable
pour la plateforme source.
2. La restructuration et la concr\'{e}tisation du mod\`{e}le abstrait suivant les r\`{e}gles pr\'{e}sent\'{e}es dans ce
chapitre : cette \'{e}tape comporte d’abord une phase enti\`{e}rement automatique qui propose une version de l’UI aux concepteurs, ensuite une phase de personnalisation qui consiste \`{a} positionner
les \'{e}l\'{e}ments graphiques, \`{a} faire le mapping avec les nouvelles ressources, \`{a} associer les objets
tangibles aux objets virtuels ou aux fonctionnalit\'{e}s, etc.
3. La g\'{e}n\'{e}ration l’UI cible finale avec le NF source : l’application r\'{e}sultante est ex\'{e}cutable sur la
table interactive et cette \'{e}tape est automatisable pour la plateforme cible.
Dans le chapitre suivant nous validons les m\'{e}canismes pr\'{e}sent\'{e}s dans ce chapitre en pr\'{e}sentant
notre prototype de migration d’une UI desktop vers la table Microsoft PixelSense.
134
CHAPITRE 6. MÉCANISMES DE MIGRATION DES UI
Quatri\`{e}me partie
Exp\'{e}rimentations
135
CHAPITRE 7
Validations du processus de migration
assist\'{e}e des UI vers les tables interactives
Sommaire
7.1
Introduction . . . . . . . . . . . . . . . . . . . . . . . . . . . . . . . . . . . . . 137
7.2
Impl\'{e}mentation du processus de migration . . . . . . . . . . . . . . . . . . . . 138
7.2.1
Abstraction de l’UI source . . . . . . . . . . . . . . . . . . . . . . . . . . 139
7.2.2
Proposition d’une version des UI migr\'{e}es . . . . . . . . . . . . . . . . . . 143
7.2.3
Personnalisation des UI propos\'{e}es . . . . . . . . . . . . . . . . . . . . . . 145
7.2.4
G\'{e}n\'{e}ration de l’UI finale . . . . . . . . . . . . . . . . . . . . . . . . . . . 148
7.2.5
Remarques . . . . . . . . . . . . . . . . . . . . . . . . . . . . . . . . . . 149
7.3
Évaluation . . . . . . . . . . . . . . . . . . . . . . . . . . . . . . . . . . . . . . 149
7.3.1
Crit\`{e}res d’\'{e}valuation de l’approche propos\'{e}e . . . . . . . . . . . . . . . . 150
7.3.2
R\'{e}sultats
. . . . . . . . . . . . . . . . . . . . . . . . . . . . . . . . . . . 151
7.3.3
Interpr\'{e}tation des r\'{e}sultats . . . . . . . . . . . . . . . . . . . . . . . . . . 154
7.1
Introduction
Nous proposons un processus semi automatique de migration assist\'{e}e des UI vers les tables interactives qui est bas\'{e} sur les PI 45 et un mod\`{e}le de structure de l’UI. L’impl\'{e}mentation de notre solution
de migration permet de d\'{e}crire un processus comportant plusieurs \'{e}tapes.
La premi\`{e}re \'{e}tape consiste \`{a} abstraire le mod\`{e}le d’une UI 46 \`{a} partir d’une repr\'{e}sentation
concr\`{e}te des applications \`{a} migrer. Dans notre cas les codes sources des UI des applications \`{a} migrer constituent des \'{e}l\'{e}ments concrets de d\'{e}part et les applications sont d\'{e}crites selon un mod\`{e}le
d’architecture qui permet une s\'{e}paration entre l’UI et le NF. La deuxi\`{e}me \'{e}tape consiste \`{a} proposer
une version de l’UI cible \`{a} l’utilisateur. L’UI propos\'{e}e est constitu\'{e}e des interacteurs \'{e}quivalents s\'{e}lectionn\'{e}s automatiquement par les r\`{e}gles de substitution d\'{e}crites \`{a} la section 6.3. La troisi\`{e}me \'{e}tape
est une personnalisation de l’UI propos\'{e}e dans le but de d\'{e}crire manuellement les aspects qui ne sont
pas pris en compte par le processus automatique. En effet, le positionnement, la taille, la police et la
couleur des \'{e}l\'{e}ments concr\'{e}tis\'{e}s sont \`{a} d\'{e}crire manuellement par le concepteur pendant la personnalisation. Cette phase offre la possibilit\'{e} aux concepteurs de modifier les interacteurs de l’UI propos\'{e}e
en les rempla\c{c}ant par les r\`{e}gles de substitution. Enfin la g\'{e}n\'{e}ration de l’UI finale est une \'{e}tape qui
consiste \`{a} g\'{e}n\'{e}rer le code source ex\'{e}cutable de l’UI finale en faisant le lien avec le NF de d\'{e}part.
Pour valider ce processus de migration, nous proposons un atelier de migration assist\'{e}e des UI
vers les tables interactives qui est compos\'{e} d’un m\'{e}canisme d’abstraction des UI de d\'{e}part, un \'{e}diteur
des UI propos\'{e}es et un g\'{e}n\'{e}rateur des applications pour la cible.
45. Elles constituent un mod\`{e}le pour exprimer les dialogues entre les utilisateurs et les syst\`{e}mes ind\'{e}pendamment des
modalit\'{e}s d’interactions
46. Le mod\`{e}le d’une UI que nous utilisons d\'{e}crit la structure et les interactions (cf chapitre 5)
137
138
CHAPITRE 7. VALIDATIONS DU PROCESSUS
Dans ce chapitre nous pr\'{e}sentons \`{a} la section 7.2 une impl\'{e}mentation du processus de migration
des UI vers les tables interactives pour valider notre approche. La section 7.3 est une \'{e}valuation des
r\`{e}gles de transformation selon diff\'{e}rents types d’applications. L’\'{e}valuation de l’approche a pour but
d’\'{e}valuer les interpr\'{e}tations des guidelines dans les r\`{e}gles de transformation. Nous nous basons sur le
regroupement, l’accessibilit\'{e} des \'{e}l\'{e}ments graphiques et l’utilisation des objets tangibles.
7.2
Impl\'{e}mentation du processus de migration
FIGURE 7.1 – Processus semi automatique de migration d’une UI vers les tables interactives
Nous illustrons les diff\'{e}rents m\'{e}canismes de notre processus de migration des UI par le diagramme de flux (cf figure 7.1).
La phase d’abstraction des UI est d\'{e}crite \`{a} la section 7.2.1. La phase de proposition d’une UI est
d\'{e}crite \`{a} la section 7.2.2, la phase de personnalisation des UI propos\'{e}es est d\'{e}crite \`{a} la section 7.2.3
et la section 7.2.4 d\'{e}crit la phase de g\'{e}n\'{e}ration de l’UI finale.
7.2. IMPLÉMENTATION DU PROCESSUS DE MIGRATION
139
7.2.1
Abstraction de l’UI source
Cette phase a pour objectif de d\'{e}crire l’UI de l’application de d\'{e}part avec notre mod\`{e}le de l’UI.
Avant de pr\'{e}senter l’algorithme d’abstraction, nous d\'{e}crivons \`{a} la sous section 7.2.1.1 l’architecture
des applications de la source. La sous section 7.2.1.2 pr\'{e}sente les algorithmes d’abstractions et nous
pr\'{e}sentons un exemple de mod\`{e}le de l’UI \`{a} la sous section 7.2.1.3
7.2.1.1
Architecture des applications \`{a} migrer
Les applications \`{a} migrer sont d\'{e}crites selon un mod\`{e}le d’architecture MVC [KP+88]. Nous choisissons ce mod\`{e}le car il permet une s\'{e}paration entre l’UI et le NF d’une application et par rapport \`{a}
ARCH [Dev92] ou \`{a} PAC-Amodeus [Nig94], l’architecture MVC est impl\'{e}ment\'{e}e \`{a} l’aide de plusieurs technologies largement r\'{e}pandues (ASP .Net, JSF, Struts, Spring MVC, etc.).
FIGURE 7.2 – Architecture des applications migr\'{e}es vers les tables interactives
L’architecture MVC que nous d\'{e}crivons \`{a} la figure 7.2 est un choix de l’impl\'{e}mentation de notre
prototype qui permet de d\'{e}crire cadre pour les applications \`{a} migrer. L’objectif de cette architecture
est de faciliter la migration de la vue sans modifier le contr\^{o}leur et le mod\`{e}le de l’UI de d\'{e}part.
Le mod\`{e}le est repr\'{e}sent\'{e} sur la figure 7.2 par le cadre vert intitul\'{e} Model et comporte des composants logiciels qui peuvent faire appel \`{a} la vue pour le mettre \`{a} jour. Les fl\`{e}ches \'{e}tiquet\'{e}es «updateView» constituent des interactions en sortie.
Le contr\^{o}leur est illustr\'{e} par le cadre vert intitul\'{e} Controler. Le contr\^{o}leur fait appel au mod\`{e}le par des interactions en entr\'{e}e illustr\'{e}es par les fl\`{e}ches \'{e}tiquet\'{e}es «inputModel». Par ailleurs, les
contr\^{o}leurs font appel aussi \`{a} la vue gr\^{a}ce aux fl\`{e}ches \'{e}tiquet\'{e}es «updateView» pour mettre \`{a} jour des
donn\'{e}es ou pour afficher des \'{e}l\'{e}ments graphiques par exemple.
La vue est repr\'{e}sent\'{e}e par un cadre intitul\'{e} View, elle est constitu\'{e}e de deux parties : la partie
d\'{e}crivant la structure, le positionnement et le style (UIStructure) et la partie d\'{e}crivant les comporte140
CHAPITRE 7. VALIDATIONS DU PROCESSUS
ments (UIBehavior) repr\'{e}sentant les interactions en entr\'{e}e et en sortie. La partie UIBehavior de la vue
dispose d’une interface (AbstractView) qui d\'{e}crit les m\'{e}thodes de mise \`{a} jour de la vue appel\'{e}es par
le mod\`{e}le ou le contr\^{o}leur. Cette interface est impl\'{e}ment\'{e}e par des classes de type ConcretView qui
utilisent les \'{e}l\'{e}ments de la biblioth\`{e}que graphique de d\'{e}part. La migration de l’UI dans ces conditions
consiste \`{a} remplacer les classes concr\`{e}tes, les \'{e}v\'{e}nements et la structure de l’UI de d\'{e}part par les
\'{e}l\'{e}ments de la biblioth\`{e}que d’arriv\'{e}e.
Les rectangles en vert repr\'{e}sentent les \'{e}l\'{e}ments des applications qui ne sont pas modifi\'{e}s pendant
la migration. Les \'{e}l\'{e}ments du mod\`{e}le et du contr\^{o}leur des applications de d\'{e}part sont conserv\'{e}s et
r\'{e}utilis\'{e}s pour les applications cibles.
UIStructure
Cette partie de la vue est constitu\'{e}e des instances des \'{e}l\'{e}ments de la biblioth\`{e}que
graphique de la plateforme de d\'{e}part. Notre impl\'{e}mentation consid\`{e}re les applications WPF, dans ce
cas la structure, le positionnement et le style sont repr\'{e}sent\'{e}s par des \'{e}l\'{e}ments d\'{e}crits en XAML.
Les fl\`{e}ches \'{e}tiquet\'{e}es «event» de notre architecture (cf la figure 7.2) repr\'{e}sentent les interactions
en entr\'{e}e sur chaque composant graphique. Les «event» repr\'{e}sentent les \'{e}v\'{e}nements impl\'{e}ment\'{e}s
(ImplementedEvents) d’un mod\`{e}le de la structure d’une instance des UI (cf figure 5.5). Les PI effectives sont identifi\'{e}es \`{a} partir des \'{e}v\'{e}nements impl\'{e}ment\'{e}s selon les r\`{e}gles d\'{e}crites \`{a} la section 5.3.2.6.
UIBehaviour
Cette partie de la vue est constitu\'{e}e des handlers des \'{e}v\'{e}nements d\'{e}clench\'{e}s par les
\'{e}l\'{e}ments de la structure et par les m\'{e}thodes de mise \`{a} jour des donn\'{e}es d’applications pr\'{e}sent\'{e}es
par la structure. Les interactions d’activation ou d’entr\'{e}e vers le NF sont repr\'{e}sent\'{e}es par les fl\`{e}ches
«inputControler», ces interactions sont repr\'{e}sent\'{e}es par la PI Activation dans notre mod\`{e}le. En effet
les appels de m\'{e}thodes appartenant aux contr\^{o}leurs constituent des interactions en entr\'{e}e. La section A.2.6 de l’annexe A d\'{e}crit le comportement de l’UI de l’application CBA.
Les interactions en sortie du NF sont repr\'{e}sent\'{e}es par les fl\`{e}ches «updateView» (de l’architecture
d\'{e}crite \`{a} la figure 7.2), ces interactions sont abstraites en PI Data Display et Widget Display dans
notre mod\`{e}le. La notification de la vue peut \^{e}tre faite par le mod\`{e}le ou par le contr\^{o}leur ; dans les
deux cas, il existe une vue abstraite (AbstractView) qui d\'{e}crit les m\'{e}thodes de notification de la vue
ind\'{e}pendamment des composants graphiques.
L’interface AbstractView fait partie des \'{e}l\'{e}ments de la vue qui ne sont pas modifi\'{e}s pendant la
migration de l’UI. Elle d\'{e}crit les m\'{e}thodes appel\'{e}es par le mod\`{e}le ou le contr\^{o}leur. Elle permet donc
de faire lien entre l’UI et le NF. Le listing A.3 de l’annexe A pr\'{e}sente un exemple de vue abstraite.
Les fl\`{e}ches \'{e}tiquet\'{e}es «updateWidget» de l’architecture d\'{e}crite \`{a} la figure 7.2 repr\'{e}sentent les
impl\'{e}mentations des m\'{e}thodes de l’interface AbstractView. Ces interactions repr\'{e}sentent les modifications des donn\'{e}es d’application ou les changements des propri\'{e}t\'{e}s graphiques d’une UI \`{a} partir du
contr\^{o}leur ou du mod\`{e}le. La section A.2.7 de l’annexe A pr\'{e}sente un exemple d’AbstractView.
Remarques
– L’architecture MVC que nous avons pr\'{e}sent\'{e}e nous permet de distinguer les \'{e}l\'{e}ments \`{a} modifier
pendant la migration des \'{e}l\'{e}ments \`{a} conserver des UI de la source.
– Les algorithmes d’abstraction parcourent les \'{e}l\'{e}ments \`{a} modifier et particuli\`{e}rement les \'{e}l\'{e}ments de type UIStructure, UIBehavior et les impl\'{e}mentations des vues abstraites pour extraire
la structure et les PI effectives.
– On consid\`{e}re qu’une UIStructure est une structure analysable qui est constitu\'{e}e d’un ensemble
7.2. IMPLÉMENTATION DU PROCESSUS DE MIGRATION
141
d’arbres (Tree) qui d\'{e}crit les diff\'{e}rentes fen\^{e}tres de l’UI \`{a} migrer.
UIStructure =
N
[
i=0
Treei, Treei = ⟨root, Node⟩
7.2.1.2
M\'{e}canismes d’abstraction
L’algorithme 2 parcourt toutes les structures repr\'{e}sent\'{e}es par des arbres et applique un algorithme
DFSAbstraction pour parcourir l’arbre en profondeur.
Algorithme 1 : DFSAbstraction
Donn\'{e}es : WidgetName : ensemble des noms des widgets d’une biblioth\`{e}que graphique,
TabImpEvent : Tableaux de correspondances des ImplementedEvent, EffPIIdRule :
R\`{e}gles d’identification des PI effectives
Entr\'{e}es : Tree, root
Sorties : interactor : Interactor
si root.child! = Null ∧∃ci ∈root.child ∧ci.name ∈WidgetName alors
pour ci ∈root.child faire
si NotMarked(ci) alors
DFSAbstraction(Tree,ci)
fin
fin
//La racine de l’arbre est un container si les fils sont des Widgets //Initialiser interactor avec
un Container \`{a} partir de root et des tableaux de correspondances de la biblioth\`{e}que cible;
interactor ←InitializeContainer(root,TaImpEvent);
sinon
//La racine est un interacteur simple sans fils, initialiser interactor avec un //UIComponent
\`{a} partir de root et des tableaux de correspondances de la biblioth\`{e}que cible;
interactor ←InitializeUIComponent(root,TabImpEvent);
fin
si interactor != Null alors
//Appliquer les r\`{e}gles d’identification des PI effectives (cf section 5.3.2.6);
interactor.E f fectivePI.add(FoundE f fectivePI(interactor,E f fPIIdRule))
fin
Mark(root);
L’algorithme DFSAbstraction parcourt en profondeur la structure de l’arbre repr\'{e}sentant l’UI de
d\'{e}part. Cet algorithme marque en premier lieu tous les fils d’un container dans le but de d\'{e}terminer
son type (Homog\`{e}ne, H\'{e}t\'{e}rog\`{e}ne, R\'{e}cursif ou Racine).
142
CHAPITRE 7. VALIDATIONS DU PROCESSUS
Algorithme 2 : Abstraction
Donn\'{e}es : AUISource : AUIStructure (Mod\`{e}le de structure de l’instance de l’UI source)
Entr\'{e}es : UIStructure : Structure de l’UI de d\'{e}part
Sorties : ModelSource : Mod\`{e}le de l’UI source (fichier XMI d\'{e}crivant le mod\`{e}le)
pour Treei ∈UIStructure faire
//Parcourir chaque nœud racine de la structure de l’UI source
pour rooti ∈Treei faire
AUISource.contains ←DFSAbstraction(Treei, rooti);
fin
fin
ModelSource ←SerializeXMI(AUISource)
7.2.1.3
Applications
Nous avons impl\'{e}ment\'{e} les mod\`{e}les des Widgets et des Interactors et les PI \`{a} l’aide du framework
de mod\'{e}lisation EMF [Fou13]. EMF offre la possibilit\'{e} de g\'{e}n\'{e}rer les codes de manipulation de ces
mod\`{e}les. Dans notre cas, les manipulations du mod\`{e}le de Widgets consistent \`{a} :
– v\'{e}rifier si un \'{e}l\'{e}ment du code source appartient \`{a} une biblioth\`{e}que graphique,
– instancier un composant graphique dans une UI finale \`{a} partir d’un Widget \'{e}quivalent.
Cette derni\`{e}re manipulation est utilis\'{e}e pour concr\'{e}tiser un composant graphique \`{a} partir de son type.
Les manipulations concernant les Interactors consistent \`{a} :
– instancier un Container ou un UIComponent \`{a} partir d’une UI finale, en pr\'{e}cisant les valeurs
des attributs.
– remplacer un interacteur par un autre dans la structure de l’UI en appliquant les r\`{e}gles de substitution.
– comparer un Interacteur avec un Widget en prenant en compte les PI et les types de donn\'{e}es
pendant les recherches d’\'{e}quivalences.
Nous avons appliqu\'{e} l’Algorithme 2 Abstraction sur l’UI de l’application CBA. Le mod\`{e}le d’une
UI r\'{e}sultant est d\'{e}crit par le listing 7.1 repr\'{e}sentant le mod\`{e}le au format XMI.
Listing 7.1 – Mod\`{e}le d’une UI
1 <?xml version="1.0" encoding="ASCII"?>
<AUIStructure name="SelecteurContenuWPF">
<contains name="MainWindow">
<Container id="1" name="Window0" containerType="Window" >
<Container id="11" name="MenuPane">
<EffectivePI>
<PrimitiveInteractions name="WidgetNavigation"/>
<PrimitiveInteractions name="WidgetSelection"/>
</EffectivePI>
<Container id="111" name="menuFile">
11
<Content propertyName="Header" value="File" />
<UIComponent id="1111" name="Open" >
<PrimitiveInteractions name="WidgetNavigation"/>
<PrimitiveInteractions name="WidgetSelection"/>
<PrimitiveInteractions name="Activation"/>
<events>
<ImplementedEvent name="Open\textbackslash{}_Click" type="Call" >
<inputDeviceType="DirectManipulationType" >
<inputDeviceType="SequentialManipulationType" >
</ImplementedEvent>
21
</events>
</UIComponent>
<... />
7.2. IMPLÉMENTATION DU PROCESSUS DE MIGRATION
143
</Container>
</Container>
<... />
<Container id="4" name="WorkspacePane" >
<Container id="44" name="Canvas1">
<PrimitiveInteractions name="WidgetNavigation"/>
<PrimitiveInteractions name="WidgetSelection"/>
31
<PrimitiveInteractions name="DataMoveIn"/>
<PrimitiveInteractions name="Activation"/>
<events>
<ImplementedEvent name="Canvas\textbackslash{}_Drop" type="Change" property="
Content" >
<inputDeviceType="DirectManipulationType" >
</ImplementedEvent>
</events>
</Container>
<... />
41
</Container>
</Container>
</contains>
</AUIStructure>
Le mod\`{e}le d’une UI obtenu apr\`{e}s abstraction est transform\'{e} pour proposer une premi\`{e}re version de
l’UI pour la cible. Dans la section suivante nous pr\'{e}sentons les diff\'{e}rents modules de notre processus
de migration.
7.2.2
Proposition d’une version des UI migr\'{e}es
Le m\'{e}canisme de proposition d’une premi\`{e}re version de l’UI se base sur un algorithme de g\'{e}n\'{e}ration de mod\`{e}le d’une UI qui utilise les r\`{e}gles de substitution et l’algorithme de classement des
Widgets (cf section 6.4). Ce m\'{e}canisme est automatique et prend en entr\'{e}e un mod\`{e}le d’une UI.
L’algorithme 3 DFSProposition parcourt une structure arborescente de l’UI de d\'{e}part afin de g\'{e}n\'{e}rer une structure de l’UI cible en utilisant les r\`{e}gles de substitution. Cet algorithme est utilis\'{e} par
l’algorithme 4 de proposition de mod\`{e}le pour parcourir toutes les structures d’un mod\`{e}le source.
144
CHAPITRE 7. VALIDATIONS DU PROCESSUS
Algorithme 3 : DFSProposition
Donn\'{e}es : EquivWidget : Dictionnaire pour contenir les Widgets \'{e}quivalents, TargetLibrary :
Ensemble des Widgets de la cible, SetofSRule :Ensemble des r\`{e}gles de substitution
R\'{e}sultat : EquivWidget
Entr\'{e}es : AUISource : AUIStructure, intSource : Interactor, nodeTarget : Container, subRule :
SubstitutionRule
Sorties : intTarget :Interactor
si nodeTarget = Null alors
nodeTarget ←intTarget
fin
si intSource.contains ̸= Null ∧∃ci ∈intSource.contains alors
//intSource est un Container
pour ci ∈intSource.contains faire
si NotMarked(ci) alors
targetParent.contains ←DFSProposition(AUISource,ci,targetParent)
fin
fin
fin
//initialisation de intTarget \`{a} partir des widgets \'{e}quivalent \`{a} intSource
si intSource /∈EquivWidget.keySet() alors
EquiWidget.put(intSource,Null);
∀widgeti ∈TargetLibrary;
intFromWidgeti ←IntializeFromWidget(widgeti);
si (∃sr ∈SetofSRule, intSource ∈sr.le ftMember, intFromWidgeti ∈sr.rightMember) ∧
(∃op ∈sr.operators, intSourceopintFromWidgeti) alors
EquiWidget.get(intSource).add(widgeti);
fin
fin
intTarget ←IntializeFromWidget(BestEquivalentWidget(EquivWidget(intSource))) ;
Mark(intSource);
Algorithme 4 : Proposition du Mod\`{e}le UI
Entr\'{e}es : ModelSource (fichier au format XMI)
Sorties : ModelTarget (fichier au format XMI)
AUISource ←DeserializeXMI(ModelSource) ;
AUITarget ←InitilizeModel() ;
pour Treei ∈AUISource.contains faire
//Parcourir chaque nœud racine de la structure de l’UI source
pour rooti ∈Treei faire
AUITarget.contains.add(DFSProposition(AUISource, rooti,Null)) ;
ModelTarget ←SerializeXMI(AUITarget) ;
7.2. IMPLÉMENTATION DU PROCESSUS DE MIGRATION
145
7.2.3
Personnalisation des UI propos\'{e}es
La phase de personnalisation des UI propos\'{e}es est constitu\'{e}e des algorithmes et des processus
suivants :
– L’afficheur et l’\'{e}diteur de l’UI propos\'{e}e permettent de modifier les substitutions effectu\'{e}es
automatiquement et de d\'{e}finir le positionnement et le style des UI propos\'{e}es. Les modifications
qui concernent le style et le positionnement sont appliqu\'{e}es uniquement sur la structure XAML
de l’UI au niveau de l’afficheur, elles n’impactent pas le mod\`{e}le de l’UI.
– Le m\'{e}canisme de modification du mod\`{e}le d’une UI cible prend en compte les actions de l’\'{e}diteur li\'{e}es \`{a} la structure ou aux PI.
– Les algorithmes de g\'{e}n\'{e}ration de la structure de l’UI cible utilisent sa biblioth\`{e}que graphique.
Ils sont constitu\'{e}s de l’algorithme 5 et de l’algorithme 6
Algorithme 5 : Concr\'{e}tisationXAML
Entr\'{e}es : ModelTarget : FichierXMI, subRule : SubstitutionRule
Sorties : UITarget : SN
i=0 Treei
AUITarget ←Deserialize(ModelTaget);
pour ∀window ∈AUITarget.contains faire
pour ∀interactori ∈window.contains faire
Treei ←DFSConcretisation(interactori)
Algorithme 6 : DFSConcretisation
Donn\'{e}es : EquivWidget : Dictionnaire pour contenir les Widgets \'{e}quivalents,
SetofCRule :Ensemble des r\`{e}gles de concr\'{e}tisation, nodeXaml : repr\'{e}sente un
composant graphique de l’UI finale
Entr\'{e}es : window : Container,
Sorties : nodeXaml : structure des \'{e}l\'{e}ments en XAML repr\'{e}sentant le container fourni en
entr\'{e}e
si window.contains ̸= Null alors
//intSource est un Container
pour interactori ∈window.contains faire
si NotMarked(interactori) alors
nodeXaml.child.add(DFSConcretisation(interactori))
fin
fin
fin
//initialisation de nodeXAML
si ∃rule ∈SetofCRule,window ∈rule.le ftMember alors
nodeXaml ←IntializeFromInteractor(window,rule)
fin
Mark(window);
146
CHAPITRE 7. VALIDATIONS DU PROCESSUS
7.2.3.1
Afficheur et \'{e}diteur des UI propos\'{e}es
– Ce m\'{e}canisme affiche la structure XAML de l’UI propos\'{e}e en instanciant les composants graphiques avec des tailles pr\'{e}d\'{e}finies. Le positionnement des \'{e}l\'{e}ments de l’UI propos\'{e}e est d\'{e}fini
par la personne en charge de la migration.
– L’\'{e}diteur des UI propos\'{e}es permet d’effectuer les actions suivantes :
- S\'{e}lectionner et positionner des \'{e}l\'{e}ments graphiques
- D\'{e}finir les tailles, les couleurs et les polices des \'{e}l\'{e}ments graphiques
- Substituer les \'{e}l\'{e}ments graphiques en utilisant des Widgets \'{e}quivalents
- Supprimer des \'{e}l\'{e}ments graphiques
- Associer les \'{e}l\'{e}ments graphiques aux objets tangibles en utilisant des Tags
- Concr\'{e}tiser l’UI personnalis\'{e}e.
Impl\'{e}mentation de la s\'{e}lection et du d\'{e}placement des \'{e}l\'{e}ments graphiques
Nous avons impl\'{e}ment\'{e} cette fonctionnalit\'{e} en d\'{e}crivant un cadre pour chaque composant graphique affich\'{e} par notre
\'{e}diteur. Les cadres sont aussi des composants graphiques qui impl\'{e}mentent des \'{e}v\'{e}nements pour
s\'{e}lectionner un composant graphique de l’UI \`{a} migrer. Ces cadres permettent aussi de modifier la
position d’un \'{e}l\'{e}ment graphique en changeant la valeur des propri\'{e}t\'{e}s de positionnement des \'{e}l\'{e}ments qu’elles contiennent. La figure 7.3 pr\'{e}sente une illustration d’un cadre contenant un bouton
s\'{e}lectionn\'{e}.
FIGURE 7.3 – Exemple de cadre de s\'{e}lection des \'{e}l\'{e}ments graphique dans l’\'{e}diteur du prototype
Impl\'{e}mentation de la modification du style
Cette fonctionnalit\'{e} permet \`{a} la personne en charge de
la migration de modifier la taille (largeur et hauteur) d’un \'{e}l\'{e}ment graphique apr\`{e}s l’avoir s\'{e}lectionner.
Les cadres qui contiennent les \'{e}l\'{e}ments graphiques \`{a} \'{e}diter permettent \`{a} la personne en charge de la
migration de redimensionner les composants graphiques par un geste de d\'{e}placement d’un coin gris
du cadre s\'{e}lectionn\'{e}. Cette fonctionnalit\'{e} permet de d\'{e}finir une couleur pour les \'{e}l\'{e}ments graphiques
s\'{e}lectionn\'{e}s.
Impl\'{e}mentation de la substitution d’un \'{e}l\'{e}ment graphique
Cette fonctionnalit\'{e} permet de remplacer un composant graphique par les \'{e}l\'{e}ments de sa classe d’\'{e}quivalence. La substitution se fait
aux niveau des instances (structure XAML de l’UI) et elle est r\'{e}percut\'{e}e au niveau du mod\`{e}le de
l’UI. Nous illustrons \`{a} la figure 7.4 un exemple de substitution d’un LibraryBar avec un composant
graphique SurfaceListBox.
Impl\'{e}mentation de la suppression d’un \'{e}l\'{e}ment graphique
Cette fonctionnalit\'{e} a pour objectif de
permettre aux concepteurs de modifier l’UI d’arriver en supprimant les composants graphiques qu’ils
7.2. IMPLÉMENTATION DU PROCESSUS DE MIGRATION
147
FIGURE 7.4 – Exemple de substitution de LibraryBar par SurfaceListBox
jugent non indispensables sur les tables interactives. Par exemple une barre de menus qui comportent
des raccourcis peut \^{e}tre supprimer si le menu principal est plus accessible apr\`{e}s la migration.
Impl\'{e}mentation de l’ajout d’un objet tangible
Cette fonctionnalit\'{e} permet \`{a} la personne en charge
d’associer un tag pour afficher/cacher un groupe d’\'{e}l\'{e}ments graphiques (menu, formulaire, etc). L’\'{e}diteur permet d’abord aux concepteurs de s\'{e}lectionner le type de tags \`{a} associer \`{a} un groupe d’\'{e}l\'{e}ments
graphiques. Ensuite il pr\'{e}cise le num\'{e}ro de s\'{e}rie et le syst\`{e}me se charge de v\'{e}rifier si le tag associ\'{e} est
d\'{e}j\`{a} utilis\'{e} dans l’UI en cours de migration. Cette v\'{e}rification se fait en parcourant la structure XAML
pr\'{e}sente dans l’\'{e}diteur. En effet la biblioth\`{e}que graphique de la table PixelSense permet gr\^{a}ce au composant graphique TagVisualizer d’associer un num\'{e}ro de tag \`{a} un groupe d’\'{e}l\'{e}ments graphiques. La
figure 7.5 pr\'{e}sente le menu d’association d’un tag \`{a} un container par exemple. Cette fonctionnalit\'{e}
ajoute un TagVisualizer dans la structure XAML affich\'{e}e par l’\'{e}diteur.
FIGURE 7.5 – Menu d’association d’un tag avec un container
Impl\'{e}mentation de la concr\'{e}tisation de la structure de l’UI finale
Cette fonctionnalit\'{e} s\'{e}rialise la
structure XAML de l’UI personnalis\'{e}e \`{a} l’aide de l’\'{e}diteur. Le m\'{e}canisme en charge de la s\'{e}rialisation
enl\`{e}ve les cadres des composants graphiques et g\'{e}n\`{e}re un fichier XAML qui comporte la position et le
style de l’UI cible. Le lien avec les contr\^{o}leurs et le mod\`{e}le est effectu\'{e} pendant la phase de g\'{e}n\'{e}ration
de l’UI finale.
Éditeur graphique du prototype
La figure 7.6 ci-dessous pr\'{e}sente une capture d’\'{e}cran de l’\'{e}diteur
graphique pour personnaliser l’UI propos\'{e}e. Cette figure pr\'{e}sente les composants graphiques de l’UI
148
CHAPITRE 7. VALIDATIONS DU PROCESSUS
CBA migr\'{e}e pour les tables interactives avant l’introduction du layout et du style.
FIGURE 7.6 – Éditeur graphique du prototype
7.2.3.2
M\'{e}canisme de modification du mod\`{e}le d’une UI
– Ce m\'{e}canisme est repr\'{e}sent\'{e} par le processus “Personnalisation du Mod\`{e}le d’UI” du diagramme de flux de la figure 7.1. Il a pour objectif de garder une coh\'{e}rence entre le mod\`{e}le
d’une UI et sa repr\'{e}sentation graphique dans l’afficheur et l’\'{e}diteur.
– Les op\'{e}rations de substitution ou de suppression des \'{e}l\'{e}ments graphiques modifient la structure
et les interactions du mod\`{e}le d’une UI.
– Nous avons impl\'{e}ment\'{e} ce m\'{e}canisme \`{a} l’aide du pattern Observer-Observable. Les actions
des utilisateurs dans l’\'{e}diteur d’une UI sont toutes observables. Le m\'{e}canisme de modification
du mod\`{e}le d’une UI dispose d’un Observer abonn\'{e} aux actions de l’\'{e}diteur.
7.2.4
G\'{e}n\'{e}ration de l’UI finale
La g\'{e}n\'{e}ration de l’UI finale consiste \`{a} s\'{e}rialiser la structure XAML 47 de l’UI affich\'{e}e par l’\'{e}diteur
et ensuite \`{a} g\'{e}n\'{e}rer les fichiers de la partie UIBehavior repr\'{e}sent\'{e}s par les handlers et les ConcretViews.
7.2.4.1
G\'{e}n\'{e}ration de la ConcretView
Cette g\'{e}n\'{e}ration a pour objectif d’obtenir \`{a} partir de l’UIStructure modifi\'{e}e pendant la phase de
personnalisation l’UI ex\'{e}cutable de la table interactive cibl\'{e}e.
Les classes repr\'{e}sentant les ConcretViews de l’UI de d\'{e}part impl\'{e}mentent les interfaces de type
AbstractView. Les m\'{e}thodes \'{e}tiquet\'{e}es «updateView» sont impl\'{e}ment\'{e}es en utilisant les composants
47. Ce m\'{e}canisme est d\'{e}crit \`{a} la section 7.2.3
7.3. ÉVALUATION
149
graphiques de la table interactive. Ces m\'{e}thodes repr\'{e}sentent les PI Data Display et Widget Display
aux niveaux de l’UI finale.
Exemple
Si un Interactor du mod\`{e}le de l’UI impl\'{e}mente la PI effective Data Display ou Widget
Display, alors il existe une m\'{e}thode \'{e}tiquet\'{e}e «updateView» dans la ConcretView de l’UI de d\'{e}part
qui doit \^{e}tre modifi\'{e}e en rempla\c{c}ant les instances du Widget de d\'{e}part par celles de la table interactive.
Si un Interactor repr\'{e}sentant un composant graphique ListBox est substitu\'{e} par un LibraryBar
et si l’Interactor impl\'{e}mente une PI Data Display et la ConretView de l’UI de d\'{e}part impl\'{e}mente
une m\'{e}thode nomm\'{e}e updateList (qui est \'{e}tiquet\'{e}e «updateView»), alors la m\'{e}thode updateList sera
modifi\'{e}e pour remplacer toutes les instances du ListBox par celles du LibraryBar pour g\'{e}n\'{e}rer les
ConcretViews de la cible. Cette substitution implique un changement de type de donn\'{e}es. Il est indispensable de d\'{e}crire un adaptateur qui permet de transformer les donn\'{e}es de type String en Object.
Le prototype que nous avons mis en place ne g\'{e}n\`{e}re pas automatiquement les ConcretViews. Mais
notre mod\`{e}le de l’UI permet d’identifier les classes et m\'{e}thodes de l’UI de d\'{e}part qui doivent \^{e}tre
modifi\'{e}es manuellement pour obtenir l’UI finale.
7.2.4.2
G\'{e}n\'{e}ration des handlers
Les handlers constituent une impl\'{e}mentation des PI en entr\'{e}e. Ils permettent de faire le lien entre
l’UIStructure et les contr\^{o}leurs.
La classe contenant les handlers est g\'{e}n\'{e}r\'{e}e \`{a} partir du mod\`{e}le de l’UI et en consid\'{e}rant les
ImplementedEvents de chaque Interactor du mod\`{e}le de l’UI. En effet, les eventTypes Call, Change et
Select des ImplementedEvents du mod\`{e}le de l’UI sont traduits en handlers dans un fichier C\textbackslash{}#.
Exemple
Si un Interactor repr\'{e}sentant un SurfaceListBox a un ImplementedEvent Call et Select ,
alors un \'{e}v\'{e}nement SelectionChanged peut \^{e}tre cr\'{e}\'{e} et sont contenu correspond \`{a} l’\'{e}v\'{e}nement du
Widget de l’UI de d\'{e}part. Les \'{e}v\'{e}nements font appel \`{a} des m\'{e}thodes du contr\^{o}leur.
7.2.5
Remarques
L’architecture d\'{e}taill\'{e}e des applications \`{a} migrer que nous avons pr\'{e}sent\'{e}e dans cette section
permet d’identifier les \'{e}l\'{e}ments de l’UI \`{a} abstraire et \`{a} transformer. Notre mod\`{e}le d’une UI permet
une abstraction de la partie UIStructure de la vue. La partie UIBehavior permet d’identifier les PI
effectives en se basant sur une impl\'{e}mentation les r\`{e}gles d’identification d\'{e}crites \`{a} la section 5.3.2.6.
Nous avons impl\'{e}ment\'{e} les diff\'{e}rentes r\`{e}gles de transformation pr\'{e}sent\'{e}es \`{a} la section 6.3. Les
algorithmes de proposition d’une version de l’UI et l’\'{e}diteur de l’UI propos\'{e}e impl\'{e}mentent les r\`{e}gles
de substitution et les r\`{e}gles de concr\'{e}tisation.
La g\'{e}n\'{e}ration automatique des ConcretViews et des handlers constituent des perspectives d’impl\'{e}mentations \`{a} moyen terme pour am\'{e}liorer notre prototype de migration.
7.3
Évaluation
L’\'{e}valuation pr\'{e}sent\'{e}e dans ce manuscrit a pour objectif d’\'{e}tudier la pertinence des r\`{e}gles de transformations pr\'{e}sent\'{e}es au chapitre 6. Pour ce faire, nous avons consid\'{e}r\'{e} quatre applications desktop
avec des structures d’une UI diff\'{e}rentes. Le choix de ces applications est motiv\'{e} par le but d’\'{e}valuer le
regroupement, l’accessibilit\'{e} des \'{e}l\'{e}ments migr\'{e}s et aussi d’\'{e}valuer l’utilisation des objets tangibles.
150
CHAPITRE 7. VALIDATIONS DU PROCESSUS
– Une application d’agenda de consultation des contacts ; nous consid\'{e}rons cette application car
elle comporte des formulaires et des tableaux. Les r\`{e}gles de transformations utilis\'{e}es par notre
processus peuvent identifier et migrer les tableaux et les formulaires de cette application tout
en garantissant leur accessibilit\'{e}.
– Une application d’album photo qui permet de consulter d’un ensemble d’images. Nous consid\'{e}rons cette application car elle d\'{e}crit des DisplayGroups qui contiennent des photos. L’UI
produite permet des dialogues multi-utilisateurs.
– Une application de dessin qui offre la possibilit\'{e} de dessiner \`{a} l’aide d’une souris en choisissant
la couleur du pinceau. La migration de cette application permet de v\'{e}rifier que les dialogues de
l’application source restent coh\'{e}rents dans un contexte multi-utilisateurs.
– Une application calculatrice ; la migration de cette application permet de v\'{e}rifier si les regroupements des \'{e}l\'{e}ments graphiques en se basant sur les PI et la structure est pertinents.
Nous caract\'{e}risons \`{a} la section 7.3.1 les crit\`{e}res de validations. La section 7.3.2 pr\'{e}sente les
r\'{e}sultats de la validation des quatre application migr\'{e}es avec notre approche. La section 7.3.3 est une
interpr\'{e}tation de ces r\'{e}sultats.
7.3.1
Crit\`{e}res d’\'{e}valuation de l’approche propos\'{e}e
Les crit\`{e}res sont le regroupement, l’accessibilit\'{e} des \'{e}l\'{e}ments graphiques et l’utilisation des objets
tangibles. Ces crit\`{e}res permettent d’\'{e}valuer uniquement le comportement de l’approche propos\'{e}e par
rapport \`{a} des applications diff\'{e}rentes r\`{e}gles de transformations.
Regroupement des \'{e}l\'{e}ments graphiques
Notre processus se base sur des containers et sur les interactions des \'{e}l\'{e}ments des containers pour d\'{e}crire un regroupement des \'{e}l\'{e}ments pendant la migration.
Ce crit\`{e}re permet de tester la pertinence de notre regroupement pour des cas diff\'{e}rents. Nous pensons
que les diff\'{e}rents types de regroupement peuvent \^{e}tre qualifi\'{e}s de trois mani\`{e}res :
– Pertinent : si tous les regroupements propos\'{e}s et les substitutions possibles sont conformes aux
attentes des concepteurs.
– Acceptable : si au moins la moiti\'{e} des regroupements propos\'{e}s et les substitutions sont
conformes aux attentes des concepteurs.
– Inexploitable : si aucune des propositions ou des substitutions est conforme aux attentes des
concepteurs.
Accessibilit\'{e} des groupes
Nous pensons que l’accessibilit\'{e} des diff\'{e}rents types de groupes 48 peut
\^{e}tre qualifi\'{e}e de trois mani\`{e}res :
– Totale : si tous les groupes identifi\'{e}s ont un mouvement de rotation et de d\'{e}placement.
– Partielle : si tous les ControlGroups au moins sont accessibles \`{a} tous les utilisateurs.
– Aucune : si aucun des groupes migr\'{e}s sont accessibles.
Objets Tangibles
Ce crit\`{e}re permet d’\'{e}valuer le m\'{e}canisme d’association des objets tangibles aux
fonctionnalit\'{e}s et aux menus. Nous qualifions ce crit\`{e}re selon la pertinence des associations propos\'{e}es.
– Pertinent : si toutes les associations propos\'{e}es peuvent \^{e}tre r\'{e}alis\'{e}es sans alt\'{e}rer l’utilisabilit\'{e}
de l’UI.
– Appr\'{e}ciable : s’il existe des associations propos\'{e}es qui peuvent \^{e}tre r\'{e}alis\'{e}es \`{a} la discr\'{e}tion du
concepteur.
– Inutile : Si la r\'{e}alisation des associations propos\'{e}es alt\`{e}re l’utilisabilit\'{e} de l’UI finale.
48. ControlGroup, DisplayGroup, UpdateGroup cf section 6.3
7.3. ÉVALUATION
151
Nous avons \'{e}valu\'{e} la migration de ces quatre applications en utilisant la fiche d’\'{e}valuation d\'{e}crite
\`{a} la figure 7.7. Nous d\'{e}crivons les r\'{e}sultats de cette \'{e}tude \`{a} la section 7.3.2.
FIGURE 7.7 – Fiche d’\'{e}valuation de la migration des applications
7.3.2
R\'{e}sultats
Dans cette section nous pr\'{e}sentons les r\'{e}sultats de l’\'{e}valuation de la migration de quatre applications.
7.3.2.1
Application Agenda
FIGURE 7.8 – Application agenda
Pour la migration de cette UI (cf figure 7.8), le processus de migration propose quatre groupes
d’\'{e}l\'{e}ments graphiques pour cette UI.
– Le menu principal constitue un ControlGroup. Les r\`{e}gles de transformation permettent de choisir un ElementMenu par exemple qui sera affich\'{e} par un tag li\'{e} \`{a} un objet tangible. Ce regroupement est pertinent car il favorise l’accessibilit\'{e} des menus.
152
CHAPITRE 7. VALIDATIONS DU PROCESSUS
– Le tableau pr\'{e}sentant les contacts constitue un DisplayGroup. Les r\`{e}gles de transformation
permettent d’utiliser un LibraryBar ou un LibraryContainer si l’on souhaite garder une pr\'{e}sentation structur\'{e}e des contacts mais avec des donn\'{e}es de type Image. Cependant, ce groupe
permet aussi de repr\'{e}senter tous les contacts dans un ScatterView de mani\`{e}re dispers\'{e}e ; dans
ce cas le panel de recherche et la liste des cat\'{e}gories peuvent \^{e}tre supprim\'{e}s. Ce regroupement
est pertinent car il permet dans les deux cas pr\'{e}sent\'{e}s d’avoir une pr\'{e}sentation plus claire des
contacts.
– Le panel de recherche des contacts et la liste des cat\'{e}gories constituent des UpdateGroups car ils
disposent des interactions en entr\'{e}e. Le panel de recherche est utile dans une pr\'{e}sentation structur\'{e}e des contacts, il facilite l’acc\`{e}s aux donn\'{e}es mais dans le cas d’une utilisation \`{a} plusieurs
il constitue un handicap pour les autres utilisateurs qui doivent attendre. Les regroupements des
UpdateGroups propos\'{e}s dans ce cas sont pertinents car ils permettent de les identifier afin de
les rendre plus accessibles ou pour les supprimer dans certains cas.
– Le MixedGroup repr\'{e}sente la fen\^{e}tre principale de l’UI. Son identification est pertinente car
elle permet de savoir les principaux \'{e}l\'{e}ments \`{a} afficher sur la table interactive.
7.3.2.2
Application Album Photo
FIGURE 7.9 – Regroupement des \'{e}l\'{e}ments de l’application Album Photo
Pour la migration de cette UI (cf figure 7.9), le processus identifie deux UpdateGroups, un
ControlGroup et un MixedGroup d’\'{e}l\'{e}ments graphiques en se basant sur la structure et les PI.
Le ControlGroup repr\'{e}sente le m\'{e}nu principal de l’application. Ce regroupement est pertinent car
il favorise l’accessibilit\'{e} des menus.
Le premier UpdateGroup est constitu\'{e} de la liste des images \`{a} afficher. Cette liste est charg\'{e}e en
s\'{e}lectionnant un r\'{e}pertoire qui contient des images. La s\'{e}lection d’un \'{e}l\'{e}ment de cette liste permet
d’afficher une image. L’identification et les \'{e}l\'{e}ments \'{e}quivalents pour sa substitution qui sont propos\'{e}s par le processus sont pertinents. En effet il est possible de substituer la liste d’images par un
LibraryBar pour afficher les images dans un carrousel. Un ScatterView par exemple permet d’afficher
toutes les images de mani\`{e}re dispers\'{e}e. L’accessibilit\'{e} des Images est accrue dans le cas d’un ScatterView car plusieurs personnes peuvent consulter les images. Le groupe d’affichage de l’image peut
\^{e}tre supprim\'{e} si l’on choisi un ScatterView par exemple.
7.3. ÉVALUATION
153
Le deuxi\`{e}me UpdateGroup est constitu\'{e} de l’image affich\'{e}e et d’un curseur pour modifier sa
taille. L’identification de ce groupe est pertinente car elle permet de le supprimer si nous utilisons
un ScatterView par exemple. Cependant les options substitution pr\'{e}sent\'{e}es sont acceptables car elles
permettent d’avoir un groupe d\'{e}pla\c{c}able mais ne peuvent pas afficher plusieurs images au m\^{e}me
moment.
Le MixedGroup repr\'{e}sente la fen\^{e}tre principale de l’UI. Son identification est pertinente car elle
permet de savoir les principaux \'{e}l\'{e}ments \`{a} afficher sur la table interactive.
7.3.2.3
Application de dessin
FIGURE 7.10 – Regroupement des \'{e}l\'{e}ments de l’application de dessin
Le processus de migration de l’UI de l’application de dessin sur la table Microsoft PixelSense (cf
figure 7.10) identifie trois groupes d’\'{e}l\'{e}ments graphiques.
Un ControlGroup qui repr\'{e}sente un panel de couleurs, un crayon et une gomme. Ce regroupement
et sa migration sur la table PixelSense en ajoutant les PI de rotation et de d\'{e}pla\c{c}ant sont pertinents
pour cette application. L’utilisation d’un objet tangible par un tag pour afficher le panel est pertinent
dans cette migration.
Un MixedGroup qui repr\'{e}sente la fen\^{e}tre principale de l’UI, son identification est pertinente car
elle permet de savoir les principaux \'{e}l\'{e}ments \`{a} afficher sur la table interactive.
Un UpdateGroup qui repr\'{e}sente la zone de dessin de l’application. Les options de substitution
propos\'{e}es pour ce groupe sont acceptables. En effet il est possible de le substituer avec un Grid en
le pla\c{c}ant au centre de l’\'{e}cran sans possibilit\'{e} de d\'{e}placer cette zone de dessin. Par ailleurs il est
possible de placer la zone de dessin dans un ScatterView pour permettre son d\'{e}placement. Cependant,
l’application dessin que nous avons migr\'{e}e sur la table PixelSense pr\'{e}sente un probl\`{e}me d’utilisabilit\'{e}
dans le cadre multi-utilisateurs. En effet, la zone de dessin permet \`{a} plusieurs personnes d’\'{e}crire, mais
avec un crayon de la m\^{e}me couleur. Un changement de couleur par un utilisateur implique aussi un
changement pour les autres utilisateurs.
Cette limite peut \^{e}tre combl\'{e}e si la table permet d’identifier les utilisateurs comme la table DiamondTouch [DL01] (cf section 3.1.1.1). En effet, l’identification permet d’associer \`{a} chaque utilisa154
CHAPITRE 7. VALIDATIONS DU PROCESSUS
teur un ControlGroup constitu\'{e} d’une palette de couleurs, d’un crayon et d’une gomme. Les changements de couleurs d’un utilisateur n’impacteront pas les autres utilisateurs dans ce cas. Par ailleurs,
cette limite s’explique car notre processus de migration ne modifie pas le NF de l’application de
d\'{e}part.
7.3.2.4
Application Calculatrice
FIGURE 7.11 – Regroupement des \'{e}l\'{e}ments de l’application de calculatrice
Le processus de migration de l’UI de cette application de calculatrice sur la table Microsoft PixelSense (cf figure 7.11) identifie quatre groupes d’\'{e}l\'{e}ments graphiques.
– Un MixedGroup repr\'{e}sentant la fen\^{e}tre principale de l’UI, son identification est pertinente car
elle permet de savoir les principaux \'{e}l\'{e}ments \`{a} afficher sur la table interactive.
– Un DisplayGroup repr\'{e}sentant l’\'{e}cran d’affichage des r\'{e}sultats.
– Deux ControlGroups repr\'{e}sentant respectivement des chiffres num\'{e}riques et les op\'{e}rations de
la calculatrice.
Les diff\'{e}rents groupes identifi\'{e}s pour cette application permettent d’avoir une calculatrice comportant trois groupes. Les regroupements propos\'{e}s dans ce cas ne sont pas pertinents, mais les substitutions permettent une reconstitution d’un groupe. L’accessibilit\'{e} des \'{e}l\'{e}ments est totale en permettant
une rotation et un d\'{e}placement de la calculatrice. L’utilisation des objets tangibles est inutile dans cette
migration car le nombre de ControlGroups n’accro\^{i}t pas la charge de travail et les ControlGroups sont
utilis\'{e}s \`{a} chaque op\'{e}ration de calcul.
7.3.3
Interpr\'{e}tation des r\'{e}sultats
Les validations que nous avons effectu\'{e}es de notre processus de migration assist\'{e}e des UI vers
les tables interactives, nous permettent d’affirmer que les regroupements et les options de substitution
propos\'{e}es ont pour objectif d’accro\^{i}tre l’accessibilit\'{e} des \'{e}l\'{e}ments graphiques pour favoriser un travail
collaboratif.
En effet la figure 7.12 ci-dessous pr\'{e}sente les \'{e}valuations de la migration des quatre applications
pr\'{e}sent\'{e}es \`{a} la section 7.3.2. Nous remarquons que les regroupements propos\'{e}s pour la migration
des applications Agenda, Album Photo et Dessin sont pertinents et ils favorisent l’accessibilit\'{e} des
\'{e}l\'{e}ments graphiques. Cependant dans le cas de l’application calculatrice les propositions de groupes
de d\'{e}part ne sont pas pertinentes mais les \'{e}quivalences propos\'{e}es permettent d’avoir une UI non
dispers\'{e}e et facilement utilisable.
7.3. ÉVALUATION
155
Nous pensons que cette diff\'{e}rence par rapport \`{a} l’application calculatrice est due \`{a} sa simplicit\'{e}.
En effet une calculatrice peut \^{e}tre utilis\'{e}e par une personne sur une table interactive sans disperser les
diff\'{e}rents groupes. La flexibilit\'{e} de notre solution permet \`{a} la personne en charge de la migration de
disperser ou non les groupes.
L’association des objets tangibles avec des groupes d’\'{e}l\'{e}ments graphiques est acceptable car elle
a pour objectif d’accro\^{i}tre l’accessibilit\'{e} des \'{e}l\'{e}ments (comme les menus) et aussi de r\'{e}duire la charge
de travail dans le cas d’une UI comportant plusieurs groupes d’\'{e}l\'{e}ments graphiques de types diff\'{e}rents.
FIGURE 7.12 – Synth\`{e}se de l’\'{e}tude de la migration des quatre applications
156
CHAPITRE 7. VALIDATIONS DU PROCESSUS
Cinqui\`{e}me partie
Conclusions et Perspectives
157
CHAPITRE 8
Conclusions et perspectives
Sommaire
8.1
Synth\`{e}se
. . . . . . . . . . . . . . . . . . . . . . . . . . . . . . . . . . . . . . . 159
8.2
Perspectives
. . . . . . . . . . . . . . . . . . . . . . . . . . . . . . . . . . . . . 161
8.2.1
A moyen terme . . . . . . . . . . . . . . . . . . . . . . . . . . . . . . . . 161
8.2.2
A long terme . . . . . . . . . . . . . . . . . . . . . . . . . . . . . . . . . 161
D
ANS ce chapitre nous faisons une synth\`{e}se des travaux pr\'{e}sent\'{e}s dans cette th\`{e}se puis nous identifions quelques perspectives de poursuite de ces travaux.
8.1
Synth\`{e}se
Nous avons observ\'{e}, en \'{e}tudiant les approches de migration des UI, au chapitre 4 qu’une approche
semi automatique est \`{a} la fois r\'{e}utilisable, flexible et permet le respect des crit\`{e}res de conception.
Une approche semi automatique de migration d’UI est r\'{e}utilisable si elle dispose des m\'{e}canismes
de transformations non sp\'{e}cifiques \`{a} une application et \`{a} une plateforme. La flexibilit\'{e} est le degr\'{e}
d’intervention de l’humain pendant la migration des UI. L’avantage d’une approche flexible de migration des UI est la possibilit\'{e} pour les concepteurs de personnaliser les UI produites. Cependant
les approches tr\`{e}s flexibles impliquent des charges de travail importantes pour les concepteurs. Par
ailleurs le respect des crit\`{e}res de conception permet de produire des UI conformes aux sp\'{e}cificit\'{e}s
de la plateforme vis\'{e}e.
Dans cette th\`{e}se proposons un atelier de migration assist\'{e}e des UI desktop vers les tables interactives. Notre solution est bas\'{e}e sur une approche semi automatique de migration des UI utilisant un
mod\`{e}le d’UI. Cette approche a pour but d’adapter les UI de d\'{e}part con\c{c}ues pour une personne en UI
favorisant la collaboration et l’utilisation des objets tangibles. Dans cette optique, nous avons suppos\'{e}
que les NF des applications de d\'{e}part seront r\'{e}utilis\'{e}s sans modification sur la cible et que l’UI est
migr\'{e}e en consid\'{e}rant la dimension des dialogues, la dimension de la structure et du positionnement
et la dimension du style. Ce sont aussi trois dimensions de probl\`{e}mes li\'{e}s \`{a} la migration des UI vers
les tables interactives.
Concernant la dimension qui adresse les probl\`{e}mes li\'{e}s \`{a} la migration des dialogues de l’UI source
vers la cible, nos travaux ont pour but de d\'{e}crire des \'{e}quivalences entre les dialogues de l’UI de d\'{e}part
et de la cible. En effet les tables interactives proposent des modalit\'{e}s d’interactions diff\'{e}rentes de
celles d’un desktop par exemple.
Ensuite nous avons identifi\'{e} la dimension adressant les probl\`{e}mes li\'{e}s \`{a} la migration de la structure et du positionnement des \'{e}l\'{e}ments graphiques d’une UI source vers une table interactive. La
disposition des \'{e}l\'{e}ments graphiques sur les tables interactives doit favoriser une utilisation \`{a} plusieurs
par exemple. Concernant cette dimension, nos travaux ont pour but de d\'{e}crire une structure \'{e}quivalente \`{a} celle de d\'{e}part mais qui prenne en compte les guidelines des tables interactives.
Enfin nous avons d\'{e}crit la dimension exprimant les probl\`{e}mes li\'{e}s \`{a} la migration du style de l’UI
de d\'{e}part. En effet l’aspect visuel et la taille des \'{e}l\'{e}ments graphiques sont importants pour d\'{e}crire
159
160
CHAPITRE 8. CONCLUSIONS ET PERSPECTIVES
des interactions pour plusieurs personnes. Concernant cette dimension, les r\'{e}ponses aux probl\`{e}mes
identifi\'{e}s constituent une perspective importante de nos travaux.
Nous concluons cette synth\`{e}se en positionnant de quelle mani\`{e}re les travaux r\'{e}alis\'{e}s dans cette
th\`{e}se permettent de r\'{e}pondre aux questions soulev\'{e}es par l’espace des probl\`{e}mes li\'{e}s \`{a} la migration
des UI d\'{e}crit \`{a} la section 2.2.
En ce qui concerne la dimension des probl\`{e}mes li\'{e}s aux dialogues, les interrogations \'{e}taient :
– Comment utiliser les objets tangibles comme moyen d’interactions ?
Nous pr\'{e}conisons d’associer les objets tangibles \`{a} des groupes d’\'{e}l\'{e}ments graphiques ou \`{a} des
fonctionnalit\'{e}s. Les groupes repr\'{e}sentent des menus, des formulaires ou des panels de modification des contenus. Nous identifions ces groupes (ou les fonctionnalit\'{e}s) \`{a} l’aide des PI qui
permettent de d\'{e}terminer les interactions effectives de chaque \'{e}l\'{e}ment d’un groupe (ou d’une
fonctionnalit\'{e}). L’utilisation des objets tangibles pour afficher (ou cacher) des groupes d’\'{e}l\'{e}ments graphiques ou pour activer des fonctionnalit\'{e}s permet de d\'{e}crire des dialogues plus accessibles. Dans notre prototype, les objets tangibles sont utilis\'{e}s \`{a} l’aide des tags. L’association
avec des objets virtuels se fait manuellement pendant la phase de personnalisation des UI gr\^{a}ce
\`{a} un \'{e}diteur graphique.
– Comment transformer les dialogues pour une plateforme multi-utilisateurs et co-localis\'{e}e ?
La transformation des dialogues pendant la migration des UI desktop vers les tables interactives impactent l’UI et le NF. Notre contribution ne concerne que la transformation de la partie
UI. Nous proposons d’abord les PI qui constituent un mod\`{e}le pour d\'{e}crire les dialogues des
utilisateurs de mani\`{e}re atomique et ind\'{e}pendamment des modalit\'{e}s d’interactions. Les PI nous
permettent de d\'{e}crire des \'{e}quivalences entre les modalit\'{e}s d’interactions des diff\'{e}rentes plateformes. Elles permettent d’assurer que les dialogues de l’UI de d\'{e}part sont pr\'{e}serv\'{e}s par notre
processus de migration gr\^{a}ce aux op\'{e}rateurs d’\'{e}quivalences.
Ensuite pour favoriser la collaboration, nous avons propos\'{e} des r\`{e}gles de transformation de l’UI
qui permettent de supprimer les activit\'{e}s bloquantes (cf section 6.3). Par exemple l’affichage
des bo\^{i}tes de dialogues des UI desktop ne facilite pas la collaboration sur une table interactive.
– Comment assurer la coh\'{e}rence des dialogues apr\`{e}s la migration ?
Pour assurer la coh\'{e}rence des dialogues apr\`{e}s la migration, nous proposons de conserver les
dialogues mono-utilisateur pour les groupes qui d\'{e}crivent des interactions en entr\'{e}e. En effet la
modification d’une donn\'{e}e ou l’appel d’une fonctionnalit\'{e} (gr\^{a}ce \`{a} un menu) sont des interactions en entr\'{e}e dont la migration n\'{e}cessite une modification du NF.
En ce qui concerne la structure des UI par transformation les questions \'{e}taient :
– Comment favoriser l’accessibilit\'{e} des \'{e}l\'{e}ments graphiques ?
Nous proposons au chapitre 6 des r\`{e}gles de substitution et de concr\'{e}tisation qui permettent de
remplacer les groupes d’\'{e}l\'{e}ments par ceux qui ont la capacit\'{e} d’\^{e}tre d\'{e}pla\c{c}ables et utilisable \`{a}
360 degr\'{e}. L’utilisation des objets tangibles comme des moyens d’interactions facilitent l’acc\`{e}s
aux composants graphiques des UI migr\'{e}es sur les tables interactives.
– Comment garantir l’utilisabilit\'{e} des UI migr\'{e}es ?
Pour aborder ce probl\`{e}me, nous avons consid\'{e}r\'{e} dans cette th\`{e}se que les r\`{e}gles de transformation de l’UI de d\'{e}part doivent \^{e}tre bas\'{e}es sur les guidelines pour la migration des UI vers les
tables interactives. Nous avons pour cela affin\'{e} les guidelines \`{a} partir des crit\`{e}res ergonomiques
de conception de Scapin [Sca86] et des sp\'{e}cificit\'{e}s des tables interactives.
De mani\`{e}re concr\`{e}te, les guidelines permettent de transformer la structure des \'{e}l\'{e}ments graphiques en garantissant l’accessibilit\'{e} des groupes d’\'{e}l\'{e}ments graphiques pertinents. Dans ce
cadre les PI et le mod\`{e}le de structure qui sont propos\'{e}s, permettent d’\'{e}tablir des \'{e}quivalences
entre les \'{e}l\'{e}ments de la plateforme de d\'{e}part et ceux des tables interactives.
8.2. PERSPECTIVES
161
Le mod\`{e}le de l’UI que nous avons propos\'{e} dans ce manuscrit d\'{e}crit les PI et la structure des
UI. En ce qui concerne le positionnement et le style des \'{e}l\'{e}ments graphiques, leur migration
est manuelle et se fait par le biais d’un \'{e}diteur graphique pour les UI destin\'{e}es aux tables
interactives.
En ce qui concerne le style des UI par transformation les questions \'{e}taient :
– Comment prendre en compte la taille de la surface d’affichage d’une table interactive ?
L’affinement des crit\`{e}res de conception en consid\'{e}rant les propri\'{e}t\'{e}s li\'{e}es aux nombre d’utilisateurs et \`{a} la taille de l’\'{e}cran des tables interactives a permis d’identifier des contraintes sur
la taille des \'{e}l\'{e}ments graphiques. En effet il est indispensable de d\'{e}crire une borne de valeurs
pour les composants graphiques redimensionnables.
– Comment assurer la coh\'{e}rence globale des styles migr\'{e}s ?
Notre solution permet une introduction manuelle du style des UI de la cible. Les questions li\'{e}es
\`{a} la coh\'{e}rence globale font partie des perspectives.
8.2
Perspectives
Ce travail engendre de nombreuses perspectives \`{a} moyen et \`{a} long terme.
8.2.1
A moyen terme
Nos travaux b\'{e}n\'{e}ficieraient des extensions suivantes \`{a} moyen terme :
– Migration semi automatique du positionnement : Les questions li\'{e}es \`{a} la transformation de
la position des \'{e}l\'{e}ments graphiques sont trait\'{e}es par les personnes en charge de la migration
avec l’aide d’un \'{e}diteur graphique. Nous pensons qu’il est possible d’identifier des guidelines
pour introduire un layout pour les UI des tables interactives pendant les phases de proposition
d’UI et de personnalisation. Cette perspective fait partie des extensions pour am\'{e}liorer l’\'{e}diteur
de notre prototype de migration des UI vers les tables interactives.
– Migration semi automatique du style : elle permet d’am\'{e}liorer la flexibilit\'{e} de notre processus
de migration des UI de d\'{e}part. Nous pensons que l’utilisation d’un mod\`{e}le pour d\'{e}crire la
dimension de style permet de d\'{e}crire des UI homog\`{e}nes et r\'{e}duira davantage le travail des
personnes en charge de la migration. L’enjeu majeur de cette extension de notre processus est la
mani\`{e}re d’interpr\'{e}ter les guidelines li\'{e}es aux styles dans un mod\`{e}le de style pour permettre aux
concepteurs de d\'{e}crire des mod\`{e}les de styles conformes aux crit\`{e}res ergonomiques des tables
interactives.
– Moteur d’apprentissage des transformations : il permettra de capitaliser les choix de substitutions et de concr\'{e}tisations des personnes en charge de la migration en fonction des \'{e}l\'{e}ments
graphiques et des UI dans le but d’am\'{e}liorer les UI propos\'{e}es et de r\'{e}duire les interventions
humaines. L’enjeu de cette am\'{e}lioration est de s\'{e}lectionner les meilleurs parmi les \'{e}l\'{e}ments
\'{e}quivalents, en tenant compte des choix effectu\'{e}s dans les migrations ant\'{e}rieures et aussi en
tenant compte de la configuration de l’UI d’arriv\'{e}e. Cependant, il est possible que les substitutions et les concr\'{e}tisations propos\'{e}es ne soient pas ad\'{e}quates pour une application, dans ce cas
il faut s’assurer que les dialogues de d\'{e}part et les guidelines soient respect\'{e}es. Par ailleurs le
moteur d’apprentissage ne doit pas exclure la phase de personnalisation des UI mais la r\'{e}duire.
8.2.2
A long terme
Nos travaux pourraient b\'{e}n\'{e}ficier des am\'{e}liorations suivantes \`{a} long terme :
162
CHAPITRE 8. CONCLUSIONS ET PERSPECTIVES
– Migration des dialogues : notre solution actuelle ne transforme pas les dialogues monoutilisateur en dialogues multi-utilisateurs. Cette transformation implique une modification de
l’UI et du NF. Pour transformer le NF, il est indispensable de consid\'{e}rer et de d\'{e}crire l’impacte
des dialogues sur le NF. En effet la migration des dialogues d’un formulaire ou d’un menu d’une
UI desktop vers une table interactive implique de permettre \`{a} plusieurs personnes de les utiliser
au m\^{e}me moment. Pour ce faire, il est d’abord indispensable d’identifier les \'{e}l\'{e}ments de l’UI et
du NF qui sont concern\'{e}s et ensuite d’\'{e}crire des r\`{e}gles de transformation qui pr\'{e}servent la coh\'{e}rence des dialogues. L’enjeu de la migration des dialogues dans ce cadre consiste \`{a} modifier
aussi le NF de l’UI de d\'{e}part.
– Autres plateformes de d\'{e}part : notre solution peut \^{e}tre \'{e}tendue pour prendre en compte
d’autres plateformes de d\'{e}part comme une tablette ou un t\'{e}l\'{e}phone portable. Nous avons \'{e}tudi\'{e}
dans notre solution l’accessibilit\'{e} des \'{e}l\'{e}ments graphiques et des interactions dans le cadre de
la migration des UI desktop vers les tables interactives. En consid\'{e}rant d’autres plateformes de
d\'{e}part, comme une tablette par exemple, des enjeux concernant la r\'{e}utilisabilit\'{e} des guidelines,
des mod\`{e}les et des r\`{e}gles de transformation de notre solution actuelle peuvent \^{e}tre d\'{e}gag\'{e}s.
Concernant les guidelines pour la migration des UI, elles sont identifi\'{e}es en consid\'{e}rant les UI
desktops. Dans le cadre des UI pour les tablettes ou les smartphones, il sera indispensable de
v\'{e}rifier si les guidelines actuelles sont r\'{e}utilisables. Une autre plateforme de d\'{e}part permettra
d’\'{e}valuer et de compl\'{e}ter les guidelines que nous avons d\'{e}crites dans cette th\`{e}se.
Par ailleurs, il est int\'{e}ressant d’\'{e}valuer les r\`{e}gles de transformations et le mod\`{e}le que nous
avons propos\'{e} dans cette th\`{e}se pour d’autres plateformes de d\'{e}part. Notre mod\`{e}le d\'{e}crit la
structure et les PI, ce choix ignore le style et le positionnement car nous avons consid\'{e}r\'{e} que
leurs diff\'{e}rences par rapport \`{a} celle des tables interactives ne facilite pas une transformation.
L’extension de notre solution \`{a} d’autres plateformes de d\'{e}part confirmera cette hypoth\`{e}se et
enrichira notre mod\`{e}le d’UI.
– Une autre plateforme d’arriv\'{e}e : dans le but de d\'{e}crire un processus de migration g\'{e}n\'{e}rique
et totalement ind\'{e}pendant d’une plateforme de d\'{e}part et d’arriv\'{e}e, nous pensons que le remplacement des tables interactives par d’autres plateformes permet de d\'{e}crire et tester les processus
d’affinement et d’interpr\'{e}tation des crit\`{e}res ergonomiques en guidelines puis en r\`{e}gles de transformation.
Comme nous pouvons le constater, de nombreuses pistes, \`{a} plus ou moins long terme, peuvent
venir compl\'{e}ter les travaux r\'{e}alis\'{e}s dans cette th\`{e}se.
Bibliographie
[App95]
Apple Computer Inc. Macintosh Human Interface Guidelines. Addison-wesley Publishing, 1995.
[Bar88]
Marie F Barthet. Logiciels interactifs et ergonomie. Paris : Dunod Informatique, 1988.
[BCL+08]
Markus Bischof, Bettina Conradi, Peter Lachenmaier, Kai Linde, Max Meier, Philipp
P\"{o}tzl, and Elisabeth Andr\'{e}. Xenakis : combining tangible interaction with probabilitybased musical composition. In Proceedings of the 2nd international conference on
Tangible and embedded interaction, TEI ’08, pages 121–124, New York, NY, USA,
2008. ACM.
[BCW+06]
Doug A. Bowman, Jian Chen, Chadwick A. Wingrave, John F. Lucas, Andrew Ray,
Nicholas F. Polys, Qing Li, Yonca Haciahmetoglu, Ji-Sun Kim, Seonho Kim, Robert
Boehringer, and Tao Ni. New directions in 3d user interfaces. IJVR, 5(2) :3–14, 2006.
[Bec00]
Kent Beck. Extreme Programming Explained : Embrace Change. Addison-Wesley
Professional, 2000.
[Ber03]
F. Berard. The Magic Table : Computer Vision Based Augmentation of a Whiteboard
for Creative Meetings, 2003.
[Bes10]
Guillaume Besacier. Interactions post-WIMP et applications existantes sur une table
interactive. PhD thesis, UNIVERSITÉ PARIS-SUD 11, 2010.
[BRNB07]
Guillaume Besacier, Ga\'{e}tan Rey, Marianne Najm, and St\'{e}phanie Buisine. Paper Metaphor for Tabletop Interaction Design. In HCII’07 Human Computer Interaction International, pages 758–767, 2007.
[BS08]
Renata Bandelloni and Carmen Santoro. Reverse Engineering Cross-Modal User Interfaces for Ubiquitous Environments. Work, 2008.
[BV02]
Laurent Bouillon and Jean Vanderdonckt. Retargeting Web pages to other computing
platforms with VAQUITA. Ninth Working Conference on Reverse Engineering, 2002.
Proceedings., pages 339–348, 2002.
[CCB+02]
Ga\"{e}lle Calvary, Jo\"{e}lle Coutaz, Laurent Bouillon, Murielle Florins, Quentin Limbourg,
L. Marucci, Fabio Paternò, Carmen Santoro, N. Souchon, David Thevenin, and Jean
Vanderdonckt. CAMELEON Project. Technical report, CAMELEON Project, 2002.
[Cre01]
M. Crease. A Toolkit of Resource-sensitive, Multimodal Widgets. University of Glasgow, 2001.
[D. 06]
D. Heinemeier Hansson. World of Resource, 2006.
[Dev92]
UIMS Tool Developers. A metamodel for the runtime architecture of an interactive
system : the uims tool developers workshop. SIGCHI Bull., 24(1) :32–37, January
1992.
[DL01]
Paul Dietz and Darren Leigh. Diamondtouch : a multi-user touch technology. In Proceedings of the 14th annual ACM symposium on User interface software and technology, UIST ’01, pages 219–226, New York, NY, USA, 2001. ACM.
[FCLD12]
Paternò Fabio, Santoro Carmen, and Spano Lucio Davide. Concur Task Trees (CTT),
2012.
[Fou13]
The Eclipse Foundation. Eclipse modeling framework project (emf), 2013.
163
164
BIBLIOGRAPHIE
[FPV95]
Christelle Farenc, Philippe Palanque, and Jean Vanderdonckt. User interface evaluation : is it ever usable ? Advances in Human Factors/Ergonomics, 20 :329–334, 1995.
[FvDFH90]
James D. Foley, Andries van Dam, Steven K. Feiner, and John F. Hughes. Computer
graphics : principles and practice (2nd ed.). Addison-Wesley Longman Publishing
Co., Inc., Boston, MA, USA, 1990.
[GH95]
Hans-w Gellersen and Hans-W. Gellersen. Modality Abstraction : Capturing Logical
Interaction Design as Abstraction from "User Interfaces for All". In 1st ERCIM Workshop on "User Inter- faces for All". ERCIM, 1995.
[GMAD05]
Derek Glover, David Miller, Doug Averis, and Victoria Door. The interactive whiteboard : a literature survey. Technology, Pedagogy and Education, 14(2) :155–170,
2005.
[GPF10]
Andr\'{e} MP Grilo, Ana CR Paiva, and João Pascoal Faria. Reverse engineering of gui
models for testing. In Information Systems and Technologies (CISTI), 2010 5th Iberian
Conference on, pages 1–6. IEEE, 2010.
[GV08]
Federico Gobbo and Matteo Vaccari. The pomodoro technique for sustainable pace in
extreme programming teams. In Agile Processes in Software Engineering and Extreme
Programming, pages 180–184. Springer, 2008.
[HCT06]
P. Hutterer, B.S. Close, and B.H. Thomas. Supporting Mixed Presence Groupware in
Tabletop Applications. In First IEEE International Workshop on Horizontal Interactive
Human-Computer Systems (TABLETOP ’06), pages 63–70. IEEE, January 2006.
[IU97]
Hiroshi Ishii and Brygg Ullmer. Tangible bits. In Proceedings of the SIGCHI conference on Human factors in computing systems - CHI ’97, pages 234–241, New York,
New York, USA, March 1997. ACM Press.
[JOF11]
C\'{e}dric JOFFROY. Composition d’applications et de leurs interfaces homme-machine
dirig\'{e}e par la composition fonctionnelle. PhD thesis, UNIVERSITÉ DE NICE SOPHIA ANTIPOLIS UFR Sciences, 2011.
[KKM03]
M. Kassoff, D. Kato, and W. Mohsin. Creating GUIs for web services. IEEE Internet
Computing, 7(5) :66–73, September 2003.
[KLL+09]
S\'{e}bastien Kubicki, Sophie Lepreux, Yoann Lebrun, Philippe Santos, Christophe
Kolski, and Jean Caelen. New human-computer interactions using tangible objects :
Application on a digital tabletop with rfid technology. In JulieA. Jacko, editor, HumanComputer Interaction. Ambient, Ubiquitous and Intelligent Interaction, volume 5612
of Lecture Notes in Computer Science, pages 446–455. Springer Berlin Heidelberg,
2009.
[KODPR12]
Andre Kalawa, Occello, Anne-Marie Dery-Pinna, and Michel Riveill. Reusing user
interface across devices with different design guidelines. Knowledge and Systems Engineering, International Conference on, 0 :211–216, 2012.
[KP+88]
Glenn E Krasner, Stephen T Pope, et al. A description of the model-view-controller
user interface paradigm in the smalltalk-80 system. Journal of object oriented programming, 1(3) :26–49, 1988.
[KSM99a]
L Kong, E Stroulia, and B Matichuk. Legacy interface migration : A task-centered
approach. . . . of the 8th International Conference on . . . , 1999.
[KSM99b]
Lanyan Kong, Eleni Stroulia, and Bruce Matichuk. Legacy interface migration : A
task-centered approach. In Proc. of 8 th int. conf. on Human-Computer Interaction
HCI International\^{a}C™99 (Munich, pages 1167–1171, 1999.
BIBLIOGRAPHIE
165
[LL09]
Kee Yong Lim and John B Long. The MUSE method for usability engineering, volume 8. Cambridge University Press, 2009.
[Lon10]
Nguyen Hoang Long. Web Visualization of Trajectory Data using Web Open Source
Visualization Libraries Web Visualization of Trajectory Data using Web Open Source
Visualization Library.
PhD thesis, INTERNATIONAL INSTITUTE FOR GEOINFORMATION SCIENCE AND EARTH OBSERVATION ENSCHEDE, THE NETHERLANDS, 2010.
[LVMB05]
Quentin Limbourg, Jean Vanderdonckt, Benjamin Michotte, and Laurent Bouillon.
USIXML : A Language Supporting Multi-path Development of User Interfaces. In
Ifip International Federation For Information Processing, pages 200–220, 2005.
[McC92]
Carma McClure. The three Rs of software automation : re-engineering, repository,
reusability. Prentice-Hall, Inc., 1992.
[Mic09]
Microsoft. Microsoft Surface User Experience Guidelines, 2009.
[Mic11]
Microsoft. Microsoft Surface 2 Design and Interaction Guide. Technical Report July,
Microsoft, 2011.
[Mic12a]
Microsoft. MSDN XAML, 2012.
[Mic12b]
Microsoft WPF. WPF, 2012.
[Mic12c]
Microsoft XNA. XNA, 2012.
[Mit]
Mitsubishi Electric Research Laboratoriesb. MERL .
[MMP00]
Nikunj R. Mehta, Nenad Medvidovic, and Sandeep Phadke. Towards a taxonomy of
software connectors. In Proceedings of the 22nd international conference on Software
engineering - ICSE ’00, pages 178–187, New York, New York, USA, June 2000. ACM
Press.
[Moo95]
James D. Mooney. Portability and reusability. In Proceedings of the 1995 ACM 23rd
annual conference on Computer science - CSC ’95, pages 150–156, New York, New
York, USA, February 1995. ACM Press.
[Moo96]
Melody M Moore. Rule-based detection for reverse engineering user interfaces. In
Reverse Engineering, 1996., Proceedings of the Third Working Conference on, pages
42–48. IEEE, 1996.
[MPV11]
Gerrit Meixner, Fabio Paternò, and Jean Vanderdonckt. Past , Present , and Future of
Model-Based User Interface Development. i-com, 10 :2–11, 2011.
[MR97]
Melody Moore and Spencer Rugaber. Using knowledge representation to understand
interactive systems. In Program Comprehension, 1997. IWPC’97. Proceedings., Fifth
Iternational Workshop on, pages 60–67. IEEE, 1997.
[MVG06]
Tom Mens and Pieter Van Gorp. A taxonomy of model transformation. Electron. Notes
Theor. Comput. Sci., 152 :125–142, March 2006.
[MVL06]
Jos\'{e} Pascual Molina Massó, Jean Vanderdonckt, and Pascual González López. Direct
manipulation of user interfaces for migration. In Proceedings of the 11th international
conference on Intelligent user interfaces - IUI ’06, page 140, New York, New York,
USA, January 2006. ACM Press.
[NC93]
Laurence Nigay and Jo\"{e}lle Coutaz. A design space for multimodal systems. In Proceedings of the SIGCHI conference on Human factors in computing systems - CHI ’93,
pages 172–178, New York, New York, USA, May 1993. ACM Press.
166
BIBLIOGRAPHIE
[Nig94]
Laurence Nigay. Conception et mod\'{e}lisation logicielles des syst\`{e}mes interactifs : application aux interfaces multimodales. PhD thesis, UNIVERSITÉ JOSEPH FOURIER
- GRENOBLE 1, 1994.
[NM90]
Jakob Nielsen and Rolf Molich. Heuristic evaluation of user interfaces. In Proceedings
of the SIGCHI Conference on Human Factors in Computing Systems, CHI ’90, pages
249–256, New York, NY, USA, 1990. ACM.
[Pat11]
Fabio Paternn¸. Migratory interactive applications for ubiquitous environments. Springer, 2011.
[PMM97]
Fabio Paternò, Cristiano Mancini, and Silvia Meniconi. Concurtasktrees : A diagrammatic notation for specifying task models. In Proceedings of the IFIP TC13 Interantional Conference on Human-Computer Interaction, volume 96, pages 362–369, 1997.
[PRI02]
James Patten, Ben Recht, and Hiroshi Ishii. Audiopad : a tag-based interface for musical performance. In Proceedings of the 2002 conference on New interfaces for musical
expression, NIME ’02, pages 1–6, Singapore, Singapore, 2002. National University of
Singapore.
[PSS09]
Fabio Paternò, Carmen Santoro, and Lucio Davide Spano.
MARIA : A Universal,
Declarative, Multiple Abstraction-Level Language for Service-Oriented Applications
in Ubiquitous Environments.
ACM Transactions on Computer-Human Interaction,
16(4) :1–30, November 2009.
[PZ10]
Fabio Paternò and Giuseppe Zichittella. Desktop-to-mobile web adaptation through
customizable two-dimensional semantic redesign. In Human-Centred Software Engineering, pages 79–94. Springer, 2010.
[Que01]
Whitney Quesenbery. Building a Better Style Guide. In Proceedings of the Usability
Professionals Association, pages 1–10, 2001.
[RBB+95]
L Alperin Resnick, Alex Borgida, Ronald J Brachman, Deborah L McGuinness, Peter F Patel-Schneider, C Isbell, and K Zalondek. CLASSIC Description and Reference
Manual For the COMMON LISP Implementation. AI Principles Research Department,
AT\textbackslash{}&T Bell Laboratories, 1995.
[RCPsC05]
Nicolas Roussel, Olivier Chapuis, Universit\'{e} Paris-sud, and Orsay Cedex. Metisse :
un syst\`{e}me de fen\^{e}trage hautement configurable et utilisable au quotidien. In IHM
’05 : Proceedings of the 17th international conference of the Association Francophone
d’Interaction Homme-Machine, 2005.
[Res13]
Microsoft Research. NUI : Natural User Interface, 2013.
[RFJ08]
Daniel Ratiu, Martin Feilkas, and Jan Jurjens. Extracting Domain Ontologies from
Domain Specific APIs. 2008 12th European Conference on Software Maintenance and
Reengineering, 1 :203–212, 2008.
[RSP11]
Yvonne Rogers, Helen Sharp, and Jenny Preece. Interaction design : beyond humancomputer interaction. Wiley, 2011.
[SC12]
Carlos Eduardo Silva and Jos\'{e} Creissac Campos. Can GUI Implementation Markup
Languages Be Used for Modelling ? In Human-Centered Software Engineering, pages
112–129, 2012.
[Sca86]
Dominique L. Scapin.
Guide ergonomique de conception des interfaces hommeordinateur.
Technical report, Institut National de Recherche en Informatique et en
Automatique, 1986.
BIBLIOGRAPHIE
167
[SVFR04]
Chia Shen, Fr\'{e}d\'{e}ric D. Vernier, Clifton Forlines, and Meredith Ringel. DiamondSpin :
an extensible toolkit for around-the-table interaction. In Proceedings of the 2004 conference on Human factors in computing systems - CHI ’04, pages 167–174, New York,
New York, USA, April 2004. ACM Press.
[SWND03]
Dave Shreiner, Mason Woo, Jackie Neider, and Tom Davis. OpenGL programming
guide : the official guide to learning openGL, version 1.4. Addison-Wesley Professional, 4th ed edition, 2003.
[TC02]
David Thevenin and Jo\"{e}lle Coutaz. Adaptation des IHM : Taxonomies et Archi. Logicielle. In IHM 2002, pages 207–210, 2002.
[TKB78]
Andrew Tanenbaum S., Paul Klint, and Wim Bohm. Guidelines for Software Portability. Software-Practice And Experience, 8(6) :681–698, November 1978.
[TKR10]
Thiago Tonelli, Krzysztof, and Ralf. Swing to swt and back : Patterns for api migration
by wrapping. In Proceedings of the 2010 IEEE International Conference on Software
Maintenance, ICSM ’10, pages 1–10, Washington, DC, USA, 2010. IEEE Computer
Society.
[TS99]
K. Tucker and R. E. Kurt Stirewalt. Model based user-interface reengineering.
In
Proceedings of the Sixth Working Conference on Reverse Engineering, WCRE ’99,
pages 56–, Washington, DC, USA, 1999. IEEE Computer Society.
[UI97]
Brygg Ullmer and Hiroshi Ishii. The metaDESK : Models and Prototypes for Tangible
User Interfaces. In UIST ’97 Proceedings of the 10th annual ACM symposium on User
interface software and technolog, pages 223 – 232, 1997.
[USJ+08]
Brygg Ullmer, Rajesh Sankaran, Srikanth Jandhyala, Blake Tregre, Cornelius Toole,
Karun Kallakuri, Christopher Laan, Matthew Hess, Farid Harhad, Urban Wiggins, and
Shining Sun. Tangible Menus and Interaction Trays : Core tangibles for common physical / digital activities. In TEI ’08, pages 209–212. ACM, 2008.
[Van97]
Jean Vanderdonckt. Conception assist\'{e}e de la pr\'{e}sentation d’une interface hommemachine ergonomique pour une application de gestion hautement interactive. Doctoral
disertation. Facult\'{e}s Universitaire Notre-Dame de la Paix, Namur, 1997.
[vD97]
Andries van Dam.
Post-WIMP user interfaces.
Communications of the ACM,
40(2) :63–67, February 1997.
[vdVNHF11] Bram JJ van der Vlist, Gerrit Niezen, Jun Hu, and Loe MG Feijs. Interaction primitives : Describing interaction capabilities of smart objects in ubiquitous computing
environments. In AFRICON, 2011, pages 1–6. IEEE, 2011.
[VLM+04]
Jean Vanderdonckt, Quentin Limbourg, Benjamin Michotte, Laurent Bouillon, Daniela
Trevisan, and Murielle Florins.
USIXML : a User Interface Description Language
for Specifying Multimodal User Interfaces The Reference Framework used for MultiDirectional UI Development. Language, pages 19–20, 2004.
[W3C03]
W3C. W3C Multimodal Interaction Framework, 2003.
[WB03]
Mike Wu and Ravin Balakrishnan. Multi-finger and whole hand gestural interaction
techniques for multi-user tabletop displays. In Proceedings of the 16th annual ACM
symposium on User interface software and technology - UIST ’03, pages 193–202, New
York, New York, USA, November 2003. ACM Press.
[Weg97]
Peter Wegner. Why Interaction Is More Powerful Than Algorithms, 1997.
168
BIBLIOGRAPHIE
[WGM08]
Xin Wang, Yaser Ghanam, and Frank Maurer. From Desktop to Tabletop : Migrating
the User Interface of AgilePlanner. In Engineering Interactive Systems 2008, pages
263–270, 2008.
Appendices
169
ANNEXE A
Application des mod\`{e}les de l’UI et des PI
Cette annexe a pour objectif de pr\'{e}senter d’abord les tableaux de correspondances pour d\'{e}crire les
\'{e}l\'{e}ments d’une biblioth\`{e}que graphique \`{a} l’aide du mod\`{e}le de type de composants graphiques pr\'{e}sent\'{e}
au chapitre 5. Elle pr\'{e}sente ensuite les ressources utilis\'{e}es par les m\'{e}canismes d’abstraction d’une UI
finale vers les mod\`{e}les d’instance. Enfin, cette annexe pr\'{e}sente l’application des r\`{e}gles d’identification
des PI pour les biblioth\`{e}ques XAML et XAMLSurface.
A.1
Comment d\'{e}crire une instance du mod\`{e}le de type de composants
graphiques ?
Nous d\'{e}crivons dans cette section le processus qui nous a permis d’instancier les Widgets des
biblioth\`{e}ques graphiques XAML Desktop et XAML Surface pour Microsoft PixelSense. Dans l’environnement .Net, les composants graphiques sont identifi\'{e}s \`{a} partir des Assemblies qui appartiennent \`{a}
des DLL (Dynamic Library Link).
La classe Widget du mod\`{e}le de type de composant, correspond \`{a} toutes les classes de l’Assembly
PresentationFramework pour XAML et de l’Assembly Microsoft.Surface.Presentation qui respectent
les conditions suivantes :
– les composants graphiques sont des classes publiques de l’Assembly,
– ils h\'{e}ritent de System.Windows.UIElement.
A.1.1
Identification d’un Widget \`{a} partir d’un Assembly
Les attributs de la classe Widget sont identifi\'{e}s \`{a} partir des composants graphiques identifi\'{e}s dans
un Assembly.
– Widget.name correspond au nom du type du composant graphique : par exemple pour la
classe Microsoft.Surface.Presentation.Controls.LibraryBar de l’Assembly PresentationFramework, Widget.name=’LibraryBar’.
– Widget.cardinality est d\'{e}termin\'{e} \`{a} l’aide du tableau d’identification des cardinalit\'{e}s (cf tableau A.2)
– Widget.contentType est d\'{e}termin\'{e} \`{a} l’aide du tableau d’identification de type de donn\'{e}es (cf
tableau A.1)
A.1.2
Identification des comportements
Les attributs de la classe Event des Widgets sont identifi\'{e}s comme suit :
– l’attribut name
– l’attribut eventType est d\'{e}termin\'{e} en parcourant tous les \'{e}v\'{e}nements et les attributs d’un composant graphique et en utilisant le tableau d’identification des types d’\'{e}v\'{e}nements (cf tableau A.3)
– l’attribut propertyType est identifi\'{e} en recherchant parmi les propri\'{e}t\'{e}s d’un composant graphique et en utilisant le tableau d’identification des types de propri\'{e}t\'{e} (cf tableau A.4)
171
172
ANNEXE A. APPLICATION DES MODÈLES DE L’UI ET DES PI
DataType
XAML et XAML Surface
Boolean
Si le composant graphique impl\'{e}mente l’attribut Checked de type bool\'{e}en
Integer
Si le composant graphique impl\'{e}mente l’attribut Value
String
Si le composant graphique impl\'{e}mente les attributs Text, Password, SelectedValue
ou la m\'{e}thode AddText
Image
Si le composant graphique impl\'{e}mente l’attribut ImageSource
MediaElement
Si le composant graphique impl\'{e}mente un attribut de nom Uri
Object
Si le composant graphique impl\'{e}mente un attribut de nom Content
Widget
Si le composant graphique impl\'{e}mente un attribut de nom Children
Null
Si aucun des cas pr\'{e}c\'{e}dent n’est satisfait
TABLE A.1 – Table d’identification de type de donn\'{e}es du mod\`{e}le
Cardinality
XAML Desktop et XAML Surface
0
Si le type de donn\'{e}es du Widget est Null
1
Si le composant graphique impl\'{e}mente l’attribut Child
N
Si le composant graphique impl\'{e}mente la m\'{e}thode AddChild
N,M
Si le composant graphique impl\'{e}mente les attributs Columns,Rows
TABLE A.2 – Table d’identification des cardinalit\'{e}s
– l’attribut inputDeviceType est identifi\'{e} en utilisant le tableau d’identification des types de propri\'{e}t\'{e} (cf tableau A.5).
A.1.3
Remarque
Cette approche a permis d’instancier les mod\`{e}les de type pour les biblioth\`{e}ques graphiques
XAML Desktop pour la source et XAML Surface pour la cible.
A.2
Abstraction des UI finales
Nous pr\'{e}sentons dans cette section les m\'{e}canismes d’abstraction d’une UI finale XAML vers le
mod\`{e}le d’instance de l’UI d\'{e}crit au chapitre 5.
A.2.1
Identification des UIComponent \`{a} partir d’un fichier XAML
La classe UIComponent est identifi\'{e}e \`{a} partir d’un code source XAML en consid\'{e}rant toutes les
balises qui repr\'{e}sentent les Widgets et qui ne contiennent pas d’autres balises de type Widget.
eventType
XAML Desktop et XAML Surface
Call
Si l’attribut suivant existe : AddHandler
Select
Selected, SelectedText, SelectItem
Change
AllowDrop, TextChanged, OnValueChanged, EditingMode, CanRotate,
LocationChanged, CanMove, ResizeMode, CanScal
TABLE A.3 – Table d’identification des types de comportements
A.2. ABSTRACTION DES UI FINALES
173
propertyType
XAML Desktop et XAML Surface
Content
Si les attributs suivant existent : TextChanged, OnValueChanged,
EditingMode, Selected, SelectedText, SelectItem, Checked, AllowDrop
Size
Si les attributs suivant existent : ResizeMode, CanScale
Position
Si les attributs suivant existent LocationChanged, CanMove
Orientation
Si l’ attribut suivant existe CanRotate
WidgetStructure
Si l’ attribut suivant existe Visible
TABLE A.4 – Table d’identification des types de propri\'{e}t\'{e}s
inputDeviceType
XAML Desktop et XAML Surface
DirectManipulationType
Si les \'{e}v\'{e}nements de la souris, de l’\'{e}cran tactile, du stylo
existent (MouseDown, TouchDown, ContactDown, StylusDown
, etc.)
SequantialManipulationType
Si les \'{e}v\'{e}nements du clavier existent (Key)
Null
Si aucun \'{e}v\'{e}nement li\'{e} \`{a} une dispositif d’entr\'{e}e n’existent
TABLE A.5 – Table d’identification des types de propri\'{e}t\'{e}s
Une balise est de type Widget si son nom correspond \`{a} celui d’un Widget de la biblioth\`{e}que
graphique de d\'{e}part.
Les attributs de la classe container sont identifi\'{e}es comme suit :
– id est un identifiant unique g\'{e}n\'{e}r\'{e} par le processus d’abstraction
– name correspond \`{a} la propri\'{e}t\'{e} Name de la balise relative \`{a} un composant graphique simple
– type correspond au Widget du m\^{e}me nom que la balise
A.2.2
Identification des Containers \`{a} partir d’un fichier XAML
La classe Container est identifi\'{e}e \`{a} partir d’un code source XAML en consid\'{e}rant toutes les balises
qui repr\'{e}sentent les Widgets et qui contiennent d’autres balises de type Widget.
Une balise est de type Widget si son nom correspond \`{a} celui d’un Widget de la biblioth\`{e}que
graphique de d\'{e}part.
Les attributs de la classe container sont identifi\'{e}s comme suit :
– id est un identifiant unique g\'{e}n\'{e}r\'{e} par le processus d’abstraction
– name correspond \`{a} la propri\'{e}t\'{e} Name de la balise relative \`{a} un container
– typeContainer est identifi\'{e} \`{a} partir des \'{e}l\'{e}ments du container
– type correspond au Widget du m\^{e}me nom que la balise.
A.2.3
Identification de la classe Content \`{a} partir d’un fichier XAML
La classe Content, correspond aux donn\'{e}es graphiques d’une UI \`{a} migrer. L’attribut value correspond \`{a} la valeur de la donn\'{e}e graphique. Par exemple pour un bouton “Valider” en XAML, value=’Valider’ et propertyName=’Content’. L’attribut propertyName correspond au nom de la propri\'{e}t\'{e}
correspondante \`{a} la donn\'{e}e. Pour les biblioth\`{e}ques graphiques comme XAML et XAML Surface, le
tableau A.6 permet d’identifier les noms de ces propri\'{e}t\'{e}s.
174
ANNEXE A. APPLICATION DES MODÈLES DE L’UI ET DES PI
XAML Desktop et XAML Surface
propertyName
Content, Value, Text
TABLE A.6 – Table d’identification des propri\'{e}t\'{e}s de contenus
eventType
XAML Desktop et XAML Surface
Change
AllowDrop, TextChanged, OnValueChanged, EditingMode, CanRotate,
LocationChanged, CanMove, ResizeMode, CanScal
Select
Selected, SelectedText, SelectItem
Call
Si les attributs suivant existent : AddHandler
Delete
Methode Remove de la propri\'{e}t\'{e} Items
Update
Methode Add de la propri\'{e}t\'{e} Items, M\'{e}thodes AddChild, AddText de ItemsControl,
Propri\'{e}t\'{e}s Content, Text, DataSource,
TABLE A.7 – Table d’identification des types de comportements
A.2.4
Identification de la classe ImplementedEvent \`{a} partir d’un fichier XAML
Les attributs de la classe Event des Widgets sont identifi\'{e}s comme suit :
– l’attribut name correspond au nom d’un \'{e}v\'{e}nement.
– l’attribut type est d\'{e}termin\'{e} en parcourant tous les \'{e}v\'{e}nements et les attributs d’un composant
graphique et en utilisant le tableau d’identification des types d’\'{e}v\'{e}nements (cf tableau A.7).
– l’attribut property est identifi\'{e} en recherchant parmi les propri\'{e}t\'{e}s d’un composant graphique
et en utilisant le tableau d’identification des types de propri\'{e}t\'{e} (cf tableau A.8).
– l’attribut inputDeviceType est identifi\'{e} en utilisant le tableau d’identification des types de propri\'{e}t\'{e} (cf tableau A.9).
A.2.5
Exemple d’UIStructure XAML
Le listing A.1 d\'{e}crit la structure et le positionnement des \'{e}l\'{e}ments de l’UI de l’application
CBA. Les «event» correspondent \`{a} l’\'{e}v\'{e}nement Click (cf les lignes 19, 44, 47, 50 du listing A.1
par exemple).
Listing A.1 – Exemple UIStructure
<Window x:Class="ExempleApplicationTools2012.Window1"
xmlns="http://schemas.microsoft.com/winfx/2006/xaml/presentation"
xmlns:x="http://schemas.microsoft.com/winfx/2006/xaml"
Title="Window1" Height="326" Width="831">
<Grid >
propertyType
XAML Desktop et XAML Surface
Content
Si les attributs suivants existent : TextChanged, OnValueChanged,
EditingMode, Selected, SelectedText, SelectItem, Checked, AllowDrop
Size
Si les attributs suivants existent : ResizeMode, CanScale
Position
Si les attributs suivants existent LocationChanged, CanMove
Orientation
Si l’attribut suivant existe CanRotate
WidgetStructure
Si l’ attribut suivant existe Visible
TABLE A.8 – Table d’identification des types de propri\'{e}t\'{e}s
A.2. ABSTRACTION DES UI FINALES
175
inputDeviceType
XAML Desktop et XAML Surface
DirectManipulationType
Si les \'{e}v\'{e}nements de la souris, de l’\'{e}cran tactile, du stylo
existent (MouseDown, TouchDown, ContactDown, StylusDown,
etc.)
SequantialManipulationType
Si les \'{e}v\'{e}nements du clavier existent (Key)
Null
Si aucun \'{e}v\'{e}nement li\'{e} \`{a} un dispositif d’entr\'{e}e n’existent
TABLE A.9 – Table d’identification des types de propri\'{e}t\'{e}s
6
<Grid.RowDefinitions>
<RowDefinition Height="30"></RowDefinition>
<RowDefinition Height="258*"></RowDefinition>
</Grid.RowDefinitions>
<Grid.ColumnDefinitions>
<ColumnDefinition Width="73*"></ColumnDefinition>
<ColumnDefinition Width="610*"></ColumnDefinition>
<ColumnDefinition Width="126.03*"></ColumnDefinition>
</Grid.ColumnDefinitions>
<Grid Name="MenuPane" Grid.ColumnSpan="3" Grid.Row="0">
16
<Menu>
<MenuItem Header="File">
<MenuItem Header="Close" Click="Close\textbackslash{}_Click"></MenuItem>
</MenuItem>
<MenuItem Header="Edit"></MenuItem>
</Menu>
</Grid>
<Grid Name="RessourcePane" Grid.Column="0" Grid.Row="1" >
<Grid.RowDefinitions>
26
<RowDefinition></RowDefinition>
<RowDefinition></RowDefinition>
<RowDefinition></RowDefinition>
</Grid.RowDefinitions>
<Grid Name="ControlPane" Grid.Row="0" VerticalAlignment="Top">
<Grid.RowDefinitions>
<RowDefinition Height="20"></RowDefinition>
<RowDefinition Height="15"></RowDefinition>
<RowDefinition Height="*"></RowDefinition>
</Grid.RowDefinitions>
36
<Grid.ColumnDefinitions>
<ColumnDefinition></ColumnDefinition>
<ColumnDefinition></ColumnDefinition>
<ColumnDefinition></ColumnDefinition>
<ColumnDefinition></ColumnDefinition>
<ColumnDefinition></ColumnDefinition>
</Grid.ColumnDefinitions>
<Label Grid.Row="0" Grid.ColumnSpan="5" >Browser</Label>
<Button Name="ScreenShoot" Grid.Column="0" Grid.Row="1" Click="ScreenShoot\textbackslash{}_Click">
<Image Source="./Ressources/screenshot.png"></Image>
46
</Button>
<Button Name="Balloon" Grid.Column="1" Grid.Row="1" Click="Balloon\textbackslash{}_Click">
<Image Source="./Ressources/baloon1.png"></Image>
</Button>
<Button Name="Text" Grid.Column="2" Grid.Row="1" Click="Text\textbackslash{}_Click">
<Image Source="./Ressources/text1.png"></Image>
</Button>
<Button Name="Label" Grid.Column="3" Grid.Row="1">
<Image Source="./Ressources/info.png"></Image>
</Button>
56
<Button Name="Favoris" Grid.Column="4" Grid.Row="1">
<Image Source="./Ressources/info.png"></Image>
</Button>
176
ANNEXE A. APPLICATION DES MODÈLES DE L’UI ET DES PI
</Grid >
<Grid Margin="0,35,0,8" Grid.RowSpan="3">
<ListBox Name="RessourceList" Margin="0,0,0,−15" VerticalAlignment="Top" PreviewMouseLeftButtonDown="
ListBox\textbackslash{}_PreviewMouseLeftButtonDown" >
<ListBox.ContextMenu>
<ContextMenu>
<MenuItem Header="File"></MenuItem>
66
</ContextMenu>
</ListBox.ContextMenu>
</ListBox>
</Grid>
</Grid>
<Grid Grid.Column="1" Grid.Row="1" Background="DarkGray" >
<Grid.RowDefinitions>
<RowDefinition ></RowDefinition>
<RowDefinition ></RowDefinition>
</Grid.RowDefinitions>
76
<Grid.ColumnDefinitions>
<ColumnDefinition ></ColumnDefinition>
<ColumnDefinition ></ColumnDefinition>
</Grid.ColumnDefinitions>
<Canvas Name="canvas1" Grid.Column="0" Grid.Row="0" Margin="8,8,8,8" Background="White" AllowDrop="True"
Drop="Canvas\textbackslash{}_Drop" >
</Canvas>
<Canvas Name="canvas2" Grid.Column="1" Grid.Row="0" Margin="8,8,8,8" Background="White"></Canvas>
<Canvas Name="canvas3" Grid.Column="0" Grid.Row="1" Margin="8,8,8,8" Background="White"></Canvas>
<Canvas Name="canvas4" Grid.Column="1" Grid.Row="1" Margin="8,8,8,8" Background="White"></Canvas>
86
</Grid>
<Grid>
</Grid>
</Grid>
</Window>
A.2.6
Exemple de UIBehavior C\textbackslash{}#
Le listing A.2 d\'{e}crit le comportement des \'{e}l\'{e}ments de l’UI de l’application CBA.
Listing A.2 – Exemple UIBehavior
namespace ExempleApplicationTools2012
\textbackslash{}{
/// <summary>
/// Logique d’interaction pour Window1.xaml
/// </summary>
public partial class Window1 : Window
7
\textbackslash{}{
public Window1()
\textbackslash{}{
InitializeComponent();
\textbackslash{}_controler = new Controler.Controler();
canc = new CanvasClass(this.canvas1, this.canvas2, this.canvas3, this.canvas4);
\textbackslash{}_controler.\textbackslash{}_view.RessourceList = RessourceList;
\textbackslash{}}
private CanvasClass canc = null;
private Controler.Controler \textbackslash{}_controler;
17
private void ListBox\textbackslash{}_DragLeave(object sender, DragEventArgs e)
\textbackslash{}{ \textbackslash{}}
private void Canvas\textbackslash{}_Drop(object sender, DragEventArgs e)
\textbackslash{}{
Canvas parent = (Canvas)sender;
\textbackslash{}_controler.DropCanvas(parent);
\textbackslash{}}
private void ListBox\textbackslash{}_DragEnter(object sender, DragEventArgs e)
A.2. ABSTRACTION DES UI FINALES
177
\textbackslash{}{ \textbackslash{}}
private void ListBox\textbackslash{}_PreviewMouseLeftButtonDown(object sender, MouseButtonEventArgs e)
27
\textbackslash{}{
ListBox parent = (ListBox)sender;
\textbackslash{}_controler.DragOutList(parent, e.GetPosition((ListBox)parent));
\textbackslash{}}
private void Text\textbackslash{}_Click(object sender, RoutedEventArgs e)
\textbackslash{}{
\textbackslash{}_controler.getItemList("text");
\textbackslash{}}
private void Balloon\textbackslash{}_Click(object sender, RoutedEventArgs e)
37
\textbackslash{}{
\textbackslash{}_controler.getItemList("ballon");
\textbackslash{}}
\textbackslash{}}
\textbackslash{}}
A.2.7
Exemple d’AbstractView C\textbackslash{}#
Le listing A.3 d\'{e}crit une vue abstraite de l’application CBA. La m\'{e}thode updateRessourceList (cf
ligne 7 du listing A.3) repr\'{e}sente les fl\`{e}ches de type «updateView» pour les interactions en sortie car
elle permet une notification de la vue par le mod\`{e}le ou le contr\^{o}leur.
Listing A.3 – Exemple AbstractView
using System;
namespace ComicBook.View
\textbackslash{}{
interface IAView
\textbackslash{}{
void canvas\textbackslash{}_Drop(object parent);
void dragOutList(object parent, System.Windows.Point position);
void updateRessourceList(System.Collections.ObjectModel.ObservableCollection<System.Windows.Controls.Image> coll);
9
\textbackslash{}}
\textbackslash{}}
Le listing A.4 est une impl\'{e}mentation de la vue abstraite selon les \'{e}l\'{e}ments de la biblioth\`{e}que graphique. Cette classe est modifi\'{e}e pendant la migration car elle n’est pas ind\'{e}pendante des instruments
d’interactions.
Listing A.4 – Exemple impl\'{e}mentation d’AbstractView
namespace ComicBook.View
\textbackslash{}{
public class AView : ComicBook.View.IAView
\textbackslash{}{
//store the "source" listbox
Object data = null;
CanvasClass canc = null;
public ListBox RessourceList;
public AView()
10
\textbackslash{}{ \textbackslash{}}
public void canvas\textbackslash{}_Drop(Object parent)
\textbackslash{}{
if (data != null)
canc.addBallon((Canvas)parent, (UIElement)data);
\textbackslash{}}
public void updateRessourceList (ObservableCollection<Image> coll)
\textbackslash{}{
RessourceList.ItemsSource = coll;
\textbackslash{}}
20
public void dragOutList(Object parent,Point position)
\textbackslash{}{
178
ANNEXE A. APPLICATION DES MODÈLES DE L’UI ET DES PI
//get the object source for the selected item
object data1 = GetObjectDataFromPoint((ListBox)parent, position);
//if the data is not null then start the drag drop operation
if (data1 != null)
\textbackslash{}{
data = data1;
DragDrop.DoDragDrop((ListBox)parent, data1, DragDropEffects.Copy);
30
\textbackslash{}}
\textbackslash{}}
private static object GetObjectDataFromPoint(ListBox source, Point point)
\textbackslash{}{
UIElement element = source.InputHitTest(point) as UIElement;
if (element != null)
\textbackslash{}{
//get the object from the element
object data = DependencyProperty.UnsetValue;
while (data == DependencyProperty.UnsetValue)
40
\textbackslash{}{
// try to get the object value for the corresponding element
data = source.ItemContainerGenerator.ItemFromContainer(element);
if (data == DependencyProperty.UnsetValue)
element = VisualTreeHelper.GetParent(element) as UIElement;
if (element == source)
return null;
\textbackslash{}}
if (data != DependencyProperty.UnsetValue)
return data;
50
\textbackslash{}}
return null;
\textbackslash{}}
\textbackslash{}}
\textbackslash{}}
A.2. ABSTRACTION DES UI FINALES
179
180
ANNEXE A. APPLICATION DES MODÈLES DE L’UI ET DES PI
A.3
Primitives d’interactions
A.3.1
Liste des Widgets
A.3. PRIMITIVES D’INTERACTIONS
181
182
ANNEXE A. APPLICATION DES MODÈLES DE L’UI ET DES PI
Table des figures
2.1
Applications sur les march\'{e}s de t\'{e}l\'{e}chargement . . . . . . . . . . . . . . . . . . . .
10
2.2
Description de la fen\^{e}tre principale de l’application CBA . . . . . . . . . . . . . . .
11
2.3
Migration de la structure d’une UI pour desktop sur une table interactive . . . . . . .
14
2.4
Espace probl\`{e}mes de la migration des UI vers les tables interactives
. . . . . . . . .
16
3.1
Mod\`{e}le d’interactions instrumentale d’une table interactive . . . . . . . . . . . . . .
24
3.2
Table interactive DiamondTouch . . . . . . . . . . . . . . . . . . . . . . . . . . . .
25
3.3
Instanciation physique des \'{e}l\'{e}ments GUI dans TUI . . . . . . . . . . . . . . . . . .
26
3.4
Instances des containers DiamondSpin . . . . . . . . . . . . . . . . . . . . . . . . .
28
3.5
Exemple de ScatterView
. . . . . . . . . . . . . . . . . . . . . . . . . . . . . . . .
29
3.6
Mod\`{e}le d’interactions abstraites de Gellersen
. . . . . . . . . . . . . . . . . . . . .
32
3.7
Illustration de la propri\'{e}t\'{e} 360˚ . . . . . . . . . . . . . . . . . . . . . . . . . . . . .
39
3.8
Repr\'{e}sentation synth\'{e}tique des guidelines selon les trois dimensions des UI . . . . .
40
3.9
Repr\'{e}sentation synth\'{e}tique des guidelines suivant deux dimensions . . . . . . . . . .
42
4.1
Application de consultation des contacts . . . . . . . . . . . . . . . . . . . . . . . .
48
4.2
Synth\`{e}se de la migration manuelle des UI . . . . . . . . . . . . . . . . . . . . . . .
48
4.3
Synth\`{e}se des approches de portage des UI sur tables interactives . . . . . . . . . . .
52
4.4
Exemple de transformation USIXML
. . . . . . . . . . . . . . . . . . . . . . . . .
54
4.5
Service de migration des UI
. . . . . . . . . . . . . . . . . . . . . . . . . . . . . .
55
4.6
Synth\`{e}se des approches automatiques de migration des UI
. . . . . . . . . . . . . .
57
4.7
Processus de migration avec MORPH . . . . . . . . . . . . . . . . . . . . . . . . .
58
4.8
Synth\`{e}se de l’approche semi automatique
. . . . . . . . . . . . . . . . . . . . . . .
62
4.9
Synth\`{e}se des approches de migration des UI . . . . . . . . . . . . . . . . . . . . . .
63
5.1
Un art\'{e}fact d’une UI
. . . . . . . . . . . . . . . . . . . . . . . . . . . . . . . . . .
72
5.2
Bo\^{i}te de dialogue . . . . . . . . . . . . . . . . . . . . . . . . . . . . . . . . . . . .
73
5.3
Types et instances de composants graphiques
. . . . . . . . . . . . . . . . . . . . .
74
5.4
Mod\`{e}le de types de composants graphiques
. . . . . . . . . . . . . . . . . . . . . .
76
5.5
Mod\`{e}le de structure d’instance d’une UI . . . . . . . . . . . . . . . . . . . . . . . .
82
5.6
Illustration des types de container
. . . . . . . . . . . . . . . . . . . . . . . . . . .
85
5.7
Mod\`{e}le abstrait et UI finale . . . . . . . . . . . . . . . . . . . . . . . . . . . . . . .
90
6.1
Processus de migration assist\'{e}e . . . . . . . . . . . . . . . . . . . . . . . . . . . . .
100
6.2
Crit\`{e}res ergonomiques de conception et guidelines pour favoriser la collaboration . .
102
6.3
Crit\`{e}res ergonomiques de conception et guidelines pour des UI tangibles . . . . . . .
103
6.4
Utilisations effectives des guidelines par les m\'{e}canismes de migration des UI
. . . .
104
6.5
Mod\`{e}le de r\`{e}gles de substitution . . . . . . . . . . . . . . . . . . . . . . . . . . . .
105
6.6
Mod\`{e}le de r\`{e}gles de concr\'{e}tisation . . . . . . . . . . . . . . . . . . . . . . . . . . .
106
6.7
Exemple des groupes d’\'{e}l\'{e}ments graphiques \`{a} transformer
. . . . . . . . . . . . . .
110
6.8
Exemple de ControlGroup
. . . . . . . . . . . . . . . . . . . . . . . . . . . . . . .
112
6.9
Exemple de ControlGroup sur table interactive Micorsoft PixelSense . . . . . . . . .
112
6.10 Exemple de DisplayGroup
. . . . . . . . . . . . . . . . . . . . . . . . . . . . . . .
114
6.11 Exemple de DisplayGroup migr\'{e} sur une table interactive . . . . . . . . . . . . . . .
115
183
184
TABLE DES FIGURES
6.12 Illustration d’un UpdateGroup . . . . . . . . . . . . . . . . . . . . . . . . . . . . .
117
6.13 Exemple d’un UpdateGroup sur une table interactive
. . . . . . . . . . . . . . . . .
117
6.14 Repr\'{e}sentation de l’UI CBA pour tables interactives . . . . . . . . . . . . . . . . . .
118
6.15 Mod\`{e}le UI CBA . . . . . . . . . . . . . . . . . . . . . . . . . . . . . . . . . . . . .
130
7.1
Processus semi automatique de migration d’une UI vers les tables interactives . . . .
138
7.2
Architecture des applications migr\'{e}es vers les tables interactives . . . . . . . . . . .
139
7.3
Exemple de cadre de s\'{e}lection des \'{e}l\'{e}ments graphique dans l’\'{e}diteur du prototype . .
146
7.4
Exemple de substitution de LibraryBar par SurfaceListBox . . . . . . . . . . . . . .
147
7.5
Menu d’association d’un tag avec un container . . . . . . . . . . . . . . . . . . . . .
147
7.6
Éditeur graphique du prototype . . . . . . . . . . . . . . . . . . . . . . . . . . . . .
148
7.7
Fiche d’\'{e}valuation de la migration des applications
. . . . . . . . . . . . . . . . . .
151
7.8
Application agenda . . . . . . . . . . . . . . . . . . . . . . . . . . . . . . . . . . .
151
7.9
Regroupement des \'{e}l\'{e}ments de l’application Album Photo
. . . . . . . . . . . . . .
152
7.10 Regroupement des \'{e}l\'{e}ments de l’application de dessin
. . . . . . . . . . . . . . . .
153
7.11 Regroupement des \'{e}l\'{e}ments de l’application de calculatrice . . . . . . . . . . . . . .
154
7.12 Synth\`{e}se de l’\'{e}tude de la migration des quatre applications . . . . . . . . . . . . . .
155
Liste des tableaux
3.1
Synth\`{e}se des dispositifs physiques d’interactions pour 3 tables interactives . . . . . .
27
3.2
Type d’UI pour les tables interactives
. . . . . . . . . . . . . . . . . . . . . . . . .
35
3.3
Affinement des crit\`{e}res ergonomiques de Scapin . . . . . . . . . . . . . . . . . . . .
37
3.4
Migration de l’aspect visuel d’un formulaire . . . . . . . . . . . . . . . . . . . . . .
42
4.1
Synth\`{e}se des technologies de portage des UI sur tables interactives . . . . . . . . . .
51
4.2
Table d’\'{e}quivalences
. . . . . . . . . . . . . . . . . . . . . . . . . . . . . . . . . .
56
4.3
R\'{e}capitulatif de MORPH . . . . . . . . . . . . . . . . . . . . . . . . . . . . . . . .
61
5.1
Exemples de primitives d’interactions en entr\'{e}e . . . . . . . . . . . . . . . . . . . .
72
5.2
Exemple de primitives d’interactions en sortie . . . . . . . . . . . . . . . . . . . . .
74
5.3
Types de container
. . . . . . . . . . . . . . . . . . . . . . . . . . . . . . . . . . .
84
6.1
Caract\'{e}ristiques des ControlGroups
. . . . . . . . . . . . . . . . . . . . . . . . . .
110
6.2
R\`{e}gle de substitution des groupes de contr\^{o}le
. . . . . . . . . . . . . . . . . . . . .
111
6.3
R\`{e}gle de concr\'{e}tisation des groupes de contr\^{o}le . . . . . . . . . . . . . . . . . . . .
111
6.4
Caract\'{e}ristiques des groupes d’affichage de contenus
. . . . . . . . . . . . . . . . .
113
6.5
R\`{e}gle de substitution des groupes DisplayGroups . . . . . . . . . . . . . . . . . . .
114
6.6
Caract\'{e}ristiques des groupes de modification de contenus . . . . . . . . . . . . . . .
115
6.7
R\`{e}gle de substitution des UpdateGroups . . . . . . . . . . . . . . . . . . . . . . . .
116
6.8
Caract\'{e}ristiques des groupes mixtes
. . . . . . . . . . . . . . . . . . . . . . . . . .
117
6.9
Synth\`{e}se des probl\`{e}mes trait\'{e}s par les transformations des groupes . . . . . . . . . .
120
6.10 Caract\'{e}ristiques des interacteurs de donn\'{e}es en sortie . . . . . . . . . . . . . . . . .
121
6.11 Caract\'{e}ristiques des interacteurs de donn\'{e}es en entr\'{e}e et en sortie
. . . . . . . . . .
122
6.12 Caract\'{e}ristiques des interacteurs d’activation
. . . . . . . . . . . . . . . . . . . . .
122
6.13 Synth\`{e}se des probl\`{e}mes trait\'{e}s par les transformations des interacteurs . . . . . . . .
123
6.14 Poids des PI par rapport aux guidelines . . . . . . . . . . . . . . . . . . . . . . . . .
126
6.15 Poids des types de donn\'{e}es par rapport aux guidelines . . . . . . . . . . . . . . . . .
127
6.16 Co\^{u}t des interventions manuelles . . . . . . . . . . . . . . . . . . . . . . . . . . . .
129
6.17 Widgets \'{e}quivalents . . . . . . . . . . . . . . . . . . . . . . . . . . . . . . . . . . .
131
A.1
Table d’identification de type de donn\'{e}es du mod\`{e}le . . . . . . . . . . . . . . . . . .
172
A.2
Table d’identification des cardinalit\'{e}s
. . . . . . . . . . . . . . . . . . . . . . . . .
172
A.3
Table d’identification des types de comportements . . . . . . . . . . . . . . . . . . .
172
A.4
Table d’identification des types de propri\'{e}t\'{e}s . . . . . . . . . . . . . . . . . . . . . .
173
A.5
Table d’identification des types de propri\'{e}t\'{e}s . . . . . . . . . . . . . . . . . . . . . .
173
A.6
Table d’identification des propri\'{e}t\'{e}s de contenus . . . . . . . . . . . . . . . . . . . .
174
A.7
Table d’identification des types de comportements . . . . . . . . . . . . . . . . . . .
174
A.8
Table d’identification des types de propri\'{e}t\'{e}s . . . . . . . . . . . . . . . . . . . . . .
174
A.9
Table d’identification des types de propri\'{e}t\'{e}s . . . . . . . . . . . . . . . . . . . . . .
175
185
186
LISTE DES TABLEAUX
R\'{e}sum\'{e}
Dans le domaine du g\'{e}nie logiciel pour les Interactions Homme Machine (IHM), la migration des interfaces
utilisateurs (UI) est un moyen pour r\'{e}utiliser des applications sur des plateformes ayant des modalit\'{e}s d’interactions diff\'{e}rentes des environnements de d\'{e}part. Les approches existantes de migration des UI sont manuelles
dans le cadre des approches sp\'{e}cifiques, elles sont automatiques dans le cadre des services d’adaptation des UI
aux contextes d’usage, ou elles sont semi automatiques dans le cadre d’une migration flexible dirig\'{e}e par un
concepteur.
Dans cette th\`{e}se nous nous int\'{e}ressons \`{a} la migration semi automatique des UI vers une cible comme une
table interactive dans l’objectif de transformer des UI Desktop en UI qui favorisent la collaboration et l’utilisation des objets tangibles. Les tables interactives sont des plateformes qui disposent des instruments d’interactions permettant de d\'{e}crire des UI tangibles et multi-utilisateurs. En consid\'{e}rant que le noyau fonctionnel
(NF) des applications de d\'{e}part peut \^{e}tre r\'{e}utilis\'{e}s sur les cibles sans changement, les UI des applications sont
caract\'{e}ris\'{e}es par la dimension des dialogues entre les utilisateurs et le syst\`{e}me, la dimension de la structure et
du positionnement des \'{e}l\'{e}ments graphiques et la dimension du style des \'{e}l\'{e}ments visuels. La migration d’une
UI dans ces conditions consiste \`{a} transformer ou \`{a} recr\'{e}er les diff\'{e}rentes dimensions d’une UI de d\'{e}part pour
la cible tout en consid\'{e}rant les crit\`{e}res de conception des UI pour les tables interactives.
Nous proposons dans cette th\`{e}se un mod\`{e}le d’interactions abstraites pour \'{e}tablir les \'{e}quivalences entre les
dialogues et la structure des UI ind\'{e}pendamment des modalit\'{e}s d’interactions des plateformes source et cible.
Les primitives d’interactions et la structure des composants graphiques permettent de d\'{e}crire des op\'{e}rateurs
d’\'{e}quivalences pour retrouver et classer les \'{e}l\'{e}ments graphiques \'{e}quivalents en prenant en compte les guidelines des tables interactives. Nous proposons aussi des r\`{e}gles de substitution et de concr\'{e}tisation pour accro\^{i}tre
l’accessibilit\'{e} des \'{e}l\'{e}ments graphiques et favoriser l’utilisation des objets tangibles.
Mots cl\'{e}s :
migration des interfaces utilisateurs, \'{e}quivalences des plateformes, modalit\'{e}s d’interactions,
crit\`{e}res de conception, guidelines
Abstract
In software engineering, in the field of human computer interaction (HCI), the migration of user interface
(UI) is a way to reuse existing applications on platforms with different interactions modalities. The existing
approaches for UI migration can be manual (for specific applications), they can be automatic (for services
which adapt UI based on context aware), or they can be mix of the previous - semi automatic (providing a
flexible migration process driven by the person in charge).
This thesis proposes a semi automatic process for migration of UI from a desktop to interactive table for the
purpose of transforming the UI of desktop to support further collaboration and usage of tangible objects. The
interactive tables are platforms with interactions instruments which allow the describtion of tangible and multi
users UIs. Considering that the functional core (FC) of source applications can be reused on target platform
without transformation, any UI can be characterized with three dimensions : the first dimension concerns the
dialogues between the users and the system, the second dimension concerns the structure and the layout of
graphical components, and the third dimension concerns the visual style of graphical elements. In this context,
the problematic regarding the UI migration is how to transform or re inject these different dimensions of source
UI into the target, while considering the UI design criteria for interactive tables.
This thesis proposes an abstract interactions model for establishing equivalences (independent of modalities
of interactions) between the source and the dialogue and structure of the target. The primitives of interaction and
the structure of graphical components are used to describe equivalence operators to find and to rank equivalent
elements on interactive tables. Furthermore, this thesis proposes substitution and concretization rules to increase
the accessibility of graphical elements and to facilitate the usage of tangible objects. The ranking process and
the transformation rules are based on guidelines for UI migration to interactive tables which are interpreted
form design criteria.
Keywords :
user interface migration, \'{e}quivalence of plateform, design criteria, guidelines


\end{document}
